\documentclass[11pt,openany]{report}

\usepackage{fontspec}
\usepackage{amsmath,amssymb,amsfonts}
\usepackage{amsthm}
\usepackage[margin=1in,headheight=17pt,headsep=14pt]{geometry}
\usepackage{booktabs,longtable,array}
\usepackage{graphicx}
\usepackage{fancyhdr}
\usepackage{titlesec}
\usepackage{microtype}
\usepackage[dvipsnames]{xcolor}
\usepackage[colorlinks=true,linkcolor=NavyBlue,citecolor=NavyBlue,urlcolor=Mahogany]{hyperref}
\makeatletter
\@ifundefined{IfFontExistsTF}{\newcommand{\IfFontExistsTF}[3]{#3}}{}
\makeatother
\IfFontExistsTF{TeX Gyre Pagella}{\setmainfont{TeX Gyre Pagella}}{%
  \IfFontExistsTF{texgyrepagella-regular.otf}{\setmainfont{texgyrepagella-regular.otf}}{}}
\IfFontExistsTF{TeX Gyre Heros}{\setsansfont{TeX Gyre Heros}}{}
\IfFontExistsTF{Latin Modern Mono}{\setmonofont{Latin Modern Mono}[Scale=0.92]}{%
  \IfFontExistsTF{lmmono10-regular.otf}{\setmonofont{lmmono10-regular.otf}[Scale=0.92]}{}}
% --- amsthm setup -----------------------------------------------------------
% \theoremstyle{plain}-family style for theorem-like environments and
% \theoremstyle{definition}-family style for definition-like environments.
% The bodies are set upright (a `plain` italic body would italicise the
% corpus's extensive \text{...} proof justifications); head formatting and
% numbering follow the plain / definition styles respectively.
\newtheoremstyle{canonplain}%
  {7pt plus 2pt minus 1pt}{7pt plus 2pt minus 1pt}%
  {\normalfont}{0pt}{\bfseries}{.}%
  {0.6em plus 0.2em}%
  {\thmname{#1}\thmnumber{ #2}\thmnote{\normalfont\ (#3)}}
\newtheoremstyle{canondef}%
  {7pt plus 2pt minus 1pt}{7pt plus 2pt minus 1pt}%
  {\normalfont}{0pt}{\bfseries}{.}%
  {0.6em plus 0.2em}%
  {\thmname{#1}\thmnumber{ #2}\thmnote{\normalfont\ (#3)}}
\theoremstyle{canonplain}
\newtheorem{theorem}{Theorem}[chapter]
\newtheorem{corollary}{Corollary}[chapter]
\newtheorem{axiom}{Axiom}[chapter]
\newtheorem{convention}{Convention}[chapter]
\theoremstyle{canondef}
\newtheorem{definition}{Definition}[chapter]
\newenvironment{movedblock}
  {\par\medskip\begingroup\small\itshape\leftskip2.2em\rightskip2.2em}
  {\par\endgroup\medskip}
% run-in heading for theorem-class blocks that contain tables or verbatim
% material (longtable/verbatim cannot live inside amsthm environments); the
% printed form matches the environment heads and the counter is shared.
\newcommand{\pseudohead}[3]{\par\medskip\noindent\refstepcounter{#1}%
  \textbf{\csname the#1\endcsname\hspace{0.6em}(#2\hspace{0.35em}---\hspace{0.35em}#3).}%
  \par\nobreak\smallskip}
% ----------------------------------------------------------------------------
\pagestyle{fancy}
\fancyhf{}
\fancyhead[LE]{\small\thepage}
\fancyhead[RO]{\small\thepage}
\fancyhead[RE]{\nouppercase{\small\leftmark}}
\fancyhead[LO]{\nouppercase{\small\rightmark}}
\renewcommand{\headrulewidth}{0.4pt}
\setlength{\parskip}{2pt plus 1pt}
\renewcommand{\arraystretch}{1.08}
\fancypagestyle{plain}{\fancyhf{}\fancyfoot[C]{\thepage}}
\titleformat{\chapter}[hang]{\normalfont\Large\bfseries}{\thechapter}{1em}{}
\titleformat{\section}{\normalfont\large\bfseries}{\thesection}{1em}{}
\titleformat{\subsection}{\normalfont\normalsize\bfseries}{\thesubsection}{1em}{}
\titlespacing{\subsection}{0pt}{8pt plus 2pt minus 1pt}{3pt}
\setcounter{tocdepth}{1}
\setcounter{secnumdepth}{3}
\allowdisplaybreaks
\emergencystretch=3em
\hbadness=10000
\vbadness=10000
\hfuzz=2pt
\begin{document}


\begin{titlepage}
\centering
\vspace*{2.4cm}
{\LARGE $\langle qu(b(i(t|t))ick)\rangle$ — the four-era canon\par}
\vspace{0.8cm}
{\Huge\bfseries Appendix B\par}
\vspace{0.5cm}
{\LARGE Structural Constraints, Refutations,\\ and No-Go Theorems\par}
\vspace{1.0cm}
{\large Addendum to the four-era canon, Volume I\par}
\vspace{0.4cm}
{\large The adjudication record, the refuted hypotheses, the no-go theorems,\\ and the trap and phantom-gap register\par}
\vfill
{\large Edition of 2026-10-09\par}
\vspace{0.3cm}
{\small Entries preserve the canonical tags and the full Thm\,$\cdot$\,Req\,$\cdot$\,Pf\,$\cdot$\,Phys structure; the mathematics is transplanted unmodified.\par}
\end{titlepage}

\tableofcontents

\section*{\texorpdfstring{Scope of this Addendum}{Scope of this Addendum}}
\addcontentsline{toc}{section}{Scope of this Addendum}

This Addendum is Appendix B of the $\langle qu(b(i(t|t))ick)\rangle$ canon
(Volume I). It
collects the corpus's structural-constraints record: category-error and status
adjudications (B.1), refuted hypotheses with their refutations (B.2), no-go
theorems (B.3), and the trap and phantom-gap register (B.4). The material
falls into three kinds of entries, marked at their heads:

\begin{itemize}
\item \textbf{Record entries} reproduce the adjudication blocks moved out of
  the main text (from Parts 2--4) and the adjudication record of the archival
  ledger (HistLog \S IV: the tier adjudication summaries FT.115, FT.121,
  FT.132, FT.140 and the closure tier F.43, FT.146--FT.152). Each keeps its
  canonical tag and its full \textbf{Thm}\,$\cdot$\,\textbf{Req}\,%
  $\cdot$\,\textbf{Pf}\,$\cdot$\,\textbf{Phys} structure, so every citation of
  these tags elsewhere in the corpus resolves here. Together these records
  trace the author's exploration pathway. They are reproduced as they stand in
  the revision-1 canon; the byte-exact pre-revision originals are preserved in
  \texttt{working\_markdowns/}.
\item \textbf{Catalog entries} record the hypothesis-level ruling of a theorem
  whose derived content is consumed by later derivations and therefore remains
  in the main corpus; the entry quotes the ruling verbatim and points to the
  main-text number.
\item \textbf{Pointer entries} (B.4) cross-reference trap disclosures that
  travel inside larger records.
\end{itemize}

Nothing in this Addendum alters the mathematics: the machine-verified canon in
\texttt{canon/} (pre-revision state frozen in \texttt{working\_markdowns/}) is the
authority, and the verification suite (\texttt{verification/run\_all.py}) checks it. Where an adjudication
or no-go stood inline in the main corpus, a stub of the form ``see Appendix B,
Section B.$x$.$y$'' marks the former locus.

\section*{\texorpdfstring{B.1 Category Errors}{B.1 Category Errors}}
\addcontentsline{toc}{section}{B.1 Category Errors}
\markboth{B.1 Category Errors}{B.1 Category Errors}
\subsection*{\texorpdfstring{B.1.1\quad FT.115 — The adjudication summary}{FT.115 B.1.1}}
\addcontentsline{toc}{subsection}{B.1.1 FT.115}
\phantomsection\label{tag:FT-115}
{\small \textit{Record.}\quad Source: HistLog §IV.1 (archival ledger, preserved at \texttt{working\_markdowns/HistLog.md}); reproduced verbatim, including its historical registry tokens, which resolve via HistLog \S I.}
\medskip
\textbf{Thm.} Provenance of every quantity of Tier F.37:

%tbl-level:0
{\setlength{\tabcolsep}{4pt}
\begin{longtable}{>{\raggedright\arraybackslash}p{\dimexpr0.4412\textwidth-2\tabcolsep\relax}|>{\raggedright\arraybackslash}p{\dimexpr0.1324\textwidth-2\tabcolsep\relax}|>{\raggedright\arraybackslash}p{\dimexpr0.4265\textwidth-2\tabcolsep\relax}}
\toprule
Statement & Source & Status \\
\addlinespace[2pt]
\midrule
\endfirsthead
\toprule
Statement & Source & Status \\
\addlinespace[2pt]
\midrule
\endhead
\bottomrule
\endlastfoot
$[J_i,\allowbreak K_j]=\allowbreak i(\delta_{ij}\Sigma-E^{ij})$ & CT1.7 & Derived (frozen) \\
$\sum_i[J_i,\allowbreak K_i]=\allowbreak 2i\Sigma$ & FT.109 & Derived \\
$\dim F_{J^\pm}=\allowbreak 3$ & CT0.12 & Derived (frozen) \\
$\operatorname{tr}T(1)=\allowbreak 3$ & FT.68(2) & Derived (frozen) \\
Tensor product $\mathcal H_{\mathrm{seal}}\otimes\mathcal H_{\mathrm{poly}}$ & FT.110 & Declared (P-H16a) \\
$C_3$ spectrum form $A+\allowbreak B\cos(\theta+\allowbreak 2\pi i/3)$ & FT.111 & Derived (representation theory) \\
$A=\allowbreak 1$, $B=\allowbreak \sqrt2$ & FT.111 & Derived (normalization) \\
$\delta=\allowbreak 2/9$ & FT.112 & Derived \\
Mass ratios $r_i^2$ & FT.113 & Derived \\
Identification with lepton masses & FT.114 & NOT derived (audit only) \\
Overall scale $\mu$ & — & Free parameter (P4) \\
\end{longtable}
}

\textbf{Req.} FT.109–FT.114; CT1.7, CT0.12 (\textbf{B}); FT.68(2) (H09); FT.73(5) (H09, pattern); P1, P4, P5 (\textbf{A}); machine: \texttt{checks/\allowbreak{}verify\_\allowbreak{}H\_\allowbreak{}koide\_\allowbreak{}channel.py} §A–§H.

\textbf{Pf.}

%dsp-level:0
\[
\text{each row is the adjudication of its own tag: FT.109 (§A–§B), FT.110 (§C), FT.111 (§D), FT.112 (§E), FT.113 (§F), FT.114 (§G–§H)}\ [\text{machine: all sections green}]
\]

%dsp-level:0
\[
\text{frozen rows re-derived at the point of use, not re-opened}\ [\mathrm{CT1.7},\ \mathrm{CT0.12},\ \mathrm{FT.68}(2)];
\]

%dsp-level:0
\[
\qquad \text{the declared row is logged}\ [\text{P-H16a},\ \mathrm{H13}]
\]

%dsp-level:0
\[
\mu:\ \neg\exists\ \text{tag fixing the overall scale of}\ m_i^2
\]

%dsp-level:0
\[
= \mu\,
\]

%dsp-level:0
\[
r_i^2\ [\text{P4: free internal configuration parameter}];
\]

%dsp-level:0
\[
\qquad \text{[EXT] rows: audit-only by FT.114}\ [\text{P1},\ \text{P5}].\ \blacksquare
\]

\textbf{Phys.} The tier derives a three-line dimensionless hierarchy from three integers of the frozen algebra — $2$ (the diagonal seal coefficient), $3$ (the seal dimension), $3$ (the transfer trace) — with one declared interface (P-H16a) and zero fitted parameters. The only external data in the tier is audited, quantified by exact shortfall factors, and left unidentified; the overall scale is left free.

\textbf{Status.} Closed. Tier F.37 adds no free parameter beyond the P4-class overall scale; P-H16a is the tier's sole declared interface; no particle is identified anywhere in it. Machine: §A–§H.

\vspace{2pt}\noindent\rule{\textwidth}{0.4pt}\vspace{2pt}

\vspace{2pt}\noindent\rule{\textwidth}{0.4pt}\vspace{2pt}

\medskip
\subsection*{\texorpdfstring{B.1.2\quad FT.121 — The adjudication summary}{FT.121 B.1.2}}
\addcontentsline{toc}{subsection}{B.1.2 FT.121}
\phantomsection\label{tag:FT-121}
{\small \textit{Record.}\quad Source: HistLog §IV.1 (archival ledger, preserved at \texttt{working\_markdowns/HistLog.md}); reproduced verbatim, including its historical registry tokens, which resolve via HistLog \S I.}
\medskip
\textbf{Thm.} Provenance of every quantity of Tier F.38:

%tbl-level:1
{\small\setlength{\tabcolsep}{4pt}
\begin{longtable}{>{\raggedright\arraybackslash}p{\dimexpr0.4845\textwidth-2\tabcolsep\relax}|>{\raggedright\arraybackslash}p{\dimexpr0.1596\textwidth-2\tabcolsep\relax}|>{\raggedright\arraybackslash}p{\dimexpr0.3559\textwidth-2\tabcolsep\relax}}
\toprule
Statement & Source & Status \\
\addlinespace[2pt]
\midrule
\endfirsthead
\toprule
Statement & Source & Status \\
\addlinespace[2pt]
\midrule
\endhead
\bottomrule
\endlastfoot
$S_3\hookrightarrow U(3)$ slot relabelling & CT3.1–CT3.3 & Derived (frozen) \\
Subgroup lattice, orders $(1,\allowbreak 2,\allowbreak 2,\allowbreak 2,\allowbreak 3,\allowbreak 6)$ & FT.116(1) & Derived \\
Commutant dimensions $(9,\allowbreak 5,\allowbreak 3,\allowbreak 2)$ & FT.116(2) & Derived \\
$S_3$ two-level obstruction; reflection $\iff\allowbreak  z\in\allowbreak \mathbb R$ & FT.116(4)–(5) & Derived \\
$\varepsilon$-selection: $\operatorname{Stab}_{S_3}(\varepsilon)=\allowbreak A_3\cong C_3$ & FT.116(6) & Derived \\
Seal-side census ($A_4$ scalars, $V_4$ diagonal, $C_3$ circulant) & FT.116(7) & Derived \\
One-body blindness: baryon $3a$, meson $0$ & FT.117(1)–(3) & Derived \\
$27=\allowbreak 1\oplus8\oplus8\oplus10$; invariant lines $1$-dimensional & FT.117(5)–(6) & Derived \\
$\operatorname{tr}T(1)=\allowbreak 3$; $\operatorname{charpoly}=\allowbreak x^2(x^2-3x-3)$ & FT.68(2) & Derived (frozen) \\
Colour factor $3=\allowbreak \log_2C_{\mathrm{cap}}-1$ & CT3.20 & Derived (frozen) \\
Tensor product $\mathcal H_{\mathrm{seal}}\otimes\mathcal H_{\mathrm{poly}}\otimes\mathcal H_{\mathrm{slot}}$ & FT.118(2) & Declared (P-H17a) \\
Spectator numerator $2$ & FT.118(4)–(5) & Derived \\
Adjoint-dressed numerator $8/3$ & FT.118(6) & Computed, NOT adopted \\
$A=\allowbreak 1$, $B=\allowbreak \sqrt2$ & FT.111(3) & Derived (frozen) \\
$\delta_q=\allowbreak 2/27$; $3\delta_q=\allowbreak \delta_\ell=\allowbreak 2/9$ & FT.119(1)–(2) & Derived \\
Mass ratios $1:3.563695074946\ldots:139.3386880284\ldots$ & FT.119(5) & Derived \\
$Q=\allowbreak \tfrac23$ (phase-free invariant) & FT.119(4) & Derived \\
Heavy-triplet correspondence & FT.120 & NOT derived (audit only) \\
Hadron-level triplets & FT.120 & Excluded (range + confinement blindness) \\
Light triplet at $(A,\allowbreak B)=\allowbreak (1,\allowbreak \sqrt2)$ & FT.120(6) & Excluded (obstruction O3) \\
Identification with quark masses & FT.120 & NOT derived (audit only) \\
Overall scale $\mu$ & — & Free parameter (P4) \\
\end{longtable}
}

\textbf{Req.} FT.116–FT.120; CT3.1, CT3.2, CT3.3, CT3.13, CT3.15, CT3.17, CT3.18, CT3.20, CT3.22 (\textbf{F}); CT0.5, CT0.6, CT0.11, CT0.12 (\textbf{B}); D3.8, D3.9, P1, P4, P5 (\textbf{A}); EW.32 (\textbf{D}); FT.68(2) (H09); FT.77(3) (H10); FT.109, FT.110, FT.111, FT.112, FT.113, FT.114 (\textbf{H16}); machine: \texttt{checks/\allowbreak{}verify\_\allowbreak{}H\_\allowbreak{}colour\_\allowbreak{}channel.py} §A–§F.

\textbf{Pf.}

%dsp-level:0
\[
\text{each row is the adjudication of its own tag: FT.116}\ (\S\mathrm{A}),
\]

%dsp-level:0
\[
\ \text{FT.117}\ (\S\mathrm{B}),
\]

%dsp-level:0
\[
\ \text{FT.118}\ (\S\mathrm{C}),
\]

%dsp-level:0
\[
\ \text{FT.119}\ (\S\mathrm{D}),
\]

%dsp-level:0
\[
\ \text{FT.120}\ (\S\mathrm{E})\ [\text{machine: all sections green}]
\]

%dsp-level:0
\[
\text{frozen rows re-derived at the point of use, not re-opened}\ [\mathrm{CT3.1\text{–}CT3.3},\ \mathrm{CT3.15},\ \mathrm{CT3.20},\ \mathrm{CT0.12},\ \mathrm{FT.68}(2),\ \mathrm{FT.109},\ \mathrm{FT.110},\ \mathrm{FT.111}(3)];
\]

%dsp-level:0
\[
\qquad \text{the declared row is logged}\ [\text{P-H17a},\ \mathrm{H13}]
\]

%dsp-level:0
\[
\mu:\ \neg\exists\ \text{tag fixing the overall scale of}\ m_i
\]

%dsp-level:0
\[
= \mu\,
\]

%dsp-level:0
\[
(r^{(q)}_i)^2\ [\text{P4: free internal configuration parameter}];
\]

%dsp-level:0
\[
\qquad [\text{EXT}]\ \text{rows: audit-only by FT.120}\ [\text{P1},\ \text{P5}].\ \blacksquare
\]

\textbf{Phys.} Tier F.38 answers the docket's three questions with one new declaration and no fitted parameter. The group is not chosen: the Levi-Civita seal that makes three triplets observable preserves exactly the even relabellings, and the full relabelling group is proven incapable of carrying a three-line family at all. The hadron branch is not set aside: it is closed by a theorem — sealed singlets read the trace and nothing else — so the channel relocates to the unsealed triplet-attached chain, whose arena contributes the single extra factor of three that separates $\delta_q=2/27$ from $\delta_\ell=2/9$. The family that falls out is therefore neither a baryon multiplet nor a meson multiplet but the triplet-attachment sector itself, with two exclusions proven on the way and one external hierarchy audited at the ratio level, its shortfalls exact and nonzero. The overall scale is left free, as in Tier F.37.

\textbf{Status.} Closed. Tier F.38 adds no free parameter beyond the P4-class overall scale; P-H17a is the tier's sole declared interface; the adjoint-dressed numerator is logged and unused; no particle is identified anywhere in it. Machine: §A–§F.

\vspace{2pt}\noindent\rule{\textwidth}{0.4pt}\vspace{2pt}

\medskip
\subsection*{\texorpdfstring{B.1.3\quad FT.132 — The adjudication summary}{FT.132 B.1.3}}
\addcontentsline{toc}{subsection}{B.1.3 FT.132}
\phantomsection\label{tag:FT-132}
{\small \textit{Record.}\quad Source: HistLog §IV.1 (archival ledger, preserved at \texttt{working\_markdowns/HistLog.md}); reproduced verbatim, including its historical registry tokens, which resolve via HistLog \S I.}
\medskip
\textbf{Thm.} Provenance of every quantity of Tier F.40:

%tbl-level:1
{\small\setlength{\tabcolsep}{4pt}
\begin{longtable}{>{\raggedright\arraybackslash}p{\dimexpr0.4988\textwidth-2\tabcolsep\relax}|>{\raggedright\arraybackslash}p{\dimexpr0.1594\textwidth-2\tabcolsep\relax}|>{\raggedright\arraybackslash}p{\dimexpr0.3418\textwidth-2\tabcolsep\relax}}
\toprule
Statement & Source & Status \\
\addlinespace[2pt]
\midrule
\endfirsthead
\toprule
Statement & Source & Status \\
\addlinespace[2pt]
\midrule
\endhead
\bottomrule
\endlastfoot
Seed images $\{W_i\mapsto\tfrac12\sigma_i,\ Y\mapsto\tfrac12\sigma_0\}$; $T_a\mathsf v=\allowbreak 0$ & EW.32 & Derived (frozen) \\
Extended carrier basis; $S$-transform; partner involution $\hat\kappa$ & Def-H.34 & Constructed (P-H18a carrier) \\
Aspect census; $\mathbf F_{\mathrm{ext}}(\mathbf c)=\allowbreak h(\mathbf c)\otimes I_\omega$ & FT.127 & Derived \\
No-flip theorem at the derived locus ($F_{LR}=\allowbreak 0$) & FT.128 & Derived \\
Parity doubling of the derived spectrum & FT.128(5) & Derived \\
$\hat\kappa$ canonicity; quaternionic triple & FT.129(1)–(2) & Derived \\
Read parity split ($\omega$-odd flips, $\omega$-even shifts) & FT.129(3) & Derived \\
Read promotion $\hat Q_{\mathrm{EM}}=\allowbreak q_{\mathrm{core}}\hat\kappa$ & FT.129(4) & Declared (P-H19a) \\
Flipping block $F_{LR}=\allowbreak q\,\allowbreak I_2$; flipping set $\{F_{\mathrm{EM}}\ \text{core read}\}$ & FT.129(5) & Derived (under P-H19a) \\
Schmidt ranks $1/2$; no entangling term & FT.130(1)–(2) & Derived \\
Promoted spectrum; splitting $2|q|\in\allowbreak \{0,2\}$ & FT.130(5) & Derived (under P-H19a) \\
Scalar-shift invariance of the flipping structure & FT.130(7) & Derived \\
Flipping frequency & — & Deferred (needs emergent time, FT.131(3)) \\
Transit-chain flipping census & — & Not performed (FT.131(1)) \\
Identification with any species & — & None (barred, FT.131(4)) \\
\end{longtable}
}

\textbf{Req.} FT.127–FT.131; EW.32, EW.24 (\textbf{D}); CT0.12 (\textbf{B}); EM.9 (\textbf{E}); P-H18a, P-H19a (\textbf{H13}); machine: \texttt{checks/\allowbreak{}verify\_\allowbreak{}H\_\allowbreak{}force\_\allowbreak{}matrix.py} §A–§E.

\textbf{Pf.}

%dsp-level:0
\[
\text{each row is the adjudication of its own tag: FT.127}\ (\S\mathrm{B}),
\]

%dsp-level:0
\[
\ \text{FT.128}\ (\S\mathrm{B}),
\]

%dsp-level:0
\[
\ \text{FT.129}\ (\S\mathrm{A},\S\mathrm{C}),
\]

%dsp-level:0
\[
\ \text{FT.130}\ (\S\mathrm{D}),
\]

%dsp-level:0
\[
\ \text{FT.131}\ (\S\mathrm{E})\ [\text{machine: all sections green}]
\]

%dsp-level:0
\[
\text{declared rows logged}\ [\text{P-H19a},\ \mathrm{H13}];
\]

%dsp-level:0
\[
\qquad \text{frequency row:}\ \neg\exists\ \text{tag supplying a carrier-time parameter}\ [\mathrm{FT.131}(3):\ \mathrm{CT0.3},\ \mathrm{FT.86}(8),\ \text{item 21(b), P-H8}];
\]

%dsp-level:0
\[
\qquad \text{no external datum anywhere in this document}.\ \blacksquare
\]

\textbf{Phys.} Nine aspects enter the census; one family of seed images acts; the complete derived matrix is a single chirality-blind tensor, and the no-flip theorem is exact. The orientation doublet contributes the unique canonical flip operator — its partner involution — and the unique $\omega$-odd force read, the ternary core charge, supplies its coefficient: the static answer of the tier is one line. A charged core splits the parity doublet by exactly two read units; a neutral core does not split it; the frequency that would make this an oscillation awaits the construction of time.

\textbf{Status.} Closed. Tier F.40 adds one declared interface (P-H19a) and no free parameter; no species is identified anywhere in it; the deferral ledger names its prerequisites exactly. Machine: §A–§E.

\vspace{2pt}\noindent\rule{\textwidth}{0.4pt}\vspace{2pt}

\vspace{2pt}\noindent\rule{\textwidth}{0.4pt}\vspace{2pt}

\medskip
\subsection*{\texorpdfstring{B.1.4\quad FT.140 — The adjudication summary and the ledger}{FT.140 B.1.4}}
\addcontentsline{toc}{subsection}{B.1.4 FT.140}
\phantomsection\label{tag:FT-140}
{\small \textit{Record.}\quad Source: HistLog §IV.1 (archival ledger, preserved at \texttt{working\_markdowns/HistLog.md}); reproduced verbatim, including its historical registry tokens, which resolve via HistLog \S I.}
\medskip
\textbf{Thm.}

%dsp-level:0
\[
(1)\quad\textbf{P-H20a, adjudicated: VACUOUS as a declaration.}\ \ \text{The identification ``the A4 branch flip}\ \iota(\alpha)
\]

%dsp-level:0
\[
= \pi-\alpha\ \text{is the parity operation}\ P\ \text{on configuration space'' is frozen derived content:}\ \mathrm{EM.24}(2)\ \text{derives}\ P:\ \mathcal P
\]

%dsp-level:0
\[
\mapsto E_{\mathcal R}\setminus\mathcal P\ (\alpha_e\mapsto\pi-\alpha_e)\ \text{from the dynamic phase}\ [\mathrm{Def\text{-}EM.20},\ \mathrm{EM.20}(1)]\ \text{and names it parity inside the CPT bridge};
\]

%dsp-level:0
\[
\ \mathrm{FT.101}(3)\ \text{records}\ \iota\ \text{as that complement with the same citations}.
\]

%dsp-level:0
\[
\text{By the corpus's own rule — a declaration must add content not otherwise obtainable — no theorem of this part uses P-H20a as a premise};
\]

%dsp-level:0
\[
\ \text{the H13 entry is logged as adjudicated-vacuous and retained as ledger history}\ [\text{P-H2 pattern}];
\]

%dsp-level:0
\[
(2)\quad\textbf{the carrier identifications, adjudicated:}\ \ \hat\omega
\]

%dsp-level:0
\[
= C\ \text{and}\ \hat\kappa
\]

%dsp-level:0
\[
= P\ \text{are rejected as operator equalities}\ —\ \text{register conflations}\ [\mathrm{FT.134}(1):\ \text{ray-swapping vs ray-fixing witnesses}];
\]

%dsp-level:0
\[
\ \text{what survives is read-level covariance}\ [\mathrm{FT.134}(2)]\ \text{and the scalar-CPT theorem}\ [\mathrm{FT.134}(3)]:\ \mathrm{CPT}\big|_{\mathcal H_\omega}
\]

%dsp-level:0
\[
= \lambda\,
\]

%dsp-level:0
\[
\mathrm{id},
\]

%dsp-level:0
\[
\ \lambda
\]

%dsp-level:0
\[
= (-1)^{|E|}e^{i\pi\Theta}\ —\ \text{chirality-even};
\]

%dsp-level:0
\[
\ \text{no}\ \hat\kappa\text{-multiple can equal it}\ (\text{trace witness});
\]

%dsp-level:0
\[
(3)\quad\textbf{the ledger rows:}\ \ \text{deficit 1 (projector): status stands — closed at the projector level}\ [\mathrm{FT.125}(1),\ \text{frozen}];
\]

%dsp-level:0
\[
\ \text{this part adds no chirality content. Deficit}\ 1'\ \text{(chiral matter assignment): stays open and gains its parity-side input — the label-level branch exchange}\ [\mathrm{FT.135}],
\]

%dsp-level:0
\[
\ \text{the scalar CPT action}\ [\mathrm{FT.134}(3)]\ \text{and the flip-block character}\ [\mathrm{FT.139}(3)]\ \text{are the derived parity data any assignment rule must consume};
\]

%dsp-level:0
\[
\ \text{the carrier-level parity-eigenvalue assignment to species joins deficit}\ 1'\ \text{as missing rule content};
\]

%dsp-level:0
\[
\text{deficits}\ 2,3:\ \text{unchanged};
\]

%dsp-level:0
\[
\qquad \text{deficit}\ 5:\ \mathrm{FT.81}(4)\text{'s prerequisites stand}\ [\mathrm{FT.125}(2)];
\]

%dsp-level:0
\[
\ \mathrm{FT.135}\ \text{adds the arena-level parity non-invariance as datum};
\]

%dsp-level:0
\[
\qquad \text{deficits}\ 7\text{–}12:\ \text{unchanged}\ [\mathrm{FT.138}(4):\ \text{ceilings parity-blind}];
\]

%dsp-level:0
\[
(4)\quad\textbf{the scope bars:}\ \ \text{no particle identification};
\]

%dsp-level:0
\[
\ \text{no external data}\ ([\text{EXT}]\ \text{enters no clause of this document — this standing-rule record is its sole marker});
\]

%dsp-level:0
\[
\ \text{no Yukawa sector (deficit 2)};
\]

%dsp-level:0
\[
\ \text{no generation selection (deficit 3)};
\]

%dsp-level:0
\[
\ \text{no Hilbert-space parity operator on the pair register}\ [\mathrm{FT.135}(4)\ \text{stands}];
\]

%dsp-level:0
\[
\ \text{no emergent time — H19's deferred frequency read stays deferred}\ [\mathrm{FT.131}(3)];
\]

%dsp-level:0
\[
\ \mathrm{FT.88}(4)\text{'s CPT scope stands: configuration-level}\ |w|^2\ \text{invariance},
\]

%dsp-level:0
\[
\ \neg\ \text{a Lorentzian field-algebra CPT theorem}.
\]

%tbl-level:1
{\small\setlength{\tabcolsep}{4pt}
\begin{longtable}{>{\raggedright\arraybackslash}p{\dimexpr0.5053\textwidth-2\tabcolsep\relax}|>{\raggedright\arraybackslash}p{\dimexpr0.1895\textwidth-2\tabcolsep\relax}|>{\raggedright\arraybackslash}p{\dimexpr0.3053\textwidth-2\tabcolsep\relax}}
\toprule
Statement & Source & Status \\
\addlinespace[2pt]
\midrule
\endfirsthead
\toprule
Statement & Source & Status \\
\addlinespace[2pt]
\midrule
\endhead
\bottomrule
\endlastfoot
$P=\allowbreak \alpha$-complement $=\allowbreak \iota$ on labels & EM.24(2), FT.101(3) & Derived (frozen) \\
$\langle C,\allowbreak P,\allowbreak T\rangle=\allowbreak \mathbb Z_2^3$; commutations; $\mathbb Z_2^4$ with $\omega$ & FT.133(1)–(2) & Derived \\
Weight identities; both CPT lifts; agreement locus & FT.133(3)–(5) & Derived (EM.24(5) re-derived) \\
Read-covariance table & FT.133(6) & Derived \\
Scalar CPT on the doublet: $\lambda\,\allowbreak \mathrm{id}$ & FT.134(3) & Derived \\
$\hat\omega=\allowbreak C$, $\hat\kappa=\allowbreak P$ as operator equalities & FT.134(1) & REJECTED (register conflation) \\
$\hat Q$-family characters: C-odd, P/T-even, $\omega$-covariant & FT.134(2) & Derived (under P-H19a) \\
Weak-branch parity exchange (label-level) & FT.135 & Derived (frozen re-run) \\
Rung-one $10/6=\allowbreak \operatorname{Ad}(\mathbb S)$ grading & FT.136(1)–(2) & Derived \\
"$S$-blocks $=\allowbreak P$-sectors" as operator content & FT.136(3) & REJECTED (distinct $\mathbb Z_2$'s) \\
Cover pair exchanged by reversal, not by $P$ & FT.137(2), (5) & Derived \\
$P$ sign character on amplitudes & FT.137(3) & Derived \\
Defect multiplets $m\le\allowbreak ftrightarrow-m$ symmetric & FT.138(1)–(2) & Derived \\
Channel constants register-blind & FT.139(2) & Derived \\
Flipping block C-odd, P-even & FT.139(3) & Derived (under P-H19a) \\
P-H20a as a declaration & FT.140(1) & Adjudicated vacuous (H13 ledger) \\
Parity-eigenvalue assignment to species & — & Open (deficit 1$'$ content) \\
\end{longtable}
}

\textbf{Req.} FT.133–FT.139; EM.23, EM.24 (\textbf{E}); FT.88, FT.89 (H11); FT.101 (\textbf{H12}); FT.125 (H18); FT.128, FT.129, FT.131 (H19); FT.81 (H10); items 20/21, P-H2, P-H19a, P-H20a (\textbf{H13}); exact bookkeeping; machine: \texttt{checks/\allowbreak{}verify\_\allowbreak{}H\_\allowbreak{}parity\_\allowbreak{}identification.py} §A–§H.

\textbf{Pf.}

%dsp-level:0
\[
\text{each row is the adjudication of its own tag: FT.133}\ (\S\mathrm{A}),
\]

%dsp-level:0
\[
\ \text{FT.134}\ (\S\mathrm{B}),
\]

%dsp-level:0
\[
\ \text{FT.135}\ (\S\mathrm{C}),
\]

%dsp-level:0
\[
\ \text{FT.136}\ (\S\mathrm{D}),
\]

%dsp-level:0
\[
\ \text{FT.137}\ (\S\mathrm{E}),
\]

%dsp-level:0
\[
\ \text{FT.138}\ (\S\mathrm{F}),
\]

%dsp-level:0
\[
\ \text{FT.139}\ (\S\mathrm{G})\ [\text{machine: all sections green}]
\]

%dsp-level:0
\[
(1):\ \mathrm{EM.24}(2)\ \text{verbatim}+\mathrm{FT.101}(3)\ \text{citation chain}:\ \text{the identification adds no content beyond frozen tags}\ \blacksquare_{\text{(row)}};
\]

%dsp-level:0
\[
\qquad \text{rejection rows:}\ \S\mathrm{B}\ \text{ray census and trace witness};\ \S\mathrm{D}\ \text{register separation}
\]

%dsp-level:0
\[
\text{open row:}\ \neg\exists\ \text{tag assigning parity eigenvalues to species}\ [\text{inventory scan: the assignment rule is deficit}\ 1'\ \text{content}].\ \blacksquare
\]

\textbf{Phys.} Parity was already in the ledger: derived from the dynamic phase, named in the CPT bridge, coupled to the weak branches, and exhaustive in its footprints. What this part adds is the complete bookkeeping — which involutions act on which registers, what each fixes and what each exchanges, and where the names agree at the read level while operator equalities are barred. The CPT theorem of the corpus touches the fermion carrier as a phase and nothing more; the only parity-sensitive structure in the whole tower is the charge-odd flipping block; and the one rule still missing is the one that would tell a species which side of any of these $\mathbb{Z}$$_2$'s it sits on.

\textbf{Status.} Closed. Tier F.41 adds no declaration and no free parameter; P-H20a is adjudicated and not adopted; no species is identified anywhere in it. Machine: §A–§H.

\vspace{2pt}\noindent\rule{\textwidth}{0.4pt}\vspace{2pt}

\medskip
\subsection*{\texorpdfstring{B.1.5\quad Tier closure exports and ledger updates (Q/R tiers)}{IV.2 B.1.5}}
\addcontentsline{toc}{subsection}{B.1.5 SIV.2}
{\small \textit{Record.}\quad Source: HistLog §IV.2 (archival ledger, preserved at \texttt{working\_markdowns/HistLog.md}); reproduced verbatim, including its historical registry tokens, which resolve via HistLog \S I.}
\medskip
\textit{(from Q\_carrier\_and\_parity.md)}

\phantomsection
\section[Closure exports of H18 (consumed by H13's registrie…]{\texorpdfstring{Closure exports of H18 (consumed by H13's registries)}{Closure exports of H18 (consumed by H13's registries)}}

%dsp-level:0
\[
\begin{array}{lll}
\mathrm{Def\text{-}H.33}\ \mathcal H_\omega(c),\ \hat\omega,\ \hat\chi,\ P_L,P_R & \mathrm{FT.122}\ \text{constructibility, canonicity}\ \pm\hat\chi,\ \text{exact}\ \mathbb Q(i) & \mathrm{FT.123}\ [W_i\otimes I,\ I\otimes P_{L/R}]=0;\ \hat\omega P_L\hat\omega^{-1}=P_R\\
\mathrm{FT.124}\ \text{obstruction evasion (seven loci)} & \mathrm{FT.125}\ \text{deficit ledger: }1\to\text{closed (projector level)},\ 1'\to\text{opened},\ 5\to\text{relocated},\ 6\to\text{partial} & \mathrm{FT.126}\ \text{chain-carrier lift},\ \mathbb Z_2\times\mathbb Z_2\\
\text{naming:}\ |L\rangle\ :=\ \text{the}\ +1\ \text{ray of the boxed}\ \hat\chi\ \text{[FT.122}(6)] & \text{one new declaration: P-H18a (carrier product)} & \text{all prior obstructions retained, no tag overturned}
\end{array}
\]

\textbf{Ledger update to H13 item 20.} Deficit 1 (chirality/Weyl/projector) is \textbf{closed at the projector level under P-H18a}, with item 21(i)'s frame-section residue standing; deficit 1$'$ (chiral matter assignment) is \textbf{opened}. Deficits 2 (Yukawa), 3 (generations), 4 (CKM/PMNS) remain open, downstream of 1$'$ rather than of 1. Deficit 5 (anomalies): blocker \textbf{relocated} to 1$'$ (+2/3) — FT.81(4)'s Weyl-symbol and hypercharge prerequisites stand. Deficit 6 (the neutrino row): projector-side prerequisite discharged; species content still absent. Deficit 7 (second scale) unchanged. Deficit 8 (Higgs 4-dof) unchanged. Deficit 9 (fractional confined charges) unchanged. Deficit 10 ($|q|=2$ multi-chain read) unchanged. Deficit 11 (spin–statistics) \textbf{unchanged} — the grade's $\mathbb Z_2$ is single-chain data, not a two-chain exchange sign [FT.80(5): no exchange operator in the inventory; FT.89(2) kind discipline]. Deficit 12 (SM masses) unchanged at the ceiling level — the ceilings are properties of the pole-side and colour-side read species [FT.124(6)–(7)], not of chirality.

\textbf{Forward pointer (H21, Tier F.42, non-load-bearing).} The rung-one chirality theorem [FT.141 (H21)] situates Def-H.33's carrier grade among four coupled gradings: the force-arena $S$-grading (the $10/6$ $\operatorname{Ad}(\mathbb S)$ split — the chirality of the complete force complex, deletion-invariant) determines which observables are chiral; $\hat\chi$ determines which chain orientations are chiral; the branch-label parity [FT.135 (H20)] grades the weak support; and the Janus arrow [FT.144 (H21)] makes the rung-one materialized locus single-arrow. Under any $\hat\omega$-covariant arrow$\leftrightarrow$ray coupling the rung-one locus occupies a single chiral ray with occupancy $(1,0)$ — the complement ray is the doublet's vacuum sector at rung one; the coupling itself, and the physical naming of the rays, remain deficit 1$'$ content [FT.125(4)]. The pointer adds no load-bearing content to this tier.

$\blacksquare$

\vspace{2pt}\noindent\rule{\textwidth}{0.4pt}\vspace{2pt}

\vspace{2pt}\noindent\rule{\textwidth}{0.4pt}\vspace{2pt}

\textit{(from Q\_carrier\_and\_parity.md)}

\phantomsection
\section[Closure exports of H19 (consumed by H13's registrie…]{\texorpdfstring{Closure exports of H19 (consumed by H13's registries)}{Closure exports of H19 (consumed by H13's registries)}}

%dsp-level:0
\[
\begin{array}{lll}
\mathrm{Def\text{-}H.34}\ \mathcal H_{\mathrm{ext}}\ \text{bases},\ S\text{-transform},\ \hat\kappa\ \text{(quaternionic triple)} & \mathrm{FT.127}\ \text{nine-aspect census};\ \mathbf F_{\mathrm{ext}}(\mathbf c)=h(\mathbf c)\otimes I_\omega & \mathrm{FT.128}\ F_{LR}=0\ \text{(derived locus)};\ \text{parity doubling}\\
\mathrm{FT.129}\ \hat\kappa\ \text{canonical};\ \hat Q_{\mathrm{EM}}=q_{\mathrm{core}}\hat\kappa\ \text{(P-H19a)};\ F_{LR}=qI_2 & \mathrm{FT.130}\ \text{Schmidt}\ 1/2;\ \operatorname{spec};\ \text{splitting}\ 2|q|\in\{0,2\} & \mathrm{FT.131}\ \text{bound scope};\ \text{deferrals (emergent time prerequisite)}\\
\mathrm{FT.132}\ \text{adjudication: frequency deferred, no identification} & \text{one new declaration: P-H19a (read promotion)} & \text{no frozen tag overturned; no free parameter added}
\end{array}
\]

\textbf{Ledger update to H13.} Item 20 deficit 1$'$ gains its static datum: the flipping block is computed under P-H19a ($F_{LR}=q_{\mathrm{core}}I_2$; flipping set $\{F_{\mathrm{EM}}\ \text{core read}\}$; chiral-partner splitting $2|q|\in\{0,2\}$; the derived locus obeys the no-flip theorem unconditionally), and the flipping frequency is deferred pending emergent time — a named prerequisite at item 21(b). The assignment rule (species $\leftrightarrow$ gauge representation $\leftrightarrow$ chiral ray) remains absent: deficit 1$'$ stays open. Items 7 and 13 gain the carrier-vs-channel superselection tension of FT.129(4)/FT.131(5) as logged bridge content.

\textbf{Forward pointer (H21, Tier F.42, non-load-bearing).} The no-flip theorem (FT.128) is the force-side half of the rung-one occupancy theorem [FT.144 (H21)]: at rung one the derived static matrix is chirality-block-diagonal ($F_{LR}=0$ exactly) and the materialized locus is single-arrow, so the complementary ray is vacuum-occupied. The flipping block $F_{LR}=q\,I_2$ (FT.129(5), under P-H19a) opens exactly at EM activation — rung two, $\dim\mathcal K_2=17$ [CT0.14] — which is where the single-sector structure ends [FT.144(6), FT.145(4) (H21)]; colour adds no carrier flip (the flipping set remains the EM core read, minimal and unique). The flipping frequency remains deferred (FT.131(3)). The pointer adds no load-bearing content to this tier.

$\blacksquare$

\vspace{2pt}\noindent\rule{\textwidth}{0.4pt}\vspace{2pt}

\vspace{2pt}\noindent\rule{\textwidth}{0.4pt}\vspace{2pt}

\textit{(from Q\_carrier\_and\_parity.md)}

\phantomsection
\section[Closure exports of H20 (consumed by H13's registrie…]{\texorpdfstring{Closure exports of H20 (consumed by H13's registries)}{Closure exports of H20 (consumed by H13's registries)}}

%dsp-level:0
\[
\begin{array}{lll}
\mathrm{FT.133}\ \mathbb Z_2^4\ \text{reversal calculus};\ \text{both CPT lifts} & \mathrm{FT.134}\ \mathrm{CPT}\big|_{\mathcal H_\omega}=\lambda\,\mathrm{id};\ \text{identifications rejected} & \mathrm{FT.135}\ \text{branch exchange, label-level only}\\
\mathrm{FT.136}\ 10/6=\operatorname{Ad}(\mathbb S);\ \text{three independent}\ \mathbb Z_2\text{'s} & \mathrm{FT.137}\ \text{reversal pairs the covers};\ P\ \text{sign character} & \mathrm{FT.138}\ D(m)=D(-m);\ \text{multiplets exact}\\
\mathrm{FT.139}\ \text{channels blind};\ F_{LR}\ \text{C-odd} & \mathrm{FT.140}\ \text{P-H20a vacuous};\ \text{deficit}\ 1'\ \text{input} & \text{no new declaration; no tag overturned}
\end{array}
\]

\textbf{Ledger update to H13.} P-H20a: adjudicated vacuous — the parity identification is frozen derived content [EM.24(2), FT.101(3)]; entry retained as ledger history [P-H2 pattern]; no H20 theorem uses it. Item 20: the parity census adds derived datum to deficit 1$'$ (label-level branch exchange; scalar CPT action on the doublet; the flip block's C-odd/P-even character) — the assignment rule remains absent; deficits 2–12 statuses as adjudicated in FT.140(3). Item 21(c): the involution-algebra program gains the exact register calculus $\mathbb Z_2^4$ [FT.133(2)] and the carrier-side scalar-CPT theorem [FT.134(3)]; the projector question on the pair register remains closed-negative [FT.101(3), FT.135(4)].

$\blacksquare$

\vspace{2pt}\noindent\rule{\textwidth}{0.4pt}\vspace{2pt}

\textit{(from R\_rung\_one.md)}

\phantomsection
\section[Closure exports of H21 (consumed by H13's registrie…]{\texorpdfstring{Closure exports of H21 (consumed by H13's registries)}{Closure exports of H21 (consumed by H13's registries)}}

%dsp-level:0
\[
\begin{array}{lll}
\mathrm{FT.141}\ 10/6=\text{the complete-force chirality};\ \text{deletion-invariant} & \mathrm{FT.142}\ \text{arena}\ \mathcal K_1=16\ (\text{derived});\ \delta_\nu=\tfrac18 & \mathrm{FT.143}\ r_i^{(\nu)};\ 1:9.58\ldots:272.0\ldots;\ Q=\tfrac23;\ \text{cubic}\ \sqrt2\cos\tfrac38\\
\mathrm{FT.144}\ \text{prefix law};\ \neg\text{backward materialization};\ \text{vacuum residence} & \mathrm{FT.144}(5)\text{–}(6)\ (1,0)\ \text{occupancy};\ \text{flip opens at rung}\ 2 & \mathrm{FT.145}\ \text{completeness}\ 16=C_{\mathrm{cap}};\ m_{\mathrm{hid}}=\tfrac{2\ln2}{\mathfrak h};\ \text{bridge}\ 17\\
\text{no new declaration; imports}\ \{\text{P-H0},\ \text{P-H18a},\ \text{P-H19a}\} & \text{angle ledger}\ 16\delta_\nu=9\delta_\ell=27\delta_q=2 & \text{no frozen tag overturned; no species identified}
\end{array}
\]

\textbf{Ledger update to H13.} No new declaration is issued by Tier F.42; the import list is confined to P-H0, P-H18a, P-H19a, with P-H16a/P-H17a cited in adjudication rows and imported nowhere. Item 20: deficit 1$'$ gains the arena-side theorem — the $10/6$ $\operatorname{Ad}(\mathbb S)$ split is the chirality of the complete rung-one force complex (deletion-invariant; no $S$-neutral observable; the weak branch selectivity a state-locus shadow with $\sum_iW_i^2=E_0+\mathbb S$) — and the assignment rule must couple four gradings: the algebra $S$-grading, the branch-label parity, the carrier $\hat\chi$-grade, and the Janus-arrow chirality of the materialized locus; the rule remains absent. Deficit 3: the third family channel runs the same democratic $C_3$ mechanism on a derived arena — the family count is $3=\dim F_{J^\pm}$ in all three channels; the selection principle remains undeclared (item 21(d) stands). Deficit 6 (the $\nu$-row): channel-side content computed — $\delta_\nu=2/16=1/8$ (the first channel denominator requiring no declared interface), spectrum $r_i^{(\nu)}=1+\sqrt2\cos(\tfrac18+\tfrac{2\pi i}3)$ with ratios $m_1^2:m_2^2:m_3^2=1:9.583926992379\ldots:271.994748749143\ldots$ (scale $\mu_\nu$ undeclared, P4), $Q=\tfrac23$ exact, cubic constant $\sqrt2\cos\tfrac38$ channel-own ($3\delta_\nu\notin\{\delta_\ell,\delta_q\}$; the shared numerator $16\delta_\nu=9\delta_\ell=27\delta_q=2$ is the only cross-channel identity), and the Janus-arrow theorem: rung-one materialization is single-arrow (the prefix law, exhaustive), backward materialization is domain-excluded — $\neg$mass-suppressed — and resides as balanced formal content in the all-virtual state with no accessible pole; single-ray occupancy follows under any $\hat\omega$-covariant arrow$\leftrightarrow$ray coupling (the naming residue is deficit 1$'$ content); the force-side flip channel opens exactly at EM activation (rung two, $F_{LR}=qI_2$ under P-H19a), so the single-sector result is a rung-one theorem. Rung-one completion: $\mathcal A(1)$ accounts for all $16=C_{\mathrm{cap}}$ pair-channel dimensions; the digon reads $m_{\mathrm{hid}}=2\ln2/\mathfrak h$; the relating equation to the channel ratios remains open (item 7). Deficits 2, 4, 5, 7–12 unchanged.

%dsp-level:0
\[
\blacksquare
\]

\medskip
\subsection*{\texorpdfstring{B.1.6\quad FT.146 — The consolidation law: the letter series, the ownership partition, and the acyclic DAG}{FT.146 B.1.6}}
\addcontentsline{toc}{subsection}{B.1.6 FT.146}
\phantomsection\label{tag:FT-146}
{\small \textit{Record.}\quad Source: HistLog §IV.3 (the closure tier F.43) (archival ledger, preserved at \texttt{working\_markdowns/HistLog.md}); reproduced verbatim, including its historical registry tokens, which resolve via HistLog \S I.}
\medskip
\textbf{Thm.}

%dsp-level:0
\[
(1)\quad\textbf{the series:}
\]

%dsp-level:0
\[
\qquad \text{the consolidated fermionic layer is the letter series}
\]

%dsp-level:0
\[
\qquad \mathbf H\to\mathbf J\to\mathbf K\to\mathbf L\to\mathbf M\to\mathbf N\to\mathbf O\to\mathbf P\to\mathbf Q\to\mathbf R\to\mathbf S\to\mathbf T
\]

%dsp-level:0
\[
\qquad (\text{reading order; the letter I is not used});
\]

%dsp-level:0
\[
\text{parts H–R carry the full working-layer text, organized by relevance}\ (\text{the first editions are archived at}\ \texttt{H\_archive/H01–H21}\ \text{as historical context});
\]

%dsp-level:0
\[
\ \text{part S carries the registries and the Open Gaps Log}\ (\text{canonical ledger text});
\]

%dsp-level:0
\[
\ \text{part T (this part) carries the closure tier F.43 and the appended H-CENSUS block};
\]

%dsp-level:0
\[
\begin{array}{c|c|c|c}\text{part} & \text{theme} & \text{tiers} & \text{tags owned}\\ \hline
\mathbf H & \text{measurement primitives, edges, subnetwork force} & \text{F.1–F.11} & \text{FT.1–FT.31};\ \text{Def-H.1–Def-H.18}\\
\mathbf J & \text{floor families, defect bridges, Platonic/polygon cells} & \text{F.12–F.16} & \text{FT.32–FT.52};\ \text{Def-H.19–Def-H.20}\\
\mathbf K & \text{spin ladder, stationarity, loops, Bott correction} & \text{F.17–F.22} & \text{FT.53–FT.66};\ \text{Def-H.21–Def-H.24}\\
\mathbf L & \text{flicker, domains, exact covers} & \text{F.23–F.25} & \text{FT.67–FT.75};\ \text{Def-H.25–Def-H.27}\\
\mathbf M & \text{pole spectra, cover orbits, gravitational reads, audits} & \text{F.26–F.30} & \text{FT.76–FT.91};\ \text{Def-H.28–Def-H.30}\\
\mathbf N & \text{Clifford census, branch grading, form factors} & \text{F.31–F.34} & \text{FT.92–FT.101}\\
\mathbf O & \text{path integrals: one-loop bridge, topology-blind sum} & \text{F.35–F.36} & \text{FT.102–FT.108};\ \text{Def-H.31–Def-H.32}\\
\mathbf P & \text{the family channels} & \text{F.37–F.38} & \text{FT.109–FT.121}\\
\mathbf Q & \text{the carrier: chirality, force matrix, parity} & \text{F.39–F.41} & \text{FT.122–FT.140};\ \text{Def-H.33–Def-H.34}\\
\mathbf R & \text{rung one: complete-force chirality, Janus arrow, the sixteen} & \text{F.42} & \text{FT.141–FT.145}\\
\mathbf S & \text{registries and the Open Gaps Log (canonical ledger text)} & — & \neg\text{FT/Def tags: the P-H registry, the revision registries, gap items}\ 1\text{–}21\\
\mathbf T & \text{rung closure: audit, separation, reframe, baryonic preparation} & \text{F.43} & \text{FT.146–FT.152}\ (\text{this part})
\end{array}
\]

%dsp-level:0
\[
(2)\quad\textbf{the ownership partition:}
\]

%dsp-level:0
\[
\qquad \text{every tag FT.1–FT.145 and Def-H.1–Def-H.34 is carried in full text by exactly one of the parts H–R};
\]

%dsp-level:0
\[
\qquad 145+34=179\ \text{tags},\ \text{no residue},\ \text{no double ownership}
\]

%dsp-level:0
\[
\qquad [\text{machine: section-heading census across the series}];
\]

%dsp-level:0
\[
(3)\quad\textbf{the canon clause:}
\]

%dsp-level:0
\[
\qquad \text{the letter parts H–R and part S carry the full}\ \text{A}\ \text{§0.4 text of every owned tag and of the ledger};
\]

%dsp-level:0
\[
\ \text{the first-edition files}\ \texttt{H\_archive/H01–H21}\ \text{are the historical locus}:
\]

%dsp-level:0
\[
\text{every transferred body is byte-identical to its first edition}\ [\text{machine: verbatim body identity per source file; headers, the layer map of part H, and this tier are the only new text}];
\]

%dsp-level:0
\[
\neg\exists\ \text{a clause of any letter part that alters, re-derives, or supersedes an owned tag};
\]

%dsp-level:0
\[
\qquad \text{the latest revision is canonical};\ \text{history is in git};
\]

%dsp-level:0
\[
(4)\quad\textbf{the acyclicity law:}
\]

%dsp-level:0
\[
\qquad \text{each tag's Req citations resolve to:}\ \text{A–G tags},
\]

%dsp-level:0
\[
\ \text{the registry}\ (\text{P-H entries; archive H13}),
\]

%dsp-level:0
\[
\ \text{or tags owned by the same or an earlier letter}
\]

%dsp-level:0
\[
\qquad [\text{machine: letter-level closure re-extracted from the archive Req lines, with citation ranges expanded}];
\]

%dsp-level:0
\[
\text{the four forward citations of the archived file granularity}\ (\mathrm{FT.14}\to\mathrm{FT.22}/\mathrm{FT.24},\ \mathrm{FT.15}\to\mathrm{FT.25},\ \mathrm{FT.21}\to\mathrm{Def\text{-}H.15}\ \text{— on-shell scope conditions cited at use})\ \text{are intra-part at letter granularity}\ (\text{all four inside}\ \mathbf H):\ \text{the partition is chosen so the letter DAG is acyclic};
\]

%dsp-level:0
\[
\text{the letter DAG is the chain}\ \mathbf H\to\mathbf J\to\cdots\to\mathbf S;\qquad\text{forward-letter citations: none}\ [\text{machine}];
\]

%dsp-level:0
\[
(5)\quad\textbf{the census block:}
\]

%dsp-level:0
\[
\qquad \text{the Space/Counts/Baryonic entries of the}\ 25\ \text{Category-B tags are consolidated in this part's H-CENSUS block, appended at the end of the series}\ [\mathrm{FT.148}(4)];
\]

%dsp-level:0
\[
\ \text{part H's layer map and the FT.148 table record which letter carries which B tag};
\]

%dsp-level:0
\[
(6)\quad\textbf{the registry domicile:}
\]

%dsp-level:0
\[
\qquad \text{the P-H entry texts and the Open Gaps Log are canonical in part S};
\]

%dsp-level:0
\[
\ \text{this part carries the verdict/status layer}\ [\mathrm{FT.147},\ \mathrm{FT.149},\ \mathrm{FT.150}];
\]

%dsp-level:0
\[
\ \text{future declarations are logged in part S's registry addendum before use}\ [\text{the author's policy}];
\]

%dsp-level:0
\[
(7)\quad\textbf{the self-containment clause:}
\]

%dsp-level:0
\[
\qquad \text{the series}\ \mathbf{H}\text{–}\mathbf T\ \text{is standalone: full theorem text, the layer map and concordance in part H, the ledger in part S — no clause of the series requires the archive to be read};
\]

%dsp-level:0
\[
\ \text{the generated}\ \texttt{index/}\ \text{covers A–G and the archived first editions}\ [\text{checker scope unchanged}].
\]

\textbf{Req.} The consolidated letter parts H–S (full-text sections, parsed at the point of use); the archived first editions H\_archive/H01–H21 (the identity locus); exact set arithmetic, letter-level graph closure [math; machine: \texttt{checks/\allowbreak{}verify\_\allowbreak{}H\_\allowbreak{}rung\_\allowbreak{}closure.py} §A–§B].

\textbf{Pf.}

%dsp-level:0
\[
(1),
\]

%dsp-level:0
\[
(2):\ \text{section-heading census across the twelve parts:}\ 179+7\ \text{tag sections, each exactly once; the ownership ranges match the table}\ [\text{machine}]
\]

%dsp-level:0
\[
(3):\ \text{each source file's body (all text below its first-edition header) occurs byte-identically inside its letter part}\ [\text{machine: verbatim containment per file}];
\]

%dsp-level:0
\[
\qquad \text{no statement text is altered}
\]

%dsp-level:0
\[
(4):\ \text{per tag, the Req token set is re-extracted from the archive section and mapped through the ownership table: every FT/Def-H cite lands in the same or an earlier letter}\ [\text{machine: zero forward-letter citations}];
\]

%dsp-level:0
\[
\qquad \text{the four scope-condition citations:}\ \mathrm{FT.14},
\]

%dsp-level:0
\[
\mathrm{FT.15},
\]

%dsp-level:0
\[
\mathrm{FT.21}\in\mathbf H\ \wedge\ \mathrm{FT.22},
\]

%dsp-level:0
\[
\mathrm{FT.24},
\]

%dsp-level:0
\[
\mathrm{FT.25},
\]

%dsp-level:0
\[
\mathrm{Def\text{-}H.15}\in\mathbf H\ [\text{machine}]
\]

%dsp-level:0
\[
(5)\text{–}(7):\ \text{block scans and checker-scope record}\ [\text{machine}].\ \blacksquare
\]

\textbf{Phys.} The working layer becomes a shelf of standing books. Twelve lettered parts carry the derivation in full, in one reading order, each owning a contiguous span of it — every theorem present in its own text, byte-identical to the first edition the archive preserves. Ownership is a partition — every theorem and definition has exactly one home — and dependency is a chain: no part leans on a later one, the four scope citations that crossed file boundaries in the archive now sit inside a single part, and the graph has no cycles because it has no backward glance. The census block waits at the end of the shelf: the twenty-five cell-supplies, labeled with the spaces they enumerate and the era that consumes them.

\textbf{Status.} Closed; the consolidation organizes and does not modify — every owned tag's content, scope and declaration status are exactly as first written, and the transfer is machine-verified byte-identical. Machine: §A–§B.

\vspace{2pt}\noindent\rule{\textwidth}{0.4pt}\vspace{2pt}

\medskip
\subsection*{\texorpdfstring{B.1.7\quad FT.147 — The P-H audit and closure}{FT.147 B.1.7}}
\addcontentsline{toc}{subsection}{B.1.7 FT.147}
\phantomsection\label{tag:FT-147}
{\small \textit{Record.}\quad Source: HistLog §IV.3 (the closure tier F.43) (archival ledger, preserved at \texttt{working\_markdowns/HistLog.md}); reproduced verbatim, including its historical registry tokens, which resolve via HistLog \S I.}
\medskip
\textbf{Thm.} Every declared choice of the registry [H13] is audited; each carries exactly one verdict from the closed vocabulary $\{$RETAIN, DISCHARGED, REDUCED, ADJUDICATED VACUOUS, RETIRED$\}$:

%dsp-level:0
\[
(1)\quad\textbf{the verdict table}\ (22\ \text{entries; machine: each exactly once, anchors existent}):
\]

%tbl-level:3
{\scriptsize\setlength{\tabcolsep}{4pt}
\begin{longtable}{>{\raggedright\arraybackslash}p{\dimexpr0.0616\textwidth-2\tabcolsep\relax}|>{\raggedright\arraybackslash}p{\dimexpr0.0974\textwidth-2\tabcolsep\relax}|>{\raggedright\arraybackslash}p{\dimexpr0.8410\textwidth-2\tabcolsep\relax}}
\toprule
P-H & Verdict & Justification (anchor tags; machine-verified existent) \\
\addlinespace[2pt]
\midrule
\endfirsthead
\toprule
P-H & Verdict & Justification (anchor tags; machine-verified existent) \\
\addlinespace[2pt]
\midrule
\endhead
\bottomrule
\endlastfoot
P-H0 & RETAIN & Def-H.4's identifying content (fermion $=$ terminus of a Janus-closed bound chain); live consumers: FT.4, FT.28, FT.144(1); foundational — no derivation exists or is claimed \\
P-H1 & RETAIN & Def-H.3 ($I\equiv\mathsf E_{\mathrm{tr}}$ carried); consumers: FT.1, FT.3, FT.5; coherence-audited and retained (item 3 closure) \\
P-H2 & DISCHARGED & FT.36(2): saturating transit-content schedules exhibited at $\le\allowbreak 4$ hops/tick capacity; retired content: the sup-attainment hypothesis; nothing remains; no future adoption required; entry retained as ledger history \\
P-H3 & REDUCED & FT.41(a): the counting half derived ($N_{\mathrm{charge}}=\allowbreak 2$ [EM.9/D3.4]; $N_{\mathrm{weak}}=\allowbreak \delta h-2$ its complement); residual: the attribution clause only (interior weight bits read through the weak/spin channel, Def-H.8); full retirement would require a derived attribution rule from registration data \\
P-H4 & RETAIN & FT.17 hypothesis 2 (per-vertex caps do not further reduce channel-visible content); consumers: FT.17, FT.28(3); the general floor-uniform claim is not derivable from CT0.6/CT4.9 bounds alone \\
P-H5 & RETAIN & Def-H.16's boundary-state read ($|\mathrm{in}\rangle,\allowbreak |\mathrm{out}\rangle$ via $\iota_\Sigma$ [Def-H.13]); consumers: FT.26, FT.27, the pole-census tiers (FT.60/FT.63/FT.76 lines); CT2.5 supplies the frozen vocabulary, not the $\iota_\Sigma$ identification \\
P-H6 & RETAIN & FT.20's H-row fold (H's coarse-grained observable content read through $m_{\mathrm{hid}}$); grounds EW.13/EW.20/EW.33 packaging; the fold itself is a declared modelling choice \\
P-H7 & RETAIN & FT.33's defect transfer (pole-side mass candidate, leading-order form); consumer: FT.33; superseded in exactness by P-H7$'$ but the leading-order packaging remains declared; derivation residue tracked at item 7 \\
P-H7$'$ & RETAIN & FT.54(4)'s exact strengthening ($m_{\mathrm{pole}}^2=\allowbreak \hat\mu'^2\,\allowbreak 4\sin^2(\Theta_R/4)$, spin-connection holonomy); consumers: FT.54, FT.76, FT.87(1) (adjudicated as its leading Taylor term); same species as P-H7: declared transfer, open to derivation (items 7/19) \\
P-H8 & REDUCED & Vocabulary upgraded to Morse–Bott on derived geometry [Def-H.24 — a construction class, no new declaration]; the non-coupling clause superseded in part by FT.66(4)'s derived anti-coupling (absolute index-0 selection inadmissible against M1); item 18 closed by FT.66(5); residual declared content: the stability vocabulary as program vocabulary; full retirement would require a derived weight/ledger-to-materialization coupling — blocked by FT.61(2)–(3)'s index-blindness (exits (a$''$)/(b) named at item 13) \\
P-H9 & RETAIN & Def-H.23/FT.60's winding-multiplicity transfer ($\mathrm{Sym}^k(\mathcal H_{\mathsf v})$, diagonal action); the coupling multiplicity $=$ winding is declared and proven irreducible to subloop decomposition [FT.60(7)]; consumers: FT.60, FT.63(1); residues: derivation/replacement (item 19), $k\ge\allowbreak 3$ carrier inventory [FT.60(8)] \\
P-H10 & REDUCED & Derived away in part: the equilibrium form $p=\allowbreak x_+\allowbreak /(x_+\allowbreak +\allowbreak x_-)$ is the unique stationary measure of the per-pair Janus generator [FT.108(4): $\ker Q$ one-dimensional, span $(x_-,\allowbreak x_+\allowbreak )$; the H14/P-H14 dictionary $p=\allowbreak 1-4\varepsilon_q$ a marked specialization]; residual declared content: $r$ (equivalently $x_\pm$) as a free internal dynamics datum — the corpus supplies no $J^\pm$ rates [D1.2 is transition vocabulary; A5's draws are unrecorded] — and the per-edge-independent product ensemble scope; consumers of the residue: FT.68(5), FT.70/FT.71 (parameterized); full retirement would require a dynamics construction supplying the rates (joins item 21(b)) \\
P-H11 & RETIRED & Candidate adjudicated: retired, not adopted (clause (2)); zero live imports [machine inventory scan]; entry retained as ledger history with its costs column intact \\
P-H12 & RETAIN & The L3(4)$\le\allowbreak ftrightarrow$FT.16 bridge statement (pair-level inertial read $=$ subnetwork informational read, "the same object read differently"); a declared bridge, $\ne\allowbreak g$a derivation — no tag of A–H relates the reads quantitatively; consumers: item 7's structure, FT.104(2)'s read-coupling line; retirement would require the relating equation (item 7) \\
P-H13 & RETAIN & The fermion-mass interpretation of the nonlinear tilt (declared reading, NOT adopted into any A–G derivation), incl. the composition rule $V_{E_{\mathrm{int}}}=\allowbreak (\tfrac9{16})^{E_{\mathrm{int}}}V_0$; consumer: FT.104(2)'s composed-tilt range (marked specialization); bars stand (no calibration, no identification) \\
P-H14 & RETAIN & Def-H.31's five-component interface complex (dictionary, family/measure, operator interface, read coupling, labelling); consumers: FT.102–FT.104; boundary-recovered identically by FT.108(6); the labelling component (e) leaves item 12 untouched \\
P-H15 & RETAIN & Def-H.32's five-component interface complex (regularization, fugacity formalism, observable conditioning, mass expectation, move system); consumers: FT.105–FT.108; the derived stationary ratio of FT.108(4) enters it as specialization, not premise \\
P-H16a & RETAIN & FT.110's seal$\times$transfer arena ($3\times3=\allowbreak 9$; the tensor-product interface declared, the factors derived: $\dim F_{J^\pm}=\allowbreak 3$ [CT0.12], $\operatorname{tr}T(1)=\allowbreak 3$ [FT.68(2)]); consumers: FT.110, FT.112 ($\delta_\ell$ denominator); contrast of record: the rung-one arena required no declaration [FT.142(3): $\mathcal K_1$ derived] \\
P-H17a & RETAIN & FT.118(2)'s colour-dressed arena ($3\cdot3\cdot3=\allowbreak 27$; tension with the orthogonal-stack-block record logged, the spectator numerator derived); consumers: FT.118, FT.119 ($\delta_q$ denominator) \\
P-H18a & RETAIN & The chain-carrier tensor product $\mathcal H_{\mathrm{int}}\otimes\mathcal H_\omega$ (declared interface; the orientation space and every operator identity on it declaration-free); consumers: FT.123, FT.126–FT.132, FT.134, FT.144(5)–(6); the FT.93(3d) Schur tension is discharged by multiplicity, not violated \\
P-H19a & RETAIN & The read promotion ($\omega$-odd aspect reads act as multiplication operators on $\mathcal H_\omega$; $\hat Q_{\mathrm{EM}}=\allowbreak q_{\mathrm{core}}\hat\kappa$); consumers: FT.129(4)–(5), FT.130(5), FT.134(2), FT.139(3), FT.144(6); the carrier-vs-channel fusion tension stays logged (items 7/13) \\
P-H20a & ADJUDICATED VACUOUS & FT.140(1): the parity identification is frozen derived content [EM.24(2) derives $P$ from the dynamic phase; FT.101(3) records $\iota=\allowbreak P|_\Phi$]; retired content: all of it — no theorem uses it; future adoption: vacuous (nothing to adopt); the carrier-level parity-eigenvalue assignment is Q1$'$ content, not a declaration; entry retained as ledger history \\
\end{longtable}
}

%dsp-level:0
\[
(2)\quad\textbf{P-H11 adjudicated: RETIRED.}\ \ \text{Grounds:}\ \text{(i) zero live imports across H01–H22}\ [\text{machine inventory scan: every occurrence is registry text, the FT.91 construction record, or item}\ 21(h)];
\]

%dsp-level:0
\[
\ \text{(ii) its own costs column stands unresolved — the spectrum becomes observer-relative and no fermionic pole exists without a materialization clause the map does not carry};
\]

%dsp-level:0
\[
\text{(iii) the constructed channel interfaces (P-H14/P-H15 complexes) already supply sector spectra and expectations without an observer-channel declaration}\ [\mathrm{FT.102}\text{–}\mathrm{FT.108}];
\]

%dsp-level:0
\[
\ \text{(iv) the family channels derive their structure from arenas and seal integers alone}\ [\mathrm{FT.109}\text{–}\mathrm{FT.121},\ \mathrm{FT.141}\text{–}\mathrm{FT.145}]:\ \neg\exists\ \text{a computation routing through the bridge map};
\]

%dsp-level:0
\[
\text{future adoption would require: a declared observer channel with resolution-ladder reads}\ [\mathrm{FT.90}(2)\text{–}(3)];
\]

%dsp-level:0
\[
\ \text{a materialization/pole clause for cover-class spectra}\ [\mathrm{Def\text{-}H.16}\ \text{pattern}];
\]

%dsp-level:0
\[
\ \operatorname{Aut}\text{-invariance of the}\ \Sigma(\mathcal O)\ \text{read}\ [\mathrm{FT.91}(4):\ \text{already machine-verified}];
\]

%dsp-level:0
\[
\ \text{and a computable consequence not routable through P-H15's conditioning}\ [\text{the P-H18a supersession precedent, item}\ 21(i)];
\]

%dsp-level:0
\[
(3)\quad\textbf{the derived-away record:}\ \ \text{P-H2 by witness}\ [\mathrm{FT.36}(2)];
\]

%dsp-level:0
\[
\ \text{P-H3 in its counting half}\ [\mathrm{FT.41}(a)];
\]

%dsp-level:0
\[
\ \text{P-H8 in its non-coupling clause, in part}\ [\mathrm{FT.66}(4)];
\]

%dsp-level:0
\[
\ \text{P-H10 in its equilibrium form}\ [\mathrm{FT.108}(4)];
\]

%dsp-level:0
\[
\ \text{P-H20a in full}\ [\mathrm{FT.140}(1)]:\ \text{five entries retired or narrowed by derivation or witness — no silent retirement: each states what is retired, what remains, and what future adoption would require}\ [\text{table rows}];
\]

%dsp-level:0
\[
(4)\quad\textbf{the no-new-declaration clause:}\ \ \text{Tier F.42 issued none}\ [\mathrm{FT.145}(6)];
\]

%dsp-level:0
\[
\ \text{Tier F.43 issues none — this audit is ledger content: verdicts cite anchors, they declare nothing};
\]

%dsp-level:0
\[
\ \text{the registry after this audit:}\ 16\ \text{RETAIN},
\]

%dsp-level:0
\[
\ 1\ \text{DISCHARGED},
\]

%dsp-level:0
\[
\ 3\ \text{REDUCED},
\]

%dsp-level:0
\[
\ 1\ \text{RETIRED},
\]

%dsp-level:0
\[
\ 1\ \text{ADJUDICATED VACUOUS}\ (22\ \text{entries});
\]

%dsp-level:0
\[
(5)\quad\textbf{import discipline of this part:}\ \ \text{every P-H mention in FT.147–FT.152 is an audit citation}\ (\text{cited},\ \neg\text{imported}):\ \neg\exists\ \text{a clause of this part whose truth depends on a declared choice};
\]

%dsp-level:0
\[
\ \text{the verdicts themselves are derived from the cited anchor tags}.
\]

\textbf{Req.} The full Declared-Choice Registry (part S; first edition H13); FT.36, FT.41 (H05); FT.17, FT.28 (H02/H04); FT.20 (H03); FT.33 (H05); FT.54, FT.60, FT.61, FT.66 (H07/H08); FT.68, FT.70, FT.71, FT.74 (H09); FT.87, FT.90, FT.91 (H11); FT.102–FT.104 (H14); FT.105–FT.108 (H15); FT.110, FT.112 (H16); FT.118, FT.119 (H17); FT.123, FT.125, FT.126 (H18); FT.129, FT.130, FT.131 (H19); FT.134, FT.139, FT.140 (H20); FT.142, FT.144, FT.145 (H21); Def-H.3, Def-H.4, Def-H.8, Def-H.13, Def-H.16, Def-H.23, Def-H.24, Def-H.31, Def-H.32, Def-H.33; A5, D1.2, EM.9, EW.13, EW.20, EW.33, CT0.12, CT2.5, D2.11, D3.4, P1, P5 (\textbf{A}/\textbf{B}/\textbf{C}/\textbf{D}/\textbf{E}); inventory scans, count arithmetic [math; machine: \texttt{checks/\allowbreak{}verify\_\allowbreak{}H\_\allowbreak{}rung\_\allowbreak{}closure.py} §A].

\textbf{Pf.}

%dsp-level:0
\[
(1):\ \text{each row's anchor tags re-located in the cited documents}\ [\text{machine: existence scan}];
\]

%dsp-level:0
\[
\qquad \text{verdict vocabulary closed at five values;}\ 22\ \text{rows, each P-H entry exactly once}\ [\text{machine parse}]
\]

%dsp-level:0
\[
\text{RETAIN rows}: \neg\exists\ \text{discharging derivation}\ [\text{inventory scan of the cited consumers}];
\]

%dsp-level:0
\[
\qquad \text{DISCHARGED/REDUCED/VACUOUS/RETIRED rows}: \text{the discharging tag or the retirement grounds are cited verbatim}\ [\mathrm{FT.36}(2),\ \mathrm{FT.41}(a),\ \mathrm{FT.66}(4)\text{+}\mathrm{Def\text{-}H.24},\ \mathrm{FT.108}(4),\ \mathrm{FT.140}(1);\ \text{clause}\ (2)]
\]

%dsp-level:0
\[
(2):\ \text{grounds (i): occurrences of the P-H11 token inventoried}\ [\text{machine}];
\]

%dsp-level:0
\[
\qquad \text{(ii): the costs column re-cited from the registry entry verbatim};
\]

%dsp-level:0
\[
\qquad \text{(iii):}\ \mathrm{FT.105}\text{–}\mathrm{FT.108}'\text{s sector constructions carry no bridge-map premise}\ [\text{Req-line scan}];
\]

%dsp-level:0
\[
\qquad \text{(iv):}\ \mathrm{FT.109}\text{–}\mathrm{FT.121}/\mathrm{FT.141}\text{–}\mathrm{FT.145}\ \text{Req lines scanned:}\ \neg\text{P-H11}\ [\text{machine}]
\]

%dsp-level:0
\[
(3):\ \text{five rows, five cited anchors}\ [\text{machine: anchor existence}];
\]

%dsp-level:0
\[
\qquad \text{the no-silent-retirement clause: every non-RETAIN row carries its retired/residual/requires triad}\ [\text{machine: row-language scan}]
\]

%dsp-level:0
\[
(4):\ \mathrm{FT.145}(6)\ \text{verbatim};
\]

%dsp-level:0
\[
\qquad \text{count arithmetic}:\ 16+1+3+1+1=22\ [\text{exact integers}];
\]

%dsp-level:0
\[
\qquad (5):\ \text{the Req lines of this part name no P-H as premise}\ [\text{machine scan}].\ \blacksquare
\]

\textbf{Phys.} The registry closes its first full pass: sixteen declarations stand because nothing in the corpus derives their content, five have been paid down — one wholly by witness, one wholly by adjudication, three narrowed by derivation — and one candidate retires without ever having been adopted, its minimal form left on the shelf as computed data. The audit adds nothing and removes nothing silently: every retirement names what is retired, what remains, and the exact price of any future re-adoption. The working layer now knows precisely which of its stones are declared and which are derived.

\textbf{Status.} Closed; verdicts are ledger adjudications anchored to cited tags; no declaration issued, none imported as premise. Machine: §A.

\vspace{2pt}\noindent\rule{\textwidth}{0.4pt}\vspace{2pt}

\medskip
\subsection*{\texorpdfstring{B.1.8\quad FT.148 — The structural/census separation of the working layer}{FT.148 B.1.8}}
\addcontentsline{toc}{subsection}{B.1.8 FT.148}
\phantomsection\label{tag:FT-148}
{\small \textit{Record.}\quad Source: HistLog §IV.3 (the closure tier F.43) (archival ledger, preserved at \texttt{working\_markdowns/HistLog.md}); reproduced verbatim, including its historical registry tokens, which resolve via HistLog \S I.}
\medskip
\textbf{Thm.} Every theorem of H01–H21 is separated into exactly one of two categories:

%dsp-level:0
\[
(1)\quad\textbf{the classification rule:}
\]

%dsp-level:0
\[
\qquad \text{Category A (structural)}\iff\text{the tag's primary deliverable is an identity, mechanism, covariance law, no-go, uniqueness, or scope adjudication about the corpus's algebra, geometry or dynamics}\ (\text{even when proved by exhaustive enumeration});
\]

%dsp-level:0
\[
\text{Category B (data-table census)}\iff\text{the tag's primary deliverable is a computed table, card, ladder, orbit set, multiplet family or count over a named configuration space, and downstream consumers cite the tabulated values};
\]

%dsp-level:0
\[
\text{mixed tags are classified by primary deliverable; a census side of an A tag is cross-referenced in the H-CENSUS preamble}\ [\text{FT.68},\ \text{FT.93},\ \text{FT.103}\ \text{rows}];
\]

%dsp-level:0
\[
(2)\quad\textbf{the separation table}\ (\text{all}\ 145\ \text{theorems; machine: each tag exactly once, categories disjoint and complete}):
\]

%tbl-level:3
{\scriptsize\setlength{\tabcolsep}{4pt}
\begin{longtable}{>{\raggedright\arraybackslash}p{\dimexpr0.0624\textwidth-2\tabcolsep\relax}|>{\raggedright\arraybackslash}p{\dimexpr0.0559\textwidth-2\tabcolsep\relax}|>{\raggedright\arraybackslash}p{\dimexpr0.8817\textwidth-2\tabcolsep\relax}}
\toprule
tag & cat & justification \\
\addlinespace[2pt]
\midrule
\endfirsthead
\toprule
tag & cat & justification \\
\addlinespace[2pt]
\midrule
\endhead
\bottomrule
\endlastfoot
FT.1 & A & the flux bound — a measurement law, not a table \\
FT.2 & A & exhaustiveness of the real-edge measurement primitive — a scope law \\
FT.3 & A & unresolvability of the lock-sourced structure — a no-go \\
FT.4 & A & terminus channel characterization with the state-counting ceiling (Cor FT.4.1) — a bound law \\
FT.5 & A & the rate form of the flux bound — a law \\
FT.6 & A & conformal invariance of combinatorial flux — a covariance law \\
FT.7 & A & the universal flux-density bound — a law \\
FT.8 & A & the flux-mediated coarse-graining connection — a mechanism \\
FT.9 & A & rung stability of the EM edge-weight part — an invariance law \\
FT.10 & A & confinement as a global sum-restriction — a structural restriction \\
FT.11 & A & the complete-force charge audit — derived absence of force self-charge \\
FT.12 & A & weak has no self-charge — a no-go \\
FT.13 & A & weak $=$ spin, algebraically — an identification identity \\
FT.14 & A & the spin–charge decomposition with $N_{\mathrm{charge}}=\allowbreak 2$ derived — a decomposition theorem \\
FT.15 & A & coarse-graining invariance — an invariance law \\
FT.16 & A & genericity of the hidden-information gap; the strong signature is structural — a law plus adjudication \\
FT.17 & A & the internal-edge lemma — a mechanism \\
FT.18 & A & contraction-completeness is not $\Sigma$-local — a no-go \\
FT.19 & A & confinement is all-or-nothing — a law \\
FT.20 & A & reduction of the aspect set to proven per-aspect criteria — a mechanism \\
FT.21 & A & the charge-type rows properly derived — a derivation \\
FT.22 & A & loop quantization $\Theta(c)\in\allowbreak 2\mathbb Z$ on materialized chains — an identity \\
FT.23 & A & orientation-evenness of the line element — an identity \\
FT.24 & A & the minimal fermion pins the geometry quantum ($\kappa\in\allowbreak \mathbb Z$) — an identity \\
FT.25 & A & on-shell well-definedness of the internal holonomy — a law \\
FT.26 & A & structural transfer of the mask/resolvent machinery — a mechanism \\
FT.27 & A & single-vertex degeneration; action versus observable — a scope adjudication \\
FT.28 & B & the digon floor card — a computed cell (see H-CENSUS) \\
FT.29 & A & strong-weight factorization, no strong coupling constant, observable sealing — a mechanism plus no-go \\
FT.30 & A & the pair-channel sector quantum audit — block-wise independent-quantum adjudication \\
FT.31 & A & the continuum limit survives on-shell quantization — a scope law \\
FT.32 & B & the cycle-floor ladder $\Sigma_L$ — computed cells (see H-CENSUS) \\
FT.33 & A & the defect bridge — a mass-candidate mechanism (P-H7 anchor) \\
FT.34 & A & round-trip quantization, integrality propagation, the sealing partition rule — identities \\
FT.35 & B & the composite floor $\Sigma_{12}$ — a computed cell, first colour-active and lock-carrying (see H-CENSUS) \\
FT.36 & A & flux-ledger completion with the attainment witnesses — a law (discharges P-H2) \\
FT.37 & A & the relating equation, complete — an exact identity \\
FT.38 & A & non-geodesic tightening and on-shell rigidity — laws \\
FT.39 & A & the confinement converse, the singlet diagonal, the filter placement — adjudications \\
FT.40 & A & the closed-core charge read — a derived read formula, ternary and $\omega$-covariant \\
FT.41 & A & the per-aspect ledger trichotomy — derived counts ($N_{\mathrm{charge}}=\allowbreak 2$; the strong bit-count grammar) \\
FT.42 & B & the Platonic card template — computed cells, five solids (see H-CENSUS) \\
FT.43 & B & the winding ledger with Euler consistency — per-solid enumerations (see H-CENSUS) \\
FT.44 & B & the duality signature and floor taxonomy — a computed taxonomy table (see H-CENSUS) \\
FT.45 & B & the polygon floors $P_n$ — a computed family of cards (see H-CENSUS) \\
FT.46 & A & distinctness is unbounded — a ceiling no-go \\
FT.47 & A & no geometric ceiling: Gauss–Bonnet binds area, never $n$ — a no-go \\
FT.48 & A & stability structure of floor-to-floor channels — a variational law \\
FT.49 & B & the tetrahedron-face full card — a computed cell (see H-CENSUS) \\
FT.50 & B & composition and gluing — computed patch cells (bowtie, double-triangle) with their gluing data (see H-CENSUS) \\
FT.51 & B & the winding-ledger enumeration with dihedral degeneracy counts (see H-CENSUS) \\
FT.52 & A & spin extraction: the constraint system, uniqueness, the localized obstruction — uniqueness plus no-go \\
FT.53 & A & geodesic constancy of the stationary on-shell ledger — a conservation identity \\
FT.54 & A & the exact holonomy-defect formula — an identity (P-H7$'$ anchor) \\
FT.55 & A & the no-go and the surviving spin realization — an adjudication \\
FT.56 & A & native ladder rungs, the spin–orbit decomposition, the missing map — structure plus no-go \\
FT.57 & A & the boundary exchange algebra: root-space shape — a shape theorem \\
FT.58 & A & hosting classification and the no-hosting theorem — a no-go \\
FT.59 & A & the higher-rung specification: non-existence, persistence refuted, hosting ceiling — void/no-go adjudications \\
FT.60 & A & the multiplicity map: conditions, ladders, basis-relativity, honesty witnesses — a construction with uniqueness (P-H9 anchor) \\
FT.61 & A & the selection side: index-blindness and the two exits — a no-go \\
FT.62 & A & the bidirectional rung audit — a correction theorem \\
FT.63 & B & the Platonic physical cards — computed read columns (see H-CENSUS) \\
FT.64 & B & the filter chain and the per-solid loop lattices — computed lattices (see H-CENSUS) \\
FT.65 & B & coarse-grained degeneracies, register product, master loop list — computed enumerations (see H-CENSUS) \\
FT.66 & A & the Bott correction: witness re-evaluation, reconstructability, selection inadmissibility — a no-go plus reconstruction \\
FT.67 & A & the exact-cover theorem and the single-vertex bypass — a classification theorem \\
FT.68 & A & joint skeletons: the transfer rule $T(t)$ and the amplitude/norm layer dichotomy — mechanisms (the cover-polynomial tables are census supply, cross-referenced in the H-CENSUS preamble) \\
FT.69 & B & the Platonic amplitude census completed — computed amplitude values (see H-CENSUS) \\
FT.70 & A & the flicker phase diagram: unique crossover, exact boundaries, the polygon limit — laws \\
FT.71 & A & the flicker-dust equation of state — a derived law \\
FT.72 & A & two-fold face decomposition; the Bernoulli marginals — a decomposition \\
FT.73 & A & the bridge audit: dielectric non-existence, the field-degree no-go, surviving species — no-go cluster \\
FT.74 & A & the extended exact-cover theorems — theorems \\
FT.75 & A & domain-relativity of the amplitude layer with the invariant bridge gap — an invariance law \\
FT.76 & B & the Platonic pole-mass multiplets and the bounded-hierarchy ceiling data — tabulated multiplets (see H-CENSUS) \\
FT.77 & A & the bounded-hierarchy negative — a ceiling theorem (its external-hierarchy rows are audit-only) \\
FT.78 & B & the orbit decompositions: tetrahedron exhaustive, cube symmetry-reduced — computed orbit sets (see H-CENSUS) \\
FT.79 & A & the combinatorial ceiling and the cover-hierarchy audit — a ceiling theorem \\
FT.80 & A & the bridge census: amplitude-real, norm-real, spectrally un-indexed — a structural status theorem \\
FT.81 & B & the particle-content read-class census — a computed census (external-scoring columns retired by FT.149; see H-CENSUS) \\
FT.82 & A & gauge-block reconstruction: the algebra closes, the coupling structure does not — an algebra theorem plus negative \\
FT.83 & B & the rung-emergence ledger — a computed per-rung status table (see H-CENSUS) \\
FT.84 & B & the ceiling-taxonomy table of the mass-read species — a computed table (external scoring retired by FT.149; see H-CENSUS) \\
FT.85 & A & vacuum census equalities, constructible gravitational reads, the non-constructibility adjudication — adjudications \\
FT.86 & A & the flicker ripple: ultralocal dynamics, cap-induced alternation — dynamical laws \\
FT.87 & A & the tumbling-model adjudication, row by row — an adjudication \\
FT.88 & A & the involution algebra $\mathbb Z_2\times\mathbb Z_2$ — a structure theorem \\
FT.89 & A & the Koide/phase-boundary adjudication: exact arithmetic, the category trap — an adjudication \\
FT.90 & A & the resolution ladder: no-hair from the channel bound, the chiral merger refuted — laws plus no-go \\
FT.91 & A & the minimal spectral-bridge form — a construction (P-H11 candidate; retired by FT.147(2)) \\
FT.92 & A & convention-exactness and the rigidity of the named basis — a re-verification theorem \\
FT.93 & A & the native Clifford census and the weak-chirality pincer — a no-go (the frame catalogue is census supply, cross-referenced in the H-CENSUS preamble) \\
FT.94 & A & the chiral mass spectrum: not constructible as chiral — a no-go \\
FT.95 & A & flicker chirality: the triviality adjudication, the absent trajectory, the zero-expectation theorem — laws \\
FT.96 & B & the reversal-chirality census of the cover sector — computed orbit pairs (see H-CENSUS) \\
FT.97 & A & the form factor formalized: node ontology, reduction identities, the general ceiling — a formalization plus ceiling \\
FT.98 & B & the polygon form factors: exact census, saturating growth — computed tables (see H-CENSUS) \\
FT.99 & B & the Platonic form factors, dominance census, reduced denominators — computed tables (see H-CENSUS) \\
FT.100 & A & the alignment pattern refines the informational mass — a refinement law \\
FT.101 & A & the weak branch grading: aligned active, antialigned blind — a structural grading theorem \\
FT.102 & A & the fluctuation spectrum, the tree-level limit, the seal-softening walls — structure \\
FT.103 & A & four exact vacua and the band decomposition of the effective potential — structure (the susceptibility census is census supply, cross-referenced in the H-CENSUS preamble) \\
FT.104 & A & the range-species classification and the fourth read — a taxonomy theorem (audit rows quarantine external data) \\
FT.105 & B & the topology census from first principles: graph counts, canonical enumeration, finiteness cells (see H-CENSUS) \\
FT.106 & A & weight transport and the two-layer sector structure — structure \\
FT.107 & A & the observable sectors and the mass expectation — a formalism with the sector-degeneracy theorem \\
FT.108 & A & the Janus sum: move algebra, budget conservation, derived detailed balance — mechanisms (reduces P-H10) \\
FT.109 & A & the diagonal seal commutator $\sum_i[J_i,\allowbreak K_i]=\allowbreak 2i\Sigma$ — an exact identity \\
FT.110 & A & the bound-chain arena dimension — a derived count under the P-H16a interface \\
FT.111 & A & the $C_3$ spectrum of the covariant seal bilinear with the normalization — a representation theorem \\
FT.112 & A & the mixing angle $\delta_\ell=\allowbreak 2/9$ — an exact rational \\
FT.113 & A & the bilinear spectrum and the mass-squared ratios — exact closed forms with enclosures \\
FT.114 & B & the external-data audit rows with exact shortfall arithmetic — a quarantined data ledger (see H-CENSUS) \\
FT.115 & A & the adjudication summary — a provenance ledger theorem \\
FT.116 & A & the relabelling lattice, the two-level obstruction, the $\varepsilon$-selection of $C_3$ — a forced-group theorem \\
FT.117 & A & confinement blindness: the one-body no-go on sealed singlets — a no-go \\
FT.118 & A & the colour-dressed arena and the spectator numerator — derived counts and integers under the P-H17a interface \\
FT.119 & A & $\delta_q=\allowbreak 2/27$ with the triple-angle link — exact rationals and spectrum \\
FT.120 & B & the external-data audit ledger: scheme columns, exact shortfalls, the two exclusions — a quarantined data ledger (see H-CENSUS) \\
FT.121 & A & the adjudication summary — a provenance ledger theorem \\
FT.122 & A & the orientation doublet is corpus-constructible, with canonicity — a construction \\
FT.123 & A & weak covariance and $\omega$-covariance on the extended carrier — commutation laws \\
FT.124 & A & the seven-locus obstruction census — locus separations (a census of no-gos is structural) \\
FT.125 & A & the deficit-ledger adjudications: projector-level closure, the opened assignment row — a ledger theorem \\
FT.126 & A & weak-covariant chiral projectors with the independent $\mathbb Z_2$'s — a construction \\
FT.127 & A & the nine-aspect census with the census conclusion — a derived-action adjudication \\
FT.128 & A & the no-flip theorem at the derived locus — an identity \\
FT.129 & A & the canonical flip candidate and the read promotion — a uniqueness construction (P-H19a anchor) \\
FT.130 & A & tensor structure, Schmidt rank, the static spectrum — structure \\
FT.131 & A & scope: bound chains, no isolated charge, the deferrals — a scope theorem \\
FT.132 & A & the adjudication summary — a provenance ledger theorem \\
FT.133 & A & the configuration-space reversal calculus $\mathbb Z_2^4$ with both CPT lifts — a structure theorem \\
FT.134 & A & the carrier bookkeeping: the scalar-CPT theorem — a theorem \\
FT.135 & A & the weak branch grading under parity, label-level only — a re-derivation theorem \\
FT.136 & A & the rung-one $S$-grading and the three independent $\mathbb Z_2$'s — structure \\
FT.137 & A & cover and flicker behaviour under the reversal family — covariance laws ($v(R^{\mathrm T})=\allowbreak v(R)$; the sign character) \\
FT.138 & A & the pole spectrum under the reversal family — a symmetry law ($D(m)=\allowbreak D(-m)$) \\
FT.139 & A & the channels under the reversal family — an invariance law plus the flip-block character \\
FT.140 & A & the adjudication summary and the ledger — a ledger theorem \\
FT.141 & A & the rung-one chirality theorem: the $10/6$ grading of the complete force complex — a grading theorem \\
FT.142 & A & the rung-one arena and $\delta_\nu=\allowbreak 1/8$ — a derived adjudication \\
FT.143 & A & the rung-one family spectrum with the angle ledger — exact closed forms \\
FT.144 & A & the Janus-arrow theorem: strict alternation, no backward materialization, single-sector occupancy — a domain/occupancy theorem \\
FT.145 & A & the rung-one completion with the bridge to rung two — an assembly adjudication \\
\end{longtable}
}

%dsp-level:0
\[
(3)\quad\textbf{the census:}\qquad\#\mathrm A=120,\qquad\#\mathrm B=25,\qquad120+25=145\qquad[\text{machine parse of}\ (2)];
\]

%dsp-level:0
\[
(4)\quad\textbf{no dismissal:}\ \ \text{every Category-B tag is a load-bearing census for the baryonic era}\ [\text{the H-CENSUS block of this part states, per entry: the configuration space, the count, and the baryonic use}];
\]

%dsp-level:0
\[
\ \neg\exists\ \text{a Category-B tag classified as peripheral};
\]

%dsp-level:0
\[
(5)\quad\textbf{the definitions:}\ \ \text{Def-H.1–Def-H.34 are structural vocabulary}\ (\neg\text{censuses});
\]

%dsp-level:0
\[
\ \text{they are carried exactly once as full definition sections in the letter parts [FT.146(2)]};
\]

%dsp-level:0
\[
\ \text{corollaries}\ (\text{Cor FT.4.1},\ \text{Cor CT0.7a-type},\ \text{Cor L0a})\ \text{are classified under their parent tags}.
\]

\textbf{Req.} The full theorem inventory of H01–H21 (titles and Thm blocks at the point of classification); FT.146 (this tier, the consolidation law); the H-CENSUS block of this part; exact count arithmetic [math; machine: \texttt{checks/\allowbreak{}verify\_\allowbreak{}H\_\allowbreak{}rung\_\allowbreak{}closure.py} §B].

\textbf{Pf.}

%dsp-level:0
\[
(2):\ \text{row-wise classification by clause}\ (1)'s\ \text{rule on each tag's Thm deliverable}\ [\text{machine: table parsed; tags}\ 1..145\ \text{each exactly once}]
\]

%dsp-level:0
\[
(3):\ \text{category counts by exact integer census of the parsed rows}\ [\text{machine}];
\]

%dsp-level:0
\[
\qquad (4),(5):\ \text{block scans}\ [\mathrm{FT.146};\ \text{the H-CENSUS block};\ \text{machine}].\ \blacksquare
\]

\textbf{Phys.} The working layer splits cleanly once the question is asked of deliverables rather than proofs: one hundred twenty theorems say something structural about the corpus — a law, an identity, an impossibility, a scope — and twenty-five hand over computed cells: cards of floors, orbits of covers, multiplets of poles, ladders of winding, ledgers of audits. The second pile is not scrap. It is the configuration space the baryonic era enumerates over, and the separation makes that role explicit, tag by tag, with no residue and no double counting.

\textbf{Status.} Closed; the classification is a ledger theorem about the layer, not a modification of any tag. Machine: §B.

\vspace{2pt}\noindent\rule{\textwidth}{0.4pt}\vspace{2pt}

\medskip
\subsection*{\texorpdfstring{B.1.9\quad FT.149 — The framing retirement and the reframed question ledger}{FT.149 B.1.9}}
\addcontentsline{toc}{subsection}{B.1.9 FT.149}
\phantomsection\label{tag:FT-149}
{\small \textit{Record.}\quad Source: HistLog §IV.3 (the closure tier F.43) (archival ledger, preserved at \texttt{working\_markdowns/HistLog.md}); reproduced verbatim, including its historical registry tokens, which resolve via HistLog \S I.}
\medskip
\textbf{Thm.}

%dsp-level:0
\[
(1)\quad\textbf{the retirement adjudication}\ (\neg\text{silent}):\ \ \text{the external-reference framing is retired from the open-question ledger: the completeness-fraction scoring of the audits}\ (\text{the}\ 13/24\ \text{fraction and its per-sector rows}),
\]

%dsp-level:0
\[
\ \text{the Yukawa-as-deficit framing},
\]

%dsp-level:0
\[
\ \text{the CKM/PMNS-as-deficit framing},
\]

%dsp-level:0
\[
\ \text{the gauge-anomaly-as-deficit framing},
\]

%dsp-level:0
\[
\ \text{the four-dof-as-deficit framing},
\]

%dsp-level:0
\[
\ \text{and the species-mass-as-deficit framing};
\]

%dsp-level:0
\[
\text{the audit theorems themselves stand unmodified as derived content}\ [\mathrm{FT.77},\ \mathrm{FT.79},\ \mathrm{FT.81}\text{–}\mathrm{FT.84},\ \mathrm{FT.87},\ \mathrm{FT.89},\ \mathrm{FT.93},\ \mathrm{FT.94},\ \mathrm{FT.97}\text{–}\mathrm{FT.99},\ \mathrm{FT.104},\ \mathrm{FT.120},\ \mathrm{FT.125}]:\ \text{their negatives are consequences of derived ceilings and inventories},
\]

%dsp-level:0
\[
\ \neg\text{failures against an external target}\ [\mathrm{FT.84}(2)\ \text{synthesis, re-read corpus-internally}];
\]

%dsp-level:0
\[
\text{the corpus stands on its own; the external reference is test data under the audit discipline}\ [\mathrm{P1},\ \mathrm{P5};\ \mathrm{FT.73}(5)\ \text{pattern}]\ \text{and enters no question below};
\]

%dsp-level:0
\[
(2)\quad\textbf{the reframed question ledger}\ (\text{corpus-internal; the former deficit numbers are mapped, not erased}):
\]

%tbl-level:3
{\scriptsize\setlength{\tabcolsep}{4pt}
\begin{longtable}{>{\raggedright\arraybackslash}p{\dimexpr0.1028\textwidth-2\tabcolsep\relax}|>{\raggedright\arraybackslash}p{\dimexpr0.1239\textwidth-2\tabcolsep\relax}|>{\raggedright\arraybackslash}p{\dimexpr0.6282\textwidth-2\tabcolsep\relax}|>{\raggedright\arraybackslash}p{\dimexpr0.1450\textwidth-2\tabcolsep\relax}}
\toprule
Q & The corpus-internal question & Derived status & Open residue \\
\addlinespace[2pt]
\midrule
\endfirsthead
\toprule
Q & The corpus-internal question & Derived status & Open residue \\
\addlinespace[2pt]
\midrule
\endhead
\bottomrule
\endlastfoot
Q1 (former deficit 1) & Where does chirality live in the corpus? & Closed at three loci: the pair-register grading is a pincer no-go [FT.93(3)]; the carrier projectors are constructible under P-H18a [FT.122–FT.126]; the force-arena grading is the complete-force $10/6$ split [FT.136, FT.141] & none at the projector level; the frame-section residue sits at item 21(i) \\
Q1$'$ (root) & Which grading couples to which matter locus — the assignment rule? & The four sides are derived: the algebra $S$-grading [FT.136/FT.141], the branch-label parity [FT.101/FT.135], the carrier $\hat\chi$-grade [FT.122–FT.126], the Janus-arrow chirality [FT.144(5)] & the coupling rule itself — genuinely open; the root of Q4/Q5 and of the Q6 naming residue \\
Q2 (mass mechanism) & What are the corpus's mass reads and how do they relate? & The corpus has its own mass mechanisms, complete within their own terms: the informational read $m_{\mathrm{hid}}=\allowbreak E_{\mathrm{int}}\ln2/\mathfrak h$ [FT.16/FT.44(4)], the defect/pole reads [P-H7$'$; FT.33/FT.54/FT.76], the tilt-contraction read [L1–L3; P-H13], the sector expectations [FT.107], the 1-loop effective potential [P-H14], the channel ratios [FT.113/FT.119/FT.143]; an external Yukawa-type sector is a different mechanism, $\ne\allowbreak g$present and $\ne\allowbreak g$required by any derived tag & the relating equation across the five reads (item 7); the overall scales are P4-class \\
Q3 (family structure) & What is the corpus's family structure, and what selects its labels? & The democratic $C_3$ mechanism on three arenas: angles $2/9$, $2/27$, $1/8$; $Q=\allowbreak \tfrac23$ in every channel; family count $3=\allowbreak \dim F_{J^\pm}$ throughout [FT.109–FT.121, FT.141–FT.145]; the triple-angle link $3\delta_q=\allowbreak \delta_\ell$; the $\nu$-channel cubic channel-own & the selection/labeling principle (item 21(d)/item 12); identification barred [P1, P5] \\
Q4 (mixing structure) & What is the corpus's mixing structure? & The $C_3$-covariant seal bilinear with its channel angles and ledgers [FT.111–FT.113, FT.119, FT.143]; external CKM/PMNS-type matrices are a different structure, $\ne\allowbreak g$present & a cross-channel mixing object would require Q1$'$ and Q3 first; none is derived \\
Q5 (anomaly structure) & What is the corpus's anomaly-side structure? & The four $\mathbb Z_2$ gradings, the scalar-CPT theorem [FT.134(3)], the C-odd/P-even flip block [FT.129(5)/FT.139(3)]; external gauge-anomaly sums are chiral constraints on a declared assignment — formable only after Q1$'$ [FT.81(4), FT.125(2)], and they are test data, never inputs [P1, P5] & Q1$'$ \\
Q6 (rung-one row) & What does the rung-one channel carry? & Computed: $\delta_\nu=\allowbreak 1/8$; ratios $1:9.58\ldots:272.0\ldots$; $Q=\allowbreak \tfrac23$; single-sector occupancy with the vacuum complement; the flip channel opens at rung two [FT.141–FT.145] & the assignment rule (Q1$'$), the scale $\mu_\nu$ (P4), the relating equation (item 7) \\
Q7 (second scale) & Can a second dimensionful scale exist? & Derived no-go: one scale $\mathfrak h$ [A §0.3, P1]; the question is closed as a theorem, $\ne\allowbreak g$open & none \\
Q8 (H-block) & What is the scalar block's derived content? & $\dim F_H=\allowbreak 1=\allowbreak \operatorname{span}\{\Sigma\}$ [CT0.12]; the background $h_0$ is P4-class [EW.13, P-H6]; external multi-dof expectations are retired as framing & none — the allocation is derived and complete \\
Q9 (charge alphabet) & What charge values are readable? & The ternary core read $q\in\allowbreak \{-1,0,+1\}$ [EM.9/D3.4/FT.40]; colour is $\ne\allowbreak g$observable by construction [CT3.22]; confined fractional values are unreadable by theorem — consistent, $\ne\allowbreak g$defective & none at single-chain level; the multi-chain read is Q10 \\
Q10 (additive charge) & Does an additive multi-chain charge read exist? & FT.41(a) counts $N_{\mathrm{charge}}=\allowbreak 2$ end registers per chain; the single-chain read is derived & a multi-chain observable construction — not in the inventory; a baryonic-era candidate \\
Q11 (exchange structure) & Does a two-chain exchange structure exist? & No exchange operator in the inventory [FT.80(5)]; the single-chain grade is not an exchange sign [FT.125(5), FT.89(2) kind discipline] & a two-chain state space — baryonic-era candidate \\
Q12 (species calibration) & Which read carries which species? & Not a corpus-internal question: calibration is barred [P1, P5 — $\ne\allowbreak g\exists$ datum]; the five read species carry exact ceilings/ranges [FT.76(3), FT.79(2), FT.97(5), FT.104(2), FT.107(4)] & retired with the external framing; internally open: item 7 and Q1$'$ \\
\end{longtable}
}

%dsp-level:0
\[
(3)\quad\textbf{the genuinely open questions survive the reframe unchanged in content:}\ \ \text{the relating equation between the mass reads}\ (\text{item}\ 7);
\]

%dsp-level:0
\[
\ \text{the family selection/labeling principle}\ (\text{item}\ 12);
\]

%dsp-level:0
\[
\ \text{the on-shell completion residue}\ (\text{item}\ 13);
\]

%dsp-level:0
\[
\ \text{the momentum-space resolvent}\ (\text{item}\ 17);
\]

%dsp-level:0
\[
\ \text{the seed-tensoring derivation residue}\ (\text{item}\ 19);
\]

%dsp-level:0
\[
\ \text{the assignment rule}\ (\text{Q1}');
\]

%dsp-level:0
\[
(4)\quad\textbf{external-data quarantine stands:}\ \ \text{the supplied external hierarchies and audit columns remain confined to their audit rows}\ [\mathrm{FT.104},\ \mathrm{FT.114},\ \mathrm{FT.120}]\ \text{and to no derivation}\ [\text{machine: the H-CENSUS block and the letter registers carry no external token}].
\]

\textbf{Req.} FT.77, FT.79, FT.81–FT.84 (H10); FT.87, FT.89 (H11); FT.93, FT.94, FT.97–FT.99 (H12); FT.104 (H14); FT.114 (H16); FT.120 (H17); FT.125 (H18); FT.134, FT.139, FT.140 (H20); FT.141–FT.145 (H21); P1, P5 (\textbf{A}); item 7/12/13/17/19 (part S); token-confinement scans [math; machine: \texttt{checks/\allowbreak{}verify\_\allowbreak{}H\_\allowbreak{}rung\_\allowbreak{}closure.py} §D].

\textbf{Pf.}

%dsp-level:0
\[
(1):\ \text{each retired framing named once, with its audit-theorem lineage cited}\ [\text{machine: framing tokens confined to this tag and to H13's retirement record}];
\]

%dsp-level:0
\[
\qquad \text{the standing clauses: the cited tags re-located unmodified}\ [\text{machine existence}]
\]

%dsp-level:0
\[
(2):\ \text{each Q-row's status rows cite derived anchors}\ [\text{machine existence}];
\]

%dsp-level:0
\[
\qquad (3):\ \text{the open items re-cited verbatim from the gap log}\ [\text{H13; the row-wise audit is the next tag — a non-load-bearing pointer}];
\]

%dsp-level:0
\[
\qquad (4):\ \mathrm{FT.114}(3)/\mathrm{FT.120}\ \text{quarantine discipline re-cited};\ \text{block scans}\ [\text{machine}].\ \blacksquare
\]

\textbf{Phys.} The ledger stops asking where the corpus falls short of somebody else's particle table and starts asking what the corpus has not yet answered about itself. The retirement is not a deletion: every audit theorem stands, every ceiling still prices every read, and the negatives keep their force — they bound what the derived reads can do. What changes is the direction of the questions. Twelve rows remain, and each is now a question with a corpus-internal answer space: an assignment rule, a relating equation, a selection principle, an additive charge read, an exchange space. The external numbers stay exactly where the audit discipline put them: in the quarantine rows, tested, never trusted.

\textbf{Status.} Closed; the reframe modifies no theorem, only the ledger's question grammar; the retirement is recorded, not silent. Machine: §D.

\vspace{2pt}\noindent\rule{\textwidth}{0.4pt}\vspace{2pt}

\medskip
\subsection*{\texorpdfstring{B.1.10\quad FT.150 — The gap-closure audit}{FT.150 B.1.10}}
\addcontentsline{toc}{subsection}{B.1.10 FT.150}
\phantomsection\label{tag:FT-150}
{\small \textit{Record.}\quad Source: HistLog §IV.3 (the closure tier F.43) (archival ledger, preserved at \texttt{working\_markdowns/HistLog.md}); reproduced verbatim, including its historical registry tokens, which resolve via HistLog \S I.}
\medskip
\textbf{Thm.} Every item of the Open Gaps Log is audited; each carries exactly one verdict from $\{$CLOSED, DEFERRED, CLOSED-DEFERRED, REFRAMED, AUDITED$\}$ (corpus numbering preserved — no renumbering):

%dsp-level:0
\[
(1)\quad\textbf{the verdict table:}
\]

%tbl-level:3
{\scriptsize\setlength{\tabcolsep}{4pt}
\begin{longtable}{>{\raggedright\arraybackslash}p{\dimexpr0.0449\textwidth-2\tabcolsep\relax}|>{\raggedright\arraybackslash}p{\dimexpr0.1876\textwidth-2\tabcolsep\relax}|>{\raggedright\arraybackslash}p{\dimexpr0.1197\textwidth-2\tabcolsep\relax}|>{\raggedright\arraybackslash}p{\dimexpr0.6478\textwidth-2\tabcolsep\relax}}
\toprule
item & title & verdict & closing tag(s) / residue and requirement \\
\addlinespace[2pt]
\midrule
\endfirsthead
\toprule
item & title & verdict & closing tag(s) / residue and requirement \\
\addlinespace[2pt]
\midrule
\endhead
\bottomrule
\endlastfoot
1 & coarse-graining scale invariance & CLOSED & FT.36(1): the saturation locus is the conformally invariant channel speed; the saturated value is $\Omega$-free \\
2 & relating equation (loop-area floor $\le\allowbreak ftrightarrow$ qubit-power floor) & CLOSED & FT.37: the complete form with the exact $\bar\kappa$ classification \\
3 & Def-H.3 modeling choice (P-H1) & CLOSED & the audit: coherent with every computed card; retained (P-H1 verdict: RETAIN) \\
4 & FT.1 non-geodesic tightening & CLOSED & FT.38(1): the proper-length form on any spacelike on-shell integer-edge path \\
5 & fermionic quantum numbers below the channel ceiling & CLOSED-DEFERRED & closed: the computed ledger — cards [FT.28], ladders [FT.32], floors [FT.42–FT.51], spin under P-H9 [FT.60], channels [FT.109–FT.121, FT.141–FT.145]; residue: the physical reading of the $(L,\allowbreak \kappa,\allowbreak k)$ ladder, contingent on items 11–12-lineage questions — now Q3; requires the selection principle (item 12) \\
6 & per-aspect bit counts & CLOSED & FT.41: the trichotomy — $N_{\mathrm{charge}}=\allowbreak 2$ derived; $N_{\mathrm{strong}}$ the sealing-grouping grammar; $G/H$ ledgers not bit-type \\
7 & relating equation (geometric $\le\allowbreak ftrightarrow$ informational) & DEFERRED & five reads exist ($m_{\mathrm{hid}}$, pole/defect, tilt, sector, 1-loop) plus the channel ratios; no declared bridge (P-H12 is a bridge statement, not a derivation); requires: a quantitative derivation relating $\Delta S(q)$ to $E_{\mathrm{int}}\ln2/\mathfrak h$-type reads, or a declared adoption with a dimensional map [FT.89(2) category discipline]; baryonic-era relevance: composed objects force the reads onto common cells \\
8 & FT.19 converse & CLOSED & FT.39(1)–(2): the converse with the exact exception locus \\
9 & Def-H.9 amendment decision & CLOSED & FT.39(3): the downstream filter placement derived \\
10 & sup-attainment (P-H2) & CLOSED & FT.36(2): witnesses exhibited; P-H2 discharged \\
11 & spin-label extraction & CLOSED & FT.60 under P-H9: $j=\allowbreak |k|/2$ with $\mathrm{Sym}^k$ and the diagonal transport; uniqueness derived; residue tracked at item 19 \\
12 & generation/flavour structure & DEFERRED & the combinatorial precondition is met and exceeded (unbounded ledger families); the family count is $3=\allowbreak \dim F_{J^\pm}$ in every channel — a datum, not a selection; requires: a principle selecting/labeling families from the computed pools [item 21(d)]; no capacity obstruction [FT.90(5): $\lceil\log_23\rceil=\allowbreak 2$] \\
13 & on-shell completion (line element) & CLOSED-DEFERRED & closed: classes (i) $\kappa=\allowbreak 0$ and (ii) $\kappa\in\allowbreak \mathbb Z$ on round-trip-closed channels [FT.24/FT.34(4)]; class (iii) semiclassically excluded [FT.53(4)]; residue: finite-$\mathfrak h$ non-stationary materializations — identical to the item-18 selection residue; requires: exit (a$''$) a declared weight extension $w\mapsto\allowbreak  w(-1)^{\operatorname{ind}_B}$ (not adopted) or exit (b) a spectral-side positivity construction; baryonic-era relevance: sealed junctions are the natural site of the residual class \\
14 & "strictly new at $\Sigma$" & CLOSED & FT.41(d) for the first targets; the exchange-sector program completed by FT.57 (root-space shape), FT.58 (hosting/no-hosting), FT.59 (the higher-rung void) \\
15 & closed-core charge read & CLOSED & FT.40: the ternary $\omega$-covariant core read \\
16 & $\mathcal A(\Sigma;C)$ classification & CLOSED-DEFERRED & closed: the computed cells (floors, ladders, gluing), residue (a) by FT.59, the hosting ceiling $\mathfrak{su}(3)$, the symmetry-class pools [FT.78], the rung ladder audited [FT.83]; residues: (b) the exchange-block declaration question (joined by the $k\ge\allowbreak 3$ carrier flag [FT.60(8)]), (c) the remaining cells (multi-solid chains beyond two floors, the single-vertex ladder, separated-contact hosting maps for $N\le\allowbreak 3$), (d) the $\aleph$-domain declaration [FT.74/FT.75: both path integrals computed, pairwise inequivalent — one author's declaration], (e) the defect corner convention [FT.76(1)–(2): variants A/B priced]; each requires a logged declaration or a new cell computation; (c)/(d)/(e) are baryonic-era gates \\
17 & momentum-space $\Sigma$-resolvent & DEFERRED & blocked pending a declared choice; requires: declared plane-wave boundary-channel states on the coarse-grained channel, a positivity or weight-extension exit (item 13 lineage), and registry logging before use; baryonic-era relevance: momentum reads of sealed objects would consume it \\
18 & stability–materialization coupling & CLOSED & FT.66: reconstructability plus the selection no-go; the open-channel residue transferred to item 13 in exact form \\
19 & seed-tensoring mechanism & CLOSED-DEFERRED & closed: the map constructed under P-H9 [Def-H.23/FT.60] with uniqueness and the composition law; residues: (i) derivation or replacement of P-H9 (joins the P-H7 family as one transfer residue under item 7), (ii) the $k\ge\allowbreak 3$ carrier-inventory question [FT.60(8) — no registered aspect block carries a $4$-dim irreducible]; requires: a derivation of multiplicity$=$winding from transport, or an inventory extension declaration \\
20 & the question ledger & REFRAMED & FT.149: the external-reference framing retired; the twelve rows re-posed as corpus-internal questions Q1–Q12 with derived status; the bracketed per-tier entries retained verbatim \\
21 & declarations the audits named & AUDITED & FT.150(3): (a) deferred, (b) deferred, (c) closed-deferred, (d) deferred (folded into Q3/item 12), (e) deferred, (f) closed [FT.89(2)], (g) closed [FT.90], (h) retired-closed [FT.147(2)], (i) deferred \\
\end{longtable}
}

%dsp-level:0
\[
(2)\quad\textbf{the deferral discipline:}\ \ \text{every DEFERRED and CLOSED-DEFERRED row states its residue, why closure precedes the baryonic era in dependency order (the required construction lives in the rung-3/4 or multi-chain or declared-dynamics stratum, none of which the fermionic pretime layer supplies), and what would close it}\ [\text{machine: row-language scan}];
\]

%dsp-level:0
\[
(3)\quad\textbf{item 21's sub-item dispositions:}
\]

%dsp-level:0
\[
\qquad (a)\ \text{DEFERRED}\ (\text{rotor/embedding datum: not adopted; adoption would require a declared circumradius or moment-of-inertia read — a P-H species});
\]

%dsp-level:0
\[
\qquad (b)\ \text{DEFERRED}\ (\text{tick-rate/field coupling: prerequisite named — emergent time}\ [\mathrm{FT.131}(3)]);
\]

%dsp-level:0
\[
(c)\ \text{CLOSED-DEFERRED}\ (\text{the involution algebra is exact at register level}\ [\mathrm{FT.133}(2)]\ \text{and scalar on the doublet}\ [\mathrm{FT.134}(3)];\ \text{the pair-register projector question is closed-negative}\ [\mathrm{FT.101}(3),\ \mathrm{FT.135}(4)];\ \text{the doublet side closed-positive}\ [\mathrm{H18}];\ \text{residue: eigenspace projectors on a declared state space});
\]

%dsp-level:0
\[
\qquad (d)\ \text{DEFERRED}\ (\text{generation-count selection: folded into Q3/item}\ 12);
\]

%dsp-level:0
\[
(e)\ \text{DEFERRED}\ (\text{per-sector entropy: requires a declared}\ \Psi_X\ \text{per configuration plus per-aspect norms});
\]

%dsp-level:0
\[
\qquad (f)\ \text{CLOSED}\ (\text{category trap}\ [\mathrm{FT.89}(2)]);
\]

%dsp-level:0
\[
\qquad (g)\ \text{CLOSED}\ (\text{resolution ladder exact}\ [\mathrm{FT.90}(2)\text{–}(3)]);
\]

%dsp-level:0
\[
\qquad (h)\ \text{RETIRED-CLOSED}\ (\text{P-H11 retired}\ [\mathrm{FT.147}(2)];\ \text{the minimal form remains computed data}\ [\mathrm{FT.91}(3)]);
\]

%dsp-level:0
\[
(i)\ \text{DEFERRED}\ (\text{the frame section}\ e\mapsto\mathcal F_e,\ \text{its EW.32-transport coupling, and a global Dirac-type operator remain unconstructed; the P-H18a supersession of the 'no computable consequence' clause stands});
\]

%dsp-level:0
\[
(4)\quad\textbf{the audit adds no closure the cited tags do not carry:}\ \ \text{every CLOSED row's closing tag is machine-located in the corpus};
\]

%dsp-level:0
\[
\ \text{no item moves between the log's Closed and Open sections by this audit}\ [\text{the CLOSED verdicts confirm existing closures; CLOSED-DEFERRED and DEFERRED items remain open with sharpened residues}].
\]

\textbf{Req.} The Open Gaps Log (part S) in full; the closing tags cited row-wise (FT.36–FT.41, FT.60, FT.66, FT.141–FT.145 lines); FT.147 (this tier, the P-H11 disposition); FT.149 (the item-20 reframe); P1, P4, P5 (\textbf{A}); row-language scans [math; machine: \texttt{checks/\allowbreak{}verify\_\allowbreak{}H\_\allowbreak{}rung\_\allowbreak{}closure.py} §D].

\textbf{Pf.}

%dsp-level:0
\[
(1):\ \text{row-wise: each verdict cites its closing tag(s) or its residue clause}\ [\text{machine: items}\ 1..21\ \text{each exactly once; verdict vocabulary closed}];
\]

%dsp-level:0
\[
\qquad \text{closing tags re-located in the corpus}\ [\text{machine existence}]
\]

%dsp-level:0
\[
(2):\ \text{residue/requires language present in every DEFERRED and CLOSED-DEFERRED row}\ [\text{machine scan}];
\]

%dsp-level:0
\[
\qquad (3):\ \text{sub-item dispositions cite their anchors}\ [\text{machine existence}];
\]

%dsp-level:0
\[
\qquad (4):\ \text{section-membership scan of the gap log unchanged}\ [\text{machine}].\ \blacksquare
\]

\textbf{Phys.} The log balances. Twelve items stand closed — each with the tag that closed it named on the row; six stand deferred with the residue spelled out and the price of closure stated; the two ledger items carry their reframe and their sub-audit. Nothing is waved away: where a closure is partial, the row says which half is closed and which half waits, and every waiting row names the stratum it waits for. The fermionic layer ends its run with its own books auditable line by line.

\textbf{Status.} Closed; verdicts confirm or sharpen — no item is closed by this audit that the cited tags do not already close. Machine: §D.

\vspace{2pt}\noindent\rule{\textwidth}{0.4pt}\vspace{2pt}

\medskip
\subsection*{\texorpdfstring{B.1.11\quad FT.151 — Baryonic preparation}{FT.151 B.1.11}}
\addcontentsline{toc}{subsection}{B.1.11 FT.151}
\phantomsection\label{tag:FT-151}
{\small \textit{Record.}\quad Source: HistLog §IV.3 (the closure tier F.43) (archival ledger, preserved at \texttt{working\_markdowns/HistLog.md}); reproduced verbatim, including its historical registry tokens, which resolve via HistLog \S I.}
\medskip
\textbf{Thm.}

%dsp-level:0
\[
(1)\quad\textbf{the foundation stack:}\ \ \text{the baryonic era builds on:}\ \text{(i) the frozen colour and sealing machinery}\ [\mathrm{CT3.1}\text{–}\mathrm{CT3.22}:\ \text{slot algebras},\ \mathfrak{su}(3),\ \varepsilon/\delta\ \text{sealing},\ \text{confinement};\ \mathrm{D3.8},\ \mathrm{D3.9};\ \mathrm{CT4.1}\text{–}\mathrm{CT4.7}:\ \text{the}\ n\ge4\ \text{lock side}];
\]

%dsp-level:0
\[
\ \text{(ii) the derived sealing grammar and junction structure}\ [\mathrm{FT.29},\ \mathrm{FT.34}(5),\ \mathrm{FT.19},\ \mathrm{FT.39}];
\]

%dsp-level:0
\[
\ \text{(iii) the computed baryon-type cells}\ [\mathrm{FT.35}:\ \Sigma_{12}\ \text{first colour-active, first}\ D^-\text{-carrying};\ \mathrm{FT.49}:\ \text{the tetrahedron-face baryon-type sealing phase},\ \Theta_{\mathrm{int}}\in4\mathbb Z];
\]

%dsp-level:0
\[
\text{(iv) the channel machinery}\ [\mathrm{FT.109}\text{–}\mathrm{FT.121}:\ \text{the}\ q\text{-channel with its}\ \varepsilon\text{-selected}\ C_3\ \text{and confinement blindness};\ \mathrm{FT.141}\text{–}\mathrm{FT.145}:\ \text{the rung-one channel and the arrow law}];
\]

%dsp-level:0
\[
\ \text{(v) the carrier and parity structure}\ [\mathrm{FT.122}\text{–}\mathrm{FT.134}];
\]

%dsp-level:0
\[
\ \text{(vi) the sector formalisms}\ [\mathrm{FT.102}\text{–}\mathrm{FT.108}:\ \text{1-loop bridge and topology-blind integral}];
\]

%dsp-level:0
\[
(2)\quad\textbf{the census consumption table}\ (\text{per H-CENSUS family; what the baryonic era builds on each}):
\]

%tbl-level:3
{\scriptsize\setlength{\tabcolsep}{4pt}
\begin{longtable}{>{\raggedright\arraybackslash}p{\dimexpr0.1199\textwidth-2\tabcolsep\relax}|>{\raggedright\arraybackslash}p{\dimexpr0.1437\textwidth-2\tabcolsep\relax}|>{\raggedright\arraybackslash}p{\dimexpr0.2631\textwidth-2\tabcolsep\relax}|>{\raggedright\arraybackslash}p{\dimexpr0.4732\textwidth-2\tabcolsep\relax}}
\toprule
census family & entries & configuration space provided & what the baryonic era builds on it \\
\addlinespace[2pt]
\midrule
\endfirsthead
\toprule
census family & entries & configuration space provided & what the baryonic era builds on it \\
\addlinespace[2pt]
\midrule
\endhead
\bottomrule
\endlastfoot
floor cards and ladders & FT.28, FT.32, FT.35, FT.42, FT.43, FT.44, FT.45, FT.49, FT.50, FT.51, FT.63 & the computed fermion floors: digon, cycles, composite, Platonic, polygon, tetrahedron-face, glued patches — with their full read columns & the cell grammar of sealed composites: baryon-type junctions are gluings of these cards; the winding ledgers and duality splits are the multiplet bookkeeping; the $\Sigma_{12}$/tetrahedron-face cells carry the first colour-active and $4\mathbb Z$-phased junctions \\
loop and cover classes & FT.64, FT.65, FT.69, FT.78, FT.96 & admissible-loop lattices, coarse-grained degeneracies, amplitude censuses, $\operatorname{Aut}$-orbits of covers with chirality tags & the enumeration substrate of sealed-object path integrals: sector conditioning $(q,\allowbreak S)$ runs over these pools; the chirality-tagged orbits feed the parity reads \\
pole and mass multiplets & FT.76 & defect-angle pole reads under both convention variants, with the ceiling & the pole-side mass cells and the hard range ceiling every baryonic read inherits; the convention variants are priced, so a declaration converts them into data immediately \\
form-factor tables & FT.98, FT.99 & polygon and Platonic form factors with growth and dominance rows & the size-bound range checks for extended baryonic skeletons \\
audit ledgers & FT.81, FT.83, FT.84, FT.114, FT.120 & the read-class inventory, the rung-emergence ledger, the ceiling taxonomy, the quarantined external-comparison rows & the card grammar and the P1/P5-compliant comparison template for the era's external tests; the activation map at $n=\allowbreak 3/n\ge\allowbreak 4$; the read-species taxonomy new reads must join \\
topology-blind cells & FT.105 & graph counts and canonical skeletons from first principles, with the budget regularization & the regulated topology pool of the baryonic path integral — the space its fugacity formalism sums over \\
\end{longtable}
}

%dsp-level:0
\[
(3)\quad\textbf{the declaration-requirements ledger}\ (\text{stated},\ \neg\text{issued}):\ \ \text{the baryonic era will require, at its own points of use:}\ \text{(a) the}\ \aleph\text{-domain choice}\ [\text{item}\ 16(d):\ \text{literal vs chain-local; both path integrals computed and pairwise inequivalent}\ [\mathrm{FT.67},\ \mathrm{FT.74},\ \mathrm{FT.75}]];
\]

%dsp-level:0
\[
\text{(b) the defect corner convention}\ [\text{item}\ 16(e):\ \text{variants A/B computed}\ [\mathrm{FT.76}(1)\text{–}(2)]];
\]

%dsp-level:0
\[
\qquad \text{(c) the exchange-block question}\ [\text{item}\ 16(b):\ \text{whether a}\ \mathcal K^{(\mathrm{exch})}\ \text{is ever declared; joined by the}\ k\ge3\ \text{carrier-inventory flag}\ [\mathrm{FT.60}(8)]];
\]

%dsp-level:0
\[
\text{(d) the dynamics rates}\ [\text{P-H10's residual: a dynamics construction supplying}\ J^\pm\ \text{rates, or the free datum retained};\ \text{item}\ 21(b)'s\ \text{tick-rate coupling prerequisite}];
\]

%dsp-level:0
\[
\qquad \text{(e) an observer/spectral bridge, if cover-class spectra are read}\ [\mathrm{FT.147}(2)'s\ \text{adoption requirements}];
\]

%dsp-level:0
\[
\text{(f) a multi-chain additive charge read}\ [\text{Q10}];
\]

%dsp-level:0
\[
\qquad \text{(g) a two-chain exchange space}\ [\text{Q11}];
\]

%dsp-level:0
\[
\qquad \text{(h) an embedding datum, if rotor reads are wanted}\ [\text{item}\ 21(a)];
\]

%dsp-level:0
\[
\text{each would be logged before use}\ [\text{the registry discipline}];
\]

%dsp-level:0
\[
\ \neg\exists\ \text{a declaration issued by this part};
\]

%dsp-level:0
\[
(4)\quad\textbf{the scope bars:}\ \ \text{no particle identification}\ [\text{P1},\ \text{P5}];
\]

%dsp-level:0
\[
\ \text{no external data}\ ([\text{EXT}]\ \text{enters no clause of this document — this standing-rule record is its sole marker});
\]

%dsp-level:0
\[
\ \text{all topology labels are internal bookkeeping}\ [\mathrm{P\text{-}H15}(c)\ \text{discipline}];
\]

%dsp-level:0
\[
\ \text{no H-sector interpretation leaks into any theorem statement};
\]

%dsp-level:0
\[
\ \text{amplitude}\ne\text{mass};
\]

%dsp-level:0
\[
\ \text{rung}\ne\text{scale}.
\]

\textbf{Req.} CT3.1–CT3.22 (\textbf{F}); CT4.1–CT4.7 (\textbf{G}); D3.8, D3.9 (\textbf{A}); FT.19, FT.29, FT.34, FT.35, FT.49 (H03/H04/H06); FT.60, FT.67, FT.74, FT.75, FT.76 (H07/H09/H10); FT.102–FT.108 (H14/H15); FT.109–FT.121 (H16/H17); FT.122–FT.134 (H18/H19/H20); FT.141–FT.145 (H21); FT.147, FT.149, FT.150 (this tier); P1, P4, P5 (\textbf{A}); items 16/21 (part S); inventory scans [math; machine: \texttt{checks/\allowbreak{}verify\_\allowbreak{}H\_\allowbreak{}rung\_\allowbreak{}closure.py} §E].

\textbf{Pf.}

%dsp-level:0
\[
(1):\ \text{each foundation row's tags re-located}\ [\text{machine existence}];
\]

%dsp-level:0
\[
\qquad (2):\ \text{family rows reference the H-CENSUS entries by tag}\ [\text{machine: every B tag appears in exactly one family row}]
\]

%dsp-level:0
\[
(3):\ \text{requirement rows cite their gap/registry anchors}\ [\text{machine existence}];
\]

%dsp-level:0
\[
\qquad \text{the stated-not-issued clause:}\ \neg\exists\ \text{a new P-H token in this part}\ [\text{machine scan}];
\]

%dsp-level:0
\[
\qquad (4):\ \text{standing-rule record}. \blacksquare
\]

\textbf{Phys.} The baryonic era does not start from zero and it does not start from hope. The colour algebra, the sealing grammar, and the lock side are frozen; the first colour-active cells are already computed, with their junction phases and their four-fold winding ledgers; the channel machinery that produced three family angles runs on any arena the era can derive; and twenty-five census supplies stand labeled with the spaces they enumerate. What the era must add is short and named: a domain choice here, a convention there, rates if it wants dynamics, an observer if it wants spectra — each a logged declaration at its point of use, none issued in advance. The rung closes with its debts itemized.

\textbf{Status.} Closed; preparation content only — no construction, no declaration, no identification. Machine: §E.

\vspace{2pt}\noindent\rule{\textwidth}{0.4pt}\vspace{2pt}

\medskip
\subsection*{\texorpdfstring{B.1.12\quad FT.152 — The status of the fermionic sector after consolidation}{FT.152 B.1.12}}
\addcontentsline{toc}{subsection}{B.1.12 FT.152}
\phantomsection\label{tag:FT-152}
{\small \textit{Record.}\quad Source: HistLog §IV.3 (the closure tier F.43) (archival ledger, preserved at \texttt{working\_markdowns/HistLog.md}); reproduced verbatim, including its historical registry tokens, which resolve via HistLog \S I.}
\medskip
\textbf{Thm.}

%dsp-level:0
\[
(1)\quad\textbf{the sector status:}
\]

%dsp-level:0
\[
\qquad \text{CLOSED (derived), as organized by the letter series:}\ \text{the measurement primitives and flux laws}\ [\text{part H}];
\]

%dsp-level:0
\[
\ \text{the edge/spin/charge structure with the ternary read}\ [\text{part H}];
\]

%dsp-level:0
\[
\ \text{the aspect machinery and the}\ 16\text{-dimensional pair-channel allocation}\ [\text{part H}];
\]

%dsp-level:0
\[
\ \text{on-shell quantization, the defect and holonomy formulas, the floors' relating equation}\ [\text{parts J, K}];
\]

%dsp-level:0
\[
\text{the stability/hosting/multiplicity structure with its no-gos}\ [\text{part K}];
\]

%dsp-level:0
\[
\ \text{the flicker and path-integral dynamics with derived detailed balance}\ [\text{parts L, O}];
\]

%dsp-level:0
\[
\ \text{chirality at three loci with the parity calculus and the scalar-CPT theorem}\ [\text{parts N, Q}];
\]

%dsp-level:0
\[
\ \text{the three family channels and the rung-one closure}\ [\text{parts P, R}];
\]

%dsp-level:0
\[
\ \text{the ceilings and negatives pricing every read}\ [\text{parts M, N}];
\]

%dsp-level:0
\[
\text{DEFERRED (open): the assignment rule}\ (\text{Q1}');
\]

%dsp-level:0
\[
\ \text{the relating equation across the five mass reads}\ (\text{item}\ 7);
\]

%dsp-level:0
\[
\ \text{the family selection/labeling principle}\ (\text{Q3/item}\ 12);
\]

%dsp-level:0
\[
\ \text{the on-shell class-(iii) residue}\ (\text{item}\ 13);
\]

%dsp-level:0
\[
\ \text{the remaining}\ \mathcal A(\Sigma;C)\ \text{cells and the two convention declarations}\ (\text{item}\ 16);
\]

%dsp-level:0
\[
\ \text{the momentum-space resolvent}\ (\text{item}\ 17);
\]

%dsp-level:0
\[
\ \text{the seed-tensoring derivation residue}\ (\text{item}\ 19);
\]

%dsp-level:0
\[
\ \text{the multi-chain and exchange-space constructions}\ (\text{Q10/Q11});
\]

%dsp-level:0
\[
\text{BUILDABLE (the baryonic foundation):}\ \mathrm{FT.151}(1)\text{'s stack with}\ \mathrm{FT.151}(2)\text{'s census consumption and}\ \mathrm{FT.151}(3)\text{'s declaration ledger};
\]

%dsp-level:0
\[
(2)\quad\textbf{the P-H status row:}
\]

%dsp-level:0
\[
\qquad 16\ \text{RETAIN},
\]

%dsp-level:0
\[
\ 1\ \text{DISCHARGED}\ (\text{P-H2}),
\]

%dsp-level:0
\[
\ 3\ \text{REDUCED}\ (\text{P-H3},\ \text{P-H8},\ \text{P-H10}),
\]

%dsp-level:0
\[
\ 1\ \text{RETIRED}\ (\text{P-H11}),
\]

%dsp-level:0
\[
\ 1\ \text{ADJUDICATED VACUOUS}\ (\text{P-H20a})
\]

%dsp-level:0
\[
\qquad [\mathrm{FT.147}(4)];
\]

%dsp-level:0
\[
(3)\quad\textbf{the gap status row:}
\]

%dsp-level:0
\[
\qquad \text{items}\ 1,2,3,4,6,8,9,10,11,14,15,18:\ \text{CLOSED};
\]

%dsp-level:0
\[
\qquad \text{items}\ 5,13,16,19:\ \text{CLOSED-DEFERRED};
\]

%dsp-level:0
\[
\qquad \text{items}\ 7,12,17:\ \text{DEFERRED};
\]

%dsp-level:0
\[
\qquad \text{item}\ 20:\ \text{REFRAMED}\ [\mathrm{FT.149}];
\]

%dsp-level:0
\[
\qquad \text{item}\ 21:\ \text{AUDITED}\ [\mathrm{FT.150}(3)]
\]

%dsp-level:0
\[
\qquad [\mathrm{FT.150}(1)];
\]

%dsp-level:0
\[
(4)\quad\textbf{the adjudication ledger:}
\]

%tbl-level:3
{\scriptsize\setlength{\tabcolsep}{4pt}
\begin{longtable}{>{\raggedright\arraybackslash}p{\dimexpr0.5890\textwidth-2\tabcolsep\relax}|>{\raggedright\arraybackslash}p{\dimexpr0.1301\textwidth-2\tabcolsep\relax}|>{\raggedright\arraybackslash}p{\dimexpr0.2808\textwidth-2\tabcolsep\relax}}
\toprule
Statement & Source & Status \\
\addlinespace[2pt]
\midrule
\endfirsthead
\toprule
Statement & Source & Status \\
\addlinespace[2pt]
\midrule
\endhead
\bottomrule
\endlastfoot
The consolidation law: the letter series H–T; ownership partition; full-text identity; acyclicity & FT.146 & Derived (machine census of the series sections) \\
P-H verdict table: 22 entries, five-verdict vocabulary & FT.147(1) & Derived (audit; anchors machine-located) \\
P-H11 retired (candidate never adopted) & FT.147(2) & Adjudicated (grounds i–iv) \\
P-H10 equilibrium form derived away & FT.108(4), FT.147(1) & Reduced (residual: the free rate datum) \\
Separation: 120 structural + 25 census = 145 & FT.148 & Derived (machine parse) \\
External-reference framing retired & FT.149(1) & Adjudicated ($\neg$silent; the audits stand) \\
Reframed question ledger Q1–Q12 & FT.149(2) & Derived (status rows anchored) \\
Gap verdicts, items 1–21 & FT.150 & Derived (closing tags machine-located) \\
Baryonic foundation stack + consumption table & FT.151(1)–(2) & Derived (tags machine-located) \\
Declaration requirements for the baryonic era & FT.151(3) & Stated ($\neg$issued) \\
New declaration issued by Tier F.43 & — & NONE \\
Species identification anywhere in this part & — & NOT derived (barred: P1, P5) \\
\end{longtable}
}

%dsp-level:0
\[
(5)\quad\textbf{the closure docket:}\ \ \text{the consolidation law is issued}\ [\mathrm{FT.146}];
\]

%dsp-level:0
\[
\ \text{the P-H audit is issued}\ [\mathrm{FT.147}];
\]

%dsp-level:0
\[
\ \text{the separation is issued}\ [\mathrm{FT.148}];
\]

%dsp-level:0
\[
\ \text{the framing is retired and the ledger reframed}\ [\mathrm{FT.149}];
\]

%dsp-level:0
\[
\ \text{the gaps are audited}\ [\mathrm{FT.150}];
\]

%dsp-level:0
\[
\ \text{the baryonic era is prepared}\ [\mathrm{FT.151}];
\]

%dsp-level:0
\[
\ \text{the status is stated}\ [\text{this tag}];
\]

%dsp-level:0
\[
\ \text{the H-CENSUS block is appended at the end of the series}.
\]

\textbf{Req.} FT.146–FT.151 (this tier); the letter parts H–R (full text); part S (the ledger); the H-CENSUS block (this part); items 5/7/12/13/16/17/19/20/21; P1, P4, P5 (\textbf{A}); exact count arithmetic [math; machine: \texttt{checks/\allowbreak{}verify\_\allowbreak{}H\_\allowbreak{}rung\_\allowbreak{}closure.py} §E–§G].

\textbf{Pf.}

%dsp-level:0
\[
(1)\text{–}(3):\ \text{assembly of FT.146–FT.151's adjudicated rows}\ [\text{machine: all sections green}];
\]

%dsp-level:0
\[
\qquad (4):\ \text{each row is the adjudication of its own tag}\ [\text{machine}];
\]

%dsp-level:0
\[
\qquad (5):\ \text{the docket scan against the tier's seven tags and the appended block}\ [\text{machine}].\ \blacksquare
\]

\textbf{Phys.} The fermionic working layer closes its rung the way the frozen corpus closes its parts: with its declarations audited and priced, its theorems sorted into law and cell-supply and shelved in an ordered letter series, its ledger turned to face inward, and its debts itemized with the stratum each waits for. One hundred forty-five theorems, thirty-four definitions, twenty-two declarations of which sixteen stand and six are paid down, twenty-one gap items of which twelve are shut. What remains open is named, and what comes next has its foundation stack, its census inventory, and its short list of declarations to make at its own points of use. The rung is closed; the books are open.

\textbf{Status.} Closed. Tier F.43 adds no declaration, no free parameter, and no identification; it retires one candidate, narrows three entries by derivation, reframes one ledger, and consolidates the layer into the letter series. Machine: §A–§G.

\vspace{2pt}\noindent\rule{\textwidth}{0.4pt}\vspace{2pt}

\phantomsection
\section[Closure exports of part T (consumed by the series a…]{\texorpdfstring{Closure exports of part T (consumed by the series and the registry addendum of part S)}{Closure exports of part T (consumed by the series and the registry addendum of part S)}}

%dsp-level:0
\[
\begin{array}{lll}
\mathrm{FT.146}\ \text{the letter series H–S};\ \text{ownership partition}\ 179;\ \text{acyclic law} & \mathrm{FT.147}\ \text{P-H audit}:\ 16/1/3/1/1;\ \text{P-H11 retired} & \mathrm{FT.148}\ 120\ \mathrm A+25\ \mathrm B=145;\ \text{rule by primary deliverable}\\
\mathrm{FT.149}\ \text{framing retired};\ \text{Q1–Q12 ledger} & \mathrm{FT.150}\ \text{items}\ 1\text{–}21\ \text{verdicts};\ 21(a)\text{–}(i)\ \text{dispositions} & \mathrm{FT.151}\ \text{baryonic stack};\ 8\ \text{declaration requirements stated}\\
\mathrm{FT.152}\ \text{sector status};\ \text{docket complete} & \text{H-CENSUS}:\ 25\ \text{entries appended at the series' end} & \text{no new declaration; no tag overturned; no species identified}
\end{array}
\]

\textbf{Ledger update (to part S's registry addendum).} The consolidated letter series H–R carries the full working-layer text as standalone canon; the first-edition files are archived for historical context, the transfer machine-verified byte-identical [FT.146(3)]. The registries and the Open Gaps Log are canonical in part S. Item 20 of the gap log is reframed by FT.149 — the external-reference scoring is retired as framing (the audit theorems stand unmodified; the bracketed per-tier entries of part S are retained verbatim as the historical status ledger), and the twelve rows are re-posed as the corpus-internal questions Q1–Q12. P-H verdicts per FT.147: RETAIN for P-H0, P-H1, P-H4, P-H5, P-H6, P-H7, P-H7$'$, P-H9, P-H12, P-H13, P-H14, P-H15, P-H16a, P-H17a, P-H18a, P-H19a; DISCHARGED for P-H2; REDUCED for P-H3 (attribution residue), P-H8 (vocabulary residue), P-H10 (free-rate-datum residue; the equilibrium form derived by FT.108(4)); RETIRED for P-H11 (adoption requirements stated at FT.147(2)); ADJUDICATED VACUOUS for P-H20a. Gap verdicts per FT.150: items 1, 2, 3, 4, 6, 8, 9, 10, 11, 14, 15, 18 CLOSED; items 5, 13, 16, 19 CLOSED-DEFERRED; items 7, 12, 17 DEFERRED; item 20 REFRAMED; item 21 AUDITED. The baryonic era's declaration requirements are itemized at FT.151(3) — stated, not issued; each requires an addendum entry before use.

$\blacksquare$

\vspace{2pt}\noindent\rule{\textwidth}{0.4pt}\vspace{2pt}

\medskip
\subsection*{\texorpdfstring{B.1.13\quad FT.89 — The Koide/phase-boundary adjudication: exact arithmetic, category trap}{FT.89 B.1.13}}
\addcontentsline{toc}{subsection}{B.1.13 FT.89}
\phantomsection\label{tag:FT-89}
{\small \textit{Record.}\quad Source: Part 3 (Part3\_Fermion\_era.md); tier: Tier F.30 — The Koide Adjudication, the Resolution Ladder, and the Minimal Bridge Declaration; canonical tag \texttt{FT.89}. Reproduced from the revision-1 canon; a cross-reference stub marks the former locus in the main corpus.}
\medskip
\textbf{Thm.} Under Def-H.26/FT.70, with the Koide constant external and marked: $(1)$ \textbf{The arithmetic is exact and re-verified.} $p^*_\Gamma$ is the unique root of $B_\Gamma(p)=p^{2E}$ in $(0,1)$ [FT.70(1)]; rational bisection of width $2^{-44}$ reproduces all eight documented polygon values ($0.807901$, $0.819671$, $0.815219$, $0.817028$, $0.816269$, $0.816596$, $0.816452$, $0.816517$ for $n=3..10$); $(p^*_n)^2$ \textbf{interlaces} $\tfrac23$ (odd $n$ below, even above) with deviations $1.40\times10^{-2},\ 2.09\times10^{-3},\ 3.72\times10^{-4},\ 7.23\times10^{-5}$ at $n=3,5,7,9$; and in $\mathbb Z[\sqrt6]$

%dsp-level:0
\[
s=2+\sqrt6:\ s^2-4s-2=0\ \checkmark;
\]

%dsp-level:0
\[
\qquad p^*_\infty=\frac{s}{1+s}=\frac{(2+\sqrt6)(3-\sqrt6)}{(3+\sqrt6)(3-\sqrt6)}=\frac{\sqrt6}{3}\ \checkmark;
\]

%dsp-level:0
\[
\qquad (p^*_\infty)^2=\frac{6}{9}=\frac23\ \textbf{exactly}.
\]

So the observation's arithmetic premise is \textbf{true and exact}, not approximate. $(2)$ \textbf{Category verdict: TRAP, on four independent grounds.} \textbf{(a) Kind mismatch.} $p^*_\infty$ is a \textit{probability} — the crossing point of the flicker ensemble's bypass measure with its floor measure [Def-H.26, FT.70(1)], conditional on the Cal-F rate datum's equilibrium reading. $Q$ [EXT] is a \textit{mass-ratio functional} of three charged-lepton masses. Both are dimensionless; dimensionless is not a kind — P2/P5 keep derived dimensionless numbers, free configuration parameters and measured constants strictly separate. \textbf{(b) No referent.} The corpus has \textbf{no lepton mass read at all}: no Yukawa sector, no lepton species, no flavour index [FT.81(5); item 20's deficits 1–3]. $Q$ therefore cannot be formed inside the corpus, and no map connects the two sides. \textbf{(c) Precedent.} FT.73(5) already adjudicated this exact species — $\alpha_{\mathrm{EM}}=\pi$, a derived pure number and a \textit{phase coefficient} [EM.20], against $1/137$, a \textit{coupling} — and ruled the comparison inadmissible. The ruling transfers verbatim: a derived dimensionless corpus number is not comparable to a measured dimensionless ratio without a declared map. \textbf{(d) Cap contingency (new, machine).} $\tfrac23$ is \textbf{not universal — it is the cap-3 value}. Repeating the identical construction with cap $c=2$ gives $(p^*_\infty)^2\to0.508004\ldots$ ($P_6,P_8,P_{10}$: $0.508193$, $0.508021$, $0.508004$), and with $c=4$ every pattern is admissible so $B\equiv1$ and \textbf{no interior boundary exists} [FT.70(1)]. Hence $(p^*_\infty)^2=\tfrac23$ is a function of $C_{\mathrm{cap}}=16$ and $D_{\mathrm{tick}}(n)=2^{n+1}$ at $n=3$ [CT0.5, CT0.6, CT0.11] — a combinatorial output of the admissibility ceiling. Genericity: $\tfrac23$ is already a derived corpus value elsewhere ($\delta_v/\pi=\tfrac23$ for the octahedron [FT.63(4)]) and is the ratio of the two smallest integers $>1$. $(3)$ \textbf{Sequence-to-generations: refuted.} $(p^*_n)^2$ approaches $\tfrac23$ with \textit{shrinking}, alternating deviations; a generation hierarchy requires \textit{growing} separation. Moreover $n$ is a polygon side count (not a generation index) and the family supplies $\ge8$ values, not $3$. $(4)$ \textbf{The A3 sub-questions, answered negatively.} (a) The transfer matrix's characteristic polynomial $x^2(x^2-3x-3)$ [FT.68(2)] has roots $(3\pm\sqrt{21})/2$ — exactly \textbf{one positive} number, the other negative, plus a double zero — so it cannot supply a three-component mass vector. (b) The $45^\circ$ cone condition ($\cos45^\circ=1/\sqrt2$) and FT.86(6)'s single-edge relative ripple $\sigma_K/\langle K\rangle=1/\sqrt2$ are the same \textit{number} of different kinds (a cone angle versus a binomial standard deviation at $p=\tfrac12$): flagged, not used. $(5)$ \textbf{Verdict.} Coincidence; closed. The equality $(p^*_\infty)^2=\tfrac23$ stands as a derived exact fact about the cap-3 polygon ensemble; its identification with the Koide constant is inadmissible on grounds (a)–(d). No structural claim, no identification, nothing adopted.

\textbf{Req.} Def-H.26; FT.63, FT.68, FT.70, FT.73, FT.81, FT.86; CT0.5, CT0.6, CT0.11 (Part 1); Cal-F, P2, P5; EM.20 (Part 2); external Koide constant (marked); exact bisection over $\mathbb Q$, exact arithmetic in $\mathbb Z[\sqrt6]$, cap-parameterized transfer enumeration [math; machine].

\textbf{Pf.}

%dsp-level:0
\[
(1):\ B_\Gamma(p)-p^{2E}\ \text{strictly decreasing/increasing}\implies\text{unique root}\ [\mathrm{FT.70}(1)];
\]

%dsp-level:0
\[
\quad\text{bisection width}\ 2^{-44}\ [\text{machine}]
\]

%dsp-level:0
\[
s^2-4s-2=(2+\sqrt6)^2-4(2+\sqrt6)-2=(10+4\sqrt6)-(8+4\sqrt6)-2=0;
\]

%dsp-level:0
\[
\qquad \tfrac{s}{1+s}=\tfrac{6-\sqrt6-2\cdot 6+2\sqrt6+\cdots}{}\ \to\ \text{rationalize}:\ \tfrac{\sqrt6}{3}
\]

%dsp-level:0
\[
(2d):\ \text{cap-}c\ \text{admissibility}:\ k_v
\]

%dsp-level:0
\[
\le c:\ \text{the transfer constraint}\ \operatorname{pop}(s_i)+\operatorname{pop}(s_{i+1})
\]

%dsp-level:0
\[
\le c\ \text{changes}\ T(t)\implies\lambda_{\max},
\]

%dsp-level:0
\[
\ p^*\ [\text{machine}]
\]

%dsp-level:0
\[
c=4\ge2+2:\ \text{no pattern excluded}:\ B\equiv1:\ \neg\exists\ \text{interior crossing}\ [\mathrm{FT.70}(1)];
\]

%dsp-level:0
\[
\qquad c=3:\ \text{the documented value}\ \tfrac23
\]

%dsp-level:0
\[
(3):\ |(p^*_n)^2-\tfrac23|\ \text{decreasing and alternating}\ [\text{machine, exact intervals}];\qquad n\ne\text{generation index}
\]

%dsp-level:0
\[
(4a):\ x^2-3x-3=0\implies x=(3\pm\sqrt{21})/2:\ \sqrt{21}>3\implies\text{one positive root};\qquad(4b):\ \text{kind distinction}
\]

%dsp-level:0
\[
(5):\ \text{conjunction of}\ (2)\text{–}(4).\ \blacksquare
\]

\textbf{Phys.} A number fell out of the polygon flicker — two thirds, exactly — and the same number is famous in lepton physics. The audit says: keep the number, drop the bridge. Two thirds here is what happens when a ring of four-way vertices is capped at three crystallized edges; cap it at two and the same machinery hands back one half; cap it at four and the phase boundary disappears altogether. A constant that moves with a bookkeeping cap cannot be the same object as a ratio of particle masses, and the corpus has no lepton masses to compare with in any case.

\textbf{Rmk (supersession of the classification clause).} The classification content of clauses (2d) and (5) — $\tfrac23$ as a mere cap-combinatorial output, and the verdict \texttt{}coincidence; no structural claim'' on the Koide-form structure — is superseded by FT.158(5): $\tfrac23=\tfrac{3-1}{3}$ is the doublet/triplet decomposition ratio of the spatial seal triad under its inner $C_3$ rotation subgroup [FT.111(1); FT.158(1)–(2)], and the Koide-form channel angle $\delta_\ell=\tfrac29=(\tfrac23)/3$ [FT.112] has its provenance derived from the three spatial directions [FT.158(2)–(4)]. Clause (1)'s arithmetic, clause (3)'s refutation, clause (4)'s negatives, and the inadmissibility grounds (2a)–(2c) are unchanged and standing: the identification of derived rationals with measured mass ratios remains inadmissible without a declared map [P1, P5].

\vspace{2pt}\noindent\rule{\textwidth}{0.4pt}\vspace{2pt}

\medskip
\subsection*{\texorpdfstring{B.1.14\quad FT.160 — The 3+1 completion: adjudication and scope}{FT.160 B.1.14}}
\addcontentsline{toc}{subsection}{B.1.14 FT.160}
\phantomsection\label{tag:FT-160}
{\small \textit{Record.}\quad Source: Part 3 (Part3\_Fermion\_era.md); tier: Tier F.44 — The 3+1 Chiral Asymmetry: The Manifold Locus of Chirality, the Spatial Provenance of the Channel Angle, and the Emergent-Time Identity; canonical tag \texttt{FT.160}. Reproduced from the revision-1 canon; a cross-reference stub marks the former locus in the main corpus.}
\medskip
\textbf{Thm.}

%dsp-level:0
\[
(1)\quad\textbf{the content summary:}
\]

%dsp-level:0
\[
\qquad \text{the}\ 3{+}1\ \text{asymmetry is derived structure at three loci of one graded pair algebra}\ [\mathrm{FT.156}:\ \text{the signature datum}\ (-,+,+,+)\ \text{at the register}\ [\mathrm{A3}]\ \text{and the metric}\ [\mathrm{CT1.12}];\ \text{the}\ 7/3\ T\text{-split}\ =(1;6;3)\ \text{component split};\ \text{reversal transport};\ \text{inner triad symmetry};\ \text{strict forwardness}];
\]

%dsp-level:0
\[
\text{its chiral read}\ [\mathrm{FT.157}:\ \text{the weak block}=\text{the shift block};\ \text{the spatial branch}\ (\vee^2,\ \text{Casimir}\ 2)\ \text{weak-active},\ \text{the temporal branch}\ (\wedge^2,\ \text{Casimir}\ 0)\ \text{weak-blind};\ \text{the selectivity}\ =\ \text{the}\ 3{+}1\ \text{asymmetry at the state locus}];
\]

%dsp-level:0
\[
\text{the spatial provenance of the channel structure}\ [\mathrm{FT.158}:\ C_3=\text{the weak rotation of the triad};\ \delta_\ell=\tfrac29=\tfrac{3-1}{3\cdot3};\ \text{the Koide-form spectrum forced};\ \mathrm{FT.89}\text{'s classification clause does not stand; its bars stand}];
\]

%dsp-level:0
\[
\text{the emergent-time identity}\ [\mathrm{FT.159}:\ \text{gravitational time}=\text{CPT time}=\text{the ledger order};\ \text{the item}\ 21(\mathrm b)\ \text{blockage removed; the rate-modulation residue restated}];
\]

%dsp-level:0
\[
(2)\quad\textbf{the adjudication table:}
\]

%dsp-level:0
\[
\begin{array}{c|c|c}\text{content} & \text{tag} & \text{status}\\ \hline\text{one signature datum}\ (-,+,+,+)\ \text{at two loci} & \mathrm{FT.156}(1) & \text{Derived}\ (\mathrm{A3}\circ\mathrm{CT1.12};\ \text{machine}\ T\text{-table})\\ 7/3\ T\text{-split}\ =(1;6;3)\ \text{dictionary split};\ S\text{-odd metric-blind} & \mathrm{FT.156}(2) & \text{Derived}\ (\mathrm{CT1.1}/\mathrm{CT1.2}/\mathrm{CT1.12})\\ \text{reversal transport};\ \operatorname{sig}\ \text{invariance} & \mathrm{FT.156}(3) & \text{Derived}\ (\text{exact rational inertia})\\ S_3\ \text{inner on the triad};\ \text{weak-generated rotations} & \mathrm{FT.156}(4),\ \mathrm{FT.158}(1) & \text{Derived}\ (\text{machine conjugation tables})\\ \text{strict forwardness};\ \neg\exists\ \text{temporal advance inverse} & \mathrm{FT.156}(5) & \text{Derived}\ (\mathrm{A6}\circ\mathrm{FT.144}\ \text{re-run})\\ \text{manifold structure}\ =\ \text{chiral structure} & \mathrm{FT.156}(6) & \text{Derived}\ (\text{three loci, locus discipline})\\ \text{weak}\ =\ \text{shift block};\ \text{the spatial force} & \mathrm{FT.157}(1)\text{–}(2) & \text{Derived}\\ \text{branch selectivity}\ =\ 3{+}1\ \text{at the state locus} & \mathrm{FT.157}(3),\ (5) & \text{Derived}\ (\mathrm{FT.101}\ \text{re-derived})\\ \text{carrier-side gradings distinct};\ \text{species rule absent} & \mathrm{FT.157}(4) & \text{Open}\ (\text{deficit}\ 1',\ \text{Q1}';\ \text{input added})\\ C_3\ \text{weak-generated};\ \delta=\tfrac{3-1}{3\cdot3} & \mathrm{FT.158}(1)\text{–}(3) & \text{Derived}\\ \text{Koide-form spectrum forced} & \mathrm{FT.158}(4) & \text{Derived}\ (\mathrm{FT.113}\ \text{re-derived})\\ \mathrm{FT.89}(2\mathrm d)/(5)\ \text{classification} & \mathrm{FT.158}(5) & \textbf{Superseded}\ (\text{bars standing})\\ \text{gravitational time}\ =\ \text{CPT time}\ =\ \text{ledger order} & \mathrm{FT.159}(1)\text{–}(3) & \text{Derived}\ (\mathrm{D1.8}\ \text{read with}\ \mathrm{CT2.6}/\mathrm{CT2.8})\\ \text{item}\ 21(\mathrm b)\ \text{``no emergent time'' blockage} & \mathrm{FT.159}(4) & \textbf{Removed}\ (\text{residue restated})\\ \text{rate modulation (redshifted flicker rate)} & \mathrm{FT.159}(4) & \text{Unconstructed}\ (\neg\text{a time declaration})\\ \text{species identification} & — & \neg\text{Derived}\ (\text{barred}:\ \mathrm{P1},\ \mathrm{P5})\end{array}
\]

%dsp-level:0
\[
(3)\quad\textbf{the ledger rows:}
\]

%dsp-level:0
\[
\qquad \text{the manifold/spacetime-structure question:}\ \textbf{CLOSED}\ [\mathrm{FT.156}\text{–}\mathrm{FT.157}];
\]

%dsp-level:0
\[
\qquad \text{the Koide-form justification question:}\ \textbf{CLOSED}\ [\mathrm{FT.158}];
\]

%dsp-level:0
\[
\text{the weak-coupling selectivity question:}\ \textbf{CLOSED}\ [\mathrm{FT.157}];
\]

%dsp-level:0
\[
\qquad \text{the emergent-time blockage}:\ \textbf{REMOVED}\ [\mathrm{FT.159}];
\]

%dsp-level:0
\[
\text{open, unchanged:}\quad\text{the relating equation}\;
\]

%dsp-level:0
\[
\qquad \text{the assignment rule}\;
\]

%dsp-level:0
\[
\qquad \text{the generation selection}\ (\text{item}\ 21(\mathrm d)/\text{item}\ 12);
\]

%dsp-level:0
\[
\qquad \text{the frame-section residue}\ (\text{item}\ 21(\mathrm i));
\]

%dsp-level:0
\[
\text{the carrier-time coupling for the frequency read}\ [\mathrm{FT.131}(3)\ —\ \text{poseable in derived time}];
\]

%dsp-level:0
\[
\qquad \text{the rate-modulation construction}\ [\mathrm{FT.159}(4)];
\]

%dsp-level:0
\[
(4)\quad\textbf{the scope bars:}
\]

%dsp-level:0
\[
\qquad \text{no particle identification}\ [\mathrm{P1},\ \mathrm{P5}];
\]

%dsp-level:0
\[
\qquad \text{no external data}\ ([\text{EXT}]\ \text{enters no clause of this tier — this standing-rule record is its sole marker});
\]

%dsp-level:0
\[
\text{all topology and branch labels are internal bookkeeping};
\]

%dsp-level:0
\[
\qquad \text{no H-sector interpretation enters any theorem statement};
\]

%dsp-level:0
\[
\qquad \text{amplitude}\ne\text{mass};
\]

%dsp-level:0
\[
\qquad \text{rung}\ne\text{scale};
\]

%dsp-level:0
\[
\text{no new declaration issues: imports are confined to the frozen derived tags of Parts 1–3};
\]

%dsp-level:0
\[
\qquad \text{the interface clauses touched carry their frozen provenance}\ [\mathrm{Def\text{-}H.41}\ \text{cited as FT.110's arena provenance and imported nowhere};\ \mathrm{Def\text{-}H.33}/\mathrm{Def\text{-}H.34}/\mathrm{Def\text{-}H.43}\ \text{enter only inside the re-cited frozen clauses FT.126/FT.129/FT.131/FT.144}];
\]

%dsp-level:0
\[
\text{no free parameter is added}\ (\text{one scale}\ \mathfrak h\ [\mathrm{P1}];\ \text{the channel scales remain P4-class}\ [\text{Cal-F}]);
\]

%dsp-level:0
\[
\qquad \text{completeness: the tier adds no new dimension to any census — the counts are CT1.1's}\ (7,3,3,3)\ \text{and CT1.12's}\ 10\ \text{read at new loci}.
\]

\textbf{Req.} FT.156, FT.157, FT.158, FT.159 (this tier); FT.125(4), FT.131(3), FT.141(6), FT.145(6) (Tiers F.39–F.42); CT1.1 (Part 1); CT1.12 (Part 2); Cal-F; items 7/12/20/21, Q1$'$; P1, P4, P5 (Parts 1–2); exact integer and rational arithmetic; the machine sections §A–§D green [math; machine: \texttt{checks/\allowbreak{}verify\_\allowbreak{}H\_\allowbreak{}3p1\_\allowbreak{}chiral\_\allowbreak{}asymmetry.py} §E].

\textbf{Pf.}

%dsp-level:0
\[
(1):\ \text{assembly of FT.156–FT.159's adjudicated rows}\ [\text{machine: §A–§D green}]
\]

%dsp-level:0
\[
(2):\ \text{each row is the adjudication of its own tag}\ [\mathrm{FT.156}\ (\S\mathrm A),\ \mathrm{FT.157}\ (\S\mathrm B),\ \mathrm{FT.158}\ (\S\mathrm C),\ \mathrm{FT.159}\ (\S\mathrm D)];
\]

%dsp-level:0
\[
\qquad \text{the open rows cite their frozen tags unmodified}\ [\mathrm{FT.125}(4),\ \mathrm{FT.141}(6),\ \mathrm{FT.131}(3)]
\]

%dsp-level:0
\[
(3):\ \text{ledger pointers}\ [\text{item}\ 21(\mathrm b)\ \text{bracket};\ \text{item}\ 20\ \text{bracket};\ \text{Q-table rows}]\ —\ \text{HistLog}\ \S\mathrm{II}\ \text{carries the resolution text}
\]

%dsp-level:0
\[
(4):\ \text{standing-rule record}\ [\mathrm{FT.145}(6)\ \text{pattern}];
\]

%dsp-level:0
\[
\qquad \text{import scan: the Req lines of this tier name no interface definition as an import; Def-H.41 appears in FT.158's provenance row alone}\ [\text{machine scan of this tier's text}].\ \blacksquare
\]

\textbf{Phys.} The tier closes as a sector: one signature datum declared at the register and derived at the metric; one dictionary turning the $7/3$ grading into lapse, spatial metric, and shift; one rotation algebra — the weak — performing every symmetry of the spatial triad from inside; one arrow of time with no inverse, whose formal mirror is exactly the CPT reversal; one channel angle that is a ratio of spatial counts; and one identity removing the last objection that the corpus had no time. What stays open stays open on the record: the rule marrying gradings to species, the relating equation between mass reads, the selection of three, and the dynamical redshift of the flicker — each now with its prerequisites derived and its locus named.

\vspace{2pt}\noindent\rule{\textwidth}{0.4pt}\vspace{2pt}

\medskip
\subsection*{\texorpdfstring{B.1.15\quad BM.12 — The Noether status of the corpus: adjudication and scope}{BM.12 B.1.15}}
\addcontentsline{toc}{subsection}{B.1.15 BM.12}
\phantomsection\label{tag:BM-12}
{\small \textit{Record.}\quad Source: Part 4 (Part4\_Baryomesonic\_era.md); tier: Tier BM.4 — Noether's Theorem from the Janus Seal Topology: the Discrete Conservation Law, the Continuum Promotion, the Stress-Energy of a Sealed Composite, and the Canonical-Pair Obstructions; canonical tag \texttt{BM.12}. Reproduced from the revision-1 canon; a cross-reference stub marks the former locus in the main corpus.}
\medskip
\textbf{Thm.}

%dsp-level:0
\[
(1)\quad\textbf{the derived rows:}\qquad\begin{array}{l|l}\text{discrete conservation} & \mathrm{BM.7}:\ \#J^+=\#J^-,\ \text{the telescoping endpoint read},\ \partial^2=0\ —\ \text{exact, integer, topological}\\\text{continuum symmetries} & \mathrm{BM.8}(1),(3),(4):\ \text{hop}\to\text{translation},\ C_3\to SO(3)\ \text{as an exact one-parameter subgroup},\ \text{slot algebras natively Lie}\\\text{stress-energy tensor} & \mathrm{BM.9}(1)\text{–}(4):\ \text{the CT1.22 read, the sign reconciliation, the translation current, the dust form}\\\text{canonical pair and bounds} & \mathrm{BM.10}(2)\text{–}(6):\ \text{the Weyl pre-limit},\ \Delta x\Delta p\ge\tfrac12,\ \text{the MT and CN forms},\ \text{the width law}\\\text{continuum conservation} & \mathrm{BM.11}:\ \partial_\mu J^\mu=0\ \text{and}\ \nabla_\mu T^{\mu\nu}=0\ \text{by two agreeing routes}\end{array}
\]

%dsp-level:0
\[
(2)\quad\textbf{the conditional rows:}
\]

%dsp-level:0
\[
\qquad \text{energy conservation under BM.8(2)'s declared stationarity scope (theorem-local, L0/CT4.5a species)};
\]

%dsp-level:0
\[
\text{the}\ (t,E)\ \text{and}\ (\theta,L)\ \text{pairs in Mandelstam–Tamm and Carruthers–Nieto form},
\]

%dsp-level:0
\[
\ \neg\text{in the naive form}\ [\mathrm{BM.10}(4),(5)];
\]

%dsp-level:0
\[
\text{the smoothing scale as}\ \mathrm{CT2.6}(\mathrm{ii})\text{–}(\mathrm{iii})\ \text{hypothesis content}\ [\mathrm{BM.10}(6),\ \mathrm{BM.11}(3)];
\]

%dsp-level:0
\[
\qquad \text{Def-BM.3's charge classes pending BM.6(1)'s construction};
\]

%dsp-level:0
\[
(3)\quad\textbf{the obstructed rows:}
\]

%dsp-level:0
\[
\qquad \neg\exists\ \text{a canonical pair in the 16-matrix or in any finite spin representation}\ [\mathrm{BM.10}(1)];
\]

%dsp-level:0
\[
\neg\exists\ \text{a self-adjoint time operator}\ [\mathrm{BM.10}(4);\ \mathrm{A6}:\ \neg\exists\,\tau^{-1}];
\]

%dsp-level:0
\[
\qquad \neg\exists\ \text{an angle operator conjugate to}\ \hat L_z\ \text{at finite spin}\ [\mathrm{BM.10}(5)];
\]

%dsp-level:0
\[
\neg\exists\ \text{a nontrivial continuous}\ U(1)\ \text{action on the ternary charge alphabet}:\ q_{\mathrm{EM}}\ \text{is}\ \neg\text{a Noether charge}\ [\mathrm{BM.7}(4)];
\]

%dsp-level:0
\[
\neg\exists\ \text{an exact time-translation invariance of the microscopic pair-state dynamics}\ [\mathrm{BM.8}(2);\ \mathrm{L2}:\ \text{deficit}\ \tfrac{\gamma(8-\gamma)}{16}>0\ \forall\gamma\in(0,1]];
\]

%dsp-level:0
\[
(4)\quad\textbf{the standing rules:}
\]

%dsp-level:0
\[
\qquad \mathrm{Def\text{-}BM.3}\text{'s identification layer is declared,}\ \neg\text{derived};
\]

%dsp-level:0
\[
\qquad \mathrm{P5}\text{'s data bar stands}: \neg\exists\ \text{a measured external value in any clause of this tier}\ [\text{the single}\ \text{marker of this record:}\ \text{no}\ \text{EXT}\ \text{datum}];
\]

%dsp-level:0
\[
\neg\text{a new dimensionful scale}\ [\mathrm{P1}:\ \mathfrak h\ \text{remains the single free scale}];
\]

%dsp-level:0
\[
\qquad \neg\text{a P4 parameter}\ [\text{§Cal-BM stays empty}];
\]

%dsp-level:0
\[
\text{imports are confined to frozen derived tags of Parts 1–3 and to this part's earlier tags};
\]

%dsp-level:0
\[
\qquad \text{completeness: the tier adds no new count to any census}\ —\ \text{the dimensions used are CT1.1's}\ (7,3,3,3),
\]

%dsp-level:0
\[
\ \mathrm{CT1.12}\text{'s}\ 10,
\]

%dsp-level:0
\[
\ \mathrm{CT0.14}\text{'s}\ 16\ \text{read at new loci};
\]

%dsp-level:0
\[
(5)\quad\textbf{the open rows touched:}
\]

%dsp-level:0
\[
\qquad \mathrm{BM.6}(1)\ \text{(the additive multi-chain charge read — Def-BM.3's classes are declared pending it)};
\]

%dsp-level:0
\[
\qquad \mathrm{BM.6}(5)\ \text{(the dynamics rates:}\ J^\pm\ \text{rates would convert BM.8(2)'s scope into dynamics)};
\]

%dsp-level:0
\[
\text{HistLog §II items 13 (the amplitude}\leftrightarrow\text{spectral bridge), 17 (the momentum-space resolvent), 21(b) (the tick-rate coupling)};
\]

%dsp-level:0
\[
\qquad \text{deficit}\ 1'\ [\mathrm{FT.125}(4)]\ \text{is untouched by the naming layer and stays open};
\]

%dsp-level:0
\[
(6)\quad\textbf{the adjudication:}
\]

%dsp-level:0
\[
\qquad \text{Noether's theorem is a consequence of the promotion,}\ \neg\text{a premise of the conservation law}:
\]

%dsp-level:0
\[
\text{the corpus's conservation law is topological and exact at every scale}\ [\mathrm{BM.7}(5)];
\]

%dsp-level:0
\[
\qquad \text{Noether's machinery enters where a discrete symmetry becomes a one-parameter group}\ [\mathrm{BM.8}];
\]

%dsp-level:0
\[
\text{it yields the stress-energy tensor as the translation current}\ [\mathrm{BM.9}(3)]\ \text{and the differential law as the density read}\ [\mathrm{BM.11}(3)];
\]

%dsp-level:0
\[
\text{where the machinery has no domain}\ —\ \text{the ternary charge, the tick arrow, the finite register}\ —\ \text{the corpus records the obstruction}\ [(3)]\ \text{rather than importing the theorem}.
\]

\textbf{Req.} BM.7, BM.8, BM.9, BM.10, BM.11 (this tier); A6 (Part 1); CT0.14, CT1.1 (Part 1); CT1.12, L2 (Parts 1–2); Def-BM.3, BM.3(1), BM.6(1), BM.6(5) (this part); FT.125(4) (Part 3); P1, P4, P5 (Parts 1–2); §Cal-BM; the machine sections §A–§E green [math; machine: \texttt{checks/\allowbreak{}verify\_\allowbreak{}H\_\allowbreak{}noether.py} §F].

\textbf{Pf.}

%dsp-level:0
\[
(1):\ \text{assembly of the adjudicated rows of BM.7–BM.11}\ [\text{machine: §A–§E green}]
\]

%dsp-level:0
\[
(2)\text{–}(3):\ \text{each row is the scope or obstruction clause of its own tag, cited unmodified}\ [\mathrm{BM.8}(2),\ \mathrm{BM.10}(1),(4),(5),(6),\ \mathrm{BM.7}(4),\ \mathrm{BM.11}(3)]
\]

%dsp-level:0
\[
(4):\ \text{standing-rule record}\ [\mathrm{BM.6}\ \text{pattern}];
\]

%dsp-level:0
\[
\qquad \text{the import scan names no interface definition of Part 3 as a premise};
\]

%dsp-level:0
\[
\qquad \text{no dimensionful or P4 datum is introduced}\ [\text{§Cal-BM}]
\]

%dsp-level:0
\[
(5):\ \text{ledger pointers}\ —\ \text{HistLog §II carries the resolution text};
\]

%dsp-level:0
\[
\qquad (6):\ \text{conjunction of}\ (1)\text{–}(5).\ \blacksquare
\]

\textbf{Phys.} The tier closes as a sector with a single verdict. Charge conservation in this corpus is topology: an integer balance written into the definition of a valid universe, telescoping exactly along every materialized chain, grounded in a boundary map whose square vanishes. Noether's theorem is not its explanation and could not be — the ternary charge alphabet admits no continuous rotation to generate a current from. Where the discrete structures do become one-parameter groups, the theorem applies and pays: the hop's commutator identity hands over momentum, the cubic relation on the weak generator hands over the whole rotation group with the seal's $120^\circ$ element as one of its rational points, and the translation current is the stress-energy tensor whose divergence vanishes on-shell exactly as the contracted Bianchi identity says it must. Where the machinery has no domain, the structure yields the obstruction instead: no canonical pair in a finite register, no time operator against an irreversible tick, no angle operator at finite spin, no exact time-translation symmetry in a dynamics that contracts. The names of the baryomesonic era now attach to these tuples by declaration, and the declaration imports no number.

\vspace{2pt}\noindent\rule{\textwidth}{0.4pt}\vspace{2pt}

\medskip
\subsection*{\texorpdfstring{B.1.16\quad BM.18 — The status of the calibration lock: adjudication and scope}{BM.18 B.1.16}}
\addcontentsline{toc}{subsection}{B.1.16 BM.18}
\phantomsection\label{tag:BM-18}
{\small \textit{Record.}\quad Source: Part 4 (Part4\_Baryomesonic\_era.md); tier: Tier BM.5 — The Calibration Lock: the Born Width, the Coherence Volume, the Curvature of a Sealed Composite, and the Neutron Row; canonical tag \texttt{BM.18}. Reproduced from the revision-1 canon; a cross-reference stub marks the former locus in the main corpus.}
\medskip
\textbf{Thm.}

%dsp-level:0
\[
(1)\quad\textbf{the derived rows:}\qquad\begin{array}{l|l}\text{the Born width} & \mathrm{BM.13}(1):\ \Delta x_{\mathrm B}=\mathfrak h/(E_{\mathrm{int}}\ln2),\ \Delta x_{\mathrm B}m_{\mathrm{hid}}=1\ —\ \text{exact, from}\ \mathrm{FT.44}(4)\\\text{the coherence volume} & \mathrm{BM.14}(1):\ V_{\mathrm{eff}}=\mathfrak h^3/(E_{\mathrm{int}}\ln2)^3\\\text{the energy density} & \mathrm{BM.14}(2):\ \rho=(E_{\mathrm{int}}\ln2)^4/\mathfrak h^4,\quad[\rho]=\mathfrak q^{+2}=[T_{\mu\nu}]\\\text{the curvature} & \mathrm{BM.15}(2),(4):\ R=\kappa^2(\rho-3p)+4\Lambda,\quad\kappa^2\rho=2\pi E_{\mathrm{int}}^4(\ln2)^3/\mathfrak h^2\\\text{the lock} & \mathrm{BM.16}(1),(4):\ (\Delta x_{\mathrm B})^2\kappa^2\rho=2\pi E_{\mathrm{int}}^2\ln2,\quad\Delta x_{\mathrm B}r_s=2G_N\\\text{the neutron row} & \mathrm{BM.17}(3),(4):\ E_{\mathrm{int}}^{(n)}=6\ \text{with all six reads exact}\ [\text{Census FT.35}(2)]\end{array}
\]

%dsp-level:0
\[
(2)\quad\textbf{the conditional and declared rows:}
\]

%dsp-level:0
\[
\qquad \text{the Compton normalization}\ c
\]

%dsp-level:0
\[
= 1\ \text{of}\ \mathrm{Def\text{-}BM.4}(1)\ \text{(declared; the dimensional-uniqueness clause prices it)};
\]

%dsp-level:0
\[
\text{the geometric factor}\ \alpha\ [\mathrm{Def\text{-}BM.4}(4):\ \text{priced},\ \neg\text{adopted}];
\]

%dsp-level:0
\[
\qquad \text{the charge class of}\ \mathrm{BM.17}(2)\ [\mathrm{Def\text{-}BM.3}(4):\ \text{pending BM.6(1)}];
\]

%dsp-level:0
\[
\text{the comparison status of}\ r_s,
\]

%dsp-level:0
\[
\ \ell_G\ [\mathrm{P3}:\ G_N\ \text{is a named combination for external comparison},\ \neg\text{an input}];
\]

%dsp-level:0
\[
\qquad \text{the chart locus of the width}\ [\mathrm{Def\text{-}BM.4}(5);\ \mathrm{CT2.6}];
\]

%dsp-level:0
\[
(3)\quad\textbf{the corrected rows:}
\]

%dsp-level:0
\[
\qquad \textbf{(i)}\ \text{the scalar curvature of a sealed composite is}\ R
\]

%dsp-level:0
\[
= \kappa^2(\rho-3p)+4\Lambda
\]

%dsp-level:0
\[
= \kappa^2\rho+4\Lambda\ \text{at}\ p
\]

%dsp-level:0
\[
= 0\ [\mathrm{BM.15}(2)]:
\]

%dsp-level:0
\[
\text{the form}\ R=-\kappa^2\rho\ \text{is refuted exactly, with residual}\ 2\kappa^2\rho-3\kappa^2p+4\Lambda\ [\text{machine}];
\]

%dsp-level:0
\[
\qquad \text{the sign is CT1.23's own trace, not a convention choice};
\]

%dsp-level:0
\[
\textbf{(ii)}\ \text{the spread-momentum relation in corpus units is}\ \Delta x
\]

%dsp-level:0
\[
\ge 1/(2\Delta p)\ \text{with}\ \hbar
\]

%dsp-level:0
\[
= 1\ [\mathrm{D2.9};\ \mathrm{BM.10}(6);\ \mathrm{BM.13}(4)]:
\]

%dsp-level:0
\[
\text{the form}\ \Delta x\sim\mathfrak h/\Delta p\ \text{is dimensionally inadmissible}\ —\ [\mathfrak h/\Delta p]
\]

%dsp-level:0
\[
= \mathfrak q^{-1/2}/\mathfrak q^{+1/2}
\]

%dsp-level:0
\[
= \mathfrak q^{-1}
\]

%dsp-level:0
\[
= [\text{area}]
\]

%dsp-level:0
\[
\ne [\text{length}]\ [\mathrm{CT2.1};\ \text{machine: exponent arithmetic}];
\]

%dsp-level:0
\[
\text{the gravitational length is}\ \ell_G
\]

%dsp-level:0
\[
= \sqrt{G_N}
\]

%dsp-level:0
\[
= \mathfrak h/(2\sqrt{\ln2})\ [\mathrm{CT2.9}],
\]

%dsp-level:0
\[
\ \text{and its relation to the width is the exact order statement of}\ \mathrm{BM.16}(5),
\]

%dsp-level:0
\[
\ \neg\text{an identity of scales};
\]

%dsp-level:0
\[
(4)\quad\textbf{the open rows touched:}
\]

%dsp-level:0
\[
\qquad \mathrm{BM.6}(1)\ \text{(the additive multi-chain charge read}\ —\ \mathrm{BM.17}(2)\text{'s charge class waits on it)};
\]

%dsp-level:0
\[
\qquad \mathrm{BM.6}(5)\ \text{(the dynamics rates)};
\]

%dsp-level:0
\[
\text{the general composite additivity law for}\ E_{\mathrm{int}}\ [\mathrm{BM.17}(6);\ \mathrm{Def\text{-}BM.1}\ \text{read clause}];
\]

%dsp-level:0
\[
\qquad \text{HistLog §II items 13, 17, 21(b)};
\]

%dsp-level:0
\[
\qquad \text{deficit}\ 1'\ [\mathrm{FT.125}(4)];
\]

%dsp-level:0
\[
\mathrm{FT.90}(6)\text{'s refutation of a Compton-wavelength relation to the generation threshold stands untouched}\ [\mathrm{Def\text{-}BM.4}(6)];
\]

%dsp-level:0
\[
(5)\quad\textbf{the standing rules:}
\]

%dsp-level:0
\[
\qquad \text{particle names are confined to the Def-BM.3 layer}\ [\mathrm{Def\text{-}BM.3}(7);\ \mathrm{BM.17}(2)\ \text{reads tuples}];
\]

%dsp-level:0
\[
\neg\exists\ \text{a measured external value in any clause}\ [\mathrm{P5}]\ —\ \text{the single marker of this record: no}\ \text{EXT}\ \text{datum};
\]

%dsp-level:0
\[
\qquad \neg\exists\ \text{a new dimensionful scale}\ [\mathrm{P1}:\ \mathfrak h];
\]

%dsp-level:0
\[
\qquad \neg\exists\ \text{a P4 addition}\ [\text{§Cal-BM stays empty}];
\]

%dsp-level:0
\[
\text{imports are confined to frozen derived tags of Parts 1–3, Census cells cited row-wise, and this part's earlier tags};
\]

%dsp-level:0
\[
\textbf{completeness:}\ \ \text{the tier adds no new count to any census}\ —\ E_{\mathrm{int}}\in\{2,6,12,24,60\}\ \text{are Census FT.28/FT.35/FT.42(5) values, and the ladder reads are substitutions into them};
\]

%dsp-level:0
\[
(6)\quad\textbf{the adjudication:}
\]

%dsp-level:0
\[
\qquad \textbf{the calibration lock is closed in the dimensionless sector:}\ (\Delta x_{\mathrm B})^2\kappa^2\rho
\]

%dsp-level:0
\[
= 2\pi E_{\mathrm{int}}^2\ln2\ \text{and}\ \Delta x_{\mathrm B}r_s
\]

%dsp-level:0
\[
= 2G_N\ \text{carry neither}\ \mathfrak h\ \text{nor a P4 parameter}\ [\mathrm{BM.16}(1),(4)];
\]

%dsp-level:0
\[
\text{the absolute scale rides on the single free}\ \mathfrak h\ [\mathrm{P1},\ \mathrm{CT2.11}:\ \text{rescaling is a symmetry}];
\]

%dsp-level:0
\[
\qquad \Lambda\ \text{is the one additive freedom in}\ R\ [\mathrm{BM.15}(6)];
\]

%dsp-level:0
\[
\text{the mass is a topological count and the width, the volume, the density and the curvature are its exact monomials}\ [\mathrm{BM.16}(3)]:
\]

%dsp-level:0
\[
\qquad \neg\exists\ \text{a continuous mass variable to tune}.
\]

\textbf{Req.} D2.9 (Parts 1–2); CT1.23, CT2.1, CT2.9, CT2.11 (Part 2); P1, P3, P4, P5 (Parts 1–2); Def-BM.1, Def-BM.3, Def-BM.4, BM.6(1), BM.6(5), BM.13, BM.14, BM.15, BM.16, BM.17 (this part); FT.90(6), FT.125(4) (Part 3); Census FT.28, FT.35(2), FT.42(5); §Cal-BM; the machine sections §A–§E green [math; machine: \texttt{checks/\allowbreak{}verify\_\allowbreak{}H\_\allowbreak{}calibration\_\allowbreak{}lock.py} §F].

\textbf{Pf.}

%dsp-level:0
\[
(1):\ \text{assembly of the adjudicated rows of BM.13–BM.17}\ [\text{machine: §A–§E green}]
\]

%dsp-level:0
\[
(2)\text{–}(3):\ \text{each row is the scope, pricing or refutation clause of its own tag, cited unmodified}\ [\mathrm{Def\text{-}BM.4}(1),(4),(5),(6);\ \mathrm{BM.15}(2);\ \mathrm{BM.16}(5);\ \mathrm{BM.13}(4)];
\]

%dsp-level:0
\[
\qquad \text{the dimensional refutation}: [\mathfrak h]-[\Delta p]=-\tfrac12-\tfrac12=-1\ne-\tfrac12\ [\mathrm{CT2.1};\ \text{machine}]
\]

%dsp-level:0
\[
(4):\ \text{ledger pointers}\ —\ \text{HistLog §II carries the resolution text}
\]

%dsp-level:0
\[
(5):\ \text{standing-rule record}\ [\mathrm{BM.12}(4)\ \text{pattern}];
\]

%dsp-level:0
\[
\qquad \text{the import scan names no interface definition of Part 3 as a premise};
\]

%dsp-level:0
\[
\qquad \text{no dimensionful or P4 datum is introduced}\ [\text{§Cal-BM}]
\]

%dsp-level:0
\[
(6):\ \text{conjunction of}\ (1)\text{–}(5).\ \blacksquare
\]

\textbf{Phys.} The lock holds where a lock can hold. Two dimensionless products — the squared width times the curvature scale, and the width times the comparison radius — are pure functions of an integer edge count, with no hop length, no internal parameter and no measured number in them. Everything dimensionful scales by its exact weight under the one free scale, which is the symmetry that makes the dimensionless statements the only predictions. The mass is a count; the width, the volume, the density and the curvature are monomials in that count; there is no continuous dial. Two forms that would have been convenient are refuted rather than adopted: the scalar curvature of a dust composite is positive in the corpus's own trace convention, and a spread written as hop over momentum has the dimensions of an area, not a length. What stays open stays open on the record — the additive charge read, the general composition law for the edge count, the dynamics rates — and the neutron row is exact where a Census cell supplies its count and parametric where none does.

\vspace{2pt}\noindent\rule{\textwidth}{0.4pt}\vspace{2pt}

\medskip
\subsection*{\texorpdfstring{B.1.17\quad BM.24 — The status of the topology-blind composite: adjudication and scope}{BM.24 B.1.17}}
\addcontentsline{toc}{subsection}{B.1.17 BM.24}
\phantomsection\label{tag:BM-24}
{\small \textit{Record.}\quad Source: Part 4 (Part4\_Baryomesonic\_era.md); tier: Tier BM.6 — The Topology-Blind Neutron and the Topology-Blind + Static Proton: the Flicker Ensemble, the Path-Integrated Mass, the Binding Read, and the Forced Carrier; canonical tag \texttt{BM.24}. Reproduced from the revision-1 canon; a cross-reference stub marks the former locus in the main corpus.}
\medskip
\textbf{Thm.}

%dsp-level:0
\[
(1)\quad\textbf{the derived rows:}\qquad\begin{array}{l|l}\text{the flicker ensemble} & \mathrm{BM.19}:\ \ker Q\ \text{one-dimensional},\ p=x_+/(x_++x_-),\ p=1-4\varepsilon_q\ \text{identically},\ E_{\mathrm{int}}(R)=|R|\\\text{the mass route} & \mathrm{BM.20}(1)\text{–}(3):\ RV=\kappa^2m,\ (R-4\Lambda)V=\kappa^2m,\ \langle m\rangle=\text{the averaged flux}\\\text{the sector mass} & \mathrm{BM.20}(4)\text{–}(5):\ m^2_B(q,S)\ \text{exact rationals, charge-degenerate, bounded by}\ \tfrac32B(B-1)\\\text{the binding read} & \mathrm{BM.21}(2)\text{–}(4):\ \text{nonzero exact rationals; zero exactly in the no-flicker limit}\\\text{the two classes} & \mathrm{BM.22}(2),(4):\ F_{LR}=qI_2,\ \text{splitting}\ 2|q|\in\{0,2\}\\\text{the transition} & \mathrm{BM.23}(1)\text{–}(5),(7)\text{–}(9):\ \Delta q=\pm1,\ \text{the forced carrier},\ \Gamma\ \text{a resolvent in}\ \varepsilon_q,\ \Gamma=0\ \text{at}\ \varepsilon_q=0\end{array}
\]

%dsp-level:0
\[
(2)\quad\textbf{the declared rows:}
\]

%dsp-level:0
\[
\qquad \mathrm{Def\text{-}BM.6}\text{'s name}\leftrightarrow\text{line attachment}\ (\text{HistLog §I addendum},\ \mathrm{Def\text{-}BM.3}(7)\text{'s rule});
\]

%dsp-level:0
\[
\qquad \mathrm{Def\text{-}H.43}\text{'s read promotion}\ [\mathrm{BM.22}(5)];
\]

%dsp-level:0
\[
\mathrm{Def\text{-}BM.3}(4)\text{'s charge-class labels, pending the additive read}\ [\mathrm{BM.6}(1)];
\]

%dsp-level:0
\[
\qquad \mathrm{Def\text{-}H.32}\text{'s own constitutive clauses}\ [\text{HistLog §I}];
\]

%dsp-level:0
\[
(3)\quad\textbf{the conditional and admissible-not-forced rows:}
\]

%dsp-level:0
\[
\qquad \text{the neutral partner's emission}\ [\mathrm{BM.23}(6)];
\]

%dsp-level:0
\[
\qquad \text{the mass form}\ (\sqrt m\propto r\ \text{vs}\ m\propto r)\ [\mathrm{Def\text{-}BM.6}(5)];
\]

%dsp-level:0
\[
\text{the}\ B\to\infty\ \text{limit of every ensemble sum}\ [\mathrm{Def\text{-}H.32}(1)];
\]

%dsp-level:0
\[
\qquad \text{the per-tick reading of}\ \Gamma\ \text{as an external rate}\ [\mathrm{BM.23}(9):\ \neg\text{available},\ \neg\text{claimed}];
\]

%dsp-level:0
\[
(4)\quad\textbf{the corrected rows:}
\]

%dsp-level:0
\[
\qquad \textbf{(i)}\ \text{the binding energy is}\ \neg\text{an integer summation}: \text{the informational count is exactly additive}\ (6=2{+}2{+}2)\ \text{and its binding read vanishes}\ [\mathrm{BM.21}(1)];
\]

%dsp-level:0
\[
\ \text{the nonzero binding is the ensemble's}\ [\mathrm{BM.21}(2)];
\]

%dsp-level:0
\[
\textbf{(ii)}\ \text{the mass enters through the curvature law on the path-integrated route}\ [\mathrm{BM.20}(1),(3)],
\]

%dsp-level:0
\[
\ \neg\text{through a constituent-mass sum};
\]

%dsp-level:0
\[
\textbf{(iii)}\ \text{the decay rate is a weight per tick in the ledger order}\ [\mathrm{D1.8};\ \mathrm{BM.23}(9)],
\]

%dsp-level:0
\[
\ \neg\text{a rate in an external time unit};
\]

%dsp-level:0
\[
(5)\quad\textbf{the open rows touched:}
\]

%dsp-level:0
\[
\qquad \mathrm{BM.6}(1)\ \text{(the additive charge read)};
\]

%dsp-level:0
\[
\qquad \mathrm{BM.6}(5)\ \text{(the dynamics rates: the weights are derived from}\ \varepsilon_q,\ \text{the per-tick scale is not)};
\]

%dsp-level:0
\[
\text{HistLog §II items 13 (the amplitude}\leftrightarrow\text{spectral bridge}\ =\ \text{the mass form}),
\]

%dsp-level:0
\[
\ 17\ \text{(the momentum-space resolvent)};
\]

%dsp-level:0
\[
\qquad \text{deficit}\ 1'\ [\mathrm{FT.125}(4):\ \text{which line sits on which chiral ray and in which gauge representation}];
\]

%dsp-level:0
\[
(6)\quad\textbf{the standing rules:}
\]

%dsp-level:0
\[
\qquad \neg\exists\ \text{a measured external value in any clause}\ [\mathrm{P5}]\ —\ \text{the single marker of this record: no}\ \text{EXT}\ \text{datum};
\]

%dsp-level:0
\[
\qquad \neg\exists\ \text{a new dimensionful scale}\ [\mathrm{P1}];
\]

%dsp-level:0
\[
\qquad \neg\exists\ \text{a P4 addition}\ [\text{§Cal-BM stays empty};\ \mathrm{BM.19}(2)];
\]

%dsp-level:0
\[
\text{imports are confined to frozen derived tags of Parts 1–3, Census cells cited row-wise, and this part's earlier tags};
\]

%dsp-level:0
\[
\qquad \text{completeness: the tier adds no new count to any census}\ —\ \text{the sector values are FT.107(3)'s and the cell counts are Census FT.35/FT.42's};
\]

%dsp-level:0
\[
(7)\quad\textbf{the adjudication:}
\]

%dsp-level:0
\[
\qquad \text{the neutron and the proton are ensembles, and the difference between them is derived three ways: the flipping block}\ (2|q|\in\{0,2\})\ [\mathrm{BM.22}(4)],
\]

%dsp-level:0
\[
\ \text{the charge compensation}\ [\mathrm{BM.23}(4)],
\]

%dsp-level:0
\[
\ \text{and the carrier selection}\ [\mathrm{BM.23}(5)];
\]

%dsp-level:0
\[
\text{the mass, the binding and the rate are all reads of the same flicker sum}\ [\mathrm{BM.20}(3),\ \mathrm{BM.21}(3),\ \mathrm{BM.23}(7)]:
\]

%dsp-level:0
\[
\qquad \text{no separate dynamics is introduced, and no datum is declared to make any of them finite}.
\]

\textbf{Req.} D1.8 (Part 1); P1, P4, P5 (Parts 1–2); Def-H.32, Def-H.43, FT.125(4) (Part 3); Def-BM.3, Def-BM.5, Def-BM.6, BM.6(1), BM.6(5), BM.19, BM.20, BM.21, BM.22, BM.23 (this part); Census FT.35(2), FT.42(5), FT.107(3) (row-wise); HistLog §I/17, Q10 (pointers); §Cal-BM; the machine sections §A–§F green [math; machine: \texttt{checks/\allowbreak{}verify\_\allowbreak{}H\_\allowbreak{}topo\_\allowbreak{}blind\_\allowbreak{}neutron.py} §G].

\textbf{Pf.}

%dsp-level:0
\[
(1):\ \text{assembly of the adjudicated rows of BM.19–BM.23}\ [\text{machine: §A–§F green}]
\]

%dsp-level:0
\[
(2)\text{–}(3):\ \text{each row is the scope, declaration or admissibility clause of its own tag, cited unmodified}
\]

%dsp-level:0
\[
(4):\ \text{the two refutations are BM.21}(1)\text{'s exact vanishing and BM.20}(1)\text{'s flux identity};
\]

%dsp-level:0
\[
\ \text{the tick reading is D1.8's coordinate}\ [\text{machine rows re-cited}]
\]

%dsp-level:0
\[
(5):\ \text{ledger pointers}\ —\ \text{HistLog §II carries the resolution text}
\]

%dsp-level:0
\[
(6):\ \text{standing-rule record}\ [\mathrm{BM.12}(4),\ \mathrm{BM.18}(5)\ \text{pattern}];
\]

%dsp-level:0
\[
\qquad \text{the import scan names no Part-3 interface definition as a premise beyond the dispositioned ones}
\]

%dsp-level:0
\[
(7):\ \text{conjunction of}\ (1)\text{–}(6).\ \blacksquare
\]

\textbf{Phys.} The tier closes with a single picture, assembled rather than assumed. A sealed composite is a finite exact sum over crystallization patterns sharing its boundary data, weighted by the unique stationary measure of the Janus move generator. Its mass is the curvature flux of that sum; its binding is the failure of the flux to add up over separated constituents, and that failure vanishes precisely when the flicker does. Its electromagnetic character is one matrix block: present with weight two in the unit charge class, absent in the neutral one. Its decay is three necessities in a row — a ladder that moves the ternary read by one, a universe whose total read is zero, an inventory with exactly one colour-blind electromagnetically active channel — leaving the emitted line forced and the neutral partner merely admissible. The rate is a resolvent in the crystallization threshold, zero at the sharp value, and counted in ticks. Nothing measured enters, nothing is tuned, and the two names on the two junctions are labels the identification layer attaches to all of it.

\vspace{2pt}\noindent\rule{\textwidth}{0.4pt}\vspace{2pt}

\medskip
\subsection*{\texorpdfstring{B.1.18\quad BM.29 — The status of the CPT sectors: adjudication and scope}{BM.29 B.1.18}}
\addcontentsline{toc}{subsection}{B.1.18 BM.29}
\phantomsection\label{tag:BM-29}
{\small \textit{Record.}\quad Source: Part 4 (Part4\_Baryomesonic\_era.md); tier: Tier BM.7 — The Proton in the Positive CPT Sector and Its Mirror: the Sector Labels, the Exact Phase, and the Charge-Blind Geometry; canonical tag \texttt{BM.29}. Reproduced from the revision-1 canon; a cross-reference stub marks the former locus in the main corpus.}
\medskip
\textbf{Thm.}

%dsp-level:0
\[
(1)\quad\textbf{the derived rows:}\qquad\begin{array}{l|l}\text{the involution} & \mathrm{BM.25}(1)\text{–}(2):\ C,P,T\ \text{derived},\ \mathrm{CPT}^2=\mathrm{id},\ \Theta_{\mathfrak L}^2=\mathrm{id},\ \text{the fixed-word set}\\\text{the exact phase} & \mathrm{BM.25}(3):\ w(\Theta_{\mathfrak L}C)=w(C)^\ast,\ \lambda\in\{\pm1,\pm i\},\ |w(\mathrm{CPT}\,C)|^2=|w(C)|^2,\ \text{the on-shell refinement}\\\text{the sector census} & \mathrm{BM.25}(4)\text{–}(5):\ \text{three orbits of size two},\ q=0\ \text{the unique charge fixed point},\ Z(0,S)=2Z(\pm1,S),\ Z(+1,S)=Z(-1,S)\\\text{the character} & \mathrm{BM.25}(6):\ \lambda=(-1)^{|E|}\ \text{on the baryon-type phase class}\\\text{the rows} & \mathrm{BM.26},\ \mathrm{BM.27}:\ \text{the proton and antiproton tuples with all reads exact}\\\text{the route theorem} & \mathrm{BM.28}:\ W(0){=}20,\ W(\pm1){=}15,\ \text{ratio}\ \tfrac43;\ \text{degenerate on the normalized route, split on the total route}\end{array}
\]

%dsp-level:0
\[
(2)\quad\textbf{the declared rows:}
\]

%dsp-level:0
\[
\qquad \mathrm{Def\text{-}BM.7}(1)\text{'s sector labels and their sign convention};
\]

%dsp-level:0
\[
\qquad \mathrm{Def\text{-}BM.3}(2)\text{'s name}\leftrightarrow\text{tuple attachments (proton, antiproton)};
\]

%dsp-level:0
\[
\mathrm{Def\text{-}BM.3}(4)\text{'s charge classes pending the additive read}\ [\mathrm{BM.6}(1)];
\]

%dsp-level:0
\[
\qquad \mathrm{Def\text{-}H.43}\text{'s read promotion, without which}\ F_{LR}=0\ [\mathrm{BM.22}(5)];
\]

%dsp-level:0
\[
(3)\quad\textbf{the corrected rows:}
\]

%dsp-level:0
\[
\qquad \textbf{(i)}\ \text{``the antiproton is not a separate object'' and ``the proton and the antiproton are not the same object'' are both exact and are not in conflict}: \text{distinct configurations, one orbit, no independent datum}\ [\mathrm{BM.27}(6)];
\]

%dsp-level:0
\[
\textbf{(ii)}\ \ \text{``backward time'' is the formal mirror}\ \Theta_{\mathfrak L},
\]

%dsp-level:0
\[
\ \neg\text{a dynamical inverse}\ [\mathrm{Def\text{-}BM.7}(2);\ \mathrm{A6};\ \mathrm{FT.144}(4)(\mathrm{ii})]:\ \text{the negative sector is a bookkeeping orbit};
\]

%dsp-level:0
\[
\textbf{(iii)}\ \ \text{the proton's width and curvature equal the neutron's because no read carries}\ q\ \text{and the scar contributes}\ 0\ \text{to}\ E_{\mathrm{int}}\ [\mathrm{BM.28}(1),(3)],
\]

%dsp-level:0
\[
\ \neg\text{because the scar is ``small''};
\]

%dsp-level:0
\[
(4)\quad\textbf{the open rows touched:}
\]

%dsp-level:0
\[
\qquad \mathrm{BM.6}(1)\ \text{(the additive charge read, on which both class labels wait)};
\]

%dsp-level:0
\[
\qquad \text{deficit}\ 1'\ [\mathrm{FT.125}(4):\ \text{which line sits on which chiral ray}];
\]

%dsp-level:0
\[
\text{HistLog §II items 13 (the amplitude}\leftrightarrow\text{spectral bridge), 17 (the momentum-space resolvent)};
\]

%dsp-level:0
\[
\qquad \text{no inter-sector dynamics}\ [\mathrm{BM.27}(7)];
\]

%dsp-level:0
\[
(5)\quad\textbf{the standing rules:}
\]

%dsp-level:0
\[
\qquad \neg\exists\ \text{a measured external value in any clause}\ [\mathrm{P5}]\ —\ \text{the single marker of this record: no}\ \text{EXT}\ \text{datum};
\]

%dsp-level:0
\[
\qquad \neg\exists\ \text{a new dimensionful scale}\ [\mathrm{P1}];
\]

%dsp-level:0
\[
\qquad \neg\exists\ \text{a P4 addition}\ [\text{§Cal-BM stays empty}];
\]

%dsp-level:0
\[
\text{species names are confined to Def-BM.3/Def-BM.6/Def-BM.7's layers}\ [\mathrm{Def\text{-}BM.3}(7)];
\]

%dsp-level:0
\[
\qquad \text{completeness: the tier adds no new count to any census}\ —\ \text{the census integers are FT.106(3)'s and the cell data are Census FT.35/FT.49's};
\]

%dsp-level:0
\[
(6)\quad\textbf{the adjudication:}
\]

%dsp-level:0
\[
\qquad \text{the sector structure is the CPT involution's orbit decomposition on derived configuration data, and it is exact}: \text{three orbits, equal charged weights, a neutral class twice either, one phase of modulus one};
\]

%dsp-level:0
\[
\text{the geometry is charge-blind and the coupling is not}:
\]

%dsp-level:0
\[
\qquad \textbf{the proton and the neutron differ in one matrix block; the proton and the antiproton differ in one sign and one orientation};
\]

%dsp-level:0
\[
\text{what the corpus does not contain is a dynamics between the sectors}\ [\mathrm{BM.27}(7)]\ —\ \text{recorded as absent},
\]

%dsp-level:0
\[
\ \neg\text{assumed small}.
\]

\textbf{Req.} A6 (Part 1); P1, P5 (Parts 1–2); Def-H.43, FT.106(3), FT.125(4) (Part 3); Def-BM.3, Def-BM.7, BM.6(1), BM.22(5), BM.25, BM.26, BM.27, BM.28 (this part); Census FT.35, FT.49, FT.106(3) (row-wise); §Cal-BM; the machine sections §A–§C green [math; machine: \texttt{checks/\allowbreak{}verify\_\allowbreak{}H\_\allowbreak{}proton\_\allowbreak{}cpt.py} §D].

\textbf{Pf.}

%dsp-level:0
\[
(1):\ \text{assembly of the adjudicated rows of BM.25–BM.28}\ [\text{machine: §A–§C green}]
\]

%dsp-level:0
\[
(2)\text{–}(3):\ \text{each row is the declaration, scope or correction clause of its own tag, cited unmodified}
\]

%dsp-level:0
\[
(4):\ \text{ledger pointers}\ —\ \text{HistLog §II carries the resolution text}
\]

%dsp-level:0
\[
(5):\ \text{standing-rule record}\ [\mathrm{BM.12}(4),\ \mathrm{BM.18}(5),\ \mathrm{BM.24}(6)\ \text{pattern}];
\]

%dsp-level:0
\[
\qquad \text{the import scan names no Part-3 interface definition as a premise beyond the dispositioned ones}
\]

%dsp-level:0
\[
(6):\ \text{conjunction of}\ (1)\text{–}(5).\ \blacksquare
\]

\textbf{Phys.} The tier closes on a clean division of labour. The CPT involution is exact and cheap: three derived involutions composing to an involution, a phase that is always one of four roots of unity, an amplitude that never changes, and an orbit census in which the two charged sectors weigh identically while the neutral one weighs exactly twice either — a ratio the norm-layer factorization already carried. The geometry is blind to all of it, so the proton, the neutron and the antiproton share twelve reads at a fixed edge count and are separated by six, and the six are named. The apparent contradiction about whether the antiproton is "the same object" dissolves once configurations and structures are counted separately: distinct configurations, one orbit, no independent datum. And the honest limit is stated rather than smoothed over — the corpus has equal-amplitude mirror sectors and no process connecting them, no annihilation and no production, because nothing in the inventory builds one.

\vspace{2pt}\noindent\rule{\textwidth}{0.4pt}\vspace{2pt}

\medskip
\subsection*{\texorpdfstring{B.1.19\quad BM.33 — The status of the mesonic flicker test}{BM.33 B.1.19}}
\addcontentsline{toc}{subsection}{B.1.19 BM.33}
\phantomsection\label{tag:BM-33}
{\small \textit{Record.}\quad Source: Part 4 (Part4\_Baryomesonic\_era.md); tier: Tier BM.9 — The Mesonic Flicker Test: the Contraction Grammar, the Colour Projector, and the Protected Junction; canonical tag \texttt{BM.33}. Reproduced from the revision-1 canon; a cross-reference stub marks the former locus in the main corpus.}
\medskip
\textbf{Thm.}

%dsp-level:0
\[
(1)\quad\textbf{the answers:}\qquad\begin{array}{l|l}\text{the extended ensemble} & \mathrm{Def\text{-}BM.9}(1):\ \text{defined as a union inside the colour-blind topology space, with the projector applied downstream}\\\text{the boundary data} & \mathrm{BM.31}(3)\text{–}(5):\ \text{charge overlaps on}\ \{-2..+2\},\ \text{phase separates partially}\ (4\mathbb Z\subsetneq2\mathbb Z),\ \text{representation content does not separate}\ (1\ \text{singlet each});\ \textbf{the attachment count does}\\\text{the mesonic contribution} & \mathrm{BM.32}(1):\ Z^{(n)}_\delta=0\ \text{exactly}\ —\ \textbf{by filtration, not cancellation}\\\text{the cancellation mechanisms} & \mathrm{BM.32}(4),(5):\ \text{orthogonal to the question};\ \text{they kill pure directed cycles, and dissolved digons survive in covers}\\\text{the mass consequence} & \mathrm{BM.32}(3):\ \text{unchanged on both routes},\ \langle m\rangle_{\mathrm{ext}}=\langle m\rangle_\varepsilon\ \text{and}\ T_{\mathrm{ext}}=T_\varepsilon\ \text{exactly}\end{array}
\]

%dsp-level:0
\[
(2)\quad\textbf{the derived rows:}
\]

%dsp-level:0
\[
\qquad g(3)=1\ \wedge\ g_\delta(3)=0\ [\text{machine}];
\]

%dsp-level:0
\[
\qquad g_\delta(N)>0\ \forall N\in\{4,\dots,7\}\ [\text{machine}];
\]

%dsp-level:0
\[
\qquad g_\varepsilon(N)\ne0\iff3\mid N\ \text{on the frozen range}\ [\text{machine}];
\]

%dsp-level:0
\[
\dim(\mathbb C^3)^{\otimes3\,\mathfrak{sl}(3)}=\dim(\mathbb C^3\otimes\bar{\mathbb C}^3)^{\mathfrak{sl}(3)}=1\ [\text{machine}];
\]

%dsp-level:0
\[
\qquad W_3=(1,6,15,20,15,6,1),\ W_2=(1,4,6,4,1)\ [\text{machine}];
\]

%dsp-level:0
\[
(3)\quad\textbf{the scope rows:}
\]

%dsp-level:0
\[
\qquad \text{the protection is}\ N
\]

%dsp-level:0
\[
= 3\text{-specific}\ [\mathrm{BM.31}(2);\ \mathrm{BM.32}\ \mathrm{Rmk}]:\ \text{at}\ N
\]

%dsp-level:0
\[
\ge 4\ \text{the extended ensemble mixes and the average moves}\ [\text{machine: a witness at}\ N=4];
\]

%dsp-level:0
\[
\text{the phase separator is partial}\ [\mathrm{BM.31}(4)];
\]

%dsp-level:0
\[
\qquad \text{the charge separator fails at}\ Q\in\{-2,\dots,+2\}\ [\mathrm{BM.31}(3)];
\]

%dsp-level:0
\[
(4)\quad\textbf{the open rows touched:}
\]

%dsp-level:0
\[
\qquad \text{the junction-gluing census at}\ N
\]

%dsp-level:0
\[
\ge 4\ [\mathrm{BM.2}(3);\ \text{Census FT.50}]:\ \text{the mixed-sealing ensemble's reads are not computed here};
\]

%dsp-level:0
\[
\mathrm{BM.6}(1)\ \text{is discharged by Def-BM.8 but its dynamics coupling remains open}\ [\mathrm{BM.6}(5)];
\]

%dsp-level:0
\[
\qquad \text{item 13 (the amplitude}\leftrightarrow\text{spectral bridge)}\ \text{governs whether}\ (5)\text{'s dissolved-digon rows can carry a mass read at all};
\]

%dsp-level:0
\[
(5)\quad\textbf{the standing rules:}
\]

%dsp-level:0
\[
\qquad \neg\exists\ \text{a measured external value in any clause}\ [\mathrm{P5}];
\]

%dsp-level:0
\[
\qquad \neg\exists\ \text{a new dimensionful scale}\ [\mathrm{P1}];
\]

%dsp-level:0
\[
\qquad \neg\exists\ \text{a P4 addition}\ [\text{§Cal-BM stays empty}];
\]

%dsp-level:0
\[
\neg\exists\ \text{a new declaration}: \mathrm{Def\text{-}BM.9}\ \text{is a construction class on Def-BM.5 with}\ \mathrm{D2.14}\text{'s derived projector at its centre};
\]

%dsp-level:0
\[
\text{species names are confined to the identification layers}\ [\mathrm{Def\text{-}BM.3}(7),\ \mathrm{Def\text{-}BM.6}(7)];
\]

%dsp-level:0
\[
\qquad \text{completeness: the tier adds one integer sequence}\ (g_\delta,g_\varepsilon)\ \text{to the census vocabulary and no new count to any Census cell};
\]

%dsp-level:0
\[
(6)\quad\textbf{the adjudication:}
\]

%dsp-level:0
\[
\qquad \textbf{the baryon's flicker ensemble is protected from mesonic fluctuations, and the protection is grammatical}: \text{three attachments admit one contraction, so the mesonic subset is empty and its contribution is exactly zero};
\]

%dsp-level:0
\[
\text{the protection is a property of the number three}\ —\ \neg\text{of the dynamics,}\ \neg\text{of the amplitude layer,}\ \neg\text{of a symmetry}\ [\mathrm{BM.31}(2)];
\]

%dsp-level:0
\[
\text{consequence for the mass reads:}\ \mathrm{BM.20}(3)\text{'s path-integrated mass and BM.28}(4)\text{'s three ratios stand as derived, with}\ \langle m\rangle^{(n)}/\langle m\rangle^{(p)}
\]

%dsp-level:0
\[
= \tfrac43\ \text{on the total-deficit route and}\ 1\ \text{on the normalized route}\ [\mathrm{BM.30}(1),(2)].
\]

\textbf{Req.} D2.14 (Parts 1–2); P1, P5 (Parts 1–2); Def-H.32, FT.106(3) (Part 3); Def-BM.3(7), Def-BM.6(7), Def-BM.8, Def-BM.9, BM.2(3), BM.6(1), BM.6(5), BM.20(3), BM.28(4), BM.30(1), BM.30(2), BM.31, BM.32 (this part); Census FT.50 (row-wise); §Cal-BM; the machine sections §A–§E green [math; machine: \texttt{checks/\allowbreak{}verify\_\allowbreak{}H\_\allowbreak{}mesonic\_\allowbreak{}flicker\_\allowbreak{}test.py} §F].

\textbf{Pf.}

%dsp-level:0
\[
(1)\text{–}(3):\ \text{assembly of the adjudicated rows of BM.31 and BM.32}\ [\text{machine: §A–§E green}]
\]

%dsp-level:0
\[
(4):\ \text{ledger pointers}\ —\ \text{HistLog §II carries the resolution text}
\]

%dsp-level:0
\[
(5):\ \text{standing-rule record}\ [\mathrm{BM.12}(4),\ \mathrm{BM.18}(5),\ \mathrm{BM.24}(6),\ \mathrm{BM.29}(5)\ \text{pattern}];
\]

%dsp-level:0
\[
\qquad \mathrm{Def\text{-}BM.9}(2)\ \text{is}\ \mathrm{D2.14}\ \text{re-cited},\ \neg\text{declared}
\]

%dsp-level:0
\[
(6):\ \text{conjunction of}\ (1)\text{–}(5).\ \blacksquare
\]

\textbf{Phys.} The test returns a clean negative with an unusual provenance. Mesonic fluctuations do not enter the baryon's flicker integral, and they do not cancel inside it either — they were never admissible, because three colour attachments have exactly one way to contract and a partial pairing leaves one dangling, which the all-or-nothing law converts into a vanished block. That makes the baryon's mass reads safe: both functionals are invariant under adding zero-weight members, so the curvature-flux mass and the four-thirds splitting stand as derived. But the protection is arithmetic, not dynamical — it belongs to the number three. From four attachments up, the groupings mix, fifteen of twenty-five at six, and the extended ensemble moves the average. The corpus records that boundary instead of blurring it.

\vspace{2pt}\noindent\rule{\textwidth}{0.4pt}\vspace{2pt}

\medskip
\subsection*{\texorpdfstring{B.1.20\quad BM.38 — The status of the metric-energy tier: adjudication and scope}{BM.38 B.1.20}}
\addcontentsline{toc}{subsection}{B.1.20 BM.38}
\phantomsection\label{tag:BM-38}
{\small \textit{Record.}\quad Source: Part 4 (Part4\_Baryomesonic\_era.md); tier: Tier BM.10 — The Kinetic Energy from the Metric and the Vacuum-Balanced Entropy: the Dispersion Read of a Sealed Composite, the Deficit Reframing of the Rest Mass, the Force-Structure Ladder, and the Route-Artifact Adjudication; canonical tag \texttt{BM.38}. Reproduced from the revision-1 canon; a cross-reference stub marks the former locus in the main corpus.}
\medskip
\textbf{Thm.}

%dsp-level:0
\[
(1)\quad\textbf{the answers:}\qquad\begin{array}{l|l}\text{the kinetic energy} & \mathrm{BM.34}:\ \textbf{derived from the emergent metric}\ —\ \text{the ADM split by Schur, the master mass-shell identity, the Eulerian energy}\\ & E=N\sqrt{m^2+\gamma^{ij}p_ip_j}-N^ip_i,\ \text{the NR series with exact rational coefficients, the zero-point bounds}\ \bigl(\tfrac74,\ \tfrac38,\ \tfrac{\sqrt7}2\bigr);\\ & \text{a chart-side read under}\ \mathrm{CT2.6}\text{'s hypotheses at the point of use};\ \mathrm{BM.10}(1)\text{'s register obstruction stands untouched}\\\text{the vacuum-balanced entropy} & \mathrm{BM.35}:\ \text{the rest mass is the ledger deficit against the all-virtual vacuum}\ —\ m_{\mathrm{rest}}=\Delta S_{\mathfrak L}/\mathfrak h=m_{\mathrm{hid}};\\ & \text{the balance identity}\ \ln4=2\ln2\ \text{per bidirectional channel};\ \text{a naming,}\ \neg\text{a substance};\ \text{three deficit currencies, one reference point, no exchange rate}\\\text{the inter-family ratios} & \mathrm{BM.36}:\ \textbf{constrained,}\ \neg\textbf{derived}\ —\ \text{the flag}\ (6,16,17,25),\ \text{the transport channel, the Gram spectrum, the exact massless line},\\ & \text{the sandwich, the ceilings};\ \text{the Yukawa-like coupling is an open gap}\ \\\text{the counting trap} & \mathrm{BM.37}:\ \textbf{closed}\ —\ \text{routes}\ (\mathrm i)\text{--}(\mathrm{iii})\ \text{are measure/amplitude functionals with input alphabets disjoint from the dispersion energy};\\ & \text{their ratios are barred as physical masses and stand as exact measure prices}\end{array}
\]

%dsp-level:0
\[
(2)\quad\textbf{the derived inventory:}
\]

%dsp-level:0
\[
\qquad N^2\det\gamma=-\det g;
\]

%dsp-level:0
\[
\qquad \bigl(p_0-N^ip_i\bigr)^2=N^2\bigl(m^2+\gamma^{ij}p_ip_j\bigr);
\]

%dsp-level:0
\[
\qquad g^{00}=-\tfrac1{N^2},\ \ g^{0i}=\tfrac{N^i}{N^2},\ \ g^{ij}=\gamma^{ij}-\tfrac{N^iN^j}{N^2};
\]

%dsp-level:0
\[
\sqrt{1+x}:\ \ \text{the coefficients}\ \bigl(1,\tfrac12,-\tfrac18,\tfrac1{16},-\tfrac5{128},\tfrac7{256}\bigr);
\]

%dsp-level:0
\[
\qquad E^2\ge\tfrac74m_{\mathrm{hid}}^2,\quad\mathrm{KE}_{\mathrm{NR}}\ge\tfrac38m_{\mathrm{hid}},\quad E\ge\tfrac{\sqrt7}2m_{\mathrm{hid}};
\]

%dsp-level:0
\[
\ln4=2\ln2;
\]

%dsp-level:0
\[
\qquad \Delta S_{\mathfrak L}=E_{\mathrm{int}}\ln2;
\]

%dsp-level:0
\[
\qquad m_{\mathrm{rest}}=m_{\mathrm{hid}};
\]

%dsp-level:0
\[
\qquad \text{the Gibbs census}\ (165\ \text{compositions, one equality case});
\]

%dsp-level:0
\[
\dim\mathcal K_n=(6,16,17,25);
\]

%dsp-level:0
\[
\qquad \text{the ratios}\ \bigl(\tfrac13,8,\tfrac1{25}\bigr);
\]

%dsp-level:0
\[
\qquad \text{the Gram spectrum}\ \bigl\{\tfrac{g^2}4,\tfrac{g^2}4,0,\tfrac{g^2+g'^2}4\bigr\};
\]

%dsp-level:0
\[
\qquad \frac{M_Z^2}{M_W^2}=\frac{g^2+g'^2}{g^2};
\]

%dsp-level:0
\[
W(0){=}20,\ W(\pm1){=}15;
\]

%dsp-level:0
\[
\qquad 1\ne\tfrac43\ne\tfrac{58}{13}\ne\tfrac{2398}{401};
\]

%dsp-level:0
\[
\qquad Q=\tfrac23\ \text{re-verified phase-independently in}\ \mathbb Q[c,s,r,u]/(c^2{+}s^2{-}1,\ r^2{-}3,\ u^2{-}2)
\]

%dsp-level:0
\[
\qquad [\text{machine: all rows exact}];
\]

%dsp-level:0
\[
(3)\quad\textbf{the open rows touched, not closed:}
\]

%dsp-level:0
\[
\qquad \text{the relating equation across the mass reads: the three deficit currencies remain unconverted}\ [\mathrm{BM.35}\ \mathrm{Rmk}];
\]

%dsp-level:0
\[
\text{the amplitude}\leftrightarrow\text{spectral/mass bridge: route}\ (\mathrm{iii})\text{'s and}\ Q=\tfrac23\text{'s mass forms};
\]

%dsp-level:0
\[
\qquad \text{the mass-mechanism inventory: the dispersion read joins the read-class taxonomy}\ [\mathrm{FT.84}(1)\ \text{pattern}];
\]

%dsp-level:0
\[
\text{deficit}\ 1'\ —\ \text{the gauge coupling of}\ Q\ [\mathrm{Def\text{-}BM.8}(8)\ \text{stands}];
\]

%dsp-level:0
\[
\qquad \mathrm{FT.159}\text{'s rate modulation untouched};
\]

%dsp-level:0
\[
\textbf{the family-profile map}\ [\mathrm{BM.36}(6)]:\ \text{the construction locus is named}\ —\ \text{the seed transport}\ [\mathrm{EW.32},\ \mathrm{EW.33},\ \mathrm{Def\text{-}H.37}(1)];
\]

%dsp-level:0
\[
\qquad \text{the construction itself issues from no clause of this tier};
\]

%dsp-level:0
\[
(4)\quad\textbf{the standing rules:}
\]

%dsp-level:0
\[
\qquad \neg\text{a new dimensionful scale}\ [\mathrm{P1}];
\]

%dsp-level:0
\[
\qquad \neg\text{a measured value}\ [\mathrm{P5};\ \text{machine: zero external tokens, zero decimal numerals in this tier}];
\]

%dsp-level:0
\[
\neg\text{a P4 addition}\ [\text{§Cal-BM stays empty}];
\]

%dsp-level:0
\[
\qquad \neg\text{a new declaration}\ [\mathrm{Def\text{-}BM.10}\ \text{is a construction: its centre is}\ \mathrm{CT1.12}\text{'s dictionary}\ +\ \mathrm{FT.44}(4)\text{'s read}];
\]

%dsp-level:0
\[
\neg\exists\ \text{an}\ [\mathrm{EXT}]\ \text{datum in any derivation};
\]

%dsp-level:0
\[
\qquad \text{the extra-axiomatic list}\ [\text{A §0.2a}]\ \text{gains no member:}\ \mathrm{CT2.6}\text{'s hypotheses and}\ \mathrm{BM.8}(2)\text{'s stationarity scope are inherited at the point of use, both pre-existing};
\]

%dsp-level:0
\[
\text{completeness: the tier adds no new count to any Census cell};
\]

%dsp-level:0
\[
\qquad \text{the new exact vocabulary: the NR series coefficients, the zero-point constants}\ \bigl(\tfrac74,\ \tfrac38,\ \tfrac{\sqrt7}2\bigr),
\]

%dsp-level:0
\[
\ \text{the Gram-spectrum rows, the dimension ratios}\ \bigl(\tfrac13,8,\tfrac1{25}\bigr);
\]

%dsp-level:0
\[
(5)\quad\textbf{the adjudication:}
\]

%dsp-level:0
\[
\qquad \boxed{\ \text{the dynamical energy of a sealed composite is derived inside the corpus, end to end}\ }
\]

%dsp-level:0
\[
\text{the vacuum supplies the entropy-maximal flat reference}\ [\mathrm{BM.35}(1),(2)];
\]

%dsp-level:0
\[
\qquad \text{crystallization books the deficit whose ledger read is the rest mass}\ [\mathrm{BM.35}(3),(4)];
\]

%dsp-level:0
\[
\text{the metric dictionary turns the same sector data into lapse, shift and spatial metric}\ [\mathrm{Def\text{-}BM.10}(2);\ \mathrm{BM.34}(1)];
\]

%dsp-level:0
\[
\qquad \text{the mass shell turns those into the dispersion relation whose motion term is the kinetic energy}\ [\mathrm{BM.34}(4),(5)];
\]

%dsp-level:0
\[
\text{the three ensemble mass routes are reclassified as static measure functionals}\ —\ \text{the trap that read their ratios as masses is closed}\ [\mathrm{BM.37}(2),(6)];
\]

%dsp-level:0
\[
\text{the open face is the family, not the energy:}\ \neg\exists\ \text{a derived map assigning transport profiles to family lines}\ [\mathrm{BM.36}(6)].
\]

\textbf{Req.} P1, P4, P5 (Parts 1–2); D2.9 (Parts 1–2); CT0.2 (Part 1); CT1.12, CT2.6 (Part 2); EW.32, EW.33 (Part 2); Def-H.37(1), FT.84(1), FT.159 (Part 3); Census FT.44(4) (row-wise); Def-BM.8, Def-BM.10, BM.34, BM.35, BM.36, BM.37 (this part); the machine sections §A–§E green [math; machine: \texttt{checks/\allowbreak{}verify\_\allowbreak{}H\_\allowbreak{}kinetic\_\allowbreak{}energy\_\allowbreak{}metric.py} §E].

\textbf{Pf.}

%dsp-level:0
\[
(1)\text{--}(4):\ \text{assembly of the adjudicated rows of BM.34--BM.37}\ [\text{machine: §A--§E green}]
\]

%dsp-level:0
\[
(5):\ \text{conjunction of}\ (1)\text{--}(4).\ \blacksquare
\]

\textbf{Phys.} The tier's verdict in one sentence: a sealed composite's energy is now a derived object from end to end, and its mass ratios across families are not. The vacuum is the reference — maximal in entropy, flat in geometry; crystallization books a deficit against it, and the ledger's bits, divided by the single free scale, are the rest mass, identical to the informational mass the floors have always carried. The same sector data that built the metric build the motion: lapse, shift, spatial metric, mass shell, dispersion, kinetic energy, with a zero-point bound that is exact arithmetic on the Born width. The three ensemble routes keep every rational they ever produced, reclassified: prices of the measure, barred as masses, the counting trap closed by disjoint input alphabets rather than by anyone's preference. What stays open is named as precisely as what is closed: no theorem converts one deficit currency into another, no bridge carries an amplitude into a mass, and no map assigns transport profiles to family lines — the Yukawa-like coupling is the tier's named open construction, with its locus in the seed transport. No new scale, no new parameter, no declaration, no measured value: the extra-axiomatic surface of the corpus is exactly as wide as it was, and the energy of the composite is one derivation narrower than the open ledger.

\medskip
\subsection*{\texorpdfstring{B.1.21\quad BM.50 — The status of the EM self-interaction correction: adjudication and scope}{BM.50 B.1.21}}
\addcontentsline{toc}{subsection}{B.1.21 BM.50}
\phantomsection\label{tag:BM-50}
{\small \textit{Record.}\quad Source: Part 4 (Part4\_Baryomesonic\_era.md); tier: Tier BM.12 — The EM Self-Interaction Correction: the Scar's Field Energy, the Infinite-Range Seal Asymmetry, the Mass-Splitting Theorem, and the Beta-Spectrum Adjudication; canonical tag \texttt{BM.50}. Reproduced from the revision-1 canon; a cross-reference stub marks the former locus in the main corpus.}
\medskip
\textbf{Thm.}

%dsp-level:0
\[
(1)\quad\textbf{the answers:}\qquad\begin{array}{l|l}\text{the scar's field energy} & \mathrm{BM.46}:\ \textbf{derived}\ —\ E_{\mathrm{EM}}^{\mathrm{self}}=c\,\varepsilon_{\mathrm{EM}}q^2m_{\mathrm{hid}},\ \text{two agreeing routes, both distributions priced}\ \bigl(\tfrac12,\ \tfrac35\bigr),\ e\ \text{symbolic}\\\text{the mass splitting} & \mathrm{BM.46}(4):\ \textbf{derived and strictly positive}\ —\ \Delta m_{\mathrm{EM}}=\tfrac{3ce^2\ln2}{2\pi\mathfrak h}>0;\ \textbf{the}\ \mathrm{BM.28}(1)\ \textbf{degeneracy is broken};\ \text{the charged member is the heavier one}\\\text{the infinite range} & \mathrm{BM.47}:\ \textbf{handled exactly}\ —\ \text{UV-finite at the Born envelope},\ \text{IR-finite per member and per universe}\ [\mathrm{EM.11}],\ \text{the seal asymmetry derived (colour local, EM global)},\ \text{the nearby fraction}\ 1-\tfrac1k\\\text{the binding potential} & \mathrm{BM.48}:\ \textbf{derived as a two-body read}\ (V=-\varepsilon_{\mathrm{EM}}/d\ \text{for the forced pair});\ \text{the bound channel}\ \neg\text{constructed};\ \text{the}\ m_{\mathrm{eff}}-E_{\mathrm{bind}}\ \text{form refuted as a category error}\\\text{the beta decay} & \mathrm{BM.49}:\ \textbf{remains kinematically closed, strictly deeper}\ —\ Q_{\mathrm{corr}}=-(\Delta m_{\mathrm{EM}}+m_\ell+m_\nu)<0;\ \text{the ``correction makes the decay possible'' claim is refuted by sign arithmetic}\\\text{the}\ E_{\max}:\langle E\rangle\ \text{proposal} & \mathrm{BM.49}(3)\text{--}(4):\ \textbf{refuted}\ —\ \text{the derived conditional ratio is}\ 2{:}1,\ \text{scale-invariant};\ \tfrac43\ \text{is a measure price}\ [\mathrm{BM.28}(2);\ \mathrm{BM.37}]\end{array}
\]

%dsp-level:0
\[
(2)\quad\textbf{the derived inventory:}
\]

%dsp-level:0
\[
\qquad E_{\mathrm{EM}}^{\mathrm{self}}\ \text{coefficients}\ \bigl\{\tfrac1{8\pi},\tfrac1{40\pi},\tfrac3{20\pi}\bigr\}\cdot e^2q^2/\Delta x_{\mathrm B};
\]

%dsp-level:0
\[
\qquad \Delta m_{\mathrm{EM}}=\tfrac{3ce^2\ln2}{2\pi\mathfrak h};
\]

%dsp-level:0
\[
\qquad \frac{\Delta m_{\mathrm{EM}}}{m_{\mathrm{hid}}}=c\,\varepsilon_{\mathrm{EM}}\in\mathfrak q^0;
\]

%dsp-level:0
\[
f(k)=1-\tfrac1k;
\]

%dsp-level:0
\[
\qquad V(d)=-\tfrac{\varepsilon_{\mathrm{EM}}}d\ (\text{the forced pair});
\]

%dsp-level:0
\[
\qquad E_n=-\tfrac{\varepsilon_{\mathrm{EM}}^2m_r}{2n^2},\ \tfrac{E_2}{E_1}=\tfrac14;
\]

%dsp-level:0
\[
\qquad \frac{\langle E\rangle}{E_{\max}}=\tfrac12,\quad E_{\max}:\langle E\rangle=2:1;
\]

%dsp-level:0
\[
\qquad 2\ne\tfrac43
\]

%dsp-level:0
\[
\qquad [\text{machine: all rows exact}];
\]

%dsp-level:0
\[
(3)\quad\textbf{the open rows touched, not closed:}
\]

%dsp-level:0
\[
\qquad \text{the numerical value of}\ e\ [\mathrm{EM.25}(7):\ \text{a calibration step requiring one datum};\ \mathrm{P5}];
\]

%dsp-level:0
\[
\qquad \text{the dynamics of}\ Q\ [\mathrm{BM.6}(5)];
\]

%dsp-level:0
\[
\qquad \text{an EM bound channel}\ [\mathrm{BM.48}(3)];
\]

%dsp-level:0
\[
\text{the internal-multipole resolution of the neutral class}\ [\mathrm{BM.46}(2)];
\]

%dsp-level:0
\[
\qquad \text{the lepton mass reads}\ [\mathrm{BM.43}(2)];
\]

%dsp-level:0
\[
\qquad \text{the route choice}\ [\mathrm{BM.30}(5)];
\]

%dsp-level:0
\[
\qquad \text{deficit}\ 1'\ \text{(the chiral assignment)}\ [\mathrm{FT.125}(4)];
\]

%dsp-level:0
\[
\text{radiation, back-reaction and the reaction's field reconfiguration during the transition}\ [\mathrm{Def\text{-}BM.12}(7)\ \text{scope}];
\]

%dsp-level:0
\[
(4)\quad\textbf{the standing rules:}
\]

%dsp-level:0
\[
\qquad \neg\text{a P4 addition}\ (e\ \text{is EW.35's logged combination};\ \text{§Cal-BM stays empty});
\]

%dsp-level:0
\[
\qquad \neg\text{a new dimensionful scale}\ [\mathrm{P1}];
\]

%dsp-level:0
\[
\qquad \neg\text{a measured value and}\ \neg\text{a numerical coupling}\ [\mathrm{P5};\ \mathrm{EM.25}(7),(9);\ \text{machine: zero external tokens, zero decimal numerals in this tier}];
\]

%dsp-level:0
\[
\neg\exists\ \text{an}\ [\mathrm{EXT}]\ \text{datum};
\]

%dsp-level:0
\[
\qquad \neg\text{a new declaration}\ [\mathrm{Def\text{-}BM.12}\ \text{is a construction; the distribution factor is priced},\ \neg\text{adopted}];
\]

%dsp-level:0
\[
\qquad \text{species names confined to the identification layers}\ [\mathrm{Def\text{-}BM.3}(7),\ \mathrm{Def\text{-}BM.6}(7)];
\]

%dsp-level:0
\[
\text{amplitude}\ne\text{mass}\ [\mathrm{FT.79}(5)];
\]

%dsp-level:0
\[
\qquad \text{rung}\ne\text{scale}\ [\mathrm{D0.8}];
\]

%dsp-level:0
\[
\qquad \text{topology labels are internal bookkeeping}\ [\mathrm{Def\text{-}H.32}(4)];
\]

%dsp-level:0
\[
\qquad \neg\text{a HistLog disposition imported as premise};
\]

%dsp-level:0
\[
\text{the extra-axiomatic list}\ [\text{A §0.2a}]\ \text{gains no member};
\]

%dsp-level:0
\[
\qquad \text{completeness: the tier adds no count to any Census cell};
\]

%dsp-level:0
\[
\ \text{the new exact vocabulary: the field-energy coefficients, the splitting monomial,}\ f(k),
\]

%dsp-level:0
\[
\ \text{the spectral moment ratio}\ 2;
\]

%dsp-level:0
\[
(5)\quad\textbf{the adjudication:}
\]

%dsp-level:0
\[
\qquad \boxed{\ \text{the EM self-interaction is derived, and it splits the baryon rows in the direction that keeps the decay shut}\ }
\]

%dsp-level:0
\[
\text{the scar's field energy is a finite, positive, charge-selected transit-channel correction to the crystallized-count read}\ —\ \text{the degeneracy}\ \mathrm{BM.28}(1)\ \text{records is broken by}\ \Delta m_{\mathrm{EM}}>0\ \text{with the proton heavier};
\]

%dsp-level:0
\[
\text{the infinite range is paid for exactly}\ —\ \text{the Born envelope in the ultraviolet,}\ \mathrm{EM.11}\text{'s universe-scale neutrality in the infrared}\ —\ \text{and the seal asymmetry with colour is a theorem};
\]

%dsp-level:0
\[
\text{the binding energy stays where it belongs}\ —\ \text{with the unconstructed bound system}\ —\ \text{and no free member's mass carries it};
\]

%dsp-level:0
\[
\text{the beta-decay vertex remains grammatically complete and kinematically closed on every adopted read, deeper by the splitting};
\]

%dsp-level:0
\[
\qquad \text{the}\ \tfrac43\ \text{measure price does not drive any spectral moment}\ —\ \text{the derived conditional ratio is}\ 2{:}1\ \text{and scale-invariant}.
\]

\textbf{Req.} P1, P4, P5 (Parts 1–2); D0.8, D2.9 (Parts 1–2); EM.6, EM.11, EM.15, EM.25(7), EM.25(9), EW.22, EW.33, EW.35 (Part 2); FT.79(5), FT.125(4), Def-H.32(4) (Part 3); Def-BM.3(7), Def-BM.4, Def-BM.6(7), Def-BM.8, Def-BM.10, Def-BM.12, BM.6(5), BM.7, BM.9(4), BM.11, BM.13, BM.17, BM.23(4), BM.26, BM.28(1), BM.28(2), BM.28(7), BM.30(5), BM.37(2), BM.37(3), BM.39(1), BM.41(3), BM.41(4), BM.41(5), BM.43(2), BM.46, BM.47, BM.48, BM.49 (this part); the machine sections §A–§E green [math; machine: \texttt{checks/\allowbreak{}verify\_\allowbreak{}H\_\allowbreak{}em\_\allowbreak{}self\_\allowbreak{}interaction.py} §E].

\textbf{Pf.}

%dsp-level:0
\[
(1)\text{--}(4):\ \text{assembly of the adjudicated rows of BM.46--BM.49}\ [\text{machine: §A--§D green}]
\]

%dsp-level:0
\[
(5):\ \text{conjunction of}\ (1)\text{--}(4).\ \blacksquare
\]

\textbf{Phys.} The tier's verdict in one sentence: the proton's scar costs it mass, exactly and positively, and the neutron's neutrality costs it nothing — so the two baryon rows split, with the charged one heavier, and the beta decay the corpus assembled stays shut, now by a wider symbolic margin. Everything numerical is rational; everything coupled is symbolic; nothing measured enters anywhere. The infinite range is handled rather than feared: the envelope regularizes the near field, the universe's own neutrality regularizes the far one, and the asymmetry with colour — a local projector against a global census — is proved rather than asserted. The binding energy is drawn as a potential and barred as a mass term, because the atom it would belong to is a construction the corpus has not made. And the four-thirds keeps the only job the adjudications ever gave it: pricing the charge filter on an ensemble measure, two-thirds of a spectrum ratio away from anything an emitted electron could measure.

\medskip
\subsection*{\texorpdfstring{B.1.22\quad BM.56 — The status of the full-scope mass splitting: adjudication and scope}{BM.56 B.1.22}}
\addcontentsline{toc}{subsection}{B.1.22 BM.56}
\phantomsection\label{tag:BM-56}
{\small \textit{Record.}\quad Source: Part 4 (Part4\_Baryomesonic\_era.md); tier: Tier BM.13 — The Full-Scope Mass Splitting: the Lyapunov Foundation, the Complete-Force Census, the Non-Uniform Clock, the Binding Zeros, and the Exact Ratio; canonical tag \texttt{BM.56}. Reproduced from the revision-1 canon; a cross-reference stub marks the former locus in the main corpus.}
\medskip
\textbf{Thm.}

%dsp-level:0
\[
(1)\quad\textbf{the answers:}\qquad\begin{array}{l|l}\text{the Lyapunov functional} & \mathrm{BM.51}:\ \textbf{the foundation}\ —\ V=8\Delta S_{\mathrm{pair}}+O(\delta\rho^3)\ [\mathrm{L3}(2)],\ \text{the inertial read}\ m_i^2(q)=\Delta S(q)\ [\mathrm{L3}(4)],\ \text{the iteration}\ (1-\tfrac\gamma4)^{2n},\\ & \bigl(\tfrac9{16}\bigr)^6\ \text{at the sharp point};\ \text{the class contribution}\ \Delta m_{\mathrm{Lyapunov}}=0,\ \text{species-independently}\\\text{the complete force operator} & \mathrm{BM.52}:\ \textbf{the sum of all sectors, censused}\ —\ J^\pm,G,H,W,\text{colour},D^\pm\ \text{contribute}\ 0\ \text{each by frozen theorem};\ \mathrm{EM}\ \text{is the only nonzero sector}:\\ & \text{the flipping block}\ 2|q|\in\{0,2\}\ \text{(dimensionless, conversion deferred}\ [\mathrm{FT.131}(3)])\ \text{and the field monomial}\ \Delta m_{\mathrm{EM}}=\tfrac{3ce^2\ln2}{2\pi\mathfrak h}\\\text{the metric} & \mathrm{BM.53}(2):\ \textbf{the emergent background with zero splitting contribution}\ —\ \text{the seven charge-blind monomials re-censused}\\\text{the non-uniform time} & \mathrm{BM.53}(3)\text{--}(5):\ \textbf{handled exactly}\ —\ \text{the contraction law and its logarithm},\ \neg\text{a uniform-}\tau\ \text{exponential},\ \text{MT with the saturated witness},\\ & \text{the resolvability bound}\ \Delta t\ge\pi\mathfrak h/(3ce^2\ln2)\\\text{the binding energies} & \mathrm{BM.54}:\ \textbf{zero between the rows on every constructed read}\ (count-additive, ensemble-degenerate);\ \text{negative where a pair system would exist}\ —\ \text{none does}\ [\mathrm{BM.48}(3)\text{--}(4)]\\\text{the self-interactions} & \mathrm{BM.46}/\mathrm{BM.52}(3):\ \textbf{positive for the charged member, zero for the neutral at monopole order};\ \text{the other sectors blind or degenerate between the rows}\\\text{the mass splitting} & \mathrm{BM.55}:\ \textbf{derived from the corpus's own structures}\ —\ \Delta m=-\tfrac{3ce^2\ln2}{2\pi\mathfrak h}<0,\quad\frac{m_n}{m_p}=\frac{4\pi}{4\pi+ce^2}\in(0,1),\quad\text{neutron lighter on every constructed read}\end{array}
\]

%dsp-level:0
\[
(2)\quad\textbf{the derived inventory:}
\]

%dsp-level:0
\[
\qquad \bigl(1-\tfrac\gamma4\bigr)^2+\tfrac{\gamma(8-\gamma)}{16}=1;
\]

%dsp-level:0
\[
\qquad \bigl(\tfrac9{16}\bigr)^6=\tfrac{531441}{16777216};
\]

%dsp-level:0
\[
\qquad V=8\Delta S_{\mathrm{pair}}\ \text{at both branch witnesses}\ \bigl(16\cdot\tfrac34=12,\ 16\cdot\tfrac1{12}=\tfrac43\bigr);
\]

%dsp-level:0
\[
2|q|\in\{0,2\};
\]

%dsp-level:0
\[
\qquad \Delta m=-\tfrac{3ce^2\ln2}{2\pi\mathfrak h};
\]

%dsp-level:0
\[
\qquad \frac{m_n}{m_p}=\frac{4\pi}{4\pi+ce^2};
\]

%dsp-level:0
\[
\qquad \Delta t\ge\frac{\pi\mathfrak h}{3ce^2\ln2};
\]

%dsp-level:0
\[
\qquad E^{\mathrm{info}}_{\mathrm{bind}}=0;
\]

%dsp-level:0
\[
\qquad \mathcal B^{(n)}=\mathcal B^{(p)}
\]

%dsp-level:0
\[
\qquad [\text{machine: all rows exact}];
\]

%dsp-level:0
\[
(3)\quad\textbf{the open rows touched, not closed:}
\]

%dsp-level:0
\[
\qquad \text{the numerical value of}\ e\ [\mathrm{EM.25}(7):\ \text{one calibration datum};\ \mathrm{P5}];
\]

%dsp-level:0
\[
\qquad \text{the deficit/inertial-to-informational conversions}\ [\mathrm{BM.51}(5)];
\]

%dsp-level:0
\[
\mathrm{FT.131}(3)\text{'s deferred frequency read}\ (\text{the}\ 2|q|\ \text{splitting's conversion});
\]

%dsp-level:0
\[
\qquad \mathrm{FT.159}\text{'s rate modulation};
\]

%dsp-level:0
\[
\qquad \text{the constituent-transport residue that would reverse the order}\ [\mathrm{BM.55}(4)];
\]

%dsp-level:0
\[
\text{the dynamics of}\ Q\ [\mathrm{BM.6}(5)];
\]

%dsp-level:0
\[
\qquad \text{the EM bound channel}\ [\mathrm{BM.48}(3)];
\]

%dsp-level:0
\[
\qquad \text{the internal-multipole resolution of the neutral class}\ [\mathrm{BM.46}(2)];
\]

%dsp-level:0
\[
\qquad \text{the route choice}\ [\mathrm{BM.30}(5)];
\]

%dsp-level:0
\[
(4)\quad\textbf{the standing rules:}
\]

%dsp-level:0
\[
\qquad \neg\text{a P4 addition}\ (e\ \text{is EW.35's logged combination};\ \text{§Cal-BM stays empty});
\]

%dsp-level:0
\[
\qquad \neg\text{a new dimensionful scale}\ [\mathrm{P1}];
\]

%dsp-level:0
\[
\qquad \neg\text{a measured value}\ [\mathrm{P5};\ \text{machine: zero external tokens, zero decimal numerals in this tier}];
\]

%dsp-level:0
\[
\neg\exists\ \text{an}\ [\mathrm{EXT}]\ \text{datum};
\]

%dsp-level:0
\[
\qquad \neg\text{a new declaration and}\ \neg\text{a new definition}\ (\text{this tier adds no Def}: \text{every object is prior construction});
\]

%dsp-level:0
\[
\qquad \text{species names confined to the identification layer}\ [\mathrm{Def\text{-}BM.3}(7),\ \mathrm{Def\text{-}BM.6}(7)];
\]

%dsp-level:0
\[
\text{amplitude}\ne\text{mass}\ [\mathrm{FT.79}(5)];
\]

%dsp-level:0
\[
\qquad \text{rung}\ne\text{scale}\ [\mathrm{D0.8}];
\]

%dsp-level:0
\[
\qquad \text{topology labels are internal bookkeeping}\ [\mathrm{Def\text{-}H.32}(4)];
\]

%dsp-level:0
\[
\qquad \neg\text{a HistLog disposition imported as premise}\ (\text{item 7 is a pointer});
\]

%dsp-level:0
\[
\text{the extra-axiomatic list}\ [\text{A §0.2a}]\ \text{gains no member}\ (\mathrm{L0}\text{'s dynamical scope,}\ \mathrm{BM.8}(2)\text{'s stationarity scope and}\ \mathrm{CT2.6}\text{'s hypotheses are inherited, all pre-existing});
\]

%dsp-level:0
\[
\qquad \text{completeness: no new count enters any Census cell};
\]

%dsp-level:0
\[
(5)\quad\textbf{the adjudication:}
\]

%dsp-level:0
\[
\qquad \boxed{\ \text{the full-scope mass splitting is derived: four contributions are exact zeros, the fifth is the EM field monomial, and the ratio is}\ \dfrac{4\pi}{4\pi+ce^2}\ }
\]

%dsp-level:0
\[
\text{the Lyapunov foundation splits nothing between the two baryon rows};
\]

%dsp-level:0
\[
\qquad \text{the complete force operator splits them by exactly one sector};
\]

%dsp-level:0
\[
\qquad \text{the metric and the binding split nothing};
\]

%dsp-level:0
\[
\text{the non-uniform clock bounds the splitting's resolvability without shifting its value};
\]

%dsp-level:0
\[
\qquad \text{the neutron is lighter than the proton on every constructed read, and the}\ \beta\text{-decay vertex stays kinematically closed};
\]

%dsp-level:0
\[
\text{what the corpus does not derive is the world's order}\ —\ \text{that residue is the constituent-transport difference}\ \mathrm{BM.55}(4)\ \text{names},
\]

%dsp-level:0
\[
\ \text{and no clause of this tier invents it}.
\]

\textbf{Req.} P1, P4, P5 (Parts 1–2); D0.8, D2.9 (Parts 1–2); EM.25(7) (Part 2); FT.79(5), FT.131(3), FT.159, Def-H.32(4) (Part 3); Def-BM.3(7), Def-BM.6(7), Def-BM.12, BM.6(5), BM.30(5), BM.46(2), BM.48(3), BM.51, BM.52, BM.53, BM.54, BM.55 (this part); the machine sections §A–§F green [math; machine: \texttt{checks/\allowbreak{}verify\_\allowbreak{}H\_\allowbreak{}full\_\allowbreak{}scope\_\allowbreak{}mass\_\allowbreak{}splitting.py} §F].

\textbf{Pf.}

%dsp-level:0
\[
(1)\text{--}(4):\ \text{assembly of the adjudicated rows of BM.51--BM.55}\ [\text{machine: §A--§E green}]
\]

%dsp-level:0
\[
(5):\ \text{conjunction of}\ (1)\text{--}(4).\ \blacksquare
\]

\textbf{Phys.} The verdict of the full scope is compact enough to state in one line of algebra. The Lyapunov tower — the functional, its contraction, its deficit links, its inertial read — gives the neutron and the proton identical masses because it gives them identical inputs. The complete force operator, censused over all six sectors with every zero backed by a frozen theorem, splits them by exactly one term: the electromagnetic field energy of the scar. The metric splits nothing, the binding splits nothing, and the non-uniform clock shifts no value while bounding how finely the splitting can be resolved. The ratio is then arithmetic on monomials — four pi over four pi plus $ce^2$ — exact, symbolic, and strictly below one: on every read the corpus can construct, the neutron is the lighter baryon and the beta decay stays shut. The world's opposite order needs a constituent-level difference the inventory does not contain, and the tier's last act of full scope is to say so precisely, to the monomial, without inventing the missing piece.

\medskip
\subsection*{\texorpdfstring{B.1.23\quad BM.64 — The status of the constituent transport splitting: adjudication and scope}{BM.64 B.1.23}}
\addcontentsline{toc}{subsection}{B.1.23 BM.64}
\phantomsection\label{tag:BM-64}
{\small \textit{Record.}\quad Source: Part 4 (Part4\_Baryomesonic\_era.md); tier: Tier BM.14 — The Constituent Transport Splitting: the Vertex Asymmetry, the Observation Frames, the Six-State Quark Assembly, and the Neutron Row; canonical tag \texttt{BM.64}. Reproduced from the revision-1 canon; a cross-reference stub marks the former locus in the main corpus.}
\medskip
\textbf{Thm.}

%dsp-level:0
\[
(1)\quad\textbf{the answers:}\qquad\begin{array}{l|l}\text{the vertex asymmetry} & \mathrm{BM.57}:\ \textbf{the count derived, the mechanism refuted}\ —\ |A_3|=3\ (2\ \text{nontrivial})\ \text{against}\ |S_1\times S_1|=1\ (0\ \text{nontrivial}),\ \text{asymmetry}\ 2\ \text{by both conventions};\\ & \Delta m_{\mathrm{vertex}}=0\ \text{exactly, on three independent annihilation rows};\ \text{the residual content is the family triplet}\ (3\ \text{lines})\\\text{the observation frame} & \mathrm{BM.58}:\ \textbf{refuted}\ —\ \text{zero charge tokens in the metric blocks},\ \text{charge exponent}\ 0\ \text{in all seven monomials},\ \text{masks act on sectors};\\ & \Delta m_{\mathrm{metric}}=0;\ \text{the frame's real content is the seal asymmetry, whose mass row is}\ \mathrm{BM.46}(4)\ \text{and points the other way}\\\text{the constituent transport splitting} & \mathrm{BM.59}:\ \textbf{derived, exact, strictly positive}\ —\ \Sigma_{\mathrm{cl}}(Q)=\ln W(Q),\ \Sigma_{\mathrm B}=\ln\tfrac43,\ \Delta m_{\mathrm{transport}}=\ln(4/3)/\mathfrak h>0;\\ & \text{the constituent row}\ \ln2/\mathfrak h;\ \mathrm{BM.55}(4)\text{'s named residue is filled, on all four of its grounds}\\\text{the quark generation structure} & \mathrm{BM.60}:\ \textbf{the}\ 3\times2=6\ \textbf{assembly derived}\ —\ \text{three floors}\ 2E\in\{12,24,60\},\ \text{the branch split on every defect datum},\\ & \text{six pairwise distinct}\ D_S\ \text{in four fields},\ \text{the branch order uniform and exact},\ \text{the generation map the unique order-isomorphism};\\ & \textbf{open:}\ \text{the flavour composition of a sealed composite}\ [\mathrm{BM.28}(7)],\ \text{the currency relating equation}\ \\\alpha_B & \mathrm{BM.61}:\ \textbf{not derivable as}\ f(\alpha_{\mathrm{EM}},\alpha_{\mathrm{Strong}})\ —\ \text{the obstruction named on four grounds}\ (\mathrm O1)\text{--}(\mathrm O4);\\ & \textbf{not added as a parameter};\ \text{the exact object at the slot is}\ \Sigma_{\mathrm B}=\ln\tfrac43,\ \text{proved coupling-blind}\\\text{the full mass splitting} & \mathrm{BM.62}:\ \textbf{derived}\ —\ \Delta m=\bigl[\ln(4/3)-6c\varepsilon_{\mathrm{EM}}\ln2\bigr]/\mathfrak h;\ \varepsilon^{\mathrm{crit}}(c)=\Sigma_{\mathrm B}/(6c\ln2);\\ & \tfrac19<\varepsilon^{\mathrm{crit}}(\tfrac35)<\tfrac18<\varepsilon^{\mathrm{crit}}(\tfrac12)<\tfrac17;\qquad\textbf{the neutron is heavier on the sub-critical branch}\\\text{the decay} & \mathrm{BM.63}:\ \textbf{closed unconditionally}\ —\ Q_\beta<-\ln3/\mathfrak h<0\ \text{at every admissible coupling};\ \text{the ordering inverts, the closure does not}\end{array}
\]

%dsp-level:0
\[
(2)\quad\textbf{the derived inventory:}
\]

%dsp-level:0
\[
\qquad |A_3|=3,\ |S_1\times S_1|=1,\ \text{asymmetry}\ 2;
\]

%dsp-level:0
\[
\qquad m_q=(2,1,1);
\]

%dsp-level:0
\[
\qquad W=(1,6,15,20,15,6,1),\ \textstyle\sum_QW=64;
\]

%dsp-level:0
\[
\Sigma_{\mathrm B}=\ln\tfrac43=2\ln2-\ln3;
\]

%dsp-level:0
\[
\qquad \Delta m_{\mathrm{transport}}=\frac{\ln(4/3)}{\mathfrak h};
\]

%dsp-level:0
\[
\qquad \Delta m=\frac{\ln(4/3)-6c\varepsilon_{\mathrm{EM}}\ln2}{\mathfrak h};
\]

%dsp-level:0
\[
\qquad \frac{m_n}{m_p}=\frac{\ln1280}{\ln960+6c\varepsilon_{\mathrm{EM}}\ln2};
\]

%dsp-level:0
\[
\varepsilon^{\mathrm{crit}}(c)=\frac{2-\log_23}{6c};
\]

%dsp-level:0
\[
\qquad \Delta_1=2\hat\mu'^2,\ \Delta_2=(\sqrt2-1)\hat\mu'^2,\ \Delta_3=(2\cos\tfrac\pi{10}-\sqrt3)\hat\mu'^2;
\]

%dsp-level:0
\[
\qquad 2E\in\{12,24,60\},\ 1:2:5;
\]

%dsp-level:0
\[
D_S\in\{2,\,4,\,2-\sqrt2,\,1,\,2-2\cos\tfrac\pi{10},\,2-\sqrt3\};
\]

%dsp-level:0
\[
\qquad V\delta_v=4\pi;
\]

%dsp-level:0
\[
\qquad Q_\beta<-\frac{\ln3}{\mathfrak h}
\]

%dsp-level:0
\[
\qquad [\text{machine: every row exact}];
\]

%dsp-level:0
\[
(3)\quad\textbf{the widened scope rows (recorded, not silent):}
\]

%dsp-level:0
\[
\qquad \mathrm{BM.52}(3)\text{'s sector census is extended by one row}\ —\ \text{the transport register, a}\ \sigma\text{-dependent functional of the seal sector's boundary data outside}\ \Delta S\text{'s alphabet}\ [\mathrm{BM.59}(8);\ \mathrm{BM.62}(11)];
\]

%dsp-level:0
\[
\ \text{its five zeros stand};
\]

%dsp-level:0
\[
\mathrm{BM.55}(1)\text{'s decomposition gains its sixth term};
\]

%dsp-level:0
\[
\qquad \mathrm{BM.55}(2)\text{'s ratio is amended to}\ (5)\ \text{of}\ \mathrm{BM.62};
\]

%dsp-level:0
\[
\qquad \mathrm{BM.55}(3)\text{'s strict order is scoped to the super-critical branch}\ [\mathrm{BM.62}(9)];
\]

%dsp-level:0
\[
\mathrm{BM.55}(4)\text{'s residue is filled}\ [\mathrm{BM.59}(1),(1\mathrm a)];
\]

%dsp-level:0
\[
\qquad \mathrm{BM.41}(4)(\mathrm e)\text{'s verdict is scoped to the reads}\ (\mathrm a)\text{--}(\mathrm d)\ \text{it enumerates};
\]

%dsp-level:0
\[
\qquad \textbf{no clause of}\ \mathrm{BM.51}\text{--}\mathrm{BM.56}\ \textbf{is retracted, and no clause of Parts 1--3 is modified};
\]

%dsp-level:0
\[
(4)\quad\textbf{the standing rules:}
\]

%dsp-level:0
\[
\qquad \neg\text{a P4 addition}\ (\text{§Cal-BM stays empty};\ e\ \text{is}\ \mathrm{EW.35}\text{'s logged combination};\ \hat\mu'^2,\mu\in\mathrm{Cal\text{-}F}\ \text{are logged});
\]

%dsp-level:0
\[
\qquad \neg\text{a new dimensionful scale}\ [\mathrm{P1}];
\]

%dsp-level:0
\[
\neg\text{a measured value and}\ \neg\text{a numerical coupling}\ [\mathrm{P5};\ \mathrm{EM.25}(7),(9);\ \text{machine: zero external tokens, zero decimal numerals in this tier}];
\]

%dsp-level:0
\[
\qquad \neg\exists\ \text{an}\ [\mathrm{EXT}]\ \text{datum in any derivation};
\]

%dsp-level:0
\[
\text{species names confined to the identification layers}\ [\mathrm{Def\text{-}BM.3}(7),\ \mathrm{Def\text{-}BM.6}(7),\ \mathrm{Def\text{-}BM.14}(8)];
\]

%dsp-level:0
\[
\qquad \text{topology labels are internal bookkeeping}\ [\mathrm{Def\text{-}H.32}(4)];
\]

%dsp-level:0
\[
\qquad \text{amplitude}\ne\text{mass}\ [\mathrm{FT.79}(5)];
\]

%dsp-level:0
\[
\qquad \text{rung}\ne\text{scale}\ [\mathrm{D0.8}];
\]

%dsp-level:0
\[
\neg\text{a HistLog disposition imported as premise}\ (\text{items}\ 7,\ 13,\ \mathrm{Q2},\ \text{deficits}\ 1',\ 4\ \text{are pointers});
\]

%dsp-level:0
\[
\qquad \text{the extra-axiomatic list}\ [\text{A §0.2a}]\ \text{gains no member};
\]

%dsp-level:0
\[
\textbf{the declarations this tier adds:}
\]

%dsp-level:0
\[
\qquad \text{exactly one name attachment}\ [\mathrm{Def\text{-}BM.14}(8):\ \text{the branch labels};\ \text{logged as the §II.4 addendum}\ \mathrm{A\text{-}BM.6}\ \text{before use},\ \text{as}\ \mathrm{Def\text{-}BM.3}(7)\ \text{requires}];
\]

%dsp-level:0
\[
\qquad \textbf{the constructions:}\ \mathrm{Def\text{-}BM.13}\ \text{and}\ \mathrm{Def\text{-}BM.14},
\]

%dsp-level:0
\[
\ \text{both abbreviations of derived objects with their non-derived clauses marked inline};
\]

%dsp-level:0
\[
(5)\quad\textbf{the adjudication:}
\]

%dsp-level:0
\[
\qquad \boxed{\ \text{the constituent transport splitting is derived, and it reverses the baryon order on an exact sub-critical branch}\ }
\]

%dsp-level:0
\[
\text{the two proposed mechanisms are refuted with their obstruction named}\ —\ \text{the vertex asymmetry is annihilated by class-blindness, seed colour-blindness and the trace law; the observation frame is annihilated by the charge exponent}\ 0;
\]

%dsp-level:0
\[
\text{the mechanism that survives is the one both directions select: the EM end-register census, whose logarithm is a class-dependent entropy in the ledger's own currency};
\]

%dsp-level:0
\[
\text{the six quark states are assembled from three floors and a branch on each, with both hierarchies derived exactly and the two currencies' inversion recorded};
\]

%dsp-level:0
\[
\qquad \alpha_B\ \text{is not derivable and is not added};
\]

%dsp-level:0
\[
\textbf{the neutron is heavier than the proton, and the beta decay stays closed, with an exact margin of one}\ \ln3.
\]

\textbf{Req.} P1, P4, P5 (Parts 1–2); D0.8, D2.9 (Parts 1–2); EM.9, EM.25(7), EM.25(9), EW.33, EW.35 (Part 2); Def-H.32(4), FT.79(5), FT.125(4) (Part 3); Def-BM.3(7), Def-BM.6(7), Def-BM.13, Def-BM.14, BM.41(4), BM.51, BM.52(3), BM.55(1)–(5), BM.57, BM.58, BM.59, BM.60, BM.61, BM.62, BM.63 (this part); the machine sections §A–§G green [math; machine: \texttt{checks/\allowbreak{}verify\_\allowbreak{}H\_\allowbreak{}constituent\_\allowbreak{}transport\_\allowbreak{}splitting.py} §G].

\textbf{Pf.}

%dsp-level:0
\[
(1)\text{--}(4):\ \text{assembly of the adjudicated rows of}\ \mathrm{BM.57}\text{--}\mathrm{BM.63}\ \text{and the two definitions}\ [\text{machine: §A--§F green}];
\]

%dsp-level:0
\[
\qquad \text{the scope-widening scan against}\ \mathrm{BM.41},\ \mathrm{BM.51}\text{--}\mathrm{BM.56}\ [\text{machine}]
\]

%dsp-level:0
\[
(5):\ \text{conjunction of}\ (1)\text{--}(4). \ \blacksquare
\]

\textbf{Phys.} The tier's verdict in one sentence: both proposed mechanisms die exactly, and the one that survives was selected independently by the force structure and by the mass spectrum before either was told about the other. The vertex asymmetry is real — three continuations against one, an asymmetry of two by either counting convention — and it splits nothing, because the two baryon rows are the same vertex, because the transported seed is colour-blind, and because a sealed singlet reads only traces. The observation frame is real too — colour projects, electromagnetism censuses — and it splits nothing geometrically, because not one of the seven locked monomials carries a charge exponent; what it does produce is the scar's field energy, which points the wrong way. Between them they exhaust the force side and leave a single survivor in the inventory: the count of admissible end-sign registers, twenty against fifteen. Read as a logarithm against the vacuum, in the ledger's own nats, through the rest slot the field energy already uses, that census is an energy, and it is positive for the neutral class. The neutron becomes the heavier baryon on an exact branch of the coupling, bounded by four integer power comparisons, and the beta decay stays shut anyway — because the transport excess is a fraction of a bit and the cheapest daughter costs two. Six quark states come out of the same pass: three Platonic mass floors climbing one to two to five, each split into a lighter and a heavier branch by the conical defect at its own trivalent vertex, with the heavier branch heavier for the same exact reason in all three generations. And $\alpha_B$ is refused: not because the arithmetic is hard, but because there is no strong coupling to take as an argument, no count-to-strength theorem to build one from, and one exact charge-class number whose derivative in every coupling the corpus owns is zero.

\medskip
\subsection*{\texorpdfstring{B.1.24\quad BM.72 — The status of the organic bridge: adjudication and scope}{BM.72 B.1.24}}
\addcontentsline{toc}{subsection}{B.1.24 BM.72}
\phantomsection\label{tag:BM-72}
{\small \textit{Record.}\quad Source: Part 4 (Part4\_Baryomesonic\_era.md); tier: Tier BM.15 — The Organic Bridge: the Operator EM Lift, the Derived-Holonomy Attempt, the Constrained Joint Entropy, and the Beta-Decay Energy Audit; canonical tag \texttt{BM.72}. Reproduced from the revision-1 canon; a cross-reference stub marks the former locus in the main corpus.}
\medskip
\textbf{Thm.}

%dsp-level:0
\[
(1)\quad\textbf{the answers:}\qquad\begin{array}{l|l}\text{the operator EM lift} & \mathrm{BM.66}:\ \textbf{derived}\ —\ \operatorname{Tr}\widehat F=2c_0\ \text{charge-blind};\ \ I_F=4\widehat Q^2;\ \ \tfrac{E_{\mathrm{int}}}4c\,\varepsilon_{\mathrm{EM}}\ln2\,I_F\ \text{reproduces}\ \mathrm{Def\text{-}BM.13}(5)\ \text{exactly on}\ Q\in\{0,\pm1\};\\ & \text{three routes to one monomial};\ \neg\text{a new parameter};\ c\ \text{priced};\ \mathrm{BM.58}(6)\text{'s uniqueness row stands}\\\text{the holonomy lift} & \mathrm{BM.67}:\ \textbf{obstructed}\ (\text{the obstruction outcome}),\ \textbf{with the partial lift derived}\ —\ \text{the}\ \mathbb Z_2\ \text{Wilson sector},\ \text{the parity lock}\ \sum\aleph\equiv|Q|\ (2),\ \text{the integer winding/sealing sectors};\\ & \text{the force census angle-free};\ \text{the two-torsion image theorem}\ (2\ \text{of}\ 6\ \text{states reachable});\ \text{the calibration-weight no-go re-cited};\ \text{the structured record issued};\ \neg\text{a declared angle}\\\text{the joint entropy} & \mathrm{BM.68}:\ \textbf{derived as a fiber cardinality}\ —\ \Omega=\Omega_{\mathrm{count}}\Omega_{\mathrm{charge}}\Omega_{\mathrm{geom}};\ \ 2^E\ \text{and}\ W(Q)\ \text{exact};\ \Omega_{\mathrm{geom}}^{\mathrm{stat}}=1\ (S_{\mathrm{geom}}=0\ \text{derived},\ \neg\text{declared});\\ & \text{the general}\ \Sigma_{12}\ \text{fiber is the named missing enumeration};\ \text{all five bridge species tested and refuted or landed at item 7}\\\text{the bridge} & \mathrm{BM.69}:\ \textbf{partially derived}\ (\text{the partial-bridge outcome})\ —\ \text{the identity holds up to the one undeclared dimensionless ratio}\ \rho_{\mathrm B}=\hat\mu'^2\mathfrak h^2/(\ln2)^2\ (\mathrm{Cal\text{-}F},\ \neg\text{added});\\ & \text{the non-derivation theorem on four grounds};\ \text{Gap 7 reduced to}\ (\mathrm R1)\text{--}(\mathrm R3)\\\text{the beta decay} & \mathrm{BM.70},\ \mathrm{BM.71}:\ \textbf{no-go}\ (\text{the no-go outcome})\ —\ F_{\ell\nu}\ge2\ln2\ \text{theorem on the informational species};\ \text{the formed-closed dichotomy};\\ & \text{the margin ladder}\ (\ln3,\ \ln12,\ \ln24);\ \text{the five-route census exhausted: three barred, one}\ \neg\text{constructed and sign-barred at the constructed order, one specification act not performed};\ \neg\text{an external}\ Q\\\text{the neutron row} & \mathrm{BM.71}(3),(4):\ \ \Delta m_{\mathrm{bridge}}=0\ \text{in currency};\ \text{the sign branch-conditional};\ \text{the decay closed};\ \neg\text{a tuning}\end{array}
\]

%dsp-level:0
\[
(2)\quad\textbf{the status adjudication:}
\]

%dsp-level:0
\[
\qquad \boxed{\ \textbf{the partial-bridge outcome and the no-go outcome are conjoined}\ }
\]

%dsp-level:0
\[
\textbf{the full-bridge outcome is}\ \neg\textbf{achieved:}\ \ \text{the exact operator identity relating}\ m_{\mathrm{pole}},
\]

%dsp-level:0
\[
\ m_{\mathrm{hid}},
\]

%dsp-level:0
\[
\ W(Q),
\]

%dsp-level:0
\[
\ \Theta_R,
\]

%dsp-level:0
\[
\ \operatorname{Tr}\widehat F^2\ \text{requires the relating equation};
\]

%dsp-level:0
\[
\ \ \text{its obstruction is exact}\ \bigl(\mathrm{BM.69}(4)\text{'s four grounds}\bigr)\ \text{and its residue is minimal and priced}\ \bigl(\mathrm{BM.69}(5)\bigr);
\]

%dsp-level:0
\[
\textbf{the opened-decay outcome is}\ \neg\textbf{achieved:}\ \ \neg\exists\ \text{a derived ledger effect opening the decay}\ —\ \text{the route census is complete and every route carries an exact bar}\ \bigl(\mathrm{BM.71}(1),(2)\bigr);
\]

%dsp-level:0
\[
\textbf{the obstruction-theorem outcome is realized inside the tier:}
\]

%dsp-level:0
\[
\qquad \mathrm{BM.67}(6),(8),(10)\ \ \text{and}\ \ \mathrm{BM.69}(4)\ \ \text{are derived obstruction theorems,}\ \neg\text{rhetorical records};
\]

%dsp-level:0
\[
(3)\quad\textbf{the derived inventory (the new exact vocabulary):}
\]

%dsp-level:0
\[
\qquad \operatorname{Tr}\widehat F=2c_0;\quad I_F=4\widehat Q^2;\quad\tfrac32=\tfrac{E_{\mathrm{int}}}4;\quad U(\gamma)=(-1)^{N_\pi(\gamma)};\quad\operatorname{im}(\text{lift})\subseteq\{0,\pi\};
\]

%dsp-level:0
\[
\text{the rung law}\ \ m\,\delta_v\in\pi\mathbb Z\iff D_S(m)\in\{0,2,4\};
\]

%dsp-level:0
\[
\qquad \textstyle\sum\aleph\equiv|Q|\ (2);
\]

%dsp-level:0
\[
\qquad \Omega_{\mathrm{geom}}^{\mathrm{stat}}=1;
\]

%dsp-level:0
\[
\qquad \text{the strict-fiber counts}\ \binom{K-1}2\ \text{with orbits}\ (1,2,4,6)\ \text{at}\ K=(3,4,5,6);
\]

%dsp-level:0
\[
F_{\ell\nu}\ge2\ln2;
\]

%dsp-level:0
\[
\qquad \text{the margin ladder}\ -(\ln3,\ \ln12,\ \ln24)\ \ \text{on the certificates}\ \ 4<12<48<96;
\]

%dsp-level:0
\[
\qquad \text{the floor ladder}\ (2,\,4,\,5)\ln2;
\]

%dsp-level:0
\[
\qquad \rho_{\mathrm B}=\tfrac{\hat\mu'^2\mathfrak h^2}{(\ln2)^2}\ \ \text{named as the residue}
\]

%dsp-level:0
\[
\qquad [\text{machine: every row exact}];
\]

%dsp-level:0
\[
(4)\quad\textbf{the scope rows (recorded,}\ \neg\textbf{silent):}
\]

%dsp-level:0
\[
\qquad \mathrm{BM.63}(2)\text{'s closure is audited and stands, with its species conditionality on its face}\ [\mathrm{BM.70}(4)];
\]

%dsp-level:0
\[
\mathrm{BM.58}(6)\text{'s uniqueness table gains no row}\ [\mathrm{BM.66}(7)];
\]

%dsp-level:0
\[
\qquad \mathrm{BM.62}(1)\text{'s decomposition gains no term}\ —\ \Delta m_{\mathrm{bridge}}=0\ \text{in currency}\ [\mathrm{BM.71}(3)];
\]

%dsp-level:0
\[
\qquad \mathrm{BM.52}(3)\text{'s sector census is refined at the operator level, its five zeros standing}\ [\mathrm{BM.66}(2),(6)];
\]

%dsp-level:0
\[
\mathrm{BM.53}(5)\text{'s MT row is re-cited as a resolvability bound,}\ \neg\text{a budget}\ [\mathrm{BM.71}(1)\ \text{row}\ 1,\ (2)];
\]

%dsp-level:0
\[
\textbf{no clause of}\ \mathrm{BM.57}\text{--}\mathrm{BM.64}\ \textbf{is retracted or reopened; no clause of Parts 1--3 is modified};
\]

%dsp-level:0
\[
\qquad \text{the conditionalities}\ (\mathrm{Def\text{-}H.37}(2),\ \mathrm{Def\text{-}H.38},\ \mathrm{Def\text{-}H.43},\ \mathrm{M1}/\mathrm{M2},\ \mathrm{CT2.6},\ \mathrm{BM.8}(2))\ \text{are inherited at the points of use};
\]

%dsp-level:0
\[
(5)\quad\textbf{the standing rules:}
\]

%dsp-level:0
\[
\qquad \neg\text{a P4 addition}\ —\ \text{§Cal-BM stays empty};
\]

%dsp-level:0
\[
\ \ \rho_{\mathrm B}\ \text{is a Cal-F packaging restated as a residue,}\ \neg\text{a new member};
\]

%dsp-level:0
\[
\qquad \neg\text{a new dimensionful scale}\ [\mathrm{P1}];
\]

%dsp-level:0
\[
\neg\text{a measured value and}\ \neg\text{a numerical coupling}\ [\mathrm{P5};\ \mathrm{EM.25}(7),(9);\ \text{machine: zero external tokens, zero decimal numerals in this tier}];
\]

%dsp-level:0
\[
\qquad \neg\exists\ \text{an}\ [\mathrm{EXT}]\ \text{datum in any derivation};
\]

%dsp-level:0
\[
\text{species names confined to the identification layers}\ [\mathrm{Def\text{-}BM.3}(7),\ \mathrm{Def\text{-}BM.6}(7),\ \mathrm{Def\text{-}BM.14}(8)];
\]

%dsp-level:0
\[
\qquad \text{topology labels are internal bookkeeping}\ [\mathrm{Def\text{-}H.32}(4)];
\]

%dsp-level:0
\[
\qquad \text{amplitude}\ne\text{mass}\ [\mathrm{FT.79}(5)];
\]

%dsp-level:0
\[
\qquad \text{rung}\ne\text{scale}\ [\mathrm{D0.8}];
\]


%dsp-level:0
\[
\neg\text{a HistLog disposition imported as premise}\ (\text{items}\ 7,\ 13,\ 19,\ \mathrm{Q2},\ \text{deficits}\ 1',\ 4\ \text{are pointers});
\]
%dsp-level:0
\[
\qquad \text{the extra-axiomatic list}\ [\text{A §0.2a}]\ \text{gains no member};
\]
%dsp-level:0
\[
\qquad \neg\text{a Yukawa-type object exists or is added}\ [\mathrm{EW.15};\ \mathrm{FT.145}(3)];
\]


%dsp-level:0
\[
\textbf{the declarations this tier adds: zero}\ —\ \neg\text{a new Def is issued}: I_F,
\]

%dsp-level:0
\[
\ F_{\ell\nu},
\]

%dsp-level:0
\[
\ \Omega,
\]

%dsp-level:0
\[
\ S_{\mathrm{joint}},
\]

%dsp-level:0
\[
\ \rho_{\mathrm B},
\]

%dsp-level:0
\[
\ \Phi(c),
\]

%dsp-level:0
\[
\ U(c)\ \text{are in-theorem abbreviations of derived objects};
\]

%dsp-level:0
\[
\alpha_B:\quad\neg\text{added},\quad\neg\text{derived},\quad\text{barred}\ —\ \mathrm{BM.61}(8)\text{'s}\ (\mathrm O1)\text{--}(\mathrm O4)\ \text{stand}
\]

%dsp-level:0
\[
\qquad [\text{machine: the token scan of this tier}\ —\ \text{every occurrence a barred-object row}];
\]

%dsp-level:0
\[
(6)\quad\textbf{the adjudication:}
\]

%dsp-level:0
\[
\qquad \boxed{\ \text{the organic bridge is derived up to one named scale ratio and two named attachments; the EM correction is an operator identity; the holonomy lift is derived as a}\ \mathbb Z_2\ \text{sector and obstructed exactly; the joint entropy is a fiber cardinality with a stationary geometric factor of one; the beta vertex is closed on every formed read}\ }
\]

%dsp-level:0
\[
\text{its opening is a specification act, which the corpus prices, audits, and}\ \neg\text{performs};
\]

%dsp-level:0
\[
\qquad \text{the ledger's open items stay open with sharper names}\ —\ (\mathrm R1)\text{--}(\mathrm R3),
\]

%dsp-level:0
\[
\ \text{the regime-scoped fiber, the lepton floors, the numerical value of}\ e.
\]

\textbf{Req.} P1, P4, P5 (Parts 1–2); D0.8 (Parts 1–2); CT2.6 (Part 2); EM.25(7), EM.25(9), EW.15 (Part 2); Def-H.20, Def-H.32(4), Def-H.37(2), Def-H.38, Def-H.43, FT.55(1), FT.73(2), FT.73(3), FT.79(5), FT.125(4), FT.145(3) (Part 3); Census FT.35(2), FT.63(6) (row-wise); Def-BM.3(7), Def-BM.6(7), Def-BM.12(4), Def-BM.13(5), Def-BM.14(5), Def-BM.14(8), BM.8(2), BM.36(3), BM.43(2), BM.52(3), BM.53(5), BM.57(8), BM.58(4), BM.58(6), BM.59(9), BM.61(8), BM.62(1), BM.62(4), BM.63(2), BM.64(3), BM.65, BM.66, BM.67, BM.68, BM.69, BM.70, BM.71 (this part); A §E (the P4 core); Cal-F; the machine sections §A–§G green [math; machine: \texttt{checks/\allowbreak{}verify\_\allowbreak{}H\_\allowbreak{}organic\_\allowbreak{}bridge\_\allowbreak{}post\_\allowbreak{}bm14.py} §G].

\textbf{Pf.}

%dsp-level:0
\[
(1)\text{--}(3):\ \text{assembly of the adjudicated rows of}\ \mathrm{BM.65}\text{--}\mathrm{BM.71}\ [\text{machine: §A--§F green}];
\]

%dsp-level:0
\[
\qquad \text{the scope-row scan against}\ \mathrm{BM.52},\ \mathrm{BM.53},\ \mathrm{BM.57}\text{--}\mathrm{BM.64}\ [\text{machine}]
\]

%dsp-level:0
\[
(4)\text{--}(5):\ \text{the standing-rule scans: external tokens, decimal numerals, unit tokens, registry tokens, declaration lines}\ [\text{machine}];
\]

%dsp-level:0
\[
\qquad \text{the}\ \alpha_B\ \text{token scan with the barred-row classification}\ [\text{machine}];
\]

%dsp-level:0
\[
\qquad \text{the symbol census against the P4 core,}\ \mathrm{Cal\text{-}2},\ \mathrm{Cal\text{-}F},\ \text{§Cal-BM}\ [\text{machine}]
\]

%dsp-level:0
\[
(6):\ \text{conjunction of}\ (1)\text{--}(5). \ \blacksquare
\]

\textbf{Phys.} The tier's verdict is two outcomes conjoined. The bridge is partially derived: the EM correction is an operator identity with a charge-blind linear trace and a four-Q-squared invariant, the joint entropy is a fiber cardinality whose stationary geometric factor is exactly one, and the pole-to-hidden ratio is exact up to a single named dimensionless packaging — Gap 7 at its smallest, with a non-derivability theorem behind the remainder rather than a silence. The full identity is not derived, and the tier says so with four independent proofs of why. The beta decay is closed: the daughter floor is a theorem wherever the Q-value is formed, the margins are exact integers, the five candidate opening terms are barred or unconstructed one by one, and the sole surviving door is the specification act the corpus declines to perform. The holonomy is not declared: the binary Wilson sector, the parity lock and the integer winding sectors are derived, the real lift is obstructed by a two-element image against six angles and by calibration weight against the winding route, and the obstruction stands on record in structured form with its missing enumerations named. Nothing is added to the declaration surface — no definition, no parameter, no scale, no measured value — and the open items of the ledger stay open, with sharper names than they carried in.

\medskip
\subsection*{\texorpdfstring{B.1.25\quad BM.78 — The status of the chiral assignment: adjudication and scope}{BM.78 B.1.25}}
\addcontentsline{toc}{subsection}{B.1.25 BM.78}
\phantomsection\label{tag:BM-78}
{\small \textit{Record.}\quad Source: Part 4 (Part4\_Baryomesonic\_era.md); tier: Tier BM.16 — The Chiral Assignment: the Closed Deficit, the Ray-Resolved Spin Census, the Attached Quark Splitting, the Chiral Vertex Dynamics, and the Specified Beta-Decay Opening; canonical tag \texttt{BM.78}. Reproduced from the revision-1 canon; a cross-reference stub marks the former locus in the main corpus.}
\medskip
\textbf{Thm.}

%dsp-level:0
\[
(1)\quad\textbf{the answers:}\qquad\begin{array}{l|l}\text{the chiral assignment} & \mathrm{BM.73}:\ \textbf{issued and logged}\ (\text{A-BM.7})\ —\ \text{one naming bit on the derived forced shape}\ [\mathrm{FT.144}(5)]:\ \text{the arrow-occupied ray is the left ray}\ (\hat\chi=+1),\ \text{weak-active label attached by consistency},\ \text{the right ray the}\ \hat\omega\text{-image};\\ & \textbf{deficit}\ 1'/\mathrm{Q1}'\ \textbf{closed at the assignment level};\ \text{the occupancy}\ (1,0)\ \text{named};\ \text{the}\ \omega/\chi\ \text{matching rows re-certified};\ \text{the four gradings stay pairwise distinct}\\\text{the spin} & \mathrm{BM.74}:\ \textbf{the ray-resolved census derived}\ —\ j=|k|/2\ \text{rows inherited from}\ \mathrm{Def\text{-}H.38}\ (\text{item 19 untouched});\ \text{the vertex table with opposite rays for the two emitted members};\\ & \text{the angular-momentum budget unchanged};\ \neg\text{a helicity row},\ \neg\text{a spin--statistics row}\ [\mathrm{FT.125}(5)\ \text{stands}]\\\text{the quark mass splitting} & \mathrm{BM.75}:\ \textbf{attached by the unique bijection}\ —\ \text{weight}\ 0\leftrightarrow q_{\mathrm{core}}\,0\leftrightarrow\text{down-type}\leftrightarrow\beta{=}{+};\ \ \text{weight}\ {+}1\leftrightarrow q_{\mathrm{core}}\,{+}1\leftrightarrow\text{up-type}\leftrightarrow\beta{=}{-};\\ & \Delta_g=(2,\ \sqrt2{-}1,\ 2\cos\tfrac\pi{10}{-}\sqrt3)\,\hat\mu'^2>0\ \text{per generation};\ \textbf{the sign-coherence theorem: the down-type constituent is heavier in every formed currency};\ \text{item 7 untouched}\\\text{the dynamics} & \mathrm{BM.76}:\ \textbf{constructed at the vertex level}\ —\ \widehat V=E_+\otimes P_L,\ \ \widehat V^2=0\ (\text{the alphabet overflow as an operator identity}),\ \operatorname{rk}\widehat V=1,\ \operatorname{im}\widehat V=\text{the up-type weight on the occupied ray},\ [\,W_3\otimes I,\widehat V\,]=\widehat V;\\ & \text{the}\ \omega\text{- and }\hat\kappa\text{-conjugates are the image vertex};\ \text{the flip, rate and arrow rows frozen};\ \text{the field-level dynamics of}\ Q\ \text{stays}\ \neg\text{constructed}\ [\mathrm{BM.6}(5);\ \mathrm{Def\text{-}BM.12}(7)]\\\text{the beta decay} & \mathrm{BM.77}:\ \textbf{specified}\ (\text{A-BM.8})\ \textbf{and allowed on an exact non-empty region}\ —\ Q_\beta^{\mathrm{spec}}\mathfrak h=\Delta T(S)-\rho_\ell\sqrt{G_\ell}-\rho_\nu\sqrt{G_\nu};\ \ \Delta T=5\,m^2_B(\cdot,S);\ \ \rho^{\mathrm{crit}}\in(\tfrac{273}{23},\tfrac{42197850}{791648})\ \text{at the witness couplings};\\ & \text{the trichotomy with integer-certified witnesses at every frozen cell};\ \text{the conditional spectrum becomes branch content};\ \text{the rescope rows recorded},\ \neg\text{retracted}\end{array}
\]

%dsp-level:0
\[
(2)\quad\textbf{the status adjudication:}
\]

%dsp-level:0
\[
\qquad \boxed{\ \textbf{deficit}\ 1'\ \textbf{is closed by the issued assignment, and the beta vertex is kinematically allowed on the specified read's exact opening region}\ }
\]

%dsp-level:0
\[
\textbf{the closures are level-exact:}
\]

%dsp-level:0
\[
\qquad 1'\ \text{closes at the assignment level}\ —\ \text{the coupling rule exists};
\]

%dsp-level:0
\[
\ \text{the frequency read}\ [\mathrm{FT.131}(3)],
\]

%dsp-level:0
\[
\ \text{the anomaly formation}\ [\mathrm{FT.125}(2):\ \text{formable on the}\ 1'\ \text{side},\ \text{the}\ +\tfrac23\ \text{side structurally unreachable}\ [\mathrm{BM.28}(7)]],
\]

%dsp-level:0
\[
\ \text{and the mixing object}\ \ \text{stay open at their own loci};
\]

%dsp-level:0
\[
\textbf{the decay row is read-relative and complete:}
\]

%dsp-level:0
\[
\qquad \text{closed on every pre-specification formed read}\ [\mathrm{BM.70}(4);\ \mathrm{BM.71}(2)\ —\ \text{standing}],
\]

%dsp-level:0
\[
\qquad \text{allowed on}\ \mathcal O_S,\ \text{closed on}\ \mathcal C_S,\ \text{degenerate on}\ \partial\mathcal O_S\ \ \text{on the specified read}\ [\mathrm{BM.77}(5),(6)];
\]

%dsp-level:0
\[
\neg\text{a frozen refutation is reopened};\qquad
\]

%dsp-level:0
\[
(3)\quad\textbf{the derived inventory (the new exact vocabulary):}
\]

%dsp-level:0
\[
\qquad \widehat V=E_+\otimes P_L;\quad\widehat V^2=0;\quad\operatorname{rk}\widehat V=1;\quad\bigl[W_3\otimes I_\omega,\widehat V\bigr]=\widehat V;
\]

%dsp-level:0
\[
\quad (I\otimes\hat\omega)\widehat V(I\otimes\hat\omega)=E_+\otimes P_R;
\]

%dsp-level:0
\[
\quad (I\otimes\hat\kappa)\widehat V(I\otimes\hat\kappa)=E_+\otimes P_R;
\]

%dsp-level:0
\[
\text{the occupancy}\ (1,0);
\]

%dsp-level:0
\[
\qquad \text{the attachment bijection}\ \{0\leftrightarrow\beta{=}{+},\ {+}1\leftrightarrow\beta{=}{-}\};
\]

%dsp-level:0
\[
\qquad \Delta T(S)=5\,m^2_B(\cdot,S)=\Bigl(\tfrac{4219785}{98956},\tfrac{9608310}{393679},\tfrac{455}{23},\tfrac{68565}{3088}\Bigr);
\]

%dsp-level:0
\[
\text{the universal witnesses}\ \rho=\tfrac{\Delta T}{4\sqrt{A_0+B_0}}\in\mathcal O_S,\quad\rho=\tfrac{\Delta T}{\sqrt{A_0-B_0}}\in\mathcal C_S;
\]

%dsp-level:0
\[
\qquad \text{the certificates}\ 11<16,
\]

%dsp-level:0
\[
\quad 828100>5819,
\]

%dsp-level:0
\[
\quad 4219785<4255108,
\]

%dsp-level:0
\[
\quad 7500>1849,
\]

%dsp-level:0
\[
\quad 99<100,
\]

%dsp-level:0
\[
\quad 75>64,
\]

%dsp-level:0
\[
\quad 273>253,
\]

%dsp-level:0
\[
\quad 42197850<42748992,
\]

%dsp-level:0
\[
\quad 3312400>5819;
\]

%dsp-level:0
\[
A_0\pm B_0\big|_{g=g'=1}=\tfrac{11}{16},\ \tfrac3{16}\qquad[\text{machine: every row exact}];
\]

%dsp-level:0
\[
(4)\quad\textbf{the scope rows (recorded,}\ \neg\textbf{silent):}
\]

%dsp-level:0
\[
\qquad \mathrm{BM.41}(3)\text{'s forced-neutral-member row stands on its equal-mass locus and is read-relative beyond it}\ [\mathrm{BM.77}(6)];
\]

%dsp-level:0
\[
\qquad \mathrm{BM.41}(4)(\mathrm e),\ \mathrm{BM.49},\ \mathrm{BM.63},\ \mathrm{BM.70}(4),\ \mathrm{BM.71}(2)\ \text{rescoped,}\ \neg\text{retracted}\ [\mathrm{BM.77}(7)];
\]

%dsp-level:0
\[
\mathrm{BM.52},\ \mathrm{BM.57},\ \mathrm{BM.58},\ \mathrm{BM.61},\ \mathrm{BM.66}\text{--}\mathrm{BM.69}\ \text{untouched};
\]

%dsp-level:0
\[
\qquad \text{A-BM.5's antineutrino row is ray-resolved}\ [\mathrm{BM.74}(4)];
\]

%dsp-level:0
\[
\qquad \mathrm{Def\text{-}BM.14}(8)/\text{A-BM.6's}\ \neg\text{declared attachment bar is discharged by derivation}\ [\mathrm{BM.75}(3)];
\]

%dsp-level:0
\[
\text{the conditionalities}\ (\mathrm{Def\text{-}H.38},\ \mathrm{Def\text{-}H.37},\ \mathrm{Def\text{-}H.43},\ \mathrm{M1}/\mathrm{M2},\ \mathrm{CT2.6},\ \mathrm{BM.8}(2))\ \text{are inherited at the points of use}\ [\text{A §0.2a}];
\]

%dsp-level:0
\[
(5)\quad\textbf{the standing rules:}
\]

%dsp-level:0
\[
\qquad \textbf{the declarations this tier adds: exactly two, both logged before use}\ —\ \text{A-BM.7}\ (\text{the chiral assignment}),
\]

%dsp-level:0
\[
\ \text{A-BM.8}\ (\text{the vertex read specification});
\]

%dsp-level:0
\[
\qquad \neg\text{a P4 addition}\ —\ \text{§Cal-BM stays empty};\ \ \rho_\ell,\ \rho_\nu\ \text{are restatements of Cal-F/P4 data};
\]

%dsp-level:0
\[
\neg\text{a new dimensionful scale}\ [\mathrm{P1}];
\]

%dsp-level:0
\[
\qquad \neg\text{a measured value and}\ \neg\text{a numerical coupling}\ [\mathrm{P5};\ \mathrm{EM.25}(7),(9);\ \text{machine: zero external tokens, zero decimal numerals in this tier}];
\]

%dsp-level:0
\[
\qquad \neg\exists\ \text{an}\ [\mathrm{EXT}]\ \text{datum in any derivation};
\]

%dsp-level:0
\[
\text{species names confined to the identification layers}\ [\mathrm{Def\text{-}BM.3}(7),\ \mathrm{Def\text{-}BM.6}(7),\ \mathrm{Def\text{-}BM.14}(8);\ \text{A-BM.7's bar extends the confinement to the left/right matter labels}];
\]

%dsp-level:0
\[
\qquad \text{topology labels are internal bookkeeping}\ [\mathrm{Def\text{-}H.32}(4)];
\]

%dsp-level:0
\[
\qquad \text{amplitude}\ne\text{mass}\ [\mathrm{FT.79}(5)];
\]

%dsp-level:0
\[
\qquad \text{rung}\ne\text{scale}\ [\mathrm{D0.8}];
\]

%dsp-level:0
\[
\neg\text{a HistLog disposition imported as premise}\ (\text{items}\ 7,\ 13,\ 19,\ 21,\ \mathrm{Q1}'\text{--}\mathrm{Q6},\ \text{deficits}\ 1',\ 4\ \text{are pointers or issued here});
\]

%dsp-level:0
\[
\qquad \text{the extra-axiomatic list}\ [\text{A §0.2a}]\ \text{gains no member};
\]

%dsp-level:0
\[
\qquad \neg\text{a Yukawa-type object exists or is added}\ [\mathrm{EW.15};\ \mathrm{FT.145}(3)];
\]

%dsp-level:0
\[
\alpha_B:\quad\neg\text{added},\quad\neg\text{derived},\quad\text{barred}\ —\ \mathrm{BM.61}(8)\text{'s}\ (\mathrm O1)\text{--}(\mathrm O4)\ \text{stand}
\]

%dsp-level:0
\[
\qquad [\text{machine: the token scan of this tier}\ —\ \text{every occurrence a barred-object row}];
\]

%dsp-level:0
\[
(6)\quad\textbf{the adjudication:}
\]

%dsp-level:0
\[
\qquad \boxed{\ \text{the left--right assignment is issued on the derived forced shape and closes deficit}\ 1';\ \ \text{the spin census is ray-resolved with every label row inherited};\ \ \text{the quark mass splitting is attached by a unique bijection and sign-coherent across the formed currencies};\ \ \text{the vertex dynamics is a nilpotent chiral step with frozen rates};\ \ \text{and the beta decay is allowed on an exact, non-empty, integer-certified region of the undeclared scale data}\ —\ \text{closed on its complement, degenerate on its exact critical locus}\ }
\]

%dsp-level:0
\[
\text{the corpus walked the one door it had left open, and every other door stayed shut}:
\]

%dsp-level:0
\[
\qquad \text{the mass-form bridge is application-layer};
\]

%dsp-level:0
\[
\ \ \text{the composition is application-layer content};
\]

%dsp-level:0
\[
\ \ \text{the scales stand undeclared};
\]

%dsp-level:0
\[
\ \ \text{and nothing measured entered anywhere}.
\]

\textbf{Req.} P1, P4, P5 (Parts 1–2); D0.8 (Parts 1–2); CT2.6 (Part 2); EM.25(7), EM.25(9), EW.15 (Part 2); Def-H.20, Def-H.32(4), Def-H.37, Def-H.38, Def-H.43, FT.79(5), FT.125(2), FT.125(5), FT.131(3), FT.144(5), FT.145(3) (Part 3); Def-BM.3(7), Def-BM.6(7), Def-BM.12(7), Def-BM.14(8), BM.6(5), BM.8(2), BM.28(7), BM.41(3), BM.41(4), BM.49, BM.52, BM.57, BM.58, BM.61(8), BM.63, BM.66(6), BM.67, BM.68, BM.69, BM.70(4), BM.71(1), BM.71(2), BM.73, BM.74, BM.75, BM.76, BM.77 (this part); A §E (the P4 core); Cal-F; the machine sections §A–§G green [math; machine: \texttt{checks/\allowbreak{}verify\_\allowbreak{}H\_\allowbreak{}chiral\_\allowbreak{}assignment\_\allowbreak{}opening.py} §G].

\textbf{Pf.}

%dsp-level:0
\[
(1)\text{--}(3):\ \text{assembly of the adjudicated rows of}\ \mathrm{BM.73}\text{--}\mathrm{BM.77}\ [\text{machine: §A--§F green}];
\]

%dsp-level:0
\[
\qquad \text{the inventory rows re-listed from their own tags}\ [\text{machine}]
\]

%dsp-level:0
\[
(4):\ \text{the rescope census against}\ \mathrm{BM.41}\text{--}\mathrm{BM.72}\ \text{clause by clause}\ [\text{machine}];
\]

%dsp-level:0
\[
\qquad \text{the locus rows re-cited}
\]

%dsp-level:0
\[
(5):\ \text{the standing-rule scans: external tokens, decimal numerals, unit tokens, registry tokens, declaration lines}\ [\text{machine}];
\]

%dsp-level:0
\[
\qquad \text{the}\ \alpha_B\ \text{token scan with the barred-row classification}\ [\text{machine}];
\]

%dsp-level:0
\[
\qquad \text{the symbol census against the P4 core}\ [\text{A §E}],\ \mathrm{Cal\text{-}2},\ \mathrm{Cal\text{-}F},\ \text{§Cal-BM}\ [\text{machine}]
\]

%dsp-level:0
\[
(6):\ \text{conjunction of}\ (1)\text{--}(5). \ \blacksquare
\]

\textbf{Phys.} The tier's verdict in one sentence: the one bit the corpus could not derive is issued on the record, and everything the ledger had waiting on that bit comes out exact. The assignment names the occupied ray left, matches it to the weak-active branch by a consistency census rather than a preference, and leaves the orientation involution to do what it always did — exchange the rays, reverse the winding, conjugate the vertex into its image. Spin resolves by ray without a single label invented: the winding rule stays a declared interface, the baryon card keeps its open rung, and helicity stays outside. The quark splitting stops floating: one of two possible gluings survives four frozen consistency rows, and the down-type constituent comes out heavier in the defect currency, the transport-register currency and the ensemble currency at once, with the electromagnetic correction the lone dissenter, dominated below its exact critical strength. The vertex becomes an operator — nilpotent, rank one, firing from the occupied ray alone — and the decay itself becomes allowed where the corpus said a specification act was the only door: on the specified read the Q-value is an exact rational minus two scale-weighted radicals, and the opening region, the closing region and the critical locus between them are all certified by integers. What did not move is named as loudly as what did: no currency was converted, no scale was fixed, no measured number entered, and every earlier verdict stands in its own species, rescoped in the open rather than quietly rewritten.

\medskip
\subsection*{\texorpdfstring{B.1.26\quad BM.86 — The status of the mass ratio census: adjudication and scope}{BM.86 B.1.26}}
\addcontentsline{toc}{subsection}{B.1.26 BM.86}
\phantomsection\label{tag:BM-86}
{\small \textit{Record.}\quad Source: Part 4 (Part4\_Baryomesonic\_era.md); tier: Tier BM.17 — The Mass Ratio Census: the Platonic Winding Ladders, the Intra- and Inter-Generational Ratios, the Six-State Bridge Census, the Four-Currency Incommensurability, and the Neutron–Proton Ratio; canonical tag \texttt{BM.86}. Reproduced from the revision-1 canon; a cross-reference stub marks the former locus in the main corpus.}
\medskip
\textbf{Thm.}

%dsp-level:0
\[
(1)\quad\textbf{the answers:}\qquad\begin{array}{l|l}\text{the quark mass ratios} & \mathrm{BM.79}\text{--}\mathrm{BM.81}:\ \textbf{derived}\ —\ \text{the ladder law}\ D(m)=4\sin^2(m\delta_v/4)=2-2T_{|m|}(\cos\tfrac{\delta_v}2)\ \text{on all five solids at both priced windings};\\ & \text{the Chebyshev tower closing at}\ T_5(\mathsf c)=0;\ \text{the intra-generational ratios as exact sine ratios}\ (1,\ 1,\ 2,\ 1{+}\sqrt3,\ \sqrt{2/(2-2\mathsf c)});\ \text{the six-state doublet ratios}\ (\sqrt2,\ \sqrt{(2{+}\sqrt2)/2},\ \sqrt{(2{-}\sqrt3)/(2{-}2\mathsf c)});\\ & \text{the splitting order}\ \Delta_1>\Delta_2>\Delta_3;\ \text{the two inverted hierarchies};\ \text{the three ladders pairwise non-proportional}\\\text{the inter-generational ratios} & \mathrm{BM.80}:\ \textbf{derived}\ —\ \text{one cubic form}\ 4z^3-24z^2+(21-12\gamma)z-(\gamma-1)^2\ \text{with three constants}\ \gamma_{\mathrm{ch}}=\sqrt2\cos3\delta_{\mathrm{ch}};\ \text{Vieta}\ \sum_iz_i=6;\ Q=\tfrac23;\\ & \text{the identity}\ r^{(\ell)}_0=1+\gamma_q;\ \text{the twelve exact}\ 10^{-12}\text{-brackets re-certifying every frozen decimal of Part 3};\ \text{the hierarchy-order theorem}\ q<\nu<\ell\\\text{the inter-family bridge} & \mathrm{BM.82},\ \mathrm{BM.83}:\ \textbf{derived as an exact factorization with one residue}\ —\ m^2_{\mathrm{pole}}/m^2_{\mathrm{hid}}=\rho_{\mathrm B}\mathcal R_s\ \text{with six exact algebraic factors};\ \text{every intra-currency ratio residue-free};\\ & \textbf{the localization:}\ \text{every cross-currency ratio is an exact expression in}\ (\rho_{\mathrm B},\sigma_\ell,\sigma_\nu)\ \text{and in no other undeclared datum};\ \neg\text{a value for any of the three is derived}\ [\mathrm{BM.69}(4)\ \text{stands}]\\\text{the neutron--proton mass ratio} & \mathrm{BM.84}:\ \textbf{derived on the sub-critical branch}\ —\ \ln1280/(\ln960+6c\,\varepsilon_{\mathrm{EM}}\ln2);\ \text{range}\ (0,\tfrac{21}{20});\ \text{irrational at every rational coupling};\ \varepsilon^{\mathrm{crit}}(c)\ \text{irrational};\\ & \text{the trichotomy total with integer-certified non-empty opening and closing regions};\ \frac{m_n}{m_p}\ne\frac43\ \text{at every admissible coupling}\\\text{the beta decay} & \mathrm{BM.85}:\ \textbf{allowed on the specified read's opening region},\ \textbf{closed on the informational currency}\ —\ \text{read-relative, complete,}\ \neg\text{retracted};\\ & \text{the ratio forms}\ Q_\beta/m_n\le-\ln3/\ln1280,\ Q_\beta/m_p\le-\ln3/\ln960;\ \text{the decay direction fixed in both currencies}\end{array}
\]

%dsp-level:0
\[
(2)\quad\textbf{the derived inventory (the new exact vocabulary of this tier):}
\]

%dsp-level:0
\[
\qquad \text{the sine-ratio law}\ r_{\mathrm{intra}}=\sin(k_f\delta_v/4)/\sin(\delta_v/4);
\]

%dsp-level:0
\[
\qquad T_5(\mathsf c)=0\iff16\mathsf c^4-20\mathsf c^2+5=0;
\]

%dsp-level:0
\[
\text{the mass-squared cubic}\ 4z^3-24z^2+(21-12\gamma)z-(\gamma-1)^2
\]

%dsp-level:0
\[
= 0\ \text{with}\ \sum_iz_i
\]

%dsp-level:0
\[
= 6,
\]

%dsp-level:0
\[
\ \ \sum_{i<j}z_iz_j
\]

%dsp-level:0
\[
= \tfrac{21-12\gamma}4,
\]

%dsp-level:0
\[
\ \ \prod_iz_i
\]

%dsp-level:0
\[
= \tfrac{(\gamma-1)^2}4;
\]

%dsp-level:0
\[
\qquad r^{(\ell)}_0=1+\gamma_q;
\]

%dsp-level:0
\[
\text{the minimal polynomials}\ x^8-24x^6+76x^4-64x^2+16\ (\text{the dodecahedral intra ratio}),
\]

%dsp-level:0
\[
\ \ x^4-24x^3+76x^2-64x+16\ (z),
\]

%dsp-level:0
\[
\ \ x^4-48x^3+304x^2-512x+256\ (\text{the ceiling}),
\]

%dsp-level:0
\[
\ \ x^8-22x^6+119x^4-38x^2+1\ (\Delta_3/\hat\mu'^2),
\]

%dsp-level:0
\[
x^8-48x^7+410x^6-1328x^5+1899x^4-1232x^3+330x^2-32x+1\ \Bigl(\frac{m^2_{\mathrm{down}}(3)}{m^2_{\mathrm{up}}(3)}\Bigr),
\]

%dsp-level:0
\[
\qquad x^4-1200x^3+190000x^2-8000000x+100000000\ \Bigl(\frac{\mathcal R_{\max}}{\mathcal R_{\min}}\Bigr);
\]

%dsp-level:0
\[
\frac{\Delta_2}{m^2_{\mathrm{up}}(2)}=\frac{\sqrt2}2;
\]

%dsp-level:0
\[
\qquad \frac{\mathcal R_{\max}}{\mathcal R_{\min}}=\frac{50}{1-\mathsf c}=25\cdot\frac{4}{2-2\mathsf c},
\]

%dsp-level:0
\[
\qquad 25=\Bigl(\frac{E(\mathcal F_3)}{E(\mathcal F_1)}\Bigr)^2;
\]

%dsp-level:0
\[
\qquad \sigma_\ell=\frac{\mu_\ell\mathfrak h}{\ln2},\quad\sigma_\nu=\frac{\mu_\nu\mathfrak h}{\ln2};
\]

%dsp-level:0
\[
\frac{\ln1280}{\ln960}\notin\mathbb Q;
\]

%dsp-level:0
\[
\qquad \varepsilon^{\mathrm{crit}}(c)\notin\mathbb Q;
\]

%dsp-level:0
\[
\qquad t\big|_{m_n=m_p}=2-\log_23\notin\mathbb Q;
\]

%dsp-level:0
\[
\qquad t\big|_{m_n/m_p=4/3}=-\,\frac{\ln405}{\ln16}<0;
\]

%dsp-level:0
\[
\textbf{the integer certificates:}
\]

%dsp-level:0
\[
\quad 2^{44}>3^{26}\!\cdot\!5,
\]

%dsp-level:0
\[
\quad 2^{34}<3^{21}\!\cdot\!5,
\]

%dsp-level:0
\[
\quad 2^{42}>3^{25}\!\cdot\!5,
\]

%dsp-level:0
\[
\quad 2^{40}<3^{24}\!\cdot\!5,
\]

%dsp-level:0
\[
\quad 3^{13}<2^{16}\!\cdot\!5^2,
\]

%dsp-level:0
\[
\quad 3^{20}>2^{24}\!\cdot\!5^3,
\]

%dsp-level:0
\[
\quad 3^{21}<2^{24}\!\cdot\!5^4,
\]

%dsp-level:0
\[
\quad 3^{16}>2^{18}\!\cdot\!5^3;
\]

%dsp-level:0
\[
720>361,
\]

%dsp-level:0
\[
\quad 12500>12321,
\]

%dsp-level:0
\[
\quad 48<49,
\]

%dsp-level:0
\[
\quad 341+140\sqrt6>0,
\]

%dsp-level:0
\[
\quad 17714961>13547520,
\]

%dsp-level:0
\[
\quad 2^{13}>3^8,
\]

%dsp-level:0
\[
\quad 2^{11}<3^7,
\]

%dsp-level:0
\[
\quad 2^8>3^5,
\]

%dsp-level:0
\[
\quad 2^{31}<3^{20}
\]

%dsp-level:0
\[
\qquad [\text{machine: every row exact}];
\]

%dsp-level:0
\[
(3)\quad\textbf{the status adjudication:}
\]

%dsp-level:0
\[
\qquad \boxed{\ \textbf{the mass-ratio census is complete inside every currency, and its entire cross-currency obstruction is three undeclared dimensionless numbers and one attachment}\ }
\]

%dsp-level:0
\[
\textbf{the level-exactness of the claim:}
\]

%dsp-level:0
\[
\qquad \text{derived: every ratio formed from}\ D_S,\ 2E,\ W(Q),\ z_i^{(\mathrm{ch})},\ \ln1280,\ \ln960,\ \Delta T(S);
\]

%dsp-level:0
\[
\qquad \neg\text{derived:}\ \rho_{\mathrm B},\ \sigma_\ell,\ \sigma_\nu,\ (\mathrm{R2}),\ (\mathrm{R3});
\]

%dsp-level:0
\[
\text{the four grounds of}\ \mathrm{BM.69}(4)\ \text{bound}\ \rho_{\mathrm B}\ \text{and leave the census untouched};
\]

%dsp-level:0
\[
\qquad \text{deficit}\ 12\text{'s two ceilings stand}\ [\text{Census FT.76}(3);\ \mathrm{FT.79}(2);\ \mathrm{FT.84}(2)];
\]

%dsp-level:0
\[
\qquad \alpha_B\ \text{stays barred on}\ \mathrm{BM.61}(8)\text{'s four grounds};
\]

%dsp-level:0
\[
(4)\quad\textbf{the scope rows (recorded,}\ \neg\textbf{silent):}
\]

%dsp-level:0
\[
\qquad \text{no clause of}\ \mathrm{BM.1}\text{--}\mathrm{BM.78}\ \text{is retracted, replaced or renumbered};
\]

%dsp-level:0
\[
\qquad \text{no clause of Parts 1--3 is modified};
\]

%dsp-level:0
\[
\qquad \text{no Census cell gains or loses a value};
\]

%dsp-level:0
\[
\mathrm{BM.69}(2)\text{'s bridge identity,}\ \mathrm{BM.62}(5),
\]

%dsp-level:0
\[
(6)\text{'s ratio and critical strength,}\ \mathrm{BM.77}(5)\text{'s trichotomy and}\ \mathrm{BM.75}(4),
\]

%dsp-level:0
\[
(5)\text{'s splitting are re-derived at their own loci and re-cited here,}\ \neg\text{superseded};
\]

%dsp-level:0
\[
\text{the conditionalities}\ (\mathrm{Def\text{-}H.37},\ \mathrm{Def\text{-}H.38},\ \mathrm{M1}/\mathrm{M2},\ \mathrm{CT2.6},\ \mathrm{BM.59}(9)\text{'s single-substitution hypothesis},\ \mathrm{BM.8}(2)\text{'s stationarity})\ \text{are inherited at every point of use}\ [\text{A §0.2a}];
\]

%dsp-level:0
\[
\text{the variant choice}\ (\mathrm A)/(\mathrm B)\ \text{of Census FT.76}(1)\ \text{is}\ \neg\text{adopted}\ —\ \text{both are reported}\ [\mathrm{BM.79}(2),(3),(7)];
\]

%dsp-level:0
\[
\qquad \text{route}\ (\mathrm{ii})\ \text{remains unadopted beyond the vertex kinematics}\ [\mathrm{BM.30}(5);\ \mathrm{BM.77}(1),(9)];
\]

%dsp-level:0
\[
(5)\quad\textbf{the standing rules:}
\]

%dsp-level:0
\[
\qquad \textbf{the declarations this tier adds: none}\ —\ \neg\text{a name attachment},
\]

%dsp-level:0
\[
\ \neg\text{a read specification},
\]

%dsp-level:0
\[
\ \neg\text{a P4 addition},
\]

%dsp-level:0
\[
\ \neg\text{a registry addendum is issued or required};
\]

%dsp-level:0
\[
\neg\text{a new dimensionful scale}\ [\mathrm{P1}];
\]

%dsp-level:0
\[
\qquad \text{§Cal-BM stays empty};
\]

%dsp-level:0
\[
\qquad \sigma_\ell,\ \sigma_\nu,\ \rho_{\mathrm B}\ \text{are restatements of Cal-F/P4 data,}\ \neg\text{new members}\ [\mathrm{BM.77}(2)\ \text{pattern}];
\]

%dsp-level:0
\[
\neg\text{a measured value and}\ \neg\text{a numerical coupling}\ [\mathrm{P5};\ \mathrm{EM.25}(7),(9);\ \text{machine: zero external tokens, zero decimal numerals, zero}\ [\mathrm{EXT}]\ \text{markers in this tier}];
\]

%dsp-level:0
\[
\text{no measured external constant is identified with any corpus quantity}\ [\mathrm{P5}];
\]

%dsp-level:0
\[
\qquad \text{species names are confined to the identification layers}\ [\mathrm{Def\text{-}BM.3}(7);\ \mathrm{Def\text{-}BM.6}(7);\ \mathrm{Def\text{-}BM.14}(8)];
\]

%dsp-level:0
\[
\text{topology labels are internal bookkeeping}\ [\mathrm{Def\text{-}H.32}(4)];
\]

%dsp-level:0
\[
\qquad \text{amplitude}\ne\text{mass}\ [\mathrm{FT.79}(5)];
\]

%dsp-level:0
\[
\qquad \text{rung}\ne\text{scale}\ [\mathrm{D0.8}];
\]

%dsp-level:0
\[
\qquad \neg\text{a HistLog disposition imported as premise};
\]

%dsp-level:0
\[
\alpha_B:\quad\neg\text{added},\quad\neg\text{derived},\quad\text{barred}\ —\ \mathrm{BM.61}(8)\text{'s}\ (\mathrm O1)\text{--}(\mathrm O4)\ \text{stand}
\]

%dsp-level:0
\[
\qquad [\text{machine: the token scan of this tier}\ —\ \text{every occurrence is a barred-object row}];
\]

%dsp-level:0
\[
\neg\text{a Yukawa-type object exists or is added}\ [\mathrm{EW.15};\ \mathrm{FT.145}(3)];
\]

%dsp-level:0
\[
\qquad \text{the extra-axiomatic list}\ [\text{A §0.2a}]\ \text{gains no member};
\]

%dsp-level:0
\[
(6)\quad\textbf{the boundary inventory (every row dispositioned in full at BM.143):}\qquad\begin{array}{l|l}(\mathrm{R1})\ \ \rho_{\mathrm B}=\hat\mu'^2\mathfrak h^2/(\ln2)^2 & \text{the one undeclared dimensionless ratio of the bridge}\ [\mathrm{BM.69}(5);\ \mathrm{BM.82}(6);\ \text{four grounds bound it}]\ \ —\ \text{item}\ 7\\ (\mathrm{R2})\ \ \text{the loop realization} & \text{a clean-enclosure vertex-figure loop as a ledger configuration}\ [\mathrm{BM.67}(9),(10);\ \text{Census FT.63}(4)\ \mathrm{Rmk}]\\ (\mathrm{R3})\ \ \text{the flavour composition} & \text{which three of the six states a sealed class occupies}\ [\mathrm{BM.28}(7);\ \mathrm{BM.75}(6)]\\ \sigma_\ell,\ \sigma_\nu,\ \mu_q & \text{the three channel scales}\ [\mathrm{P4};\ \text{Cal-F}]\\ e & \text{the numerical value of the electric coupling}\ —\ \text{one calibration datum decides the branch}\ [\mathrm{EM.25}(7);\ \mathrm{BM.62}(12);\ \mathrm{BM.84}(9)]\\ m_\nu\ \text{absolute} & \neg\text{derived}\ [\mathrm{FT.143}(3)];\ \text{the}\ \nu\text{-channel's internal ratios are derived}\ [\mathrm{BM.80}(6)]\\ \text{the}\ \nu\text{-line index}\ i_\nu & \text{unresolved}\ [\mathrm{BM.77}(2):\ \text{every row of this tier holds uniformly in}\ i_\nu]\\ \text{the mixing matrix} & \neg\text{derived}\ [\mathrm{FT.145}(3);\ \mathrm{BM.75}(6)]\\ \text{the generation count as a selection} & \text{three lines derived, three selected}\ \neg\text{derived}\ [\mathrm{BM.60}(1)]\\ \text{the two mass ceilings} & \text{deficit}\ 12\ \text{stands}\ [\text{Census FT.76}(3);\ \mathrm{FT.79}(2);\ \mathrm{FT.84}(2)]\end{array}
\]

%dsp-level:0
\[
(7)\quad\textbf{the adjudication:}
\]

%dsp-level:0
\[
\qquad \boxed{\ \text{the quark mass ratios are derived from the Platonic}\ D(m)\ \text{values and the Gram phase};\qquad\text{the inter-family bridge is derived up to three named dimensionless numbers};}
\]

%dsp-level:0
\[
\boxed{\ \text{the neutron--proton mass ratio is derived on the sub-critical branch, exact, bounded by}\ \tfrac{21}{20},\ \text{and irrational at every rational coupling};\qquad\text{the beta decay is allowed on the specified read's opening region and closed on the informational currency}\ }
\]

%dsp-level:0
\[
\textbf{the currencies are distinct and stay distinct:}
\]

%dsp-level:0
\[
\qquad \text{four reads, two formed ratios and two formed differences, one proven incommensurability, one localization, zero conversions};
\]

%dsp-level:0
\[
\text{item 7 stands with its residue named};
\]

%dsp-level:0
\[
\qquad \text{the scales stand undeclared};
\]

%dsp-level:0
\[
\qquad \text{nothing measured entered anywhere}.
\]

\textbf{Req.} D0.8 (Parts 1–2); A §0.2a, P1, P4, P5 (Parts 1–2); EM.25(7), EM.25(9), EW.15 (Part 2); Def-H.32(4), Def-H.37, Def-H.38, FT.79(2), FT.79(5), FT.84(2), FT.143(3), FT.145(3) (Part 3); Census FT.63(4), FT.76(1), FT.76(3) (row-wise); Def-BM.3(7), Def-BM.6(7), Def-BM.8(2), Def-BM.14(8), BM.8(2), BM.28(4), BM.28(7), BM.30(5), BM.59(9), BM.60(1), BM.61(8), BM.62(5), BM.62(6), BM.62(12), BM.67(9), BM.67(10), BM.69(4), BM.69(5), BM.75(4), BM.75(5), BM.75(6), BM.77(1), BM.77(2), BM.77(5), BM.77(9), BM.79, BM.80, BM.81, BM.82, BM.83, BM.84, BM.85 (this part); Cal-F; the machine sections §A–§H green [math; machine: \texttt{checks/\allowbreak{}verify\_\allowbreak{}H\_\allowbreak{}mass\_\allowbreak{}ratio\_\allowbreak{}census.py} §I].

\textbf{Pf.}

%dsp-level:0
\[
(1)\text{--}(2):\ \text{assembly of the adjudicated rows of}\ \mathrm{BM.79}\text{--}\mathrm{BM.85}\ [\text{machine: §A--§H green}];
\]

%dsp-level:0
\[
\qquad \text{the inventory rows re-listed from their own tags}\ [\text{machine}]
\]

%dsp-level:0
\[
(3):\ \text{the derived/}\neg\text{derived partition by the symbol census of the tier}\ [\text{machine}];
\]

%dsp-level:0
\[
\qquad \mathrm{BM.69}(4)\text{'s four grounds and}\ \mathrm{BM.61}(8)\text{'s obstruction re-cited}
\]

%dsp-level:0
\[
(4):\ \text{the clause census of}\ \mathrm{BM.1}\text{--}\mathrm{BM.78}\ \text{against retraction, replacement and renumbering}\ [\text{machine}];
\]

%dsp-level:0
\[
\qquad \text{the Census-cell invariance scan}\ [\text{machine}];
\]

%dsp-level:0
\[
\qquad \text{the conditionality rows re-cited at their loci}
\]

%dsp-level:0
\[
(5):\ \text{the standing-rule scans: external tokens, decimal numerals, unit tokens, registry tokens, declaration lines, the barred object}\ \alpha_B\ [\text{machine}];
\]

%dsp-level:0
\[
\qquad (6):\ \text{the ten open rows re-cited from their own tags}\ [\text{machine}]
\]

%dsp-level:0
\[
(7):\ \text{assembly}.\ \blacksquare
\]

\textbf{Phys.} The census ends where it began, with a shorter list. Every mass ratio the corpus can form — inside a Platonic ladder, across a generation, between two branches of a doublet, between a cone defect and an edge count, between two charge classes in bits, and between the parent and daughters of a decay — is now written down exactly, in its own currency, with its own minimal polynomial or its own rational bracket, and with the integer certificate that decides its sign or its order. What is left is small and named: one dimensionless packaging between the angle currency and the bit currency, two channel scales, one loop the ledger has not realized, one composition no clause attaches, and one coupling constant the corpus has never been given. The two negatives are as exact as the positives. The informational neutron-to-proton ratio is irrational at every rational coupling and never equals the ensemble's four thirds at any coupling at all, because the coupling that would make them equal is minus the logarithm of four hundred five over the logarithm of sixteen — a negative number, outside the admissible range by a proof, not by a margin. And no scale makes the floor ladder proportional to the channel ladder, or the channel ladder proportional to the transport ladder, so the generation map remains an order and never becomes a number. Three families share one cubic and differ by one cosine each; the heaviest lines of all three nearly coincide while the lightest differ by a factor of twenty-six, exactly as the Platonic ladders pile their hierarchy at the light end. Nothing was fitted, nothing was converted, nothing measured entered, and the ledger's oldest open item stands where it stood — with everything around it now counted.

\medskip
\subsection*{\texorpdfstring{B.1.27\quad BM.87 — The category adjudication: the void distribution factor, the derived exterior value, and the standing rows}{BM.87 B.1.27}}
\addcontentsline{toc}{subsection}{B.1.27 BM.87}
\phantomsection\label{tag:BM-87}
{\small \textit{Record.}\quad Source: Part 4 (Part4\_Baryomesonic\_era.md); tier: Tier BM.18 — The EM Self-Energy Correction: the Category Adjudication, the Running Coupling, the Path-Integral Self-Energy, the Branch-Blindness Theorem, and the Corrected Neutron–Proton Splitting; canonical tag \texttt{BM.87}. Reproduced from the revision-1 canon; a cross-reference stub marks the former locus in the main corpus.}
\medskip
\textbf{Thm.}

%dsp-level:0
\[
(1)\quad\textbf{the referent census (what the graph supplies against what a distribution requires):}\qquad\begin{array}{l|l}\text{a charge distribution requires} & \text{the corpus supplies}\\\hline\text{a volume form}\ dV & \neg\text{one at the graph locus: the pre-dimension graph is a time-directed graph of nodes}\ [\mathrm{D0.1};\ \mathrm{D0.10};\ \mathrm{CT0.3}];\ \text{the metric is emergent and chart-side}\ [\mathrm{CT1.12};\ \mathrm{CT2.6}\ \text{hypotheses}]\\\text{a density}\ \rho_Q(x) & \text{an integer read}\ Q=\sum_iq_{\mathrm{core}}(c_i)\in\mathbb Z\ \text{with no position argument}\ [\mathrm{Def\text{-}BM.8}(1);\ \mathrm{D3.4};\ \mathrm{EM.9}]\\\text{a support with an interior} & \text{an envelope locus}\ \Delta x_{\mathrm B},\ \textbf{a chart-side envelope,}\ \neg\textbf{a ledger distance}\ [\mathrm{Def\text{-}BM.4}(5);\ \mathrm{Def\text{-}BM.12}(4);\ \mathrm{BM.13}(1)]\\\text{a radial profile of the source} & \neg\text{one: the end signs carry no position datum, so internal multipoles are}\ \neg\text{resolved}\ [\mathrm{BM.46}(2);\ \mathrm{BM.28}(7);\ \mathrm{Def\text{-}BM.8}(5)]\end{array}
\]

%dsp-level:0
\[
[\text{machine: the token and inventory census of}\ \mathrm{BM.58}(6)\text{'s table,}\ \mathrm{Def\text{-}BM.12}\text{'s clauses and}\ \mathrm{BM.37}(2)\text{'s alphabets}];
\]

%dsp-level:0
\[
(2)\quad\textbf{the two-route decomposition of the chart-side value:}
\]

%dsp-level:0
\[
\qquad E_{\mathrm{EM}}^{\mathrm{self}}=\frac{c\,e^2q^2}{4\pi\,\Delta x_{\mathrm B}},
\]

%dsp-level:0
\[
\qquad c=\frac12+\iota,
\]

%dsp-level:0
\[
\qquad \iota:=\frac{\text{the interior term}}{e^2q^2/(4\pi\Delta x_{\mathrm B})}
\]

%dsp-level:0
\[
\textbf{the exterior term (distribution-free):}
\]

%dsp-level:0
\[
\qquad \int_{\Delta x_{\mathrm B}}^{\Lambda}\tfrac12E_r^2\,
\]

%dsp-level:0
\[
4\pi r^2dr
\]

%dsp-level:0
\[
= \frac{e^2q^2}{8\pi}\Bigl(\frac1{\Delta x_{\mathrm B}}-\frac1\Lambda\Bigr)\xrightarrow[\Lambda\to\infty]{}\frac{e^2q^2}{8\pi\Delta x_{\mathrm B}}
\]

%dsp-level:0
\[
= \frac12\cdot\frac{e^2q^2}{4\pi\Delta x_{\mathrm B}}
\]

%dsp-level:0
\[
\qquad \implies\boxed{\ c_{\mathrm{ext}}=\frac12\ }
\]

%dsp-level:0
\[
\textbf{the potential route (a point evaluation, no volume, no density):}
\]

%dsp-level:0
\[
\qquad \tfrac12\,q\,\phi(\Delta x_{\mathrm B})=\tfrac12(eq)\frac{eq}{4\pi\Delta x_{\mathrm B}}=\frac{e^2q^2}{8\pi\Delta x_{\mathrm B}}
\]

%dsp-level:0
\[
\qquad [\mathrm{BM.46}(1)\text{'s two routes, re-derived};\ \text{machine: they agree exactly}];
\]

%dsp-level:0
\[
\textbf{the interior term (distribution-dependent):}
\]

%dsp-level:0
\[
\qquad \text{shell}:\ \iota=0;
\]

%dsp-level:0
\[
\qquad \text{uniform ball}:\ \iota=\frac1{10},\quad\frac{E_{\mathrm{in}}}{E_{\mathrm{out}}}=\frac{1/40\pi}{1/8\pi}=\frac15,\quad c_{\mathrm{ball}}=\frac12+\frac1{10}=\frac35
\]

%dsp-level:0
\[
\qquad [\mathrm{BM.46}(1);\ \mathrm{Def\text{-}BM.12}(4);\ \text{machine: exact rationals}];
\]

%dsp-level:0
\[
(3)\quad\textbf{the void theorem:}
\]

%dsp-level:0
\[
\qquad \iota\ \text{is a functional of an interior charge density};
\]

%dsp-level:0
\[
\qquad \neg\exists\ \text{an interior density in the inventory}\ [(1)\ \text{row}\ 4]:
\]

%dsp-level:0
\[
\boxed{\ \textbf{the interior term is}\ \neg\textbf{formed, and the question ``shell or ball?'' has no corpus referent}\ }
\]

%dsp-level:0
\[
\qquad [\text{the category error, stated exactly}];
\]

%dsp-level:0
\[
\textbf{the formed value is therefore derived, not priced:}
\]

%dsp-level:0
\[
\qquad \boxed{\ c_{\mathrm{formed}}=\frac12\ \ \text{exactly, on two independent distribution-free routes}\ }
\]

%dsp-level:0
\[
\textbf{the ball value}\ \tfrac35\ \textbf{is void as a mass datum};
\]

%dsp-level:0
\[
\qquad \mathrm{Def\text{-}BM.12}(4)\text{'s pricing row is void with it, and the definition's clauses}\ (1),
\]

%dsp-level:0
\[
(2),
\]

%dsp-level:0
\[
(3),
\]

%dsp-level:0
\[
(5),
\]

%dsp-level:0
\[
(6)\ \text{stand unmodified};
\]

%dsp-level:0
\[
(4)\quad\textbf{what stands (re-derived,}\ \neg\textbf{retracted):}
\]

%dsp-level:0
\[
\qquad \textbf{the charge selector:}\quad E_{\mathrm{EM}}^{\mathrm{self}}\propto q^2\ \ \text{with the census}\ (1,0,1)\ \ \text{at}\ q=(\pm1,0)
\]

%dsp-level:0
\[
\qquad [\mathrm{BM.46}(2);\ \text{machine}];
\]

%dsp-level:0
\[
\textbf{the monomial shape and its dimension:}
\]

%dsp-level:0
\[
\qquad E_{\mathrm{EM}}^{\mathrm{self}}=\frac12\,\varepsilon_{\mathrm{EM}}\,q^2\,m_{\mathrm{hid}}=\frac12\,\frac{\varepsilon_{\mathrm{EM}}\,q^2\,E_{\mathrm{int}}\ln2}{\mathfrak h},
\]

%dsp-level:0
\[
\qquad [\varepsilon_{\mathrm{EM}}]=\mathfrak q^0,\quad[E_{\mathrm{EM}}^{\mathrm{self}}]=\mathfrak q^{+1/2}=[\text{energy}]
\]

%dsp-level:0
\[
[\mathrm{BM.46}(1)\text{'s unified form at}\ c=\tfrac12;\ \mathrm{Def\text{-}BM.4}(1):\ \Delta x_{\mathrm B}\cdot m_{\mathrm{hid}}=1;\ \mathrm{CT2.1};\ \text{machine: the exponent arithmetic}];
\]

%dsp-level:0
\[
\textbf{the sign theorem, which never needed}\ c\textbf{:}
\]

%dsp-level:0
\[
\qquad \boxed{\ \Delta m_{\mathrm{EM}}>0\ \ \text{for every}\ c>0\ }
\]

%dsp-level:0
\[
\qquad [\mathrm{BM.46}(4),\ \mathrm{Rmk};\ \text{machine: the positive coefficient census}];
\]

%dsp-level:0
\[
\textbf{the envelope locus, the UV/IR finiteness rows and the seal asymmetry:}
\]

%dsp-level:0
\[
\qquad \Delta x_{\mathrm B}>0,\quad E_{\mathrm{EM}}^{\mathrm{self}}<\infty,\quad f(k)=1-\tfrac1k<1\ \forall k<\infty,\quad\text{EM seals at universe scale}
\]

%dsp-level:0
\[
[\mathrm{BM.47}(2),(3),(4),(6);\ \mathrm{EM.11};\ \text{machine}];
\]

%dsp-level:0
\[
\qquad \textbf{the two-body potential, the hydrogenic rows and the binding category bar:}
\]

%dsp-level:0
\[
\quad V=-\varepsilon_{\mathrm{EM}}/d,
\]

%dsp-level:0
\[
\quad E_n=-\tfrac{\varepsilon_{\mathrm{EM}}^2m_r}{2n^2},\quad\neg\text{a bound channel}\ [\mathrm{BM.48}(1)\text{--}(4);\ \text{standing}];
\]

%dsp-level:0
\[
(5)\quad\textbf{the correct method (the relocation):}
\]

%dsp-level:0
\[
\qquad \boxed{\ E_{\mathrm{EM}}^{\mathrm{self}}\ :=\ \sum_{R\,\in\,\mathcal E(\mathcal X)}P(R)\cdot E_{\mathrm{EM}}(R)\ }
\]

%dsp-level:0
\[
\qquad [\mathrm{Def\text{-}BM.5};\ \mathrm{Def\text{-}H.32}(1);\ \mathrm{BM.19}(1)\text{--}(3):\ \text{the derived stationary measure}]:
\]

%dsp-level:0
\[
P(R)\ \text{is the ensemble weight of the configuration}\ R\ [\mathrm{BM.19}(1):\ p=x_+/(x_++x_-)=1-4\varepsilon_q];
\]

%dsp-level:0
\[
\qquad E_{\mathrm{EM}}(R)\ \text{is the formed monomial of}\ (4)\ \text{evaluated on}\ R\ \ [\mathrm{BM.89}];
\]

%dsp-level:0
\[
\textbf{the ensemble is finite at every budget and non-singleton at every nonzero threshold}
\]

%dsp-level:0
\[
\qquad [\mathrm{BM.19}(4),(6);\ \mathrm{Def\text{-}H.32}(1)];
\]

%dsp-level:0
\[
(6)\quad\textbf{the rescope census (recorded,}\ \neg\textbf{retracted):}
\]

%dsp-level:0
\[
\qquad \mathrm{BM.46}(1)\text{'s field-route volume integral and}\ \mathrm{Def\text{-}BM.12}(4)\text{'s two-value pricing:}\ \textbf{void as mass data}\ —\ \text{the arithmetic stands as the priced alternative's exact content};
\]

%dsp-level:0
\[
\mathrm{BM.46}(1)\text{'s potential route,}\ (2)\text{'s selector,}\ (3)\text{'s corrected rest read,}\ (4)\text{'s splitting theorem,}\ (5)\text{'s correction-species row:}\ \textbf{stand},
\]

%dsp-level:0
\[
\ \text{with}\ c\ \text{evaluated at}\ \tfrac12;
\]

%dsp-level:0
\[
\mathrm{BM.47}(2)\text{--}(7):\ \textbf{stand untouched}\ —\ \text{each is a flux or exponent row,}\ \neg\text{a density row};
\]

%dsp-level:0
\[
\qquad \mathrm{BM.48},\ \mathrm{BM.49}(1),(2),\ \mathrm{BM.50}:\ \textbf{stand},\ \text{with}\ c=\tfrac12\ \text{substituted where a value is displayed};
\]

%dsp-level:0
\[
\mathrm{BM.52}(4),
\]

%dsp-level:0
\[
(5),
\]

%dsp-level:0
\[
\ \mathrm{BM.55}(1),
\]

%dsp-level:0
\[
(2),
\]

%dsp-level:0
\[
\ \mathrm{BM.62}(3)\text{--}(9),
\]

%dsp-level:0
\[
\ \mathrm{BM.63},
\]

%dsp-level:0
\[
\ \mathrm{BM.65}(3)\text{--}(5),
\]

%dsp-level:0
\[
\ \mathrm{BM.66}(3),
\]

%dsp-level:0
\[
\ \mathrm{BM.70},
\]

%dsp-level:0
\[
\ \mathrm{BM.75}(4)\text{'s opposed row},
\]

%dsp-level:0
\[
\ \mathrm{BM.77}(9),
\]

%dsp-level:0
\[
\ \mathrm{BM.84}(1)\text{--}(9),
\]

%dsp-level:0
\[
\ \mathrm{BM.85}(1)\text{--}(3):
\]

%dsp-level:0
\[
\textbf{stand on their own locus}\ —\ \text{the fully crystallized, vacuum-gauge, chart-side read}\ (w{=}1,\ N{=}1,\ c\ \text{symbolic});
\]

%dsp-level:0
\[
\qquad \mathrm{BM.91}\ \text{generalizes them with two derived factors and re-derives every certificate};
\]

%dsp-level:0
\[
\neg\text{a frozen refutation is reopened};\qquad\neg\text{a Census cell changes};\qquad
\]

%dsp-level:0
\[
(7)\quad\textbf{status:}
\]

%dsp-level:0
\[
\qquad \textbf{derived:}\ (1)\text{--}(5);
\]

%dsp-level:0
\[
\qquad \textbf{recorded:}\ (6);
\]

%dsp-level:0
\[
\qquad \textbf{void:}\ \text{the distribution pricing as a mass datum}\ (3);
\]

%dsp-level:0
\[
\neg\text{a declaration},
\]

%dsp-level:0
\[
\quad\neg\text{a P4 addition},
\]

%dsp-level:0
\[
\quad\neg\text{a new dimensionful scale}\ [\mathrm{P1}],
\]

%dsp-level:0
\[
\quad\neg\text{a measured value}\ [\mathrm{P5}],
\]

%dsp-level:0
\[
\quad\text{§Cal-BM stays empty};
\]

%dsp-level:0
\[
\qquad \textbf{the residue:}\ \text{the interior multipole content}\ [\mathrm{BM.46}(2)\text{'s resolution row, named, unfilled}].
\]

\textbf{Req.} D0.1, D0.10, D2.9, D3.4 (Parts 1–2); CT0.3, CT0.8, CT1.12, CT1.22, CT2.1, CT2.6, D1.5, EM.9, EM.11, EM.15 (Part 2); P1, P4, P5 (Parts 1–2); A4 (Part 1); Def-H.32(1) (Part 3); Def-BM.4(1), Def-BM.4(5), Def-BM.5, Def-BM.8(1), Def-BM.8(5), Def-BM.12(1), Def-BM.12(2), Def-BM.12(3), Def-BM.12(4), Def-BM.12(5), Def-BM.12(6), BM.13(1), BM.19(1), BM.19(2), BM.19(3), BM.19(4), BM.19(6), BM.28(1), BM.28(7), BM.35(5), BM.37(2), BM.46(1), BM.46(2), BM.46(3), BM.46(4), BM.46(5), BM.47(2), BM.47(3), BM.47(4), BM.47(6), BM.48(1), BM.48(4), BM.49(1), BM.49(2), BM.50, BM.52(4), BM.52(5), BM.55(1), BM.55(2), BM.58(6), BM.62(3), BM.62(12), BM.63, BM.65(3), BM.65(5), BM.66(3), BM.69(4), BM.69(5), BM.70, BM.75(4), BM.77(7), BM.77(9), BM.82(4), BM.83(3), BM.83(5), BM.84, BM.85(1), BM.85(3) (this part); exact radial integrals and their limit rows, exact rational coefficient arithmetic, Laurent-monomial assembly with $\mathfrak q$-exponent arithmetic, token and alphabet censuses [math; machine: \texttt{checks/\allowbreak{}verify\_\allowbreak{}H\_\allowbreak{}em\_\allowbreak{}self\_\allowbreak{}energy\_\allowbreak{}correction.py} §A].

\textbf{Pf.}

%dsp-level:0
\[
(1):\ \text{the four requirement rows against the inventory}\ [\text{machine: the token census of}\ \mathrm{Def\text{-}BM.12},\ \mathrm{BM.37}(2)\text{'s four alphabets and}\ \mathrm{BM.58}(6)\text{'s seven rows}];
\]

%dsp-level:0
\[
\qquad \Delta x_{\mathrm B}\cdot m_{\mathrm{hid}}=1\ \text{with}\ \Delta x_{\mathrm B}\ \text{an envelope}\ [\mathrm{Def\text{-}BM.4}(5)\text{'s own scope row re-cited}]
\]

%dsp-level:0
\[
(2):\ \int_{\Delta x_{\mathrm B}}^{\Lambda}\tfrac12(eq/4\pi r^2)^2\,
\]

%dsp-level:0
\[
4\pi r^2dr
\]

%dsp-level:0
\[
= \tfrac{e^2q^2}{8\pi}(1/\Delta x_{\mathrm B}-1/\Lambda)\ [\text{machine: the exact antiderivative}\ -1/r];
\]

%dsp-level:0
\[
\qquad \Lambda\to\infty:\ 1/\Lambda\to0\ [\text{machine: the limit row}];
\]

%dsp-level:0
\[
\qquad \tfrac{e^2q^2}{8\pi\Delta x_{\mathrm B}}=\tfrac12\cdot\tfrac{e^2q^2}{4\pi\Delta x_{\mathrm B}}\ [\text{machine:}\ 1/8\pi=(1/2)\cdot1/4\pi]
\]

%dsp-level:0
\[
\tfrac12q\phi(\Delta x_{\mathrm B})\ \text{with}\ \phi
\]

%dsp-level:0
\[
= eq/(4\pi\Delta x_{\mathrm B})\ [\text{Gauss}+\text{harmonicity of}\ 1/r;\ \mathrm{Def\text{-}BM.12}(3)]
\]

%dsp-level:0
\[
= \tfrac{e^2q^2}{8\pi\Delta x_{\mathrm B}}\ [\text{machine: the two routes agree}];
\]

%dsp-level:0
\[
\qquad \text{ball}:\ E_{\mathrm{in}}/E_{\mathrm{out}}=(1/40\pi)/(1/8\pi)=1/5,\quad\tfrac12+\tfrac12\cdot\tfrac15=\tfrac35\ [\text{machine: exact Fractions}]
\]

%dsp-level:0
\[
(3):\ \iota\ \text{is defined by an integral over the source's interior support}\ [(2)];
\]

%dsp-level:0
\[
\qquad \neg\exists\ \text{a position datum on the end signs}\ [\mathrm{D3.4};\ \mathrm{BM.46}(2);\ \mathrm{BM.28}(7)]\implies\neg\exists\ \text{an interior density}\implies\iota\ \neg\text{formed}\ [\text{machine: the alphabet census of}\ \iota\text{'s defining expression against}\ \mathrm{BM.37}(2)\text{'s four alphabets}];
\]

%dsp-level:0
\[
\qquad c_{\mathrm{formed}}=c_{\mathrm{ext}}=\tfrac12\ [\text{machine}]
\]

%dsp-level:0
\[
(4):\ q^2\ \text{at}\ q\in\{-1,0,+1\}=(1,0,1)\ [\text{machine}];
\]

%dsp-level:0
\[
\qquad \tfrac12\varepsilon_{\mathrm{EM}}q^2m_{\mathrm{hid}}=\tfrac12\varepsilon_{\mathrm{EM}}q^2E_{\mathrm{int}}\ln2/\mathfrak h\ [\text{machine: one Laurent monomial}];
\]

%dsp-level:0
\[
\qquad [\varepsilon_{\mathrm{EM}}]=[e^2]=\mathfrak q^0,\ [\Delta x_{\mathrm B}]=\mathfrak q^{-1/2}\implies[E]=\mathfrak q^{+1/2}\ [\text{machine: the exponent arithmetic}];
\]

%dsp-level:0
\[
\qquad \Delta m_{\mathrm{EM}}>0\ \forall c>0\ [\text{machine: the coefficient census at}\ e^2>0,\ q^2\ge0,\ m_{\mathrm{hid}}>0]
\]

%dsp-level:0
\[
(5):\ \mathcal E(\mathcal X)\ \text{finite at each budget}\ [\mathrm{Def\text{-}H.32}(1);\ \mathrm{BM.19}(4)];
\]

%dsp-level:0
\[
\qquad p=x_+/(x_++x_-)=1-4\varepsilon_q\ [\mathrm{BM.19}(1),(2)];
\]

%dsp-level:0
\[
\qquad |\mathcal E|>1\iff\varepsilon_q>0\ [\mathrm{BM.19}(6)];
\]

%dsp-level:0
\[
\qquad \varepsilon_q=0\implies\mathcal E=\{\text{the crystallized configuration}\}\ [\mathrm{BM.19}(3);\ \mathrm{Def\text{-}BM.5}(5)]
\]

%dsp-level:0
\[
(6)\text{--}(7):\ \text{the clause-by-clause rescope census of}\ \mathrm{BM.46}\text{--}\mathrm{BM.85}\ [\text{machine}];
\]

%dsp-level:0
\[
\qquad \text{the standing-rule scans}\ [\text{machine}].\ \blacksquare
\]

\textbf{Phys.} The corpus never had a charged ball. It had an integer read on the ends of chains, a width derived from an edge count, and a field equation varied from an action it already owned. Asking whether the charge sits on a shell or fills a ball asks for a density the graph cannot supply: the end signs carry no positions, so there is nothing to smear and nothing to integrate over the interior. That single observation disposes of the priced coefficient and, at the same time, rescues the number that mattered — because the exterior field of a sealed charge is fixed by Gauss alone, and the potential route needs nothing but the potential at one radius. Both give exactly one half, twice, with no distribution in either argument. So the self-energy keeps its shape, its selector, its sign and its dimension, loses a classical artifact, and gains a locus: it is now an average over the flicker ensemble, which is where the composite actually lives once the threshold is nonzero. The interior multipole content stays open exactly where the earlier audit left it, named and unfilled, and every downstream row is re-scoped rather than erased — the chart-side numbers stand as the fully crystallized, vacuum-gauge specialization of a formula that now carries two derived factors the chart could not see.

\medskip
\subsection*{\texorpdfstring{B.1.28\quad BM.93 — The status of the EM self-energy correction: adjudication and scope}{BM.93 B.1.28}}
\addcontentsline{toc}{subsection}{B.1.28 BM.93}
\phantomsection\label{tag:BM-93}
{\small \textit{Record.}\quad Source: Part 4 (Part4\_Baryomesonic\_era.md); tier: Tier BM.18 — The EM Self-Energy Correction: the Category Adjudication, the Running Coupling, the Path-Integral Self-Energy, the Branch-Blindness Theorem, and the Corrected Neutron–Proton Splitting; canonical tag \texttt{BM.93}. Reproduced from the revision-1 canon; a cross-reference stub marks the former locus in the main corpus.}
\medskip
\textbf{Thm.}

%dsp-level:0
\[
(1)\quad\textbf{the answers:}\qquad\begin{array}{l|l}\mathrm{BM.46}\ \text{and}\ \mathrm{Def\text{-}BM.12}(4) & \mathrm{BM.87}:\ \textbf{the distribution pricing is void as a mass datum}\ —\ \text{the interior term requires a charge density, and}\ \neg\exists\ \text{one}:\\ & \text{the end signs carry no position datum};\ \text{the question ``shell or ball?'' has no corpus referent};\\ & \textbf{the formed value is derived twice:}\ c_{\mathrm{formed}}=\tfrac12,\ \text{from the exterior flux integral and from the potential at one radius};\\ & \text{the selector}\ q^2,\ \text{the sign theorem, the monomial shape, the dimension row, the envelope locus and all of}\ \mathrm{BM.47}\ \textbf{stand}\\\text{the running coupling} & \mathrm{BM.88}:\ \textbf{derived as an exact geometric profile}\ —\ e(r,E,N)=e_0\sqrt{N(E)\,f(r/\Delta x_{\mathrm B}(E))\,\mathsf g(G(E))},\\ & f(k)=1-\tfrac1k,\quad\Delta x_{\mathrm B}(E)=\mathfrak h/(E_{\mathrm{int}}(E)\ln2),\quad G=(\Delta x_{\mathrm B})^2\kappa^2\rho=2\pi E_{\mathrm{int}}^2\ln2\ \ (\mathfrak h\text{-free},\ \Lambda\text{-free},\ \mathrm{P4}\text{-free});\\ & f(1)=0:\ \text{the coupling vanishes at the member's own envelope};\ f\ \text{strictly increasing in the energy scale};\\ & \textbf{priced,}\ \neg\textbf{adopted:}\ \mathsf g\in\{1,\,1+G,\,(1+G)^{-1}\};\qquad\textbf{symbolic:}\ e_0\ [\mathrm{EM.25}(7)]\\\text{the EM self-energy} & \mathrm{BM.89}:\ \textbf{computed from the path integral}\ —\ E_{\mathrm{EM}}^{\mathrm{self}}(Q)=\tfrac12 N\,w\,\mathsf g\,\varepsilon_{\mathrm{EM}}Q^2E_{\mathrm{int}}^{\mathrm{cryst}}\ln2/\mathfrak h;\\ & \text{the derived measure}\ p=1-4\varepsilon_q;\ \ w(0)=1\ \text{exactly};\ 0<w\le1;\ \text{the Bernoulli specialization}\ w=1-4\varepsilon_q\ \text{marked};\\ & \text{the crystallized vacuum-gauge specialization reproduces}\ \mathrm{BM.87}(2)\text{'s exterior value coefficient by coefficient};\\ & \text{route-relative:}\ (\mathrm i)\ \text{charge-degenerate},\ (\mathrm{ii})\ \text{the four exact frozen values},\ (\mathrm{iii})\ \text{barred}\\\text{the quark splitting and}\ m_{\mathrm{hid}} & \mathrm{BM.90}:\ \textbf{hypothesis (W) refuted as a derived coupling}\ —\ \text{the flicker measure's alphabet is}\ \{\varepsilon_q\},\ \text{the defect values' is}\ \{\delta_v,\hat\mu'^2\},\ \text{disjoint};\\ & \text{four independent grounds}\ (\text{degree }16,\ \text{two-torsion},\ \text{calibration weight},\ \text{strength bar});\ \ w=w(\varepsilon_q)\ \textbf{is branch-blind};\\ & \text{the six states act in the pole currency alone};\ \neg\text{a count}\to\text{strength datum is issued}\\\text{the neutron--proton splitting} & \mathrm{BM.91}:\ \textbf{recomputed}\ —\ \Delta m\,\mathfrak h=\ln\tfrac43-3Nw\mathsf g\,\varepsilon_{\mathrm{EM}}\ln2;\ \ w\ \text{cancels in the ledger difference};\\ & \varepsilon^{\mathrm{crit}}=(2-\log_23)/(3Nw\mathsf g)\ \text{irrational at rational}\ (N,w,\mathsf g);\ \textbf{the parameter-free bar}\ Nw\mathsf g\,\varepsilon_{\mathrm{EM}}\le\tfrac18\implies m_n>m_p;\\ & 1<m_n/m_p<\tfrac{10}9\ \text{universally, with}\ \ln20/\ln15\in(\tfrac{11}{10},\tfrac{10}9)\ \text{the dissolved supremum};\ \tfrac43\ \text{unattainable on}\ 64<405\\\text{the beta decay} & \mathrm{BM.92}:\ \textbf{closed on the informational currency uniformly in}\ (w,N,\mathsf g,\varepsilon_{\mathrm{EM}}),\ \textbf{allowed on the specified read's}\ \mathcal O_S;\\ & \text{the specified read is untouched by both acts};\ \text{the decay direction is fixed in both currencies};\ \text{both verdicts read-relative,}\ \neg\text{retracted}\end{array}
\]

%dsp-level:0
\[
(2)\quad\textbf{the derived inventory (the new exact vocabulary of this tier):}
\]

%dsp-level:0
\[
\qquad c_{\mathrm{formed}}=\tfrac12;
\]

%dsp-level:0
\[
\qquad \iota_{\mathrm{ball}}=\tfrac1{10};
\]

%dsp-level:0
\[
\qquad \frac{E_{\mathrm{in}}}{E_{\mathrm{out}}}=\frac15;
\]

%dsp-level:0
\[
\qquad \frac{E_{\mathrm{tot}}}{E_{\mathrm{out}}}=\frac65;
\]

%dsp-level:0
\[
f(k)=1-\frac1k,\qquad k=\frac{r\,E_{\mathrm{int}}\ln2}{\mathfrak h},\qquad G=2\pi E_{\mathrm{int}}^2\ln2,\qquad N=\sqrt{\Omega^2+\gamma^{ij}N_iN_j},\qquad e=e_0\sqrt{N\,f\,\mathsf g};
\]

%dsp-level:0
\[
w=\frac{\langle E_{\mathrm{int}}\rangle_{\mathcal E}}{E_{\mathrm{int}}^{\mathrm{cryst}}},
\]

%dsp-level:0
\[
\qquad w(0)=1,
\]

%dsp-level:0
\[
\qquad w=1-4\varepsilon_q\ \ (\text{the marked specialization}),
\]

%dsp-level:0
\[
\qquad E_{\mathrm{EM}}^{\mathrm{self}}(Q)=\frac12Nw\mathsf g\,\varepsilon_{\mathrm{EM}}Q^2\frac{E_{\mathrm{int}}^{\mathrm{cryst}}\ln2}{\mathfrak h};
\]

%dsp-level:0
\[
\Delta m\,\mathfrak h=\ln\frac43-3Nw\mathsf g\,\varepsilon_{\mathrm{EM}}\ln2,
\]

%dsp-level:0
\[
\qquad \varepsilon^{\mathrm{crit}}=\frac{2-\log_23}{3Nw\mathsf g},
\]

%dsp-level:0
\[
\qquad \varepsilon^{*}_{\mathrm{EM}}\big|_{m_n/m_p=4/3}=\frac{(6-6w)\ln2-\ln405}{12Nw\mathsf g\ln2}<0;
\]

%dsp-level:0
\[
\frac{m_n}{m_p}\in\Bigl[\frac{\ln1280}{\ln960},\ \frac{\ln20}{\ln15}\Bigr),
\]

%dsp-level:0
\[
\qquad \frac{\ln20}{\ln15}\in\Bigl(\frac{11}{10},\frac{10}9\Bigr),
\]

%dsp-level:0
\[
\qquad \Delta m_{\mathrm{EM}}^{(\mathrm{ii})}
\]

%dsp-level:0
\[
= \frac{N\mathsf g\,\varepsilon_{\mathrm{EM}}}{2\mathfrak h}\Bigl(\frac{12659355}{98956},\frac{28824930}{393679},\frac{1365}{23},\frac{205695}{3088}\Bigr);
\]

%dsp-level:0
\[
\frac{Q_\beta}{m_n}\le-\frac{\ln3}{6w\ln2+\ln20},
\]

%dsp-level:0
\[
\qquad -\frac{\ln3}{\ln20}\in\Bigl(-\frac25,-\frac13\Bigr);
\]

%dsp-level:0
\[
\qquad \textbf{the certificates:}\quad8192>6561,\quad2048<2187,\quad20^{10}>15^{11},\quad20^9<15^{10},\quad27>20,\quad243<400,\quad64<405
\]

%dsp-level:0
\[
[\text{machine: every row exact}];
\]

%dsp-level:0
\[
(3)\quad\textbf{the status adjudication:}
\]

%dsp-level:0
\[
\qquad \boxed{\ \textbf{the chart-side distribution factor is void, the formed exterior value}\ \tfrac12\ \textbf{is derived, the coupling runs geometrically, and the self-energy is an ensemble average with one derived attenuation}\ }
\]

%dsp-level:0
\[
\textbf{the level-exactness:}
\]

%dsp-level:0
\[
\qquad \text{derived:}\ c_{\mathrm{formed}},
\]

%dsp-level:0
\[
\ f,
\]

%dsp-level:0
\[
\ G,
\]

%dsp-level:0
\[
\ N\text{'s ADM row},
\]

%dsp-level:0
\[
\ w(0)
\]

%dsp-level:0
\[
= 1,
\]

%dsp-level:0
\[
\ \text{the monomial shapes, every inequality and certificate};
\]

%dsp-level:0
\[
\neg\text{derived:}\ e_0\ (\text{one calibration datum}),
\]

%dsp-level:0
\[
\ \mathsf g\ (\text{the curvature composition}),
\]

%dsp-level:0
\[
\ N\ (\text{solution data}),
\]

%dsp-level:0
\[
\ w(\varepsilon_q)\ \text{beyond the specialization},
\]

%dsp-level:0
\[
\ \iota\ (\text{the interior multipole content});
\]

%dsp-level:0
\[
\textbf{the two ceilings and the obstructions stand:}
\]

%dsp-level:0
\[
\qquad \alpha_B\ \text{stays barred on}\ \mathrm{BM.61}(8)\text{'s four grounds};
\]

%dsp-level:0
\[
\qquad \text{deficit}\ 12\ \ [\text{Census FT.76}(3);\ \mathrm{FT.79}(2);\ \mathrm{FT.84}(2)];
\]

%dsp-level:0
\[
\qquad \text{item}\ 7\ \ [\mathrm{BM.69}(4),(5);\ \mathrm{BM.83}(5)];
\]

%dsp-level:0
\[
(4)\quad\textbf{the scope rows (recorded,}\ \neg\textbf{silent):}
\]

%dsp-level:0
\[
\qquad \textbf{void as mass data:}\ \mathrm{Def\text{-}BM.12}(4)\text{'s two-value pricing row and}\ \mathrm{BM.46}(1)\text{'s interior term};
\]

%dsp-level:0
\[
\textbf{standing at their own locus:}\ \mathrm{BM.46}(1)\text{'s potential route,}\ (2)\text{'s selector,}\ (3)\text{'s corrected rest read,}\ (4)\text{'s splitting theorem,}\ (5)\text{'s species row};
\]

%dsp-level:0
\[
\ \mathrm{BM.47}(1)\text{--}(7);
\]

%dsp-level:0
\[
\ \mathrm{BM.48}(1)\text{--}(4);
\]

%dsp-level:0
\[
\ \mathrm{BM.49}(1),
\]

%dsp-level:0
\[
(2);\qquad \ \mathrm{BM.50};
\]

%dsp-level:0
\[
\mathrm{BM.52}(4),
\]

%dsp-level:0
\[
(5);
\]

%dsp-level:0
\[
\ \mathrm{BM.55}(1),
\]

%dsp-level:0
\[
(2);
\]

%dsp-level:0
\[
\ \mathrm{BM.62}(1)\text{--}(12);
\]

%dsp-level:0
\[
\ \mathrm{BM.63};
\]

%dsp-level:0
\[
\ \mathrm{BM.65}(3)\text{--}(5);
\]

%dsp-level:0
\[
\ \mathrm{BM.66}(3);
\]

%dsp-level:0
\[
\ \mathrm{BM.70};
\]

%dsp-level:0
\[
\ \mathrm{BM.75}(4)\text{'s opposed row};
\]

%dsp-level:0
\[
\ \mathrm{BM.77};
\]

%dsp-level:0
\[
\ \mathrm{BM.84};
\]

%dsp-level:0
\[
\ \mathrm{BM.85}\ \ —\ \ \text{each as the}\ (w,N,\mathsf g)
\]

%dsp-level:0
\[
= (1,1,1)\ \text{specialization};
\]

%dsp-level:0
\[
\neg\text{a clause of}\ \mathrm{BM.1}\text{--}\mathrm{BM.86}\ \text{is retracted, replaced or renumbered};
\]

%dsp-level:0
\[
\qquad \neg\text{a clause of Parts 1--3 is modified};
\]

%dsp-level:0
\[
\qquad \neg\text{a Census cell gains or loses a value};
\]

%dsp-level:0
\[
\text{the conditionalities}\ (\mathrm{Def\text{-}H.37},\ \mathrm{Def\text{-}H.38},\ \mathrm{M1}/\mathrm{M2},\ \mathrm{CT2.6},\ \mathrm{CT2.7},\ \mathrm{BM.8}(2),\ \mathrm{BM.59}(9)\text{'s hypothesis}\ (\mathrm H))\ \text{are inherited at every point of use}\ [\text{A §0.2a}];
\]

%dsp-level:0
\[
(5)\quad\textbf{the standing rules:}
\]

%dsp-level:0
\[
\qquad \textbf{the declarations this tier adds: none}\ —\ \neg\text{a name attachment},
\]

%dsp-level:0
\[
\ \neg\text{a read specification},
\]

%dsp-level:0
\[
\ \neg\text{a P4 addition},
\]

%dsp-level:0
\[
\ \neg\text{a registry addendum is issued or required};
\]

%dsp-level:0
\[
\neg\text{a new dimensionful scale}\ [\mathrm{P1}];
\]

%dsp-level:0
\[
\qquad \text{§Cal-BM stays empty};
\]

%dsp-level:0
\[
\qquad N,\ w,\ \mathsf g,\ \varepsilon_{\mathrm{EM}},\ e_0\ \text{are restatements of derived or Cal-F/P4-class symbols,}\ \neg\text{new members};
\]

%dsp-level:0
\[
\neg\text{a measured value and}\ \neg\text{a numerical coupling}\ [\mathrm{P5};\ \mathrm{EM.25}(7),(9);\ \text{machine: zero external tokens, zero decimal numerals, zero}\ [\mathrm{EXT}]\ \text{markers in this tier}];
\]

%dsp-level:0
\[
\text{no measured external constant is identified with any corpus quantity}\ [\mathrm{P5}];
\]

%dsp-level:0
\[
\qquad \text{species names are confined to the identification layers}\ [\mathrm{Def\text{-}BM.3}(7);\ \mathrm{Def\text{-}BM.6}(7);\ \mathrm{Def\text{-}BM.14}(8)];
\]

%dsp-level:0
\[
\text{topology labels are internal bookkeeping}\ [\mathrm{Def\text{-}H.32}(4)];
\]

%dsp-level:0
\[
\qquad \text{amplitude}\ne\text{mass}\ [\mathrm{FT.79}(5)];
\]

%dsp-level:0
\[
\qquad \text{rung}\ne\text{scale}\ [\mathrm{D0.8}];
\]

%dsp-level:0
\[
\qquad \neg\text{a currency is converted}\ [\mathrm{BM.83}(3);\ \mathrm{BM.90}(5)];
\]

%dsp-level:0
\[
\alpha_B:\quad\neg\text{added},\quad\neg\text{derived},\quad\text{barred}\ —\ \mathrm{BM.61}(8)\text{'s}\ (\mathrm O1)\text{--}(\mathrm O4)\ \text{stand}
\]

%dsp-level:0
\[
\qquad [\text{machine: the token scan of this tier}];
\]

%dsp-level:0
\[
\qquad \neg\text{a Yukawa-type object exists or is added}\ [\mathrm{EW.15};\ \mathrm{FT.145}(3)];
\]

%dsp-level:0
\[
\neg\text{a HistLog disposition imported as premise}\ —\ \text{the flicker measure is}\ \mathrm{BM.19}(1),
\]

%dsp-level:0
\[
(2)\text{'s derived row and}\ \mathrm{A4}\text{'s threshold, and item 7 is cited as a pointer};
\]

%dsp-level:0
\[
(6)\quad\textbf{the boundary inventory (every row dispositioned in full at BM.143):}\qquad\begin{array}{l|l}(\mathrm{R1})\ \ \rho_{\mathrm B}=\hat\mu'^2\mathfrak h^2/(\ln2)^2 & \text{the bridge's single undeclared dimensionless ratio}\ [\mathrm{BM.69}(5);\ \mathrm{BM.82}(6)]\\(\mathrm{R2})\ \ \text{the loop realization} & \text{a clean-enclosure vertex-figure loop as a ledger configuration}\ [\mathrm{BM.67}(9),(10);\ \text{Census FT.63}(4)\ \mathrm{Rmk}]\\(\mathrm{R3})\ \ \text{the flavour composition} & \text{which three of the six states a sealed class occupies}\ [\mathrm{BM.28}(7);\ \mathrm{BM.75}(6)]\\(\mathrm{R4})\ \ w(\varepsilon_q)\ \text{in general} & \text{the ensemble's configuration distribution beyond the per-datum specialization, and the transport register's ensemble average}\ [\mathrm{BM.89}(8);\ \mathrm{BM.91}(2)]\\(\mathrm{R5})\ \ \iota\ \text{and the interior multipoles} & \text{the internal charge resolution the end signs cannot carry}\ [\mathrm{BM.46}(2);\ \mathrm{BM.87}(3),(7)]\\ e_0 & \text{the numerical value of the electric coupling}\ —\ \text{one calibration datum}\ [\mathrm{EM.25}(7);\ \mathrm{BM.62}(12);\ \mathrm{BM.84}(9)]\\\mathsf g,\ N & \text{the curvature composition and the lapse as solution data}\ [\mathrm{BM.88}(4),(8);\ \mathrm{CT2.7};\ \mathrm{BM.34}(2)]\\ m_\nu\ \text{absolute} & \neg\text{derived}\ [\mathrm{FT.143}(3);\ \mathrm{BM.83}(8)]\\ i_\nu,\ \text{the mixing matrix} & \text{unresolved}\ [\mathrm{BM.77}(2);\ \mathrm{BM.75}(6);\ \mathrm{FT.145}(3)]\\\text{the generation count as a selection} & \text{three lines derived, three selected}\ \neg\text{derived}\ [\mathrm{BM.60}(1)]\end{array}
\]

%dsp-level:0
\[
(7)\quad\textbf{the adjudication:}
\]

%dsp-level:0
\[
\qquad \boxed{\ \text{the EM self-energy is an ensemble average of one derived monomial, its coefficient is}\ \tfrac12\ \text{by two distribution-free routes, its coupling runs geometrically, and its classical distribution factor is void}\ }
\]

%dsp-level:0
\[
\boxed{\ \text{the neutron--proton splitting keeps its exact trichotomy, gains a parameter-free bar}\ Nw\mathsf g\,\varepsilon_{\mathrm{EM}}\le\tfrac18,\ \text{and its ratio stays below}\ \tfrac{10}9\ \text{everywhere}\ }
\]

%dsp-level:0
\[
\textbf{one classical artifact left the corpus, two derived factors entered, and nothing was fitted:}
\]

%dsp-level:0
\[
\qquad \text{the void row is a coefficient, the new rows are a lapse, a curvature lock, a nearby fraction and a crystallization weight};
\]

%dsp-level:0
\[
\text{item 7 stands with its residue named};
\]

%dsp-level:0
\[
\qquad \text{the scales stand undeclared};
\]

%dsp-level:0
\[
\qquad \text{nothing measured entered anywhere}.
\]

\textbf{Req.} D0.8 (Parts 1–2); P1, P4, P5 (Parts 1–2); A §0.2a, CT2.7, EM.25(7), EM.25(9), EW.15 (Part 2); Def-H.32(4), Def-H.37, Def-H.38, FT.79(2), FT.79(5), FT.84(2), FT.143(3), FT.145(3) (Part 3); Census FT.63(4), FT.76(3) (row-wise); Def-BM.3(7), Def-BM.6(7), Def-BM.12(4), Def-BM.14(8), BM.8(2), BM.19(1), BM.19(2), BM.28(7), BM.34(2), BM.46(1), BM.46(2), BM.46(3), BM.46(4), BM.46(5), BM.47, BM.48, BM.49(1), BM.49(2), BM.50, BM.52(4), BM.52(5), BM.55(1), BM.55(2), BM.59(9), BM.60(1), BM.61(8), BM.62, BM.63, BM.65(3), BM.65(5), BM.66(3), BM.67(9), BM.67(10), BM.69(4), BM.69(5), BM.70, BM.75(4), BM.75(6), BM.77, BM.82(6), BM.83(3), BM.83(5), BM.83(8), BM.84, BM.84(9), BM.85, BM.86, BM.87, BM.88, BM.89, BM.89(8), BM.90, BM.90(5), BM.91, BM.91(2), BM.92 (this part); Cal-F; the machine sections §A–§G green [math; machine: \texttt{checks/\allowbreak{}verify\_\allowbreak{}H\_\allowbreak{}em\_\allowbreak{}self\_\allowbreak{}energy\_\allowbreak{}correction.py} §G].

\textbf{Pf.}

%dsp-level:0
\[
(1)\text{--}(2):\ \text{assembly of the adjudicated rows of}\ \mathrm{BM.87}\text{--}\mathrm{BM.92}\ [\text{machine: §A--§F green}];
\]

%dsp-level:0
\[
\qquad \text{the inventory rows re-listed from their own tags}\ [\text{machine}]
\]

%dsp-level:0
\[
(3):\ \text{the derived/}\neg\text{derived partition by the symbol census of the tier}\ [\text{machine}];
\]

%dsp-level:0
\[
\qquad \text{the obstruction rows re-cited from}\ \mathrm{BM.61}(8),
\]

%dsp-level:0
\[
\ \mathrm{BM.69}(4),
\]

%dsp-level:0
\[
(5),
\]

%dsp-level:0
\[
\ \text{Census FT.76}(3),
\]

%dsp-level:0
\[
\ \mathrm{FT.79}(2),
\]

%dsp-level:0
\[
\ \mathrm{FT.84}(2)
\]

%dsp-level:0
\[
(4):\ \text{the clause census of}\ \mathrm{BM.1}\text{--}\mathrm{BM.86}\ \text{against retraction, replacement and renumbering}\ [\text{machine}];
\]

%dsp-level:0
\[
\qquad \text{the Census-cell invariance scan}\ [\text{machine}];
\]

%dsp-level:0
\[
\qquad \text{the conditionality rows re-cited at their loci}
\]

%dsp-level:0
\[
(5):\ \text{the standing-rule scans: external tokens, decimal numerals, unit tokens, registry tokens, declaration lines, the barred object}\ \alpha_B\ [\text{machine}];
\]

%dsp-level:0
\[
\qquad (6):\ \text{the ten open rows re-cited from their own tags}\ [\text{machine}]
\]

%dsp-level:0
\[
(7):\ \text{assembly}.\ \blacksquare
\]

\textbf{Phys.} The tier began with an accusation and ended with a smaller corpus and a sharper one. The accusation was right: a shell or a ball is a shape the graph cannot hold, because a charge here is an integer on the end of a chain and the ends carry no positions. But the accusation cost nothing that mattered, because the same value falls out twice without a distribution — once from the flux outside the envelope, which Gauss fixes from the total charge alone, and once from the potential at a single radius, which needs no volume at all. Both say one half. So the coefficient the corpus had priced twice and adopted never is now derived once, and the interior term that needed a density is named as a residue instead of being guessed. What replaces the classical integral is an average over the flicker ensemble, and the average brings three exact dimensionless companions: the lapse of the emergent metric, the nearby-space fraction one minus one over k, and the calibration lock's curvature monomial, which is two pi times the square of the edge count times the logarithm of two and carries no h, no cosmological term and no declared parameter. The coupling runs with all three, vanishes at the member's own envelope, and grows with energy for purely geometric reasons. The splitting that results is the old one with three symbols in its electromagnetic term, its critical strength pushed outward by every failure to crystallize, its ratio bounded by ten ninths at every admissible value of everything, and its trichotomy intact. The decay stays closed where a read closes it and open where a read opens it, and the one thing both reads still agree on is which baryon decays.

\medskip
\subsection*{\texorpdfstring{B.1.29\quad BM.100 — The status of the flicker integral and the rest-mass read: adjudication and scope}{BM.100 B.1.29}}
\addcontentsline{toc}{subsection}{B.1.29 BM.100}
\phantomsection\label{tag:BM-100}
{\small \textit{Record.}\quad Source: Part 4 (Part4\_Baryomesonic\_era.md); tier: Tier BM.19 — The Flicker Integral, the Rest-Mass Q-Value, and the Relating Equation: the Materialized Ensemble, the Vertex Superposition, the Rung-One $\textbackslash{}nu$ Line, and the Bridge in Four Exact Forms; canonical tag \texttt{BM.100}. Reproduced from the revision-1 canon; a cross-reference stub marks the former locus in the main corpus.}
\medskip
\textbf{Thm.}

%dsp-level:0
\[
(1)\quad\textbf{the answers:}\qquad\begin{array}{l|l}\text{the snapshot read} & \mathrm{BM.94}:\ \textbf{fixed}\ —\ E_{\mathrm{int}}=6\ \text{is one point of an exact distribution with weight}\ p^2/D<1;\ \text{the integral}\ \textstyle\sum_RP(R)E_{\mathrm{int}}(R)\ \text{replaces it};\\ & \text{the derived distribution}\ \bigl((1{-}p)^3,\,3p(1{-}p)^2,\,3p^2(1{-}p),\,p^3\bigr)\ \text{on}\ \{0,2,4,6\};\ \text{the materialization restriction}\ E\ge2\ \text{derived from}\ \mathrm{FT.144}(4)\ \text{and}\ \mathrm{Def\text{-}H.4};\\ & w_{\mathrm{mat}}=1/(1{+}u{+}u^2)\in(\tfrac13,1];\quad\langle E\rangle_{\mathrm{mat}}\in(2,6];\quad\sigma/\langle E\rangle<\tfrac12\ \text{on}\ 175<256\\\text{the vertex constituent} & \mathrm{BM.95}:\ \textbf{a superposition, derived}\ —\ \text{the fibre weights}\ \prod_im_{q_i}/W(Q)\ \text{exact rationals};\ \Sigma_{\mathrm{cl}}=H+\langle\ln\prod_im\rangle;\\ & \ln\tfrac43=\bigl[\tfrac95\ln2-\ln3\bigr]+\bigl[\tfrac15\ln2\bigr]\ \text{two positive parts on}\ 512>243;\ \max P=2/5<1/2\ \text{on}\ 16<20\\\text{the}\ \nu\ \text{line} & \mathrm{BM.96}:\ \textbf{left ray only, derived}\ —\ \text{occupancy}\ (1,0);\ \text{the complement is the vacuum sector with no accessible pole at any rung};\\ & m_\nu^{\mathrm{rest}\,2}=\hat\mu_\nu^2(\ell_c^{(\nu)})^2G(c(i_\nu))\ \text{at the derived}\ \delta_\nu=\tfrac18;\ G_2>G_1>A_0>G_0;\ m_\nu^{\mathrm{rest}}>0\ \text{strictly};\\ & \neg\text{a seesaw-type mechanism on four bars}\ (\mathrm S1)\text{--}(\mathrm S4);\ \neg\text{a right-ray matter line};\ \text{the scale undeclared}\\\text{the rest-mass Q-value} & \mathrm{BM.97}:\ \textbf{computed on rest masses}\ —\ \text{the rest mass is the pole mass};\ \text{the specified read is flicker-blind};\\ & \text{the pure-pole read formed with}\ \Delta_g^{\mathrm{crit}}=s(2M_0+s);\ \text{the trichotomy total with integer-certified witnesses on both sides}\\\text{the relating equation} & \mathrm{BM.98}:\ \textbf{closed in form}\ —\ \mathcal R=\rho_{\mathrm B}D_S/(4E^2)=(\pi\ln2/2)\rho_{\mathrm B}D_S/G=(\hat\mu'^2D_Sr^2/4)(1-f)^2;\\ & \mathcal R_{\mathrm{mat}}/\mathcal R=D^2<9;\quad G^{\mathrm{eff}}\in(8\pi\ln2,\,72\pi\ln2];\quad\mathcal R\ \text{strictly decreasing in the running coupling};\quad(\mathrm{R1})\ \text{stands}\\\text{the splitting and the decay} & \mathrm{BM.99}:\ \textbf{recomputed on the materialized read}\ —\ \text{the two-sided parameter-free bar}\ (\tfrac18,\ \tfrac5{12});\ 1<m_n/m_p<\ln80/\ln60<\tfrac{15}{14};\\ & Q_\beta\mathfrak h\le-\ln3\ \text{uniformly};\ -\ln3/\ln80\in(-\tfrac13,-\tfrac14);\ \text{the dissolution row rescoped}\end{array}
\]

%dsp-level:0
\[
(2)\quad\textbf{the derived inventory (the new exact vocabulary of this tier):}
\]

%dsp-level:0
\[
\qquad D=3-3p+p^2=1+u+u^2;
\]

%dsp-level:0
\[
\qquad w_{\mathrm{mat}}=\frac1D;
\]

%dsp-level:0
\[
\qquad \langle E\rangle_{\mathrm{mat}}=\frac6D;
\]

%dsp-level:0
\[
\qquad \operatorname{Var}_{\mathrm{mat}}(E)=\frac{12p(1-p)(3-2p)}{D^2};
\]

%dsp-level:0
\[
\frac{\operatorname{Var}_{\mathrm{mat}}}{\langle E\rangle_{\mathrm{mat}}^2}=\frac{p(1-p)(3-2p)}3\le\frac{10+7\sqrt7}{162}<\frac15;
\]

%dsp-level:0
\[
\qquad \langle E\rangle_{\mathrm{mat}}\bigl\langle E^{-1}\bigr\rangle_{\mathrm{mat}}=\frac{18-27p+11p^2}{2D^2}\in\Bigl[1,\frac65\Bigr);
\]

%dsp-level:0
\[
\qquad \Delta x_{\mathrm B}^{\mathrm{eff}}=D\,\Delta x_{\mathrm B}^{\mathrm{cryst}};
\]

%dsp-level:0
\[
P(\text{fibre}\,|\,Q)=\frac{\prod_im_{q_i}}{W(Q)};
\]

%dsp-level:0
\[
\qquad H(0)=\ln5+\frac15\ln2;
\]

%dsp-level:0
\[
\qquad H(\pm1)=\ln3+\ln5-\frac85\ln2;
\]

%dsp-level:0
\[
\qquad \ln\frac43=\Bigl(\frac95\ln2-\ln3\Bigr)+\frac15\ln2;
\]

%dsp-level:0
\[
G\bigl(c(i)\bigr)=A_0-B_0\cos\Bigl(\frac18-\frac{2\pi i}3\Bigr);\qquad G_2>G_1>A_0>G_0;\qquad\Delta_g^{\mathrm{crit}}=s(2M_0+s);\qquad s^{\mathrm{crit}}=\sqrt{M_0^2+\Delta_g}-M_0;
\]

%dsp-level:0
\[
\mathcal R
\]

%dsp-level:0
\[
= \frac{\rho_{\mathrm B}D_S}{4E^2}
\]

%dsp-level:0
\[
= \frac{\pi\ln2}2\cdot\frac{\rho_{\mathrm B}D_S}{G}
\]

%dsp-level:0
\[
= \frac{\hat\mu'^2D_Sr^2}4(1-f)^2
\]

%dsp-level:0
\[
= \frac{\hat\mu'^2D_Sr^2}4\Bigl(1-\frac{e^2}{e_0^2N\mathsf g}\Bigr)^{\!2};
\]

%dsp-level:0
\[
\qquad \frac{\mathcal R_{\mathrm{mat}}}{\mathcal R}=D^2<9;
\]

%dsp-level:0
\[
\qquad G^{\mathrm{eff}}=\frac{G}{D^2};
\]

%dsp-level:0
\[
\varepsilon^{\mathrm{crit}}_{\mathrm{mat}}=\frac{2-\log_23}{3Nw_{\mathrm{mat}}\mathsf g};
\]

%dsp-level:0
\[
\qquad N\mathsf g\,\varepsilon_{\mathrm{EM}}\le\frac18\implies\Delta m>0;
\]

%dsp-level:0
\[
\qquad N\mathsf g\,\varepsilon_{\mathrm{EM}}\ge\frac5{12}\implies\Delta m<0;
\]

%dsp-level:0
\[
\qquad \frac{m_n}{m_p}<\frac{\ln80}{\ln60}<\frac{15}{14};
\]

%dsp-level:0
\[
\textbf{the certificates:}
\]

%dsp-level:0
\[
\quad 175<256,
\]

%dsp-level:0
\[
\quad 512>243,
\]

%dsp-level:0
\[
\quad 16<20,
\]

%dsp-level:0
\[
\quad 12<15,
\]

%dsp-level:0
\[
\quad 2^{13}>3^8,
\]

%dsp-level:0
\[
\quad 3^{12}>2^{19},
\]

%dsp-level:0
\[
\quad 12<40,
\]

%dsp-level:0
\[
\quad 2^{28}>3^{16}\!\cdot\!5,
\]

%dsp-level:0
\[
\quad 2^{26}<3^{15}\!\cdot\!5,
\]

%dsp-level:0
\[
\quad 150<154,\qquad \quad 135<140,
\]

%dsp-level:0
\[
81>80,
\]

%dsp-level:0
\[
\quad 80>27,
\]

%dsp-level:0
\[
\quad 64<405,
\]

%dsp-level:0
\[
\quad 5000>101,
\]

%dsp-level:0
\[
\quad 11<16,
\]

%dsp-level:0
\[
\quad 1875>1600,
\]

%dsp-level:0
\[
\quad 75>64,
\]

%dsp-level:0
\[
\quad 6560>2,
\]

%dsp-level:0
\[
\quad 7921>7500,
\]

%dsp-level:0
\[
\quad 192<196,\qquad \quad 290<294
\]

%dsp-level:0
\[
\qquad [\text{machine: every row exact}];
\]

%dsp-level:0
\[
(3)\quad\textbf{the status adjudication:}
\]

%dsp-level:0
\[
\qquad \boxed{\ \textbf{the snapshot read is replaced by the flicker integral, whose materialized mean is bounded below by the digon's two}\ }
\]

%dsp-level:0
\[
\boxed{\ \textbf{the Q-value is read on rest masses}\ —\ \text{the pole/transport species}\ —\ \textbf{and that read is flicker-blind, so the fix and the specification are orthogonal acts}\ }
\]

%dsp-level:0
\[
\boxed{\ \textbf{the}\ \nu\ \textbf{line occupies one ray, is strictly massive in form, and admits no seesaw-type mechanism on four independent bars}\ }
\]

%dsp-level:0
\[
\boxed{\ \textbf{the relating equation is derived in four exact forms, flicker-weighted, bounded by an exact nine, and closed up to the single undeclared ratio}\ (\mathrm{R1})\ }
\]

%dsp-level:0
\[
\textbf{the level-exactness:}
\]

%dsp-level:0
\[
\qquad \text{derived:}\ w_{\mathrm{mat}}\ \text{and its bound, every moment and enclosure, the fibre weights and the entropy decomposition, the}\ \nu\ \text{line's order and positivity, the occupancy row, the four bars, the four bridge forms, the bound}\ 9,
\]

%dsp-level:0
\[
\ G^{\mathrm{eff}}\text{'s enclosure, both bars, both tightenings, every certificate};
\]

%dsp-level:0
\[
\neg\text{derived:}\ \rho_{\mathrm B}\ (\mathrm{R1}),
\]

%dsp-level:0
\[
\ \ \text{the flavour composition}\ (\mathrm{R3}),
\]

%dsp-level:0
\[
\ \ w(\varepsilon_q)\ \text{in general}\ (\mathrm{R4}),
\]

%dsp-level:0
\[
\ \ \iota\ (\mathrm{R5}),
\]

%dsp-level:0
\[
\ \ \text{the Jensen residue}\ (\mathrm{R6}),
\]

%dsp-level:0
\[
\ \ e_0,
\]

%dsp-level:0
\[
\ \mathsf g,
\]

%dsp-level:0
\[
\ N,
\]

%dsp-level:0
\[
\ \hat\mu_\nu,
\]

%dsp-level:0
\[
\ \ell_c^{(\nu)},
\]

%dsp-level:0
\[
\ \mu_\nu,
\]

%dsp-level:0
\[
\ i_\nu,
\]

%dsp-level:0
\[
\ \text{the mixing matrix};
\]

%dsp-level:0
\[
\textbf{the mass-form bridge disposition:}
\]

%dsp-level:0
\[
\qquad \textbf{the equation is closed in form and its residue is localized to one dimensionless ratio};
\]

%dsp-level:0
\[
\qquad \mathrm{BM.69}(4)\text{'s non-derivation theorem stands};
\]

%dsp-level:0
\[
\qquad \neg\text{a value for}\ \rho_{\mathrm B}\ \text{is derived, declared or fitted};
\]

%dsp-level:0
\[
\textbf{the two ceilings and the obstructions stand:}
\]

%dsp-level:0
\[
\qquad \alpha_B\ \text{barred on}\ \mathrm{BM.61}(8)\text{'s four grounds};
\]

%dsp-level:0
\[
\qquad \text{deficit}\ 12\ [\mathrm{Census\ FT.76}(3);\ \mathrm{FT.79}(2);\ \mathrm{FT.84}(2)];
\]

%dsp-level:0
\[
\qquad \text{deficits}\ 4,\ 6\ \text{unchanged};
\]

%dsp-level:0
\[
(4)\quad\textbf{the scope rows (recorded,}\ \neg\textbf{silent):}
\]

%dsp-level:0
\[
\qquad \neg\text{a clause of}\ \mathrm{BM.1}\text{--}\mathrm{BM.93}\ \text{is retracted, replaced or renumbered};
\]

%dsp-level:0
\[
\qquad \neg\text{a clause of Parts 1--3 is modified};
\]

%dsp-level:0
\[
\qquad \neg\text{a Census cell gains or loses a value};
\]

%dsp-level:0
\[
\textbf{standing at their own locus:}\ \mathrm{Def\text{-}BM.5}(4),
\]

%dsp-level:0
\[
\ \mathrm{BM.19}(1)\text{--}(6),
\]

%dsp-level:0
\[
\ \mathrm{BM.20}(3),
\]

%dsp-level:0
\[
\ \mathrm{BM.21},
\]

%dsp-level:0
\[
\ \mathrm{BM.28}(2),
\]

%dsp-level:0
\[
(4),
\]

%dsp-level:0
\[
(5),
\]

%dsp-level:0
\[
\ \mathrm{BM.39},
\]

%dsp-level:0
\[
\ \mathrm{BM.41},
\]

%dsp-level:0
\[
\ \mathrm{BM.45},
\]

%dsp-level:0
\[
\ \mathrm{BM.59}(4),
\]

%dsp-level:0
\[
(9),
\]

%dsp-level:0
\[
\ \mathrm{BM.69}(2)\text{--}(5),
\]

%dsp-level:0
\[
\ \mathrm{BM.70},
\]

%dsp-level:0
\[
\ \mathrm{BM.77},
\]

%dsp-level:0
\[
\ \mathrm{BM.84},
\]

%dsp-level:0
\[
\ \mathrm{BM.85},
\]

%dsp-level:0
\[
\ \mathrm{BM.86}(4),
\]

%dsp-level:0
\[
\mathrm{BM.88}(2)\text{--}(6),
\]

%dsp-level:0
\[
\ \mathrm{BM.89}(4),
\]

%dsp-level:0
\[
(5),
\]

%dsp-level:0
\[
(8),
\]

%dsp-level:0
\[
\ \mathrm{BM.90}(4),
\]

%dsp-level:0
\[
(5),
\]

%dsp-level:0
\[
\ \mathrm{BM.91}(2),
\]

%dsp-level:0
\[
(4),
\]

%dsp-level:0
\[
(5),
\]

%dsp-level:0
\[
(7),
\]

%dsp-level:0
\[
(8),
\]

%dsp-level:0
\[
\ \mathrm{BM.92}(1),
\]

%dsp-level:0
\[
(2),
\]

%dsp-level:0
\[
(3),
\]

%dsp-level:0
\[
(5),
\]

%dsp-level:0
\[
\ \mathrm{BM.93}\ \ —\ \ \text{each re-derived at the specialization it carries and none weakened};
\]

%dsp-level:0
\[
\mathrm{BM.91}(4)\text{'s dissolution row is rescoped at}\ \mathrm{BM.99}(7);
\]

%dsp-level:0
\[
\qquad \mathrm{BM.89}(4)\text{'s}\ w\ \text{is the unconditioned weight and}\ w_{\mathrm{mat}}\ \text{is its materialized companion},
\]

%dsp-level:0
\[
\ \neg\text{a replacement};
\]

%dsp-level:0
\[
\text{the conditionalities}\ (\mathrm{Def\text{-}H.37},\ \mathrm{Def\text{-}H.38},\ \mathrm{M1}/\mathrm{M2},\ \mathrm{CT2.6},\ \mathrm{CT2.7},\ \mathrm{BM.8}(2),\ \mathrm{BM.59}(9)\text{'s hypothesis}\ (\mathrm H))\ \text{are inherited at every point of use}\ [\text{A §0.2a}];
\]

%dsp-level:0
\[
\text{the marked specializations}\ (\text{the pair-Bernoulli distribution},\ \text{the witness scale points})\ \text{are marked},
\]

%dsp-level:0
\[
\ \neg\text{adopted as general measures or as declared scales};
\]

%dsp-level:0
\[
(5)\quad\textbf{the standing rules:}
\]

%dsp-level:0
\[
\qquad \textbf{the declarations this tier adds: none}\ —\ \neg\text{a name attachment},
\]

%dsp-level:0
\[
\ \neg\text{a read specification},
\]

%dsp-level:0
\[
\ \neg\text{a P4 addition},
\]

%dsp-level:0
\[
\ \neg\text{a registry addendum is issued or required};
\]

%dsp-level:0
\[
\neg\text{a new dimensionful scale}\ [\mathrm{P1}];
\]

%dsp-level:0
\[
\qquad \text{§Cal-BM stays empty};
\]

%dsp-level:0
\[
\qquad p,
\]

%dsp-level:0
\[
\ u,
\]

%dsp-level:0
\[
\ D,
\]

%dsp-level:0
\[
\ w_{\mathrm{mat}},
\]

%dsp-level:0
\[
\ \Delta_g^{\mathrm{crit}},
\]

%dsp-level:0
\[
\ G^{\mathrm{eff}},
\]

%dsp-level:0
\[
\ \mathcal R_{\mathrm{mat}}\ \ \text{are in-theorem abbreviations of derived objects},
\]

%dsp-level:0
\[
\ \neg\text{new members};
\]

%dsp-level:0
\[
\neg\text{a measured value and}\ \neg\text{a numerical coupling}\ [\mathrm{P5};\ \mathrm{EM.25}(7),(9);\ \text{machine: zero external tokens, zero decimal numerals, zero}\ [\mathrm{EXT}]\ \text{markers in this tier}];
\]

%dsp-level:0
\[
\text{no measured external constant is identified with any corpus quantity}\ [\mathrm{P5}];
\]

%dsp-level:0
\[
\qquad \text{species names are confined to the identification layers}\ [\mathrm{Def\text{-}BM.3}(7);\ \mathrm{Def\text{-}BM.6}(7);\ \mathrm{Def\text{-}BM.14}(8)];
\]

%dsp-level:0
\[
\text{topology labels are internal bookkeeping}\ [\mathrm{Def\text{-}H.32}(4)];
\]

%dsp-level:0
\[
\qquad \text{amplitude}\ne\text{mass}\ [\mathrm{FT.79}(5)];
\]

%dsp-level:0
\[
\qquad \text{rung}\ne\text{scale}\ [\mathrm{D0.8}];
\]

%dsp-level:0
\[
\qquad \neg\text{a currency is converted}\ [\mathrm{BM.83}(3);\ \mathrm{BM.90}(5);\ \mathrm{BM.98}(6)];
\]

%dsp-level:0
\[
\alpha_B:\quad\neg\text{added},\quad\neg\text{derived},\quad\text{barred}\ —\ \mathrm{BM.61}(8)\text{'s}\ (\mathrm O1)\text{--}(\mathrm O4)\ \text{stand}\qquad[\text{machine: the token scan of this tier}];
\]

%dsp-level:0
\[
\neg\text{a Yukawa-type object exists or is added},
\]

%dsp-level:0
\[
\qquad \neg\text{a seesaw-type mechanism is suggested},
\]

%dsp-level:0
\[
\qquad \neg\text{a right-ray matter line is added}
\]

%dsp-level:0
\[
\qquad [\mathrm{EW.15};\ \mathrm{FT.145}(3);\ \mathrm{BM.96}(4),(6);\ \text{machine: every occurrence sits in a barred-object row}];
\]

%dsp-level:0
\[
\neg\text{a HistLog disposition imported as premise}\ —\ \text{the flicker measure is}\ \mathrm{BM.19}(1),
\]

%dsp-level:0
\[
(2)\text{'s derived row and}\ \mathrm{A4}\text{'s threshold},
\]

%dsp-level:0
\[
\ \text{the materialization restriction is}\ \mathrm{FT.144}(4)\text{'s domain},
\]

%dsp-level:0
\[
\ \text{and items 7 and 20 are cited as pointers};
\]

%dsp-level:0
\[
(6)\quad\textbf{the boundary inventory (every row dispositioned in full at BM.143):}\qquad\begin{array}{l|l}(\mathrm{R1})\ \ \rho_{\mathrm B}=\hat\mu'^2\mathfrak h^2/(\ln2)^2 & \text{the bridge's single undeclared dimensionless ratio}\ [\mathrm{BM.69}(5);\ \mathrm{BM.98}(5)]\\(\mathrm{R2})\ \ \text{the loop realization} & \text{a clean-enclosure vertex-figure loop as a ledger configuration}\ [\mathrm{BM.67}(9),(10);\ \mathrm{Census\ FT.63}(4)]\\(\mathrm{R3})\ \ \text{the flavour composition} & \text{which three of the six states a sealed class occupies}\ [\mathrm{BM.28}(7)]\\(\mathrm{R4})\ \ w(\varepsilon_q)\ \text{in general} & \text{the ensemble's configuration distribution beyond the pair specialization, and the transport register's ensemble average}\ [\mathrm{BM.89}(8);\ \mathrm{BM.91}(2);\ \mathrm{BM.99}(1)]\\(\mathrm{R5})\ \ \iota\ \text{and the interior multipoles} & \text{the internal charge resolution the end signs cannot carry}\ [\mathrm{BM.46}(2);\ \mathrm{BM.87}(3)]\\(\mathrm{R6})\ \ \text{the Jensen residue} & \langle E\rangle_{\mathrm{mat}}\langle E^{-1}\rangle_{\mathrm{mat}}-1:\ \text{the two formed reads of the reciprocal register}\ [\mathrm{BM.94}(8);\ \mathrm{BM.98}(5)]\\ e_0 & \text{the numerical value of the electric coupling}\ [\mathrm{EM.25}(7);\ \mathrm{BM.84}(9)]\\ \mathsf g,\ N & \text{the curvature composition and the lapse as solution data}\ [\mathrm{BM.88}(4),(8);\ \mathrm{CT2.7}]\\ m_\nu\ \text{absolute},\ \mu_\nu & \neg\text{derived}\ [\mathrm{FT.143}(3);\ \mathrm{BM.96}(3)]\\ i_\nu,\ \text{the mixing matrix} & \text{unresolved}\ [\mathrm{BM.77}(2);\ \mathrm{FT.145}(3)]\\ \text{the generation count as a selection} & \text{three lines derived, three selected}\ \neg\text{derived}\ [\mathrm{BM.60}(1)]\\ \text{the two mass ceilings} & \text{deficit}\ 12\ \text{stands}\ [\mathrm{Census\ FT.76}(3);\ \mathrm{FT.79}(2);\ \mathrm{FT.84}(2)]\end{array}
\]

%dsp-level:0
\[
(7)\quad\textbf{the adjudication:}
\]

%dsp-level:0
\[
\qquad \boxed{\ \textbf{one read was replaced by an integral, one Q-value was moved onto the rest-mass species, one line was shown to occupy a single ray, and one ledger item was closed in form}\ }
\]

%dsp-level:0
\[
\textbf{the cost:}
\]

%dsp-level:0
\[
\qquad \text{one new named residue}\ (\mathrm{R6}),
\]

%dsp-level:0
\[
\ \ \text{one marked specialization}\ (\text{the pair distribution}),
\]

%dsp-level:0
\[
\ \ \text{zero declarations},
\]

%dsp-level:0
\[
\ \ \text{zero parameters},
\]

%dsp-level:0
\[
\ \ \text{zero scales},
\]

%dsp-level:0
\[
\ \ \text{zero measured values};
\]

%dsp-level:0
\[
\textbf{the gain:}
\]

%dsp-level:0
\[
\qquad \text{an exact distribution with bounded relative width},
\]

%dsp-level:0
\[
\ \ \text{a weight bounded below by a third},
\]

%dsp-level:0
\[
\ \ \text{two tightenings of the ratio and Q-value bounds},
\]

%dsp-level:0
\[
\ \ \text{a two-sided parameter-free bar},
\]

%dsp-level:0
\[
\ \ \text{a bridge bounded by nine and readable in the coupling};
\]

%dsp-level:0
\[
\text{item 7 stands with its residue named and its equation derived};
\]

%dsp-level:0
\[
\qquad \text{the scales stand undeclared};
\]

%dsp-level:0
\[
\qquad \text{nothing measured entered anywhere}.
\]

\textbf{Req.} D0.8 (Parts 1–2); A §0.2a, CT2.7, EM.25(7), EM.25(9), EW.15 (Part 2); P1, P4, P5 (Parts 1–2); Def-H.32(4), Def-H.37, Def-H.38, FT.73(2), FT.73(3), FT.76, FT.79(2), FT.79(5), FT.84(2), FT.143(3), FT.144(4), FT.145(3) (Part 3); Census FT.28, FT.63(4), FT.76(3) (row-wise); Def-BM.3(7), Def-BM.5(4), Def-BM.6(7), Def-BM.14(8), BM.8(2), BM.19, BM.20(3), BM.21, BM.28(7), BM.39, BM.41, BM.45, BM.46(2), BM.59(9), BM.60(1), BM.61(8), BM.67(9), BM.67(10), BM.69(2), BM.69(3), BM.69(4), BM.69(5), BM.70, BM.77, BM.84, BM.84(9), BM.85, BM.86(4), BM.87(3), BM.88(4), BM.88(8), BM.89(4), BM.89(8), BM.90(4), BM.90(5), BM.91(2), BM.92(3), BM.93, BM.94, BM.95, BM.96, BM.97, BM.98, BM.99 (this part); Cal-F; the machine sections §A–§G green [math; machine: \texttt{checks/\allowbreak{}verify\_\allowbreak{}H\_\allowbreak{}flicker\_\allowbreak{}integral\_\allowbreak{}restmass.py} §G].

\textbf{Pf.}

%dsp-level:0
\[
(1)\text{--}(2):\ \text{assembly of the adjudicated rows of}\ \mathrm{BM.94}\text{--}\mathrm{BM.99}\ [\text{machine: §A--§F green}];
\]

%dsp-level:0
\[
\qquad \text{the inventory rows re-listed from their own tags}\ [\text{machine}]
\]

%dsp-level:0
\[
(3):\ \text{the derived/}\neg\text{derived partition by the symbol census of the tier}\ [\text{machine}];
\]

%dsp-level:0
\[
\qquad \mathrm{BM.69}(4)\text{'s four grounds and}\ \mathrm{BM.61}(8)\text{'s obstruction re-cited};
\]

%dsp-level:0
\[
\qquad \text{the ceiling rows re-cited from}\ \mathrm{Census\ FT.76}(3),\ \mathrm{FT.79}(2),\ \mathrm{FT.84}(2)
\]

%dsp-level:0
\[
(4):\ \text{the clause census of}\ \mathrm{BM.1}\text{--}\mathrm{BM.93}\ \text{against retraction, replacement and renumbering}\ [\text{machine}];
\]

%dsp-level:0
\[
\qquad \text{the Census-cell invariance scan}\ [\text{machine}];
\]

%dsp-level:0
\[
\qquad \text{the conditionality rows re-cited at their loci}\ [\text{A §0.2a}]
\]

%dsp-level:0
\[
(5):\ \text{the standing-rule scans: external tokens, decimal numerals, unit tokens, registry tokens, declaration lines, the barred object}\ \alpha_B,
\]

%dsp-level:0
\[
\ \text{the seesaw-type and right-ray tokens}\ [\text{machine}];
\]

%dsp-level:0
\[
\qquad (6):\ \text{the open rows re-cited from their own tags}\ [\text{machine}]
\]

%dsp-level:0
\[
(7):\ \text{assembly}.\ \blacksquare
\]

\textbf{Phys.} Four acts, one ledger. The snapshot read is replaced by an integral over an exact distribution whose materialized mean cannot fall below the digon's two edges and whose relative width cannot exceed a half; the weight that prices the failure to crystallize is now bounded below by a third, and that one bound converts a weight-dependent inequality into a two-sided bar with no parameter in it at all. The Q-value moves onto the rest masses, which in this corpus means the pole masses, and the move turns out to be free of consequences for the flicker, because the specified read's alphabet contains no edge count and no threshold — the two acts are orthogonal, and every witness the specification certified still certifies. The neutral line occupies a single ray, is strictly massive in form from a derived channel angle, and cannot be given a partner: four independent bars, each a theorem of an earlier tier, close every route to a seesaw-type mechanism, and the smallness of the line stays a question about undeclared scale rather than an invitation to invent one. The relating equation is derived in four forms — count, curvature, coupling, flicker-weighted — with its enhancement bounded by an exact nine and its curvature argument confined between two computed cells, and it closes in the only sense this corpus allows: the equation is a theorem and its one undeclared number remains undeclared, with four grounds standing against deriving it. The cost is one named residue and one marked specialization. Nothing was declared, nothing was fitted, and nothing measured entered anywhere.

\medskip
\subsection*{\texorpdfstring{B.1.30\quad BM.107 — The status of the six resolutions: adjudication and scope}{BM.107 B.1.30}}
\addcontentsline{toc}{subsection}{B.1.30 BM.107}
\phantomsection\label{tag:BM-107}
{\small \textit{Record.}\quad Source: Part 4 (Part4\_Baryomesonic\_era.md); tier: Tier BM.20 — Six Objective Resolutions: the Energy-Scale Graph, the Dimensional Bridge, the Eigenvalue Composition, the Curvature Composition, the General Crystallization Weight, and the Kind Structure of the Arrow; canonical tag \texttt{BM.107}. Reproduced from the revision-1 canon; a cross-reference stub marks the former locus in the main corpus.}
\medskip
\textbf{Thm.}

%dsp-level:0
\[
(1)\quad\textbf{the answers:}\qquad\begin{array}{l|l}\text{the energy-scale graph} & \mathrm{BM.101}:\ \textbf{derived}\ —\ Q^{\mathrm{inv}}=-Q^{\mathrm{fwd}}\ \text{identically};\ \Delta m(E)\,\mathfrak h=\ln\tfrac43-3w\ln2\cdot\varepsilon_{\mathrm{eff}};\ \text{the crossing strength}\ \varepsilon_\ast=(2-\log_23)/(3w)\in(\tfrac18,\tfrac5{12});\\ & \text{the crossing locus is the profile's level set at}\ e_\ast^2=4\pi\varepsilon_\ast;\ \text{three profile shapes: one crossing, always one, or a window with two};\\ & P_{\max}=1-\tfrac3{2x_\ast},\ \ x_\ast\ \text{the unique root}\ >1\ \text{of}\ 2\beta x^3-3\beta x^2-1;\ \ \text{the specified and pure-pole graphs are flat in}\ E\\\text{the dimensional bridge} & \mathrm{BM.102}:\ \textbf{censused, and the two bridges separated}\ —\ \text{one dimension}\ \mathfrak q;\ \text{the seven rows derived or adjudicated};\ [\text{current}]=\mathfrak q^0\ \text{(ledger)}\ \text{and}\ \mathfrak q^{+1/2}\ \text{(chart)};\ [\text{temperature}]=\mathfrak q^{+1/2}\ \text{on}\ k_B{=}1;\\ & \textbf{the dimensional bridge is}\ \mathfrak h,\ \textbf{the numerical one is}\ \rho_{\mathrm B}\in\mathfrak q^0;\ \textbf{A-BM.9 declared}: \text{one calibration slot, value}\ \neg\text{issued};\ \textbf{item 7 closed as an equation}\\\text{the eigenvalue composition} & \mathrm{BM.103}:\ \textbf{adjudicated}\ \neg\textbf{formed}\ —\ \det(\lambda I-M_{\mathrm{ensemble}})\ \text{has no referent}: \neg\text{a three-chain state space},\ \neg\text{an exchange operator},\ \neg\text{a flavour index},\ \neg\text{a triple resolution};\\ & \textbf{the formed eigenvalue equations:}\ \text{the family cubic with}\ \sum_iz_i=6;\ \ W(Q)=\binom6{Q+3}\ \text{with one root}\ {-}1\ \text{of multiplicity six};\ \text{the octet}\ \neg\text{formed};\ (\mathrm{R3})\ \text{stands}\\\text{the curvature composition} & \mathrm{BM.104}:\ \textbf{the fixed point derived, the selection refused}\ —\ \mathsf g_{\mathrm{fp}}=1,\ \tfrac{1+G}{1-aG},\ \text{the quadratic's positive root};\ \text{the pole locus}\ \pi Nw\,\varepsilon_{\mathrm{EM}}Q^2E^2\ln2=1,\ \text{unreachable at}\ Q{=}0;\\ & \text{the loop is excluded by the lock's P4-freedom};\ \mathsf g(G)\ \neg\text{derived on four grounds};\ \textbf{the composition is observable in the graph's topology and still}\ \neg\text{selected}\\\text{the crystallization weight} & \mathrm{BM.105}:\ \textbf{generalized}\ —\ w_{\mathrm{mat}}=\tfrac13\sum_ia_i\in[\tfrac13,1]\ \text{measure-free, with both equality loci};\ \text{the family}\ S/(3(1-R))\ \text{with seven exact configuration weights};\\ & \text{AM--GM gives}\ w_{\mathrm{mat}}\le1\ \text{with equality iff every pair crystallizes};\ \textbf{hypothesis (W6) refuted};\ (\mathrm{R4})\ \text{closes at two levels and narrows to}\ (\mathrm{R4}')\\\text{spin and statistics} & \mathrm{BM.106}:\ \textbf{deficit 11 stands}\ —\ j=|k|/2\ \text{derived, the tick law derived, and}\ \kappa{=}1,2\ \text{give}\ j=\tfrac12,1\ \text{on one kind}: \neg\text{a spin-parity theorem};\\ & \neg\text{a statistics referent}: \neg\text{a multi-chain state space},\ \neg\text{an exchange operator},\ \neg\text{a permutation sign};\ \text{the four}\ \mathbb Z_2\text{'s are single-object};\ (\mathrm{R7})\ \text{named}\end{array}
\]

%dsp-level:0
\[
(2)\quad\textbf{the derived inventory (the new exact vocabulary of this tier):}
\]

%dsp-level:0
\[
\qquad \varepsilon_{\mathrm{eff}}=Nf\,\mathsf g\,\varepsilon_{\mathrm{EM}}=e^2/(4\pi);
\]

%dsp-level:0
\[
\qquad \varepsilon_\ast=\frac{2-\log_23}{3w};
\]

%dsp-level:0
\[
\qquad e_\ast^2=\frac{4\pi(2-\log_23)}{3w};
\]

%dsp-level:0
\[
\qquad x_\ast=\frac1{1-c}\ \ (\mathsf g{=}1);
\]

%dsp-level:0
\[
\beta x^3-\beta x^2+(1-c)x-1=0\ \ (\mathsf g{=}1{+}G);
\]

%dsp-level:0
\[
\qquad c\beta x^3+(c-1)x+1=0,\quad P_{\max}=1-\frac3{2x_\ast},\quad2\beta x_\ast^3-3\beta x_\ast^2=1\ \ (\mathsf g{=}(1{+}G)^{-1});
\]

%dsp-level:0
\[
\beta=\frac{2\pi\mathfrak h^2}{r^2\ln2},\qquad E_1=\frac{\mathfrak h}{r\ln2},\qquad h\bigl(\tfrac32\bigr)=-1,\qquad h(2)=4\beta-1;
\]

%dsp-level:0
\[
[\text{current}]_{\text{chart}}=\mathfrak q^{+3/2}\ \ (\text{density}),
\]

%dsp-level:0
\[
\qquad [\text{current}]_{\text{ledger}}=\mathfrak q^0,
\]

%dsp-level:0
\[
\qquad [F^{\mu\nu}]=\mathfrak q^{+1},
\]

%dsp-level:0
\[
\qquad [\rho_{\mathrm B}]=\mathfrak q^0,
\]

%dsp-level:0
\[
\qquad [\mathfrak h]=\mathfrak q^{-1/2};
\]

%dsp-level:0
\[
W(Q)=\binom6{Q+3},\qquad\bigl(2+z+z^{-1}\bigr)^3=\frac{(z+1)^6}{z^3},\qquad4z^3-24z^2+(21-12\gamma)z-(\gamma-1)^2=0,\qquad\sum_iz_i=6;
\]

%dsp-level:0
\[
G_{\mathrm{corr}}=G(1+a\mathsf g),
\]

%dsp-level:0
\[
\qquad a=\tfrac12Nw\,\varepsilon_{\mathrm{EM}}Q^2,
\]

%dsp-level:0
\[
\qquad \mathsf g_{\mathrm{fp}}=\frac{1+G}{1-aG},
\]

%dsp-level:0
\[
\qquad aG\,\mathsf g^2+(1+G)\mathsf g-1=0,
\]

%dsp-level:0
\[
\qquad aG=\pi Nw\,\varepsilon_{\mathrm{EM}}Q^2E^2\ln2;
\]

%dsp-level:0
\[
w_{\mathrm{mat}}=\frac13\sum_{i=1}^3a_i\in\Bigl[\frac13,1\Bigr],
\]

%dsp-level:0
\[
\qquad w_{\mathrm{mat}}^{\mathrm{family}}=\frac{\sum_ip_i}{3\bigl(1-\prod_iu_i\bigr)},
\]

%dsp-level:0
\[
\qquad 3(1-R)-S=\sum_iu_i-3\prod_iu_i\ge0;
\]

%dsp-level:0
\[
j=\frac{|\kappa|}2,\qquad\kappa\in\mathbb Z\setminus\{0\},\qquad\Theta_{\mathrm{int}}=2\kappa,\qquad\Bigl(\frac9{16}\Bigr)^{\!6}=\frac{531441}{16777216},\qquad16\delta_\nu=9\delta_\ell=27\delta_q=2;
\]

%dsp-level:0
\[
\textbf{the certificates:}
\]

%dsp-level:0
\[
\quad 2^{13}>3^8,
\]

%dsp-level:0
\[
\quad 3^{12}>2^{19},
\]

%dsp-level:0
\[
\quad 3^5<2^8,
\]

%dsp-level:0
\[
\quad 12<40,
\]

%dsp-level:0
\[
\quad 4<12,\quad\binom6{k}=(1,6,15,20,15,6,1),
\]

%dsp-level:0
\[
\quad 2^6=64,
\]

%dsp-level:0
\[
\quad 531441<16777216
\]

%dsp-level:0
\[
\qquad [\text{machine: every row exact}];
\]

%dsp-level:0
\[
(3)\quad\textbf{the status adjudication:}
\]

%dsp-level:0
\[
\qquad \boxed{\ \textbf{one graph derived with its crossing strength and its composition-dependent crossing count};\qquad\textbf{one dimension counted, two bridges separated, one calibration slot declared with no value issued};}
\]

%dsp-level:0
\[
\boxed{\ \textbf{one composition question adjudicated}\ \neg\textbf{formed, with the two formed eigenvalue equations derived};\qquad\textbf{one fixed point solved three times and one selection refused};}
\]

%dsp-level:0
\[
\boxed{\ \textbf{one weight generalized measure-free and one defect conditioning refuted};\qquad\textbf{one kind structure derived and one statistics}\ \neg\textbf{formed}\ }
\]

%dsp-level:0
\[
\textbf{the level-exactness:}
\]

%dsp-level:0
\[
\qquad \text{derived:}\ \varepsilon_\ast,
\]

%dsp-level:0
\[
\ e_\ast^2,
\]

%dsp-level:0
\[
\ \text{the three crossing equations},
\]

%dsp-level:0
\[
\ P_{\max},
\]

%dsp-level:0
\[
\ x_\ast\text{'s uniqueness and enclosure},
\]

%dsp-level:0
\[
\ \text{the whole exponent census},
\]

%dsp-level:0
\[
\ W(Q)
\]

%dsp-level:0
\[
= \binom6{Q+3},
\]

%dsp-level:0
\[
\ \text{the family cubic's Vieta rows},
\]

%dsp-level:0
\[
G_{\mathrm{corr}},
\]

%dsp-level:0
\[
\ \text{the three fixed points},
\]

%dsp-level:0
\[
\ \text{the pole locus},
\]

%dsp-level:0
\[
\ \text{the measure-free bounds and their equality loci},
\]

%dsp-level:0
\[
\ \text{the family form},
\]

%dsp-level:0
\[
\ \text{the AM--GM row},
\]

%dsp-level:0
\[
\ j
\]

%dsp-level:0
\[
= |\kappa|/2\ \text{and the two-parity census},
\]

%dsp-level:0
\[
\ \text{the tick and contraction rows},
\]

%dsp-level:0
\[
\ \text{the four gradings' distinctness},
\]

%dsp-level:0
\[
\ \text{the wall-parity census};
\]

%dsp-level:0
\[
\neg\textbf{derived:}\quad\mathsf g\ (\text{the curvature composition}),
\]

%dsp-level:0
\[
\quad e_0,
\]

%dsp-level:0
\[
\quad N\ (\text{solution data}),\quad\rho_{\mathrm B}\text{'s value}\ (\text{A-BM.9's slot}),
\]

%dsp-level:0
\[
\quad M_{\mathrm{ensemble}},
\]

%dsp-level:0
\[
\quad\text{a stability condition},
\]

%dsp-level:0
\[
\quad\text{the flavour composition}\ (\mathrm{R3}),
\]

%dsp-level:0
\[
\quad\text{the pair correlations}\ (\mathrm{R4}'),
\]

%dsp-level:0
\[
\iota\ (\mathrm{R5}),\quad\text{the Jensen residue}\ (\mathrm{R6}),\quad\text{an exchange structure}\ (\mathrm{R7}),\quad\hat\mu_\nu,\ \ell_c^{(\nu)},\ \mu_\nu,
\]

%dsp-level:0
\[
\quad i_\nu\ \text{and}\ \neg\text{a derived mixing matrix},\quad\text{the generation selection};
\]

%dsp-level:0
\[
(4)\quad\textbf{the declarations this tier adds: one}\ —\ \textbf{A-BM.9},
\]

%dsp-level:0
\[
\ \text{the bridge's dimensionless calibration slot}\ \rho_{\mathrm B},
\]

%dsp-level:0
\[
\ \text{logged before use, with}\ \neg\textbf{a value issued}
\]

%dsp-level:0
\[
\qquad [\mathrm{BM.102}(5)]:
\]

%dsp-level:0
\[
\text{§Cal-BM gains exactly one declared member and}\ \neg\text{a new free parameter is introduced}\ —\ \text{the slot packages}\ \hat\mu'^2\ [\mathrm{Cal\text{-}F}]\ \ \text{and}\ \ \mathfrak h\ [\mathrm{P1}];
\]

%dsp-level:0
\[
\qquad \neg\text{a new dimensionful scale}\ [\mathrm{P1}];
\]

%dsp-level:0
\[
\qquad \neg\text{a measured value}\ [\mathrm{P5}];
\]

%dsp-level:0
\[
\textbf{the declarations this tier adds beyond that slot: none}\ —\ \neg\text{a name attachment},
\]

%dsp-level:0
\[
\ \neg\text{a read specification},
\]

%dsp-level:0
\[
\ \neg\text{a P4 addition},
\]

%dsp-level:0
\[
\ \neg\text{a second registry addendum};
\]

%dsp-level:0
\[
(5)\quad\textbf{the scope rows (recorded,}\ \neg\textbf{silent):}
\]

%dsp-level:0
\[
\qquad \neg\text{a clause of}\ \mathrm{BM.1}\text{--}\mathrm{BM.100}\ \text{is retracted, replaced or renumbered};
\]

%dsp-level:0
\[
\qquad \neg\text{a clause of Parts 1--3 is modified};
\]

%dsp-level:0
\[
\qquad \neg\text{a Census cell gains or loses a value};
\]

%dsp-level:0
\[
\textbf{rescoped at their own loci:}\ \mathrm{BM.91}(2),\ \mathrm{BM.99}(1),(3)\ \text{as the far-field locus}\ f=1\ \text{of}\ \mathrm{BM.101}(2);
\]

%dsp-level:0
\[
\qquad \mathrm{BM.94}(7)\text{'s strict bound as the interior row of}\ \mathrm{BM.105}(2)\text{'s measure-free one};
\]

%dsp-level:0
\[
\qquad (\mathrm{R4})\ \text{as}\ (\mathrm{R4}')\ \ [\mathrm{BM.105}(6)];
\]

%dsp-level:0
\[
\textbf{standing untouched:}\ \mathrm{BM.35}(5),
\]

%dsp-level:0
\[
\ \mathrm{BM.61}(8),
\]

%dsp-level:0
\[
\ \mathrm{BM.69}(2)\text{--}(5),
\]

%dsp-level:0
\[
\ \mathrm{BM.77},
\]

%dsp-level:0
\[
\ \mathrm{BM.80},
\]

%dsp-level:0
\[
\ \mathrm{BM.84},
\]

%dsp-level:0
\[
\ \mathrm{BM.88},
\]

%dsp-level:0
\[
\ \mathrm{BM.89},
\]

%dsp-level:0
\[
\ \mathrm{BM.90},
\]

%dsp-level:0
\[
\ \mathrm{BM.94}\text{--}\mathrm{BM.99},
\]

%dsp-level:0
\[
\ \mathrm{BM.102}(2)\text{'s census rows},
\]

%dsp-level:0
\[
\ \mathrm{FT.125}(5),
\]

%dsp-level:0
\[
\ \mathrm{FT.66}(4);
\]

%dsp-level:0
\[
\text{the conditionalities}\ (\mathrm{Def\text{-}H.37},\ \mathrm{Def\text{-}H.38},\ \mathrm{M1}/\mathrm{M2},\ \mathrm{CT2.6},\ \mathrm{CT2.7},\ \mathrm{BM.8}(2),\ \mathrm{BM.59}(9)\text{'s hypothesis}\ (\mathrm H),\ \mathrm{BM.105}(3)\text{'s hypothesis}\ (\mathrm F))\ \text{are inherited at every point of use}\ [\text{A §0.2a}];
\]

%dsp-level:0
\[
(6)\quad\textbf{the standing rules:}
\]

%dsp-level:0
\[
\qquad \neg\text{a new dimensionful scale}\ [\mathrm{P1}];
\]

%dsp-level:0
\[
\qquad \neg\text{a measured value and}\ \neg\text{a numerical coupling}\ [\mathrm{P5};\ \mathrm{EM.25}(7),(9);\ \text{machine: zero external tokens, zero decimal numerals, zero}\ [\mathrm{EXT}]\ \text{markers in this tier}];
\]

%dsp-level:0
\[
\text{no measured external constant is identified with any corpus quantity}\ [\mathrm{P5}];
\]

%dsp-level:0
\[
\qquad \neg\text{the observed neutron--proton ratio and}\ \neg\text{the observed electric coupling are imported}\ [\mathrm{BM.99}(5);\ \mathrm{BM.101}(3),(4)];
\]

%dsp-level:0
\[
\text{species names are confined to the identification layers}\ [\mathrm{Def\text{-}BM.3}(7);\ \mathrm{Def\text{-}BM.6}(7);\ \mathrm{Def\text{-}BM.14}(8)];
\]

%dsp-level:0
\[
\qquad \text{topology labels are internal bookkeeping}\ [\mathrm{Def\text{-}H.32}(4)];
\]

%dsp-level:0
\[
\text{amplitude}\ne\text{mass}\ [\mathrm{FT.79}(5)];
\]

%dsp-level:0
\[
\qquad \text{rung}\ne\text{scale}\ [\mathrm{D0.8}];
\]

%dsp-level:0
\[
\qquad \neg\text{a currency is converted}\ [\mathrm{BM.83}(3);\ \mathrm{BM.90}(5);\ \mathrm{BM.98}(6);\ \mathrm{BM.102}(6)];
\]

%dsp-level:0
\[
\alpha_B:\quad\neg\text{added},\quad\neg\text{derived},\quad\text{barred}\ —\ \mathrm{BM.61}(8)\text{'s}\ (\mathrm O1)\text{--}(\mathrm O4)\ \text{stand}\qquad[\text{machine: the token scan of this tier}];
\]

%dsp-level:0
\[
\neg\text{a Yukawa-type object exists or is added},
\]

%dsp-level:0
\[
\qquad \neg\text{a seesaw-type mechanism is suggested},
\]

%dsp-level:0
\[
\qquad \neg\text{a right-ray matter line is added},
\]

%dsp-level:0
\[
\qquad \neg\text{an external mixing matrix is imported}
\]

%dsp-level:0
\[
\qquad [\mathrm{EW.15};\ \mathrm{FT.145}(3);\ \mathrm{BM.96}(4),(6);\ \mathrm{BM.106}(6);\ \text{machine}];
\]

%dsp-level:0
\[
\neg\text{a HistLog disposition imported as premise}\ —\ \text{items 7, 19, 20, Q1}',
\]

%dsp-level:0
\[
\ \text{Q7},
\]

%dsp-level:0
\[
\ \text{Q9},
\]

%dsp-level:0
\[
\ \text{Q11 and the deficits are cited as pointers, and A-BM.9 is issued here and logged before use};
\]

%dsp-level:0
\[
(7)\quad\textbf{the boundary inventory (every row dispositioned in full at BM.143):}\qquad\begin{array}{l|l}(\mathrm{R1})\ \ \rho_{\mathrm B}\text{'s value} & \text{the declared slot}\ \text{A-BM.9};\ \neg\text{derived on}\ \mathrm{BM.69}(4)\text{'s four grounds};\ \neg\text{fixed by measurement}\ [\mathrm{P5}]\\(\mathrm{R2})\ \ \text{the loop realization} & \text{a clean-enclosure vertex-figure loop as a ledger configuration}\ [\mathrm{BM.67}(9),(10);\ \mathrm{Census\ FT.63}(4)]\\(\mathrm{R3})\ \ \text{the flavour composition} & \text{which three of the six states a sealed class occupies}\ [\mathrm{BM.28}(7);\ \mathrm{BM.103}(4),(8)]\\(\mathrm{R4}')\ \ \text{the pair correlations} & \text{the ensemble's configuration distribution beyond per-pair independence, and the transport register's ensemble average}\ [\mathrm{BM.105}(6);\ \mathrm{BM.89}(8)]\\(\mathrm{R5})\ \ \iota\ \text{and the interior multipoles} & \text{the internal charge resolution the end signs cannot carry}\ [\mathrm{BM.46}(2);\ \mathrm{BM.87}(3)]\\(\mathrm{R6})\ \ \text{the Jensen residue} & \langle E\rangle_{\mathrm{mat}}\langle E^{-1}\rangle_{\mathrm{mat}}-1:\ \text{the two formed reads of the reciprocal register}\ [\mathrm{BM.94}(8);\ \mathrm{BM.98}(5)]\\(\mathrm{R7})\ \ \text{an exchange structure} & \text{a two-chain state space, an exchange operator, a permutation-sign read}\ [\mathrm{BM.106}(7)]\\\mathsf g,\ e_0,\ N & \text{the curvature composition, the electric coupling's value, the lapse as solution data}\ [\mathrm{BM.104}(5),(7);\ \mathrm{EM.25}(7);\ \mathrm{CT2.7}]\\m_\nu\ \text{absolute},\ \mu_\nu,\ i_\nu,\ \text{the mixing matrix} & \neg\text{derived}\ [\mathrm{BM.96}(3);\ \mathrm{BM.106}(6)]\\\text{the generation count as a selection} & \text{three lines derived, three selected}\ \neg\text{derived}\ [\mathrm{BM.60}(1)]\\\text{the two mass ceilings} & \text{deficit}\ 12\ \text{stands}\ [\mathrm{Census\ FT.76}(3);\ \mathrm{FT.79}(2);\ \mathrm{FT.84}(2)]\end{array}
\]

%dsp-level:0
\[
(8)\quad\textbf{the adjudication:}
\]

%dsp-level:0
\[
\qquad \boxed{\ \textbf{six objectives, six exact verdicts: one derived, one declared, three adjudicated}\ \neg\textbf{formed or}\ \neg\textbf{selected, one generalized}\ }
\]

%dsp-level:0
\[
\textbf{the cost:}
\]

%dsp-level:0
\[
\qquad \text{one declared calibration slot with no value, two new named residues}\ (\mathrm{R4}'),
\]

%dsp-level:0
\[
\ (\mathrm{R7}),
\]

%dsp-level:0
\[
\ \ \text{one marked hypothesis}\ (\mathrm F),
\]

%dsp-level:0
\[
\ \ \text{zero parameters fitted},
\]

%dsp-level:0
\[
\ \ \text{zero scales fixed},
\]

%dsp-level:0
\[
\ \ \text{zero measured values};
\]

%dsp-level:0
\[
\textbf{the gain:}
\]

%dsp-level:0
\[
\qquad \text{a crossing strength enclosed by two integers, a graph whose shape discriminates the priced compositions, a complete dimension census, a binomial characteristic polynomial, three exact fixed points, a measure-free weight bound, and a spin-parity refutation at the floor};
\]

%dsp-level:0
\[
\text{item 7 is closed as an equation with one declared slot};
\]

%dsp-level:0
\[
\qquad \text{deficits 4, 6, 11, 12 stand};
\]

%dsp-level:0
\[
\qquad \text{nothing measured entered anywhere}.
\]

\textbf{Req.} D0.8, D2.9 (Parts 1–2); A §0.2a, CT2.1, CT2.7, EM.25(7), EM.25(9), EW.15 (Part 2); P1, P4, P5 (Parts 1–2); Def-H.32(4), Def-H.37, Def-H.38, FT.66(4), FT.73(2), FT.73(3), FT.76, FT.79(2), FT.79(5), FT.80(5), FT.84(2), FT.125(5), FT.143(3), FT.145(3) (Part 3); Census FT.63(4), FT.76(3) (row-wise); Def-BM.3(7), Def-BM.6(7), Def-BM.14(8), BM.8(2), BM.28(7), BM.35(5), BM.46(2), BM.59(9), BM.60(1), BM.61(8), BM.67(9), BM.67(10), BM.69(2), BM.69(4), BM.69(5), BM.77, BM.80, BM.83(3), BM.84, BM.87(3), BM.88, BM.89(8), BM.90(5), BM.94, BM.94(7), BM.94(8), BM.95, BM.96, BM.96(3), BM.96(4), BM.96(6), BM.97, BM.98, BM.98(5), BM.98(6), BM.99, BM.99(1), BM.99(3), BM.99(5), BM.100, BM.101, BM.102, BM.103, BM.104, BM.105, BM.106 (this part); Cal-F; HistLog, 19, 20, 21(d), Q1', Q7, Q9, Q11, §II.4 (the addendum A-BM.9), deficits $1'$, 4, 6, 11, 12 (pointers); the machine sections §A–§G green [math; machine: \texttt{checks/\allowbreak{}verify\_\allowbreak{}H\_\allowbreak{}six\_\allowbreak{}objective\_\allowbreak{}resolution.py} §G].

\textbf{Pf.}

%dsp-level:0
\[
(1)\text{--}(2):\ \text{assembly of the adjudicated rows of}\ \mathrm{BM.101}\text{--}\mathrm{BM.106}\ [\text{machine: §A--§F green}];
\]

%dsp-level:0
\[
\qquad \text{the inventory rows re-listed from their own tags}\ [\text{machine}]
\]

%dsp-level:0
\[
(3):\ \text{the derived/}\neg\text{derived partition by the symbol census of the tier}\ [\text{machine}];
\]

%dsp-level:0
\[
\qquad \mathrm{BM.69}(4)\text{'s four grounds,}\ \mathrm{BM.61}(8)\text{'s obstruction and}\ \mathrm{FT.125}(5)\text{'s kind bar re-cited}
\]

%dsp-level:0
\[
(4):\ \text{A-BM.9's species, costs and bars re-cited from the registry entry}\ [\mathrm{BM.102}(5);\ \text{machine: the packaging census}];
\]

%dsp-level:0
\[
\qquad \neg\text{a second addendum is issued}\ [\text{machine}]
\]

%dsp-level:0
\[
(5):\ \text{the clause census of}\ \mathrm{BM.1}\text{--}\mathrm{BM.100}\ \text{against retraction, replacement and renumbering}\ [\text{machine}];
\]

%dsp-level:0
\[
\qquad \text{the Census-cell invariance scan}\ [\text{machine}];
\]

%dsp-level:0
\[
\qquad \text{the conditionality rows re-cited at their loci}\ [\text{A §0.2a}]
\]

%dsp-level:0
\[
(6):\ \text{the standing-rule scans: external tokens, decimal numerals, unit tokens, registry tokens, declaration lines, the barred object}\ \alpha_B,
\]

%dsp-level:0
\[
\ \text{the seesaw-type, right-ray and mixing-matrix rows}\ [\text{machine}]
\]

%dsp-level:0
\[
(7):\ \text{the open rows re-cited from their own tags}\ [\text{machine}];\qquad(8):\ \text{assembly}.\ \blacksquare
\]

\textbf{Phys.} Six objectives came in and six verdicts went out, and only one of them is a plain yes. The graph is the yes: the two reactions are one curve reflected, the crossing strength is a derived number enclosed between an eighth and five twelfths by two integer powers, and the curve's shape — a half-line, an always-crossing half-line, or a window that opens and shuts — is decided by which curvature composition the corpus is eventually given, so the unresolved choice has an observable consequence inside the theory even though nothing here can make it. The dimension count is a yes with a correction: one flux, seven rows derived or adjudicated, and two bridges where the question assumed one, because the scale that carries bits to mass is the free scale and the number that exchanges bits for angles is dimensionless. That number is now declared, as a slot with no value in it, which closes the relating equation as an equation and leaves the calibration where the corpus leaves every calibration. The other three objectives are exact negatives with content. There is no three-quark matrix to take eigenvalues of, because there is no state space, no exchange operator and no flavour index — but there are two eigenvalue equations that do exist, a family cubic summing to six and a binomial whose single root has multiplicity six, one per directed edge. The curvature composition solves its own consistency condition three times over and is selected by none of them, and the back-reaction that would have selected it is excluded by the lock's own freedom census. The crystallization weight generalizes past its specialization: measure-free bounds with both equality loci, a threshold family with seven exact configuration weights, and a refutation of the six-state version on the four old grounds. And spin meets statistics at a wall the corpus has already mapped — the spin label is derived, the arrow is derived, and at the floor the same kind of chain carries both parities, so what is missing is not a step but an exchange structure, now named. Two residues were added to the ledger's list, one slot was declared, and nothing was fitted.

\medskip
\subsection*{\texorpdfstring{B.1.31\quad BM.118 — The status of the Antimatter and Resonance Tier: adjudication, scope, and the open items}{BM.118 B.1.31}}
\addcontentsline{toc}{subsection}{B.1.31 BM.118}
\phantomsection\label{tag:BM-118}
{\small \textit{Record.}\quad Source: Part 4 (Part4\_Baryomesonic\_era.md); tier: Tier BM.21 — The Antimatter and Resonance Tier: the CP-Conjugate, the Meson Seal, the Exchange Operator, and the Local Resonance Spectrum; canonical tag \texttt{BM.118}. Reproduced from the revision-1 canon; a cross-reference stub marks the former locus in the main corpus.}
\medskip
\textbf{Thm.}

%dsp-level:0
\[
(1)\quad\textbf{the four objectives and their four verdicts:}\qquad\begin{array}{l|l|l}\text{objective} & \text{verdict} & \text{locus}\\\hline\text{the CP-conjugate chain} & \textbf{derived, and forced by the ledger's own balance} & [\mathrm{BM.108}(2),(5);\ \mathrm{Def\text{-}BM.15}(2)]\\\text{the meson seal} & \textbf{derived, with the two-chain state space constructed and}\ \mathrm{BM.6}(2)\ \textbf{discharged} & [\mathrm{BM.110}(1),(2);\ \mathrm{Def\text{-}BM.16}]\\\text{spin--statistics} & \textbf{derived: deficit 11 CLOSED, residue}\ (\mathrm{R7})\ \textbf{discharged at}\ n=2 & [\mathrm{BM.113}(4),(7);\ \mathrm{BM.112}(1),(2)]\\\text{the resonance spectrum} & \textbf{derived: two characteristic polynomials, the orbit law, one exact collision, the threshold and width functions} & [\mathrm{BM.114};\ \mathrm{BM.115};\ \mathrm{BM.116}]\end{array}
\]

%dsp-level:0
\[
(2)\quad\textbf{the inventory gained:}\qquad\begin{array}{l|l}\text{new object} & \text{its exact content}\\\hline\mathrm{CP},\ c^{\mathrm{CP}},\ \mu_{\mathrm{CP}} & \text{an involution, its orbit, and}\ e^{i\pi\Theta}=1\ \text{on the materialized locus}\ [\mathrm{Def\text{-}BM.15};\ \mathrm{BM.108}(3)]\\\mathcal W_{\text{pair}} & 4\varepsilon_q(1-4\varepsilon_q)\le\tfrac14,\ \text{max at}\ \varepsilon_q=\tfrac18,\ \text{zero at both endpoints}\ [\mathrm{BM.109}(2)]\\\mathcal H_2,\ \mathcal H_{\text{meson}} & \text{the first multi-chain state spaces;}\ \dim=4(|k_1|+1)(|k_2|+1)\ [\mathrm{Def\text{-}BM.16}(2)]\\\hat P_{12} & \hat P_{12}^2=\mathrm{id}\ \text{unconditionally};\ \operatorname{spec}\subseteq\{\pm1\};\ \operatorname{Tr}=d\,\mu_{\mathrm{CP}}\ [\mathrm{BM.112}(1),(2)]\\\varepsilon(c) & (-1)^{|k|}=(-1)^{2j}\ [\mathrm{BM.113}(2),(4)]\\\mathcal R,\ \lambda^{(n)},\ \mathbb M & \text{the resonance tuple, its eigenvalue, and the degree-16 eigenvalue field}\ [\mathrm{Def\text{-}BM.17}(1);\ \mathrm{BM.114}(3)]\\W_n(Q) & \binom{2n}{Q+n}\ \text{on}\ 2n+1\ \text{classes, summing to}\ 4^n\ [\mathrm{BM.114}(1)]\\\mathcal T,\ \Gamma & \text{the two exact local functions of}\ (E,N,G,\varepsilon_{\mathrm{EM}})\ [\mathrm{Def\text{-}BM.17}(3),(4);\ \mathrm{BM.115}(1);\ \mathrm{BM.116}(1)]\end{array}
\]

%dsp-level:0
\[
(3)\quad\textbf{the derived negatives (each with its ground):}\qquad\begin{array}{l|l}\text{negative} & \text{ground}\\\hline\neg\text{a dynamical inverse tick} & \mathrm{A6}:\ \neg\exists\tau^{-1};\quad\mathrm{L2}:\ V'=\tfrac9{16}V\ \text{at}\ \gamma=1;\quad\mathrm{FT.144}(4)(\mathrm{ii})\ [\mathrm{BM.108}(1)]\\\neg\text{an annihilation channel or cross-section} & \text{the inventory scan}\ [\mathrm{BM.27}(7);\ \mathrm{BM.108}(7);\ \mathrm{BM.109}(5)]\\\neg\text{a meson inside a baryon's ensemble} & \mathcal E_\delta(\mathcal X_{N=3})=\varnothing\ \text{by the grammar}\ [\mathrm{BM.31}(1);\ \mathrm{BM.110}(5)]\\\neg\text{a seal-binding term on two routes} & E^{\mathrm{info}}_{\mathrm{bind}}=0\ \text{and}\ \mathcal B_{\mathrm{pole}}=0\ \text{identically}\ [\mathrm{BM.111}(1),(2)]\\\neg\text{a meson ensemble binding} & \neg\text{a}\ 2\text{-chain sector mass exists}\ [\mathrm{BM.111}(3)]\\\neg\text{an octet or decuplet grouping} & 8\ \text{occurs in no grouping family}\ [\mathrm{BM.117}(3)]\\\neg\text{a permanently forbidden state} & \mathcal T\in(0,\infty)\implies\text{every locus is a nonempty half-line}\ [\mathrm{BM.115}(5);\ \mathrm{BM.117}(6)]\\\neg\text{a binary stable/unstable category} & \Gamma\ \text{is a function with three exact zeros and one indicator}\ [\mathrm{BM.116}(3),(4);\ \mathrm{Def\text{-}BM.17}(6)(\mathrm{ii})]\\\neg\text{a cosmological coordinate} & \text{the ledger owns}\ \tau,\ G,\ N,\ E\ \text{and no global time}\ [\mathrm{Def\text{-}BM.17}(6)(\mathrm i);\ \mathrm{A6};\ \mathrm{D1.8}]\end{array}
\]

%dsp-level:0
\[
(4)\quad\textbf{the declarations this tier adds: none.}
\]

%dsp-level:0
\[
\qquad \neg\text{a}\ \mathrm{A\text{-}BM}\ \text{addendum is issued};
\]

%dsp-level:0
\[
\qquad \neg\text{a P4 member is added};
\]

%dsp-level:0
\[
\qquad \text{§Cal-BM keeps}\ \mathrm{A\text{-}BM.9}\text{'s logged slot with its value}\ \neg\text{issued and gains nothing}\ [\mathrm{P4};\ \mathrm{P1};\ \text{machine: the packaging census}]:
\]

%dsp-level:0
\[
\textbf{the one non-derived row is a sector choice, priced and}\ \neg\textbf{declared:}
\]

%dsp-level:0
\[
\qquad \mathrm{Def\text{-}BM.16}(5)\text{'s}\ \vee^2\ \text{locus for the sign read, with the four alternatives evaluated in}\ \mathrm{BM.113}(6);
\]

%dsp-level:0
\[
(5)\quad\textbf{the rescope record (each row re-scoped at its own locus, none retracted, replaced or renumbered):}\qquad\begin{array}{l|l}\text{row} & \text{the rescope}\\\hline\mathrm{BM.6}(2) & \textbf{discharged:}\ \text{the exchange-space specification is constructed}\ [\mathrm{Def\text{-}BM.16}(6)]\\\mathrm{BM.103}(4)\ \text{rows 1--2} & \textbf{discharged at}\ n=2:\ \text{a two-chain state space and an exchange operator exist; row 3 (a flavour index), row 4 (a triple resolution) and row 5 (the 27-arena) stand}\ [\mathrm{BM.110}(2)]\\\mathrm{BM.106}(3) & \textbf{discharged on its first three rows;}\ \text{the fourth (the kind bar) and the fifth (the counting spaces) stand unchanged}\ [\mathrm{BM.110}(2);\ \mathrm{BM.112}(5)]\\\mathrm{BM.106}(7),\ \mathrm{BM.107}(7) & (\mathrm{R7})\ \textbf{discharged at}\ n=2;\ \text{the}\ n=3\ \text{exchange structure remains open}\ [\mathrm{BM.113}(7);\ \mathrm{Def\text{-}BM.16}(7)]\\\mathrm{BM.106}(1) & \textbf{stands:}\ \neg\text{a spin-parity theorem on the}\ \mathrm{D1.5}\ \text{kind;}\ \varepsilon\ \text{is a function of the winding}\ [\mathrm{BM.113}(5)]\\\mathrm{FT.80}(5),\ \mathrm{FT.125}(5) & \textbf{frozen on their own loci:}\ \neg\text{an exchange operator exists in the pre-BM.21 inventory}\ [\mathrm{BM.110}(6)]\\\mathrm{BM.32}(5),(6) & \textbf{extended to the sealed composite:}\ \text{the meson seal lives on the norm layer and the}\ \mathcal C\text{-projector}\ [\mathrm{BM.111}(4)]\\\mathrm{BM.28}(2),\ \mathrm{BM.103}(3) & \textbf{generalized to every}\ n:\ W_n(Q)=\binom{2n}{Q+n}\ [\mathrm{BM.114}(1)]\\\mathrm{BM.21}(1),\ \mathrm{BM.54}(1) & \textbf{instantiated at the pair:}\ \mathcal B_{\mathrm{info}}=0\ [\mathrm{BM.109}(1);\ \mathrm{BM.111}(1)]\\\mathrm{FT.76}(3),\ \mathrm{FT.79}(2),\ \mathrm{FT.97}(5) & \textbf{untouched:}\ \text{the ceilings are single-chain reads and a composite eigenvalue is a sum of them}\ [\mathrm{BM.117}(4),(5)]\\\text{every clause of}\ \mathrm{BM.1}\text{--}\mathrm{BM.107} & \textbf{no retraction, no replacement, no renumbering}\ [\text{machine: the clause census}]\end{array}
\]

%dsp-level:0
\[
(6)\quad\textbf{the standing rules:}
\]

%dsp-level:0
\[
\qquad \neg\text{a new dimensionful scale}\ [\mathrm{P1}];
\]

%dsp-level:0
\[
\qquad \neg\text{a measured value and}\ \neg\text{a numerical coupling}\ [\mathrm{P5};\ \mathrm{EM.25}(7),(9);\ \text{machine: zero external tokens, zero decimal numerals, zero}\ [\mathrm{EXT}]\ \text{markers in this tier}];
\]

%dsp-level:0
\[
\text{no measured external constant is identified with any corpus quantity}\ [\mathrm{P5}];
\]

%dsp-level:0
\[
\qquad \neg\text{the observed lifetime, width, mass or coupling of any object is imported}\ [\mathrm{BM.44}(5);\ \mathrm{BM.99}(5)];
\]

%dsp-level:0
\[
\text{species names are confined to the identification layers}\ [\mathrm{Def\text{-}BM.3}(7);\ \mathrm{Def\text{-}BM.6}(7);\ \mathrm{Def\text{-}BM.14}(8)];
\]

%dsp-level:0
\[
\qquad \text{topology labels are internal bookkeeping}\ [\mathrm{Def\text{-}H.32}(4)];
\]

%dsp-level:0
\[
\text{amplitude}\ne\text{mass}\ [\mathrm{FT.79}(5);\ \mathrm{BM.111}(6)];
\]

%dsp-level:0
\[
\qquad \text{rung}\ne\text{scale}\ [\mathrm{D0.8}];
\]

%dsp-level:0
\[
\qquad \neg\text{a currency is converted}\ [\mathrm{BM.83}(3);\ \mathrm{BM.102}(6);\ \mathrm{BM.117}(4)];
\]

%dsp-level:0
\[
\alpha_B:\quad\neg\text{added},\quad\neg\text{derived},\quad\text{barred}\ —\ \mathrm{BM.61}(8)\text{'s}\ (\mathrm O1)\text{--}(\mathrm O4)\ \text{stand}\qquad[\text{machine: the token scan of this tier}];
\]

%dsp-level:0
\[
\neg\text{a Yukawa-type object exists or is added},
\]

%dsp-level:0
\[
\qquad \neg\text{a seesaw-type mechanism is suggested},
\]

%dsp-level:0
\[
\qquad \neg\text{a right-ray matter line is added},
\]

%dsp-level:0
\[
\qquad \neg\text{an external mixing matrix is imported}
\]

%dsp-level:0
\[
\qquad [\mathrm{EW.15};\ \mathrm{FT.145}(3);\ \mathrm{BM.96}(4),(6);\ \mathrm{BM.106}(6);\ \text{machine}];
\]

%dsp-level:0
\[
\neg\text{a Dirac operator, a Lorentz spinor or a negative-energy spectrum is constructed}\ [\mathrm{BM.6}(8);\ \mathrm{FT.93}(3);\ \mathrm{BM.108}(6)];
\]

%dsp-level:0
\[
\neg\text{a HistLog disposition imported as premise}\ —\ \text{Q9, Q10, Q11, items 7, 13, 20 and the deficits are cited as pointers};
\]

%dsp-level:0
\[
(7)\quad\textbf{the boundary inventory (every row dispositioned in full at BM.143):}\qquad\begin{array}{l|l}(\mathrm{R1})\ \ \rho_{\mathrm B}\text{'s value} & \text{the declared slot}\ \mathrm{A\text{-}BM.9};\ \neg\text{derived on}\ \mathrm{BM.69}(4)\text{'s four grounds};\ \neg\text{fixed by measurement}\ [\mathrm{P5}]\\(\mathrm{R2})\ \ \text{the loop realization} & \text{a clean-enclosure vertex-figure loop as a ledger configuration}\ [\mathrm{BM.67}(9),(10);\ \mathrm{BM.107}(7)]\\(\mathrm{R3})\ \ \text{the flavour composition} & \text{which three of the six states a sealed class occupies}\ [\mathrm{BM.28}(7);\ \mathrm{BM.103}(8);\ \mathrm{BM.117}(6)]\\(\mathrm{R4}')\ \ \text{the pair correlations} & \text{the ensemble's configuration distribution beyond per-pair independence}\ [\mathrm{BM.105}(6);\ \mathrm{BM.89}(8)]\\(\mathrm{R5})\ \ \iota\ \text{and the interior multipoles} & \text{the internal charge resolution the end signs cannot carry}\ [\mathrm{BM.46}(2);\ \mathrm{BM.87}(3)]\\(\mathrm{R6})\ \ \text{the Jensen residue} & \langle E\rangle_{\mathrm{mat}}\langle E^{-1}\rangle_{\mathrm{mat}}-1\ [\mathrm{BM.94}(8);\ \mathrm{BM.98}(5)]\\(\mathrm{R7})\ \ \text{the exchange structure at}\ n\ge3 & \textbf{discharged at}\ n=2;\ \ \mathcal H_3\ \text{and its}\ S_3\ \text{action are}\ \neg\text{constructed}\ [\mathrm{Def\text{-}BM.16}(7);\ \mathrm{BM.113}(7)]\\(\mathrm{R8})\ \ \text{the}\ 2\text{-chain sector masses} & m^2_B(\cdot,S)\big|_{n=2}:\ \text{the meson's ensemble binding and its}\ \mathcal B\ne0\ \text{row}\ [\mathrm{BM.111}(3);\ \textbf{new}]\ \\(\mathrm{R9})\ \ \text{the annihilation channel} & \text{an inter-sector transition operator between}\ \mathrm{EM}^+\ \text{and}\ \mathrm{EM}^-:\ \neg\text{in the inventory}\ [\mathrm{BM.27}(7);\ \mathrm{BM.109}(5);\ \textbf{new}]\\\mathsf g,\ e_0,\ N & \text{the curvature composition, the electric coupling's value, the lapse as solution data}\ [\mathrm{BM.104}(5);\ \mathrm{EM.25}(7);\ \mathrm{CT2.7}]\\m_\nu\ \text{absolute},\ \mu_\nu,\ i_\nu,\ \text{the mixing matrix} & \neg\text{derived}\ [\mathrm{BM.96}(3);\ \mathrm{BM.106}(6)]\\\text{the generation count as a selection} & \text{three lines derived, three selected}\ \neg\text{derived}\ [\mathrm{BM.60}(1)]\\\text{the two mass ceilings} & \text{deficit}\ 12\ \text{stands}\ [\mathrm{FT.76}(3);\ \mathrm{FT.79}(2);\ \mathrm{FT.97}(5);\ \mathrm{BM.117}(5)]\\\text{the}\ \vee^2\ \text{sector choice} & \text{priced in}\ \mathrm{BM.113}(6),\ \neg\text{forced},\ \neg\text{declared}\ [\textbf{new}]\end{array}
\]

%dsp-level:0
\[
(8)\quad\textbf{the adjudication:}
\]

%dsp-level:0
\[
\qquad \boxed{\ \textbf{four objectives, four derived verdicts, zero declarations, two new named residues}\ }
\]

%dsp-level:0
\[
\textbf{the cost:}
\]

%dsp-level:0
\[
\qquad \text{two new residues}\ (\mathrm{R8}),\ (\mathrm{R9})\ \text{and one priced sector choice};
\]

%dsp-level:0
\[
\qquad \text{zero parameters fitted};
\]

%dsp-level:0
\[
\qquad \text{zero scales fixed};
\]

%dsp-level:0
\[
\qquad \text{zero measured values};
\]

%dsp-level:0
\[
\qquad \text{zero declarations};
\]

%dsp-level:0
\[
\textbf{the gain:}
\]

%dsp-level:0
\[
\qquad \text{a forced CP-conjugate with a mass-degeneracy theorem and a creation weight whose maximum is exactly one quarter};
\]

%dsp-level:0
\[
\ \text{the first two-chain state space with an involution exchange operator and a permutation-sign read};
\]

%dsp-level:0
\[
\text{the spin--statistics connection as}\ \varepsilon
\]

%dsp-level:0
\[
= (-1)^{|k|}
\]

%dsp-level:0
\[
= (-1)^{2j}\ \text{with deficit 11 closed};
\]

%dsp-level:0
\[
\ \text{two characteristic polynomials at every chain count};
\]

%dsp-level:0
\[
\ \text{the multiset-orbit law with stabilizer multiplicities};
\]

%dsp-level:0
\[
\ \text{exactly one exact collision in the three-chain ladder};
\]

%dsp-level:0
\[
\text{a degree-16 eigenvalue field with one generator of full degree};
\]

%dsp-level:0
\[
\ \text{threshold ladders spanning a factor of exactly thirty at every level};
\]

%dsp-level:0
\[
\ \text{a width with three exact zeros and a three-tick lifetime floor};
\]

%dsp-level:0
\[
\textbf{deficit 11 is closed};
\]

%dsp-level:0
\[
\qquad \text{residue}\ (\mathrm{R7})\ \text{is discharged at}\ n=2;
\]

%dsp-level:0
\[
\qquad \mathrm{BM.6}(2)\ \text{is discharged as written};
\]

%dsp-level:0
\[
\qquad \text{deficits}\ 1',\ 3,\ 4,\ 6,\ 12\ \text{stand};
\]

%dsp-level:0
\[
\qquad \text{nothing measured entered anywhere}.
\]

\textbf{Req.} D0.8, D1.8, D2.9 (Parts 1–2); A6, CT2.7, EM.25(7), EM.25(9), EW.15 (Part 1, Part 2); P1, P4, P5 (Parts 1–2); Def-H.32(4), FT.76(3), FT.79(2), FT.80(5), FT.89(2), FT.93(3), FT.97(5), FT.125(5), FT.145(3) (Part 3); Def-BM.3(7), Def-BM.6(7), Def-BM.14(8), Def-BM.15, Def-BM.16(5), Def-BM.16(6), Def-BM.16(7), BM.6(2), BM.6(8), BM.21(1), BM.27(7), BM.28(2), BM.28(7), BM.31(1), BM.32(5), BM.32(6), BM.44(5), BM.46(2), BM.54(1), BM.60(1), BM.61(8), BM.67(9), BM.67(10), BM.69(4), BM.83(3), BM.87(3), BM.89(8), BM.94(8), BM.96(3), BM.96(4), BM.96(6), BM.98(5), BM.99(5), BM.102(6), BM.103(3), BM.103(4), BM.103(7), BM.103(8), BM.104(5), BM.105(6), BM.106(1), BM.106(3), BM.106(4), BM.106(6), BM.106(7), BM.107(7), BM.108, BM.109, BM.110, BM.111, BM.112, BM.113, BM.114, BM.115, BM.116, BM.117, A-BM.9 (this part); the machine sections §A–§G green [math; machine: \texttt{checks/\allowbreak{}verify\_\allowbreak{}H\_\allowbreak{}antimatter\_\allowbreak{}resonance\_\allowbreak{}spectrum.py} §G].

\textbf{Pf.}

%dsp-level:0
\[
(1)\text{--}(2):\  \text{assembly of the derived rows of}\ \mathrm{BM.108}\text{--}\mathrm{BM.117}\ [\text{machine: §A--§F green}];
\]

%dsp-level:0
\[
\qquad \text{the inventory rows re-listed from their own tags}\ [\text{machine}]
\]

%dsp-level:0
\[
(3):\  \text{the nine negative rows re-cited with their grounds}\ [\text{machine}];
\]

%dsp-level:0
\[
\qquad (4):\  \text{the packaging census:}\ \neg\text{an}\ \mathrm{A\text{-}BM}\ \text{addendum is issued,}\ \neg\text{a P4 member is added}\ [\text{machine}]
\]

%dsp-level:0
\[
(5):\  \text{the clause census of}\ \mathrm{BM.1}\text{--}\mathrm{BM.107}\ \text{against retraction, replacement and renumbering}\ [\text{machine}];
\]

%dsp-level:0
\[
\qquad \text{the eleven rescope rows verified at their own loci}\ [\text{machine}]
\]

%dsp-level:0
\[
(6):\  \text{the standing-rule scans: external tokens, decimal numerals, unit tokens, registry tokens, declaration lines, the barred object}\ \alpha_B,
\]

%dsp-level:0
\[
\ \text{the Yukawa-type, seesaw-type, right-ray, mixing-matrix and Dirac-operator rows}\ [\text{machine}]
\]

%dsp-level:0
\[
(7):\  \text{the open rows re-cited from their own tags, with}\ (\mathrm{R8}),
\]

%dsp-level:0
\[
\ (\mathrm{R9})\ \text{and the sector-choice row named here}\ [\text{machine}];
\]

%dsp-level:0
\[
\qquad (8):\  \text{assembly}.\ \blacksquare
\]

\textbf{Phys.} Four objectives came in and four derived verdicts went out, with nothing declared. The conjugate is forced: total charge vanishes in every valid universe, the tick has no inverse so the compensator must be a configuration rather than a reversed evolution, and of the derived maps exactly one produces a configuration with the negated charge, the dual attachment, the opposite ray and the same edge count — which makes it mass-degenerate by structure, since not one of the seven geometric reads carries a charge exponent. The pair it forms is neutral, weighs exactly twice a member, and is created with the Janus resolvent's own weight, a parabola in the single crystallization threshold with maximum one quarter at one eighth and zeros at both ends, so pairs appear exactly where the ledger flickers. The meson is the unique seal at two attachments and it brings the corpus's first two-chain state space, which is what three successive tiers said a statistics had no referent without; the specification written as a build-spec in the open-construction list is discharged as written, and the charge side is exact — the three self-conjugate pairs saturate the neutral class at six with nothing left over, while the charged classes are made of pairs that are not conjugates. The exchange operator is the flip composed with the transport its own identification forces, which is what makes it an involution for every winding, and its sign on the symmetric sector is the fourth-root phase at the half-turn, which is the winding's parity, which is $(-1)^{2j}$: deficit eleven closes on a phase alphabet and an additivity law, with no Lorentz argument, no positivity of energy and no microcausality, and with the earlier spin-parity refutation standing untouched because that one was about the kind and this one is about the winding. The spectrum is counted rather than named: two characteristic polynomials, one binomial transport census and one defect ladder inside a degree-sixteen field with exactly one collision in two hundred sixteen triples, threshold ladders spanning a factor of thirty at every level, and a width that is a product of four frozen factors and one count with three exact zeros and a three-tick floor. Nothing is forbidden and nothing is categorically stable; what varies is where a threshold sits and what a function evaluates to. Two residues were added, one sector choice was priced, and no measured number entered anywhere.

\medskip
\subsection*{\texorpdfstring{B.1.32\quad BM.119 — The c-threshold adjudication: no velocity boundary, the exact subluminal deficit, and the register boundary that is derived}{BM.119 B.1.32}}
\addcontentsline{toc}{subsection}{B.1.32 BM.119}
\phantomsection\label{tag:BM-119}
{\small \textit{Record.}\quad Source: Part 4 (Part4\_Baryomesonic\_era.md); tier: Tier BM.22 — The Annihilation Channel, Concrete Spin, and Fermion Identification: the Janus Tunnel, the Zig-Zag World Line, and the Spin-Dependent Transport Defect; canonical tag \texttt{BM.119}. Reproduced from the revision-1 canon; a cross-reference stub marks the former locus in the main corpus.}
\medskip
\textbf{Thm.}

%dsp-level:0
\[
(1)\quad\textbf{the convention row:}\qquad c=\hbar=k_B=1\qquad[\mathrm{D2.9};\ \mathrm{BM.102}(1):\ \text{the one-dimension law}]:
\]

%dsp-level:0
\[
\boxed{\ c\ \ \textbf{is a unit convention,}\ \neg\textbf{a threshold,}\ \neg\textbf{a boundary and}\ \neg\textbf{a locus in the ledger}\ }
\]

%dsp-level:0
\[
\qquad [\text{machine: the symbol census of}\ \mathrm{D2.9}\text{'s row carries}\ \neg\text{an inequality}];
\]

%dsp-level:0
\[
(2)\quad\textbf{the exact subluminal deficit:}
\]

%dsp-level:0
\[
\qquad E=N\sqrt{m^2+p^2}-N^ip_i,
\]

%dsp-level:0
\[
\qquad \text{the vacuum gauge}\ (\Omega=1,\ w=0):\quad E^2=m^2+p^2
\]

%dsp-level:0
\[
\qquad [\mathrm{BM.34}(4);\ \mathrm{BM.34}(2)]:
\]

%dsp-level:0
\[
v\ :=\ \frac{p}{E}\ =\ \frac{p}{\sqrt{m^2+p^2}}\qquad[\text{the chart-side group velocity};\ \mathrm{CT2.6}\text{'s hypotheses at the point of use}];
\]

%dsp-level:0
\[
\boxed{\ 1-v\ =\ \frac{\sqrt{m^2+p^2}-p}{\sqrt{m^2+p^2}}\ =\ \frac{m^2}{\sqrt{m^2+p^2}\,\bigl(\sqrt{m^2+p^2}+p\bigr)}\ >\ 0\qquad\forall\,m^2>0,\ \forall\,p<\infty\ }
\]

%dsp-level:0
\[
[\text{machine: the cleared identity}\ (E-p)(E+p)=E^2-p^2=m^2\ \text{verified on exact rational witnesses}\ (m^2,p)\in\{4,\tfrac94,1,\tfrac14\}\times\{3,\tfrac52,1,1000\}];
\]

%dsp-level:0
\[
(3)\quad\textbf{every materialized chain is strictly subluminal, with a derived positive mass:}
\]

%dsp-level:0
\[
\qquad m_{\mathrm{hid}}=\frac{E_{\mathrm{int}}\ln2}{\mathfrak h}\ \ge\ \frac{2\ln2}{\mathfrak h}\ >\ 0
\]

%dsp-level:0
\[
\qquad [\mathrm{FT.44}(4);\ \mathrm{FT.28};\ \mathrm{BM.13}(3):\ E_{\mathrm{int}}\ge2;\ \mathrm{BM.70}(2):\ \text{the two-bit floor};\ \tfrac12<\ln2<1]:
\]

%dsp-level:0
\[
\boxed{\ \textbf{no materialized chain reaches}\ c\ \textbf{at any finite momentum; the deficit}\ 1-v\ \textbf{is exact and strictly positive, and it tends to}\ 0\ \textbf{only as}\ p\to\infty\ }
\]

%dsp-level:0
\[
\textbf{the monotone row:}
\]

%dsp-level:0
\[
\qquad p_1<p_2\implies v(p_1)<v(p_2)<1
\]

%dsp-level:0
\[
\qquad [\text{machine:}\ v^2=p^2/(m^2+p^2)\ \text{strictly increasing in}\ p^2\ \text{on exact rationals}];
\]

%dsp-level:0
\[
(4)\quad\textbf{the boundary that is derived is the Janus register, not a speed:}
\]

%dsp-level:0
\[
\qquad n_{\mathcal R}=0\implies\forall(p,q):\ (p,q)\in\mathcal V\ \wedge\ \rho_{pq}=\rho_\circ=\tfrac14I,
\]

%dsp-level:0
\[
\qquad S_{\mathrm{raw}}(\rho_\circ)=\ln4=\max,
\]

%dsp-level:0
\[
\qquad \arg\max=\bigl\{\tfrac14I\bigr\}
\]

%dsp-level:0
\[
[\mathrm{FT.144}(4);\ \mathrm{CT0.1};\ \mathrm{CT0.2};\ \mathrm{A5};\ \text{machine:}\ \mathfrak T\rho_\circ=1,\ \operatorname{charpoly}(\rho_\circ)=(x-\tfrac14)^4];
\]

%dsp-level:0
\[
\text{virtual pairs are unregistered}\ [\mathrm{D1.1}];\qquad\text{chains require real edges}\ [\mathrm{D1.4}:\ e_i\in E_{\mathcal R}]:
\]

%dsp-level:0
\[
\boxed{\ |E_{\mathcal R}|=0\implies\text{the closed-chain count is}\ 0\implies\neg\exists\ \text{a fermion terminus}\ }
\]

%dsp-level:0
\[
\qquad [\mathrm{Def\text{-}H.4};\ \mathrm{FT.144}(4);\ \text{machine}];
\]

%dsp-level:0
\[
\therefore\ \textbf{the materialized sector's boundary is the}\ \mathcal R/\mathcal V\ \textbf{partition of the pair register, and crossing it is the move}\ \hat J_\pm,
\]

%dsp-level:0
\[
\ \neg\textbf{an acceleration};
\]

%dsp-level:0
\[
(5)\quad\textbf{the crossing threshold is a ledger threshold:}
\]

%dsp-level:0
\[
\qquad \text{the dissolution weight is}\ x_-=4\varepsilon_q\ \ \text{and the crystallization weight is}\ x_+=1-4\varepsilon_q
\]

%dsp-level:0
\[
\qquad [\mathrm{BM.19}(2);\ \mathrm{Def\text{-}BM.5}(3)]:
\]

%dsp-level:0
\[
\boxed{\ x_-=0\iff\varepsilon_q=0\iff\neg\exists\ \text{a dissolution move}\iff\text{the ensemble collapses to its crystallized member}\ }
\]

%dsp-level:0
\[
\qquad [\mathrm{BM.19}(3);\ \mathrm{BM.23}(8);\ \text{machine}];
\]

%dsp-level:0
\[
\textbf{the threshold parameter is}\ \varepsilon_q
\]

%dsp-level:0
\[
\in [0,\tfrac14)\ \textbf{— the}\ \mathrm{A4}\ \textbf{crystallization threshold,}\ \neg\textbf{a velocity and}\ \neg\textbf{a P-class datum}\ [\mathrm{D0.4};\ \mathrm{CT0.8};\ \mathrm{BM.19}(2)];
\]

%dsp-level:0
\[
(6)\quad\textbf{the adjudication:}\qquad\begin{array}{l|l}\text{the claim} & \text{the verdict}\\\hline\text{``the chain is accelerated to}\ c\text{''} & \neg\textbf{formed:}\ (2),(3)\ \text{give a strict inequality with an exact positive deficit at every finite}\ p\\\text{``}c\ \text{is the boundary of the materialized sector''} & \neg\textbf{formed:}\ c\ \text{is a unit convention}\ (1);\ \textbf{the derived boundary is the}\ \mathcal R/\mathcal V\ \textbf{partition}\ (4)\\\text{``the chain crosses via a velocity threshold''} & \neg\textbf{formed:}\ \textbf{the crossing is the move}\ \hat J_\pm\ \textbf{with weight}\ x_\pm\ (5)\\\text{``the crossing is a Janus tunnel''} & \textbf{derived, on the register:}\ \mathrm{BM.120}\ \text{constructs}\ \hat J_\pm,\ Q,\ \ker Q,\ \operatorname{spec}Q\end{array}
\]

%dsp-level:0
\[
(7)\quad\textbf{status:}
\]

%dsp-level:0
\[
\qquad \textbf{derived:}\ (2),\ (3),\ (4),\ (5);
\]

%dsp-level:0
\[
\qquad \textbf{refuted as formed:}\ (1)\text{'s and}\ (6)\text{'s three velocity readings};
\]

%dsp-level:0
\[
\neg\text{a declaration},
\]

%dsp-level:0
\[
\quad\neg\text{a P4 addition},
\]

%dsp-level:0
\[
\quad\neg\text{a new dimensionful scale}\ [\mathrm{P1}],
\]

%dsp-level:0
\[
\quad\neg\text{a measured value}\ [\mathrm{P5}],
\]

%dsp-level:0
\[
\quad\neg\text{a species identification}\ [\mathrm{Def\text{-}BM.3}(7)],
\]

%dsp-level:0
\[
\quad\neg\text{a cosmological coordinate}\ [\mathrm{Def\text{-}BM.17}(6)(\mathrm i)];
\]

%dsp-level:0
\[
\textbf{the chart-side scope:}\ (2),
\]

%dsp-level:0
\[
(3)\ \text{are continuum-chart reads under}\ \mathrm{CT2.6}\text{'s hypotheses at the point of use}\ [\mathrm{BM.34}\ \mathrm{Rmk};\ \mathrm{BM.10}(1):\ \neg\text{a canonical pair in the finite register}];
\]

%dsp-level:0
\[
\neg\text{a HistLog disposition imported as premise}.
\]

\textbf{Rmk (why the refutation is the useful half).} Clause (2) is a one-line algebraic identity, and it does more work than it looks: it says the corpus cannot describe a massive chain arriving at $c$, because at every finite momentum there is an exact positive gap whose size is the mass squared over a product of two energies. So a mechanism that needs "reaching $c$" has no referent, and the tier supplies the referent that does exist — a two-state register with a crossing weight that is a polynomial in one threshold. The gap between the two is not a gap in the physics but a difference in kind: one is a limit of a chart-side dispersion relation, the other is a move in a finite ensemble.

\textbf{Req.} A4, A5, A6, D0.4, D0.8, D1.1, D1.4, D1.8, D2.9 (Part 1); CT0.1, CT0.2, CT0.8, CT2.1, CT2.6 (Part 1, Part 2); P1, P4, P5 (Parts 1–2); Def-H.4, Def-H.32, FT.22, FT.28, FT.44(4), FT.108(4), FT.144(4) (Part 3, Census); Def-BM.3(7), Def-BM.5(3), Def-BM.17(6), Def-BM.18, BM.10(1), BM.13(3), BM.19(2), BM.19(3), BM.23(8), BM.34(2), BM.34(4), BM.70(2), BM.102(1), BM.115(1) (this part); exact rational dispersion identities, monotonicity on exact rationals, exact rational $\ln2$ enclosures, symbol censuses [math; machine: \texttt{checks/\allowbreak{}verify\_\allowbreak{}H\_\allowbreak{}janus\_\allowbreak{}tunnel\_\allowbreak{}concrete\_\allowbreak{}spin.py} §A].

\textbf{Pf.}

%dsp-level:0
\[
(1):\ \mathrm{D2.9},\ \mathrm{BM.102}(1)\ \text{re-cited};
\]

%dsp-level:0
\[
\qquad \text{the row}\ c=\hbar=k_B=1\ \text{is an equality of units and carries}\ \neg\text{an inequality}\ [\text{machine}]
\]

%dsp-level:0
\[
(2):\ E^2=m^2+p^2\ [\mathrm{BM.34}(4)\ \text{in the vacuum gauge of}\ \mathrm{BM.34}(2)];\qquad v=p/E;\qquad1-v=(E-p)/E
\]

%dsp-level:0
\[
(E-p)(E+p)
\]

%dsp-level:0
\[
= E^2-p^2
\]

%dsp-level:0
\[
= (m^2+p^2)-p^2
\]

%dsp-level:0
\[
= m^2
\]

%dsp-level:0
\[
\implies E-p
\]

%dsp-level:0
\[
= \frac{m^2}{E+p}
\]

%dsp-level:0
\[
\implies 1-v
\]

%dsp-level:0
\[
= \frac{m^2}{E(E+p)}\ [\text{machine: the cleared identity on exact rational witnesses}]
\]

%dsp-level:0
\[
m^2>0\wedge E>0\wedge E+p>0\implies1-v>0\ [\text{machine: sign census}];\qquad v<1\iff p<E\iff p^2<m^2+p^2\iff m^2>0\ [\text{machine}]
\]

%dsp-level:0
\[
(3):\ m_{\mathrm{hid}}=E_{\mathrm{int}}\ln2/\mathfrak h\ge2\ln2/\mathfrak h\ [\mathrm{FT.44}(4);\ \mathrm{BM.13}(3)];
\]

%dsp-level:0
\[
\qquad \tfrac12<\ln2<1\implies2\ln2>1>0\ [\mathrm{BM.13}(3)\text{'s exact enclosure};\ \text{machine}]
\]

%dsp-level:0
\[
v^2=\frac{p^2}{m^2+p^2};
\]

%dsp-level:0
\[
\qquad p_1^2<p_2^2\implies\frac{p_1^2}{m^2+p_1^2}<\frac{p_2^2}{m^2+p_2^2}
\]

%dsp-level:0
\[
\iff p_1^2(m^2+p_2^2)<p_2^2(m^2+p_1^2)
\]

%dsp-level:0
\[
\iff p_1^2m^2<p_2^2m^2
\]

%dsp-level:0
\[
\iff p_1^2<p_2^2\ [\text{machine}]
\]

%dsp-level:0
\[
(4):\ \mathrm{FT.144}(4)\text{'s}\ \mathcal V\text{-sector rows re-cited (frozen, machine-verified at Tier F.41)};
\]

%dsp-level:0
\[
\qquad |E_{\mathcal R}|=0\implies\neg\exists\ \text{a closed chain of real edges}\ [\mathrm{D1.4}]\implies\neg\exists\ \text{a terminus}\ [\mathrm{Def\text{-}H.4}]
\]

%dsp-level:0
\[
(5):\ x_-=4\varepsilon_q,\ x_+=1-4\varepsilon_q\ [\mathrm{BM.19}(2);\ \text{machine}];
\]

%dsp-level:0
\[
\qquad \varepsilon_q=0\implies x_-=0\implies\neg\exists\ \text{a dissolution generator}\ [\mathrm{BM.19}(3);\ \mathrm{D1.2}]
\]

%dsp-level:0
\[
(6),(7):\ \text{the four-row adjudication assembled from}\ (1)\text{--}(5);
\]

%dsp-level:0
\[
\qquad \text{the standing-rule scans}\ [\text{machine}].\ \blacksquare
\]

\textbf{Phys.} Three of the four claims in the question have no referent, and the tier says which and why. The speed of light is a unit convention here, not a threshold; a massive chain never reaches it, because the dispersion gives an exact positive gap equal to the mass squared over a product of two energies, and every materialized chain has a strictly positive informational mass from the two-bit floor upward. So nothing crosses by accelerating. What the corpus does derive is a boundary of a different kind: when no edge is registered the pair sits in the virtual sector at the maximally raw-entropic state, one quarter of the identity, which is the unique argmax, and with no real edges there are no closed chains and hence no fermion termini at all. That partition — registered against unregistered — is the materialized sector's boundary, and crossing it is a move with a weight, not a limit of a velocity. The weight is a polynomial in the single crystallization threshold: four epsilon for dissolution, one minus four epsilon for crystallization, and at the sharp threshold the dissolution weight is exactly zero so nothing crosses back. The refutation is the useful half, because it relocates the question onto an object that exists and is finite, exact, and already carries a derived stationary measure.

\medskip
\subsection*{\texorpdfstring{B.1.33\quad BM.130 — The status of the Annihilation, Concrete Spin and Fermion Identification Tier: adjudication, scope, and the open items}{BM.130 B.1.33}}
\addcontentsline{toc}{subsection}{B.1.33 BM.130}
\phantomsection\label{tag:BM-130}
{\small \textit{Record.}\quad Source: Part 4 (Part4\_Baryomesonic\_era.md); tier: Tier BM.22 — The Annihilation Channel, Concrete Spin, and Fermion Identification: the Janus Tunnel, the Zig-Zag World Line, and the Spin-Dependent Transport Defect; canonical tag \texttt{BM.130}. Reproduced from the revision-1 canon; a cross-reference stub marks the former locus in the main corpus.}
\medskip
\textbf{Thm.}

%dsp-level:0
\[
(1)\quad\textbf{the objectives and their verdicts:}\qquad\begin{array}{l|l|l}\text{objective} & \text{verdict} & \text{locus}\\\hline\text{the Janus tunnel} & \textbf{constructed, on the register}\ —\ \text{with the velocity reading refuted as formed} & [\mathrm{Def\text{-}BM.18};\ \mathrm{BM.119}(6);\ \mathrm{BM.120}]\\\text{the zig-zag world line} & \textbf{derived}\ —\ \text{the even resolvent, every leg forward} & [\mathrm{BM.121}(2),(3),(6)]\\\text{pair production} & \textbf{the channel derived; the amplitude}\ \neg\textbf{formed on four grounds} & [\mathrm{BM.122}(1)\text{--}(5),\ (6)]\\\text{annihilation} & \textbf{the channel derived;}\ (\mathrm{R9})\ \textbf{closed at the channel level and narrowed to}\ (\mathrm{R9}'),(\mathrm{R9}'') & [\mathrm{BM.123}(1)\text{--}(6)]\\\text{concrete spin} & \textbf{derived}\ —\ \Theta=4j,\ \text{the multiplet ladder, the standing zero, the ceiling} & [\mathrm{BM.124};\ \mathrm{BM.125}]\\\text{the spin-dependent transport defect}\ \Delta(j) & \neg\textbf{formed as a single value: it is a ladder of}\ 2j+1\ \text{values, and}\ \Theta_R\ne2\pi k\ \text{by dimension} & [\mathrm{BM.124}(3),(4),(5)]\\\text{the spin-dependent Janus crossing} & \neg\textbf{formed: hypothesis (W7) refuted on}\ (\mathrm O1)\text{--}(\mathrm O4) & [\mathrm{BM.126}(1)\text{--}(3),\ (7)]\\\text{the fermion/boson identification} & \textbf{derived as a ledger read, with the two reads separated} & [\mathrm{BM.127}(1),(2),(4),(5)]\\\text{the meson is a boson} & \textbf{derived unconditionally, and robust under the orientation-sign row} & [\mathrm{BM.129}(1),(2)]\\\text{the baryon is a fermion} & \textbf{derived as an exact criterion; its value per named class awaits}\ (\mathrm{R3}) & [\mathrm{BM.129}(4),(5),(6)]\\\text{the three-chain exchange structure} & \textbf{constructed;}\ (\mathrm{R7})\ \textbf{discharged at}\ n=3,\ \text{narrowed to}\ (\mathrm{R7}') & [\mathrm{BM.128}(1)\text{--}(8)]\end{array}
\]

%dsp-level:0
\[
(2)\quad\textbf{the inventory gained:}\qquad\begin{array}{l|l}\text{new object} & \text{its exact content}\\\hline\mathcal J,\ \hat J_\pm,\ Q & \text{the pair register, two nilpotent moves anticommuting to the identity, and}\ Q=x_+(\hat J_+-P_{\mathcal V})+x_-(\hat J_--P_{\mathcal R})\ \text{with}\ \operatorname{spec}\{0,-(x_++x_-)\},\ \operatorname{rk}1\ [\mathrm{Def\text{-}BM.18};\ \mathrm{BM.120}]\\\mathcal T_{\mathcal J}(c) & \mathcal J\otimes\mathcal H_\omega\otimes\Psi,\qquad\dim=4(|k|+1)\ [\mathrm{Def\text{-}BM.18}(4)]\\\text{the zig-zag} & \text{the even resolvent}\ 1/(1-w^2x_+x_-),\ \ \mathcal W_m=(x_+x_-)^m\le4^{-m},\ \text{pole at}\ |w|\ge2\ [\mathrm{BM.121}]\\\mathcal W_{\text{channel}},\ \mathcal W_{\text{ann}},\ \mathcal W_{\text{round}} & x_+^2\chi,\quad x_-^2\chi,\quad x_+x_-\le\tfrac14;\qquad\text{the ratio}\ (x_-/x_+)^2=1\iff\varepsilon_q=\tfrac18\ [\mathrm{BM.122}(3);\ \mathrm{BM.123}(1),(2)]\\1-v & \frac{m^2}{E(E+p)}>0\ \ \text{at every finite}\ p\ [\mathrm{BM.119}(2)]\\\Theta=4j & \text{the exact spin--winding relation on the cone}\ [\mathrm{BM.124}(1)]\\\Delta(m) & 4\sin^2\frac{m\Theta_R}4=2-2T_{|m|}(\cos\frac{\Theta_R}2),\qquad\#\{m\}=2j+1,\qquad\Delta(0)=0\ \text{for even}\ k\ [\mathrm{BM.124}(3),(6)]\\\chi(\sigma)=d^{\,c(\sigma)} & \text{the}\ S_n\ \text{character on}\ (\mathbb C^d)^{\otimes n}\ [\mathrm{BM.128}(3)]\\\mathcal H_3 & \dim=8\prod_a(|k_a|+1);\qquad(64,\ 216,\ 96,\ 192)\ \text{at the four anchors}\ [\mathrm{BM.128}(1)]\\\dim\wedge^3 & \binom23=0\ \text{on the carrier against}\ \binom33=1\ \text{on colour}\ [\mathrm{BM.128}(5)]\\\varepsilon(\Gamma)=\prod_a\varepsilon(c_a) & \text{the multiplicative exchange law at an open connector}\ [\mathrm{BM.128}(7)]\\\varepsilon_M,\ \varepsilon_B & +1\ \text{unconditionally};\qquad(-1)^{k_1+k_2+k_3}\ \text{with the odd-count criterion}\ [\mathrm{BM.129}(1),(4)]\end{array}
\]

%dsp-level:0
\[
(3)\quad\textbf{the derived negatives (each with its ground):}\qquad\begin{array}{l|l}\text{negative} & \text{ground}\\\hline\neg\text{a velocity threshold and}\ \neg\text{an acceleration to}\ c & c=1\ \text{is a convention};\ 1-v=m^2/(E(E+p))>0\ \text{exactly}\ [\mathrm{BM.119}(1),(2),(3)]\\\neg\text{a crossing operator on}\ \mathcal H_\omega & \hat\omega^2=\mathrm{id}\ \text{against}\ \hat J_\pm^2=0:\ \neg\text{an algebra isomorphism}\ [\mathrm{BM.120}(6)]\\\neg\text{a backward leg in the zig-zag} & \neg\exists\tau^{-1};\ \text{the pure reversal is inadmissible};\ \Theta_{\mathfrak L}\ \text{fixes every completed word}\ [\mathrm{BM.121}(6)]\\\neg\text{a two-transit-quantum supply vertex} & (\mathrm O1)\text{--}(\mathrm O4):\ \neg\text{a gauged chiral coupling, deficit}\ 1',\ \neg\text{a momentum resolvent,}\ \neg\text{a winding carrier on transit}\ [\mathrm{BM.122}(6)]\\\neg\text{a cross-section, an S-matrix element or a phase space} & \mathrm{BM.6}(6);\ \mathrm{Def\text{-}BM.18}(6);\ \text{item 17}\\\neg\text{a direct}\ \mathrm{EM}^+\to\mathrm{EM}^-\ \text{operator} & \text{the sectors are amplitude-equal orbits; the channel is mediated through}\ \mathcal V\ [\mathrm{BM.123}(5)]\\\neg\text{a single-valued}\ \Delta(j) & \#\{m\}=2j+1\ge2\ [\mathrm{BM.124}(4)]\\\Theta_R=2\pi k\ \text{is impossible} & [\Theta_R]=\mathfrak q^0\ \text{against}\ [k]=\mathfrak q^{+1}\ [\mathrm{FT.55}(1);\ \mathrm{BM.124}(5)]\\\neg\text{a spin-gated tunnel} & \text{hypothesis (W7), the four grounds}\ (\mathrm O1)\text{--}(\mathrm O4)\ [\mathrm{BM.126}(2),(3)]\\\neg\text{a product mass form}\ \hat\mu'\ell_c\sqrt{G\Delta} & \text{three grounds: no single}\ \Delta,\ \text{two reads not two factors, dimensions}\ [\mathrm{BM.125}(4)]\\\neg\text{a totally antisymmetric carrier state at}\ n=3 & \dim\wedge^3(\mathbb C^2)=0\ [\mathrm{BM.128}(5)]\\\neg\text{a flavour}\ \mathsf{SU}(3),\ \neg\text{an octet,}\ \neg\text{a decuplet} & S_3\ \text{permutes slots};\ 8\ \text{occurs in no grouping family}\ [\mathrm{BM.117}(3);\ \mathrm{BM.128}(9);\ \mathrm{BM.129}(7)]\\\neg\text{a massless}\ \nu\ \text{line} & \mathrm{BM.96}(3),(4)\text{'s two SOS sandwiches stand}\ [\mathrm{BM.125}(5)]\\\neg\text{a cosmological coordinate} & \text{the ledger owns}\ \tau,\ G,\ N,\ E\ \text{and no global time}\ [\mathrm{Def\text{-}BM.17}(6)(\mathrm i)]\end{array}
\]

%dsp-level:0
\[
(4)\quad\textbf{the declarations this tier adds: none.}
\]

%dsp-level:0
\[
\qquad \neg\text{an}\ \mathrm{A\text{-}BM}\ \text{addendum is issued};
\]

%dsp-level:0
\[
\qquad \neg\text{a P4 member is added};
\]

%dsp-level:0
\[
\qquad \text{§Cal-BM keeps}\ \mathrm{A\text{-}BM.9}\text{'s logged slot with its value}\ \neg\text{issued and gains nothing}\ [\mathrm{P4};\ \mathrm{P1};\ \text{machine: the packaging census}]:
\]

%dsp-level:0
\[
\textbf{the constituent clauses are constructions on derived objects and read specifications, each priced:}
\]

%dsp-level:0
\[
\qquad \mathrm{Def\text{-}BM.18}(1),(2),(5)\text{'s register, moves and three names};
\]

%dsp-level:0
\[
\qquad \mathrm{Def\text{-}BM.19}(3)\text{'s locus of the spin-dependence, forced by}\ \mathrm{FT.55}(1)\text{'s no-go};
\]

%dsp-level:0
\[
(5)\quad\textbf{the rescope record (each row re-scoped at its own locus, none retracted, replaced or renumbered):}\qquad\begin{array}{l|l}\text{row} & \text{the rescope}\\\hline\mathrm{BM.118}(7)\text{'s}\ (\mathrm{R9}) & \textbf{closed at the channel level and narrowed to}\ (\mathrm{R9}'),(\mathrm{R9}'')\ [\mathrm{BM.123}(6)]\\\mathrm{BM.112}(7),\ \mathrm{Def\text{-}BM.16}(7)\text{'s}\ (\mathrm{R7}) & \textbf{discharged at}\ n=3\ \text{and narrowed to}\ (\mathrm{R7}')\ \text{for}\ n\ge4\ [\mathrm{BM.128}(8)]\\\mathrm{BM.109}(2)\text{'s}\ \mathcal W_{\text{pair}} & \textbf{sharpened: it is the round-trip weight}\ x_+x_-;\ \text{the one-way weights are}\ x_+\ \text{and}\ x_-\ [\mathrm{BM.121}(4);\ \mathrm{BM.120}(5)]\\\mathrm{BM.27}(7),\ \mathrm{BM.29}(5),\ \mathrm{Def\text{-}BM.15}(5),\ \mathrm{BM.108}(7),\ \mathrm{BM.109}(5) & \textbf{frozen on their own loci:}\ \neg\text{a}\ \textbf{direct}\ \text{inter-sector operator exists; the mediated channel is new}\ [\mathrm{BM.123}(5)]\\\mathrm{BM.114}(2)\text{'s}\ \text{colour-antisymmetry row} & \textbf{sharpened by the wedge dimensions}\ (0\ \text{against}\ 1)\ [\mathrm{BM.128}(6)]\\\mathrm{BM.114}(4)\text{'s}\ \text{orbit-count rows} & \textbf{stand and are}\ \neg\text{promoted to flavour multiplets}\ [\mathrm{BM.129}(7)]\\\mathrm{BM.106}(1)\text{'s}\ \text{spin-parity refutation} & \textbf{stands:}\ \varepsilon\ \text{is a function of}\ k,\ \neg\text{of the}\ \mathrm{D1.5}\ \text{kind}\ [\mathrm{BM.127}(4)]\\\mathrm{BM.90}(5),\ \mathrm{BM.105}(7) & \textbf{extended by one hypothesis:}\ (\mathrm{W7})\ \text{falls to the same four grounds}\ [\mathrm{BM.126}(3)]\\\mathrm{BM.126}(5)\text{'s}\ \text{domain rows} & \textbf{new:}\ \text{the bound/transit scope of}\ \varepsilon\ \text{and of the register's moves}\\\mathrm{FT.76}(3),\ \mathrm{FT.79}(2),\ \mathrm{FT.97}(5) & \textbf{untouched:}\ \lambda^{(n)}\le4n\ \text{is a sum of single-chain ceilings}\ [\mathrm{BM.125}(2);\ \mathrm{BM.117}(4)]\\\text{every clause of}\ \mathrm{BM.1}\text{--}\mathrm{BM.118} & \textbf{no retraction, no replacement, no renumbering}\ [\text{machine: the clause census}]\end{array}
\]

%dsp-level:0
\[
(6)\quad\textbf{the standing rules:}
\]

%dsp-level:0
\[
\qquad \neg\text{a new dimensionful scale}\ [\mathrm{P1}];
\]

%dsp-level:0
\[
\qquad \neg\text{a measured value and}\ \neg\text{a numerical coupling}\ [\mathrm{P5};\ \mathrm{EM.25}(7),(9);\ \text{machine: zero external tokens, zero decimal numerals, zero}\ [\mathrm{EXT}]\ \text{markers in this tier}];
\]

%dsp-level:0
\[
\text{no measured external constant is identified with any corpus quantity}\ [\mathrm{P5}];
\]

%dsp-level:0
\[
\qquad \neg\text{an observed lifetime, width, cross-section, branching fraction or mass is imported}\ [\mathrm{BM.44}(5);\ \mathrm{BM.99}(5);\ \mathrm{BM.118}(6)];
\]

%dsp-level:0
\[
\text{species names are confined to the identification layers}\ [\mathrm{Def\text{-}BM.3}(7);\ \mathrm{Def\text{-}BM.6}(7);\ \mathrm{Def\text{-}BM.14}(8);\ \text{machine: the}\ \textbf{Thm}\text{-clause scan}];
\]

%dsp-level:0
\[
\qquad \text{topology labels are internal bookkeeping}\ [\mathrm{Def\text{-}H.32}(4)];
\]

%dsp-level:0
\[
\text{amplitude}\ne\text{mass}\ [\mathrm{FT.79}(5);\ \mathrm{BM.125}(6)];
\]

%dsp-level:0
\[
\qquad \text{rung}\ne\text{scale}\ [\mathrm{D0.8}];
\]

%dsp-level:0
\[
\qquad \neg\text{a currency is converted}\ [\mathrm{BM.83}(3);\ \mathrm{BM.102}(6);\ \mathrm{BM.125}(6)];
\]

%dsp-level:0
\[
\alpha_B:\quad\neg\text{added},
\]

%dsp-level:0
\[
\quad\neg\text{derived},
\]

%dsp-level:0
\[
\quad\text{barred}\ —\ \mathrm{BM.61}(8)\text{'s}\ (\mathrm O1)\text{--}(\mathrm O4)\ \text{stand, and they are the same four grounds that refute}\ (\mathrm{W7})
\]

%dsp-level:0
\[
\qquad [\text{machine: the token scan of this tier}];
\]

%dsp-level:0
\[
\neg\text{a Yukawa-type object exists or is added},
\]

%dsp-level:0
\[
\qquad \neg\text{a seesaw-type mechanism is suggested},
\]

%dsp-level:0
\[
\qquad \neg\text{a right-ray matter line is added},
\]

%dsp-level:0
\[
\qquad \neg\text{an external mixing matrix is imported}
\]

%dsp-level:0
\[
\qquad [\mathrm{EW.15};\ \mathrm{FT.145}(3);\ \mathrm{BM.96}(4),(6);\ \mathrm{BM.106}(6);\ \mathrm{BM.122}(8);\ \mathrm{BM.125}(5);\ \text{machine}];
\]

%dsp-level:0
\[
\neg\text{a cosmological coordinate is used and}\ \neg\text{a black-hole physics is derived:}\quad\text{the accelerator and deep-gravity readings are evaluations of one threshold function at local}\ (E,N,G)\ [\mathrm{BM.115}(6);\ \mathrm{Def\text{-}BM.17}(6)(\mathrm i);\ \mathrm{P3}:\ r_s\ \text{quarantined}];
\]

%dsp-level:0
\[
\neg\text{a HistLog disposition imported as premise}\ —\ \text{Q7, Q9, Q11, items 7, 13, 16(b), 17, 19, 20 and the deficits are cited as pointers};
\]

%dsp-level:0
\[
(7)\quad\textbf{the boundary inventory (every row dispositioned in full at BM.143):}\qquad\begin{array}{l|l}(\mathrm{R1})\ \ \rho_{\mathrm B}\text{'s value} & \text{the declared slot}\ \mathrm{A\text{-}BM.9};\ \neg\text{derived on}\ \mathrm{BM.69}(4)\text{'s four grounds};\ \neg\text{fixed by measurement}\ [\mathrm{P5}]\\(\mathrm{R2})\ \ \text{the loop realization} & \text{the vertex-figure locus}\ \Theta_R=\delta_v\ \text{as a ledger configuration}\ [\mathrm{BM.67}(9),(10);\ \mathrm{BM.82}(6);\ \mathrm{Def\text{-}BM.19}(7);\ \textbf{inherited by every row of}\ \mathrm{BM.124}\text{--}\mathrm{BM.125}]\\(\mathrm{R3})\ \ \text{the flavour composition} & \text{which three of the six states a class occupies, and which winding each carries}\ [\mathrm{BM.28}(7);\ \mathrm{BM.103}(8);\ \mathrm{BM.127}(6);\ \mathrm{BM.129}(6)]\\(\mathrm{R4}')\ \ \text{the pair correlations} & \text{the ensemble's configuration distribution beyond per-pair independence}\ [\mathrm{BM.105}(6);\ \mathrm{BM.89}(8)]\\(\mathrm{R5})\ \ \iota\ \text{and the interior multipoles} & \text{the internal charge resolution the end signs cannot carry}\ [\mathrm{BM.46}(2);\ \mathrm{BM.87}(3)]\\(\mathrm{R6})\ \ \text{the Jensen residue} & \langle E\rangle_{\mathrm{mat}}\langle E^{-1}\rangle_{\mathrm{mat}}-1\ [\mathrm{BM.94}(8);\ \mathrm{BM.98}(5)]\\(\mathrm{R7}')\ \ \text{the exchange structure at}\ n\ge4 & \mathcal H_n\ \text{and its}\ S_n\ \text{action}\ [\mathrm{BM.128}(8);\ \textbf{new}];\ \ \text{the grammar mixes}\ \delta\ \text{and}\ \varepsilon\ \text{from}\ N=4\ [\mathrm{BM.31}(2)]\\(\mathrm{R8})\ \ \text{the}\ 2\text{-chain sector masses} & m^2_B(\cdot,S)\big|_{n=2}\ [\mathrm{BM.111}(3)]\\(\mathrm{R9}')\ \ \text{the annihilation cross-section} & \text{requires the momentum-space}\ \Sigma\text{-resolvent}\ [\mathrm{BM.6}(6)]\ \textbf{and}\ \text{the gauge coupling of the charge read}\ [\mathrm{Def\text{-}BM.8}(8);\ \textbf{new}]\\(\mathrm{R9}'')\ \ \text{the radiation legs} & \neg\text{a transit-quantum supply vertex}\ [\mathrm{BM.122}(6)(\mathrm O1)\text{--}(\mathrm O4);\ \textbf{new}]\\(\mathrm{R10})\ \ \text{the connector winding}\ d & \Theta(P)\ \text{per sealed class; it alone fixes}\ \varepsilon_M\ \text{if the connector is closed}\ [\mathrm{BM.2}(2);\ \mathrm{BM.129}(3);\ \textbf{new}]\\\mathsf g,\ e_0,\ N & \text{the curvature composition, the electric coupling's value, the lapse as solution data}\ [\mathrm{BM.104}(5);\ \mathrm{EM.25}(7);\ \mathrm{CT2.7}]\\m_\nu\ \text{absolute},\ \mu_\nu,\ i_\nu,\ \text{the mixing matrix} & \neg\text{derived}\ [\mathrm{BM.96}(3);\ \mathrm{BM.106}(6)]\\\text{the generation count as a selection} & \text{three lines derived, three selected}\ \neg\text{derived}\ [\mathrm{BM.60}(1)]\\\text{the two mass ceilings} & \text{deficit}\ 12\ \text{stands}\ [\mathrm{FT.76}(3);\ \mathrm{FT.79}(2);\ \mathrm{FT.97}(5);\ \mathrm{BM.125}(2)]\end{array}
\]

%dsp-level:0
\[
(8)\quad\textbf{the adjudication:}
\]

%dsp-level:0
\[
\qquad \boxed{\ \textbf{four directives, eleven derived or adjudicated verdicts, zero declarations, four new named residues}\ }
\]

%dsp-level:0
\[
\textbf{the cost:}
\]

%dsp-level:0
\[
\qquad \text{four new residues}\ (\mathrm{R7}'),\ (\mathrm{R9}'),\ (\mathrm{R9}''),\ (\mathrm{R10}),\ \text{two narrowed from prior ones};
\]

%dsp-level:0
\[
\qquad \text{zero parameters fitted};
\]

%dsp-level:0
\[
\qquad \text{zero scales fixed};
\]

%dsp-level:0
\[
\qquad \text{zero measured values};
\]

%dsp-level:0
\[
\qquad \text{zero declarations};
\]

%dsp-level:0
\[
\textbf{the gain:}
\]

%dsp-level:0
\[
\qquad \text{a crossing operator with an exact spectrum and a derived stationary measure};
\]

%dsp-level:0
\[
\ \text{the zig-zag as the resolvent's evenness with every leg forward and no inverse tick};
\]

%dsp-level:0
\[
\text{a pair channel with an exact threshold, weight and charge balance};
\]

%dsp-level:0
\[
\ \text{an annihilation channel with an exact ratio, a detailed-balance indicator and an entropy destination};
\]

%dsp-level:0
\[
\ (\mathrm{R9})\ \text{closed at the channel level};
\]

%dsp-level:0
\[
\text{the exact spin--winding relation}\ \Theta
\]

%dsp-level:0
\[
= 4j;
\]

%dsp-level:0
\[
\ \text{the multiplet defect ladder with its Chebyshev form, its standing zero and its ceiling};
\]

%dsp-level:0
\[
\ \text{the dimensional bar against fusing the two angles};
\]

%dsp-level:0
\[
\text{the}\ S_n\ \text{character formula}\ d^{\,c(\sigma)};
\]

%dsp-level:0
\[
\ \mathcal H_3\ \text{and}\ (\mathrm{R7})\ \text{discharged at}\ n
\]

%dsp-level:0
\[
= 3;
\]

%dsp-level:0
\[
\ \text{the vanishing carrier wedge against the one-dimensional colour wedge};
\]

%dsp-level:0
\[
\text{the multiplicative exchange law};
\]

%dsp-level:0
\[
\ \text{the meson's unconditional boson sign};
\]

%dsp-level:0
\[
\ \text{the baryon's exact fermion criterion};
\]

%dsp-level:0
\[
\textbf{deficit 11 stays closed};
\]

%dsp-level:0
\[
\qquad (\mathrm{R7})\ \text{and}\ (\mathrm{R9})\ \text{are discharged and narrowed};
\]

%dsp-level:0
\[
\qquad \text{deficits}\ 1',\ 2,\ 3,\ 4,\ 6,\ 12\ \text{stand};
\]

%dsp-level:0
\[
\qquad \text{nothing measured entered anywhere}.
\]

\textbf{Req.} D0.8, D1.4, D1.5, D2.9 (Parts 1–2); CT2.7, CT3.22, EM.25(7), EM.25(9), EW.15 (Part 1, Part 2); P1, P3, P4, P5 (Parts 1–2); Def-H.32(4), Def-H.38, FT.76(3), FT.79(2), FT.79(5), FT.89(2), FT.97(5), FT.125(5), FT.145(3) (Part 3); Def-BM.3(7), Def-BM.6(7), Def-BM.8(8), Def-BM.14(8), Def-BM.15(5), Def-BM.16(7), Def-BM.17(6), Def-BM.18(6), Def-BM.19(7), BM.2(2), BM.6(6), BM.27(7), BM.28(7), BM.29(5), BM.31(2), BM.44(5), BM.46(2), BM.60(1), BM.61(8), BM.67(9), BM.67(10), BM.69(4), BM.82(6), BM.83(3), BM.87(3), BM.89(8), BM.90(5), BM.94(8), BM.96(3), BM.96(4), BM.96(6), BM.98(5), BM.99(5), BM.102(6), BM.103(8), BM.104(5), BM.105(6), BM.105(7), BM.106(1), BM.106(6), BM.108(7), BM.109(2), BM.109(5), BM.111(3), BM.112(7), BM.113, BM.114(2), BM.114(4), BM.115(4), BM.115(6), BM.117(3), BM.117(4), BM.118(6), BM.118(7), BM.119, BM.120, BM.121, BM.122, BM.123, BM.124, BM.125, BM.126, BM.127, BM.128, BM.129, A-BM.9 (this part); the machine sections §A–§F green [math; machine: \texttt{checks/\allowbreak{}verify\_\allowbreak{}H\_\allowbreak{}janus\_\allowbreak{}tunnel\_\allowbreak{}concrete\_\allowbreak{}spin.py} §F].

\textbf{Pf.}

%dsp-level:0
\[
(1)\text{--}(2):\ \text{assembly of the derived and adjudicated rows of}\ \mathrm{BM.119}\text{--}\mathrm{BM.129}\ [\text{machine: §A--§E green}];
\]

%dsp-level:0
\[
\qquad \text{the inventory rows re-listed from their own tags}\ [\text{machine}]
\]

%dsp-level:0
\[
(3):\ \text{the fourteen negative rows re-cited with their grounds}\ [\text{machine}];
\]

%dsp-level:0
\[
\qquad (4):\ \text{the packaging census:}\ \neg\text{an}\ \mathrm{A\text{-}BM}\ \text{addendum is issued,}\ \neg\text{a P4 member is added}\ [\text{machine}]
\]

%dsp-level:0
\[
(5):\ \text{the clause census of}\ \mathrm{BM.1}\text{--}\mathrm{BM.118}\ \text{against retraction, replacement and renumbering}\ [\text{machine}];
\]

%dsp-level:0
\[
\qquad \text{the eleven rescope rows verified at their own loci}\ [\text{machine}]
\]

%dsp-level:0
\[
(6):\ \text{the standing-rule scans: external tokens, decimal numerals, unit tokens, registry tokens, declaration lines, the barred object}\ \alpha_B,
\]

%dsp-level:0
\[
\ \text{the Yukawa-type, seesaw-type, right-ray, mixing-matrix, cosmological and black-hole rows, and the}\ \textbf{Thm}\text{-clause species scan}\ [\text{machine}]
\]

%dsp-level:0
\[
(7):\ \text{the open rows re-cited from their own tags, with}\ (\mathrm{R7}'),
\]

%dsp-level:0
\[
\ (\mathrm{R9}'),
\]

%dsp-level:0
\[
\ (\mathrm{R9}''),
\]

%dsp-level:0
\[
\ (\mathrm{R10})\ \text{named here}\ [\text{machine}];
\]

%dsp-level:0
\[
\qquad (8):\ \text{assembly}.\ \blacksquare
\]

\textbf{Phys.} Four directives came in and eleven verdicts went out, of which six are plain derivations, four are refutations with grounds, and one is a closure with a narrowed successor. The tunnel exists, but not where the question put it: nothing reaches the speed of light, because the dispersion leaves an exact positive gap at every finite momentum and every materialized chain has positive mass from the two-bit floor up, so the crossing is not an acceleration but a move on a two-state register — a register whose generator has one stationary mode and one relaxation mode, whose kernel gives the stationary measure the corpus had already derived, and whose two moves are nilpotent and anticommute to the identity. The zig-zag is that register's resolvent, and the resolvent is even: every odd matrix element vanishes, so each term is a create-propagate-annihilate excursion weighted by the parabola to a power, with all legs forward and three independent reasons why no leg runs backward — the sharpest being that the corpus's own record involution already fixes every completed alternating word. Pair production and annihilation are channels with exact weights, an exact threshold and an exact charge balance, and the ratio between the two directions is one precisely where the round-trip weight peaks, which is the detailed-balance point. What is not formed is the supply vertex and the cross-section: four grounds bar the first and two open interfaces bar the second, so residue (R9) closes at the channel level and narrows into two named successors rather than disappearing. On the spin side the relation is exact — the winding is four times the spin — but the defect is a ladder and not a number, because the winding and the holonomy angle carry different dimensions and fusing them is barred by a no-go the corpus proved earlier; the ladder has two-j-plus-one rungs, an integer-polynomial form, a hard ceiling, and an exact massless middle weight precisely when the winding is even, which is precisely when the exchange sign is plus one. A spin-gated tunnel is refuted on the same four grounds that killed two earlier state-indexed thresholds, and what replaces it is a domain statement: the sign is defined on what crossed, not on the rate. The identification that follows is a ledger read and nothing more, kept strictly apart from the bound-or-transit label that cannot carry a statistics. Then three chains get their space, their group and their character formula, and one structural fact settles where antisymmetry lives: three copies of the two-dimensional carrier have no antisymmetric state at all while three colour slots have exactly one, so the baryon's antisymmetry is colour's and its statistics is the winding's, multiplied over constituents. The meson comes out a boson unconditionally, robust even against the corpus's unresolved orientation-sign convention; the baryon comes out fermionic exactly when an odd number of its constituents wind oddly, which the minimal floor satisfies. Four residues were added, two prior ones discharged and narrowed, and no measured number entered anywhere.

\medskip
\subsection*{\texorpdfstring{B.1.34\quad BM.144 — The status of the Closing Tier and the closure of the Baryomesonic era}{BM.144 B.1.34}}
\addcontentsline{toc}{subsection}{B.1.34 BM.144}
\phantomsection\label{tag:BM-144}
{\small \textit{Record.}\quad Source: Part 4 (Part4\_Baryomesonic\_era.md); tier: Tier BM.23 — The Closing Tier: the Haar–Poisson Duality and the Lapse, the Selected Curvature Composition, the Momentum-Space Resolvent, the Charge and Statistics Censuses, the Two-Chain Ensemble, the Seed-Tensoring Mechanism, the Convention Declarations, the Ceiling Study, the Jensen Residue, the Pair Amplitude, and the Era Closure; canonical tag \texttt{BM.144}. Reproduced from the revision-1 canon; a cross-reference stub marks the former locus in the main corpus.}
\medskip
\textbf{Thm.}

%dsp-level:0
\[
(1)\quad\textbf{the objectives and their verdicts:}\qquad\begin{array}{l|l|l}\text{objective} & \text{verdict} & \text{locus}\\\hline\text{the Haar--Poisson duality} & \textbf{derived}\ —\ \text{the twirl fixed point, the complementary factorization, the unique common fixed point} & [\mathrm{BM.131}(1)\text{--}(6)]\\\hline\text{the lapse}\ N(E) & \textbf{derived}\ —\ \text{the exact function, the formed domain, the envelope row; the gap closed} & [\mathrm{BM.132}]\\\hline\text{the curvature composition}\ \mathsf g(G) & \textbf{selected}\ —\ \frac1{1+G}\ \text{by two derived criteria; the running coupling complete} & [\mathrm{BM.133}]\\\hline\text{the momentum-space resolvent } & \textbf{constructed}\ —\ \text{the form uniqueness, the pole reads, the phase space; item 17 closed} & [\mathrm{Def\text{-}BM.20};\ \mathrm{BM.134}]\\\hline\text{the multi-chain charge read } & \textbf{censused and gauged}\ —\ \text{deficits 10 closed, 9 re-certified; Q10 closed} & [\mathrm{BM.135}]\\\hline\text{the}\ S_3\ \text{action and the statistics per class} & \textbf{derived}\ —\ \text{the colour-vertex equivariance, the exchange structure at every}\ n;\ (\mathrm{R7}')\ \text{closed} & [\mathrm{BM.136}]\\\hline\text{the 2-chain pole masses}\ (\mathrm{R8}) & \textbf{computed}\ —\ \text{the pincer zero, the formed ensemble binding, the 21-value spectrum} & [\mathrm{BM.137}]\\\hline\text{the seed-tensoring } & \textbf{derived}\ —\ \text{the twirl mechanism, the transport and composition closure; item 19 closed} & [\mathrm{BM.138}]\\\hline\text{the conventions } & \textbf{declared}\ —\ \text{four addenda; item 16 closed} & [\mathrm{BM.139}]\\\hline\text{deficit 12's ceilings} & \textbf{derived and explained}\ —\ \text{the dodecahedron selection, the five-currency census; deficit 12 closed as adjudicated} & [\mathrm{BM.140}]\\\hline\text{the Jensen residue}\ (\mathrm{R6}) & \textbf{computed}\ —\ \text{the closed form, the supremum enclosure, the companion} & [\mathrm{BM.141}]\\\hline\text{the pair amplitude and the cross-sections} & \textbf{constructed}\ —\ (\mathrm{R9}'),(\mathrm{R9}'')\ \text{closed; the four grounds resolved} & [\mathrm{BM.142}]\\\hline\text{the house cleaning} & \textbf{executed}\ —\ \text{the void closed, the phantoms reclassified, the register consolidated, Part 5 named} & [\mathrm{BM.143}]\end{array}
\]

%dsp-level:0
\[
(2)\quad\textbf{the inventory gained:}\qquad\begin{array}{l|l}\text{new object} & \text{its exact content}\\\hline\text{the duality rows} & \mathrm{tw}(\rho)=\frac{\mathfrak T(\Pi_{\mathcal R}\rho)}3\Pi_{\mathcal R}+\mathfrak T(\Pi_{\mathcal V}\rho)\Pi_{\mathcal V};\qquad\delta K=\gamma\rho+\log Z_\gamma I;\qquad\{\text{common fixed points}\}=\{\tfrac14I\}\ [\mathrm{BM.131}]\\\hline N(k,E) & \sqrt{1-\frac{G(E)}{4\pi k}},\qquad\text{formed}\iff k\ge\frac{G}{4\pi}=\frac{E_{\mathrm{int}}^2\ln2}2\ [\mathrm{BM.132}]\\\hline\mathsf g(G) & \frac1{1+G};\qquad e(r,E)=e_0\bigl[Nf/(1+G)\bigr]^{1/2}\ [\mathrm{BM.133}]\\\hline R(z;\mathbf p),\ d\mu_{\mathrm{shell}},\ \Phi_2,\ F & \text{the resolvent, the on-shell measure}\ \frac{d^3p}{2E},\ \ \Phi_2=\frac{\sqrt{z^2-4N^2m^2}}{8\pi N^3z},\ \ F_{\gamma\gamma}=\frac{2z^2}{N^2}\ [\mathrm{BM.134}]\\\hline w^{(n)}_{\mathrm{mat}},\ \mathcal B^{(n)}_{\mathrm{ens}} & \frac1{1+u+\cdots+u^{n-1}};\qquad-\frac{2nuD_{n-1}}{D_n}\cdot\frac{\ln2}{\mathfrak h}\ [\mathrm{BM.137}]\\\hline m_2^2(q,S) & 0\ \ \text{on every formed sector, by the pincer}\ [\mathrm{BM.137}(4)]\\\hline\mathcal H_n,\ \chi(\sigma) & 2^n\prod(|k_a|{+}1);\qquad d^{\,c(\sigma)}\ \ \text{at every}\ n\ [\mathrm{BM.136}(3)]\\\hline\Pi_k & \text{the}\ S_k\text{-twirl},\ \ \operatorname{rk}=k{+}1,\ \ \text{the carrier mechanism}\ [\mathrm{BM.138}]\\\hline J_3,\ J_2 & \frac{p(1-p)(9-10p+2p^2)}{2D^2};\qquad\frac{p(1-p)}{(2-p)^2}\le\frac18\ [\mathrm{BM.141}]\\\hline\widehat{\mathcal V}_{\mathrm{sup}},\ \mathcal M,\ \sigma & e^2(\hat J_+\otimes\hat J_+),\ \ \text{nilpotent};\qquad\sigma_{\mathrm{pair}},\ \sigma_{\mathrm{ann}}\ \text{exact symbolic},\ [\sigma]=\mathfrak q^{-1}\ [\mathrm{BM.142}]\end{array}
\]

%dsp-level:0
\[
(3)\quad\textbf{the derived negatives (each with its ground):}\qquad\begin{array}{l|l}\text{negative} & \text{ground}\\\hline\neg\text{a formed self-lapse at the envelope} & 1-\frac{E_{\mathrm{int}}^2\ln2}2<0\ \ \text{on every computed cell}\ [\mathrm{BM.132}(3)]\\\hline\neg\text{a self-consistency loop on}\ \mathsf g & \mathrm{BM.104}(4)\text{'s freedom census stands; the selection runs on}\ (\mathrm S1),(\mathrm S2)\ [\mathrm{BM.133}(5)]\\\hline\neg\text{a four-fold colour singlet} & \dim\operatorname{Inv}\bigl((\mathbb C^3)^{\otimes4}\bigr)=0\ [\mathrm{BM.136}(3)]\\\hline\neg\text{an octet, a nonet, a flavour multiplet} & 8\ \ \text{occurs in no grouping family}\ [\mathrm{BM.117}(3);\ \mathrm{BM.137}(6)]\\\hline\neg\text{a collision in the}\ n{=}2\ \text{pole ladder} & \text{the}\ 21\ \text{sums are distinct in}\ \mathbb M\ [\mathrm{BM.137}(5)]\\\hline\neg\text{a chart-frequency width at}\ \mathcal V{=}1 & [\Gamma^{\mathrm{chart}}]=\mathfrak q^{-1/2}\ \ \text{is a time,}\ \neg\text{a rate; the dressing is}\ \neg\text{declared}\ [\mathrm{BM.142}(7)]\\\hline\neg\text{a redshifted self-modulating rate} & \mathrm{BM.131}(7)\ \text{is a reparametrization,}\ \neg\text{a feedback loop}\ [\mathrm{BM.104}(4)\ \text{pattern}]\\\hline\neg\text{a numerical identification of the two deficit species} & \text{disjoint alphabets}\ [\mathrm{BM.37}(2);\ \mathrm{BM.131}(7)]\\\hline\neg\text{black-hole physics at the lapse boundary} & r_s\ \ \text{quarantined}\ [\mathrm{P3};\ \mathrm{BM.132}(6)]\\\hline\neg\text{a}\ \Theta_R{=}2\pi k\ \text{fusion},\ \neg\text{a product mass form} & \mathrm{BM.124}(5),\ \mathrm{BM.125}(4)\ \text{re-cited, untouched}\\\hline\neg\text{an internal multipole read} & \text{the trace-only law and the void theorem}\ [\mathrm{BM.143}(1)]\\\hline\neg\text{a W-species chiral gauging} & \mathrm{BM.142}(1)(\mathrm O1)\text{'s recorded bar}\\\hline\neg\text{a cosmological coordinate} & \text{the ledger owns}\ \tau,G,N,E\ \text{and no global time}\ [\mathrm{Def\text{-}BM.17}(6)(\mathrm i)]\end{array}
\]

%dsp-level:0
\[
(4)\quad\textbf{the declarations this tier adds: four.}
\]

%dsp-level:0
\[
\qquad \text{A-BM.10}\ (\aleph\text{-domain, pointwise}),
\]

%dsp-level:0
\[
\quad\text{A-BM.11}\ (\text{corner convention, variant B}),
\]

%dsp-level:0
\[
\quad\text{A-BM.12}\ (\text{resolvent interface}),
\]

%dsp-level:0
\[
\quad\text{A-BM.13}\ (\vee^2\ \text{locus})
\]

%dsp-level:0
\[
\qquad [\mathrm{BM.139}(7);\ \mathrm{Def\text{-}BM.20};\ \text{all four logged in §II.4 before use}];
\]

%dsp-level:0
\[
\neg\text{a P4 member is added by any of them};
\]

%dsp-level:0
\[
\qquad \neg\text{a dimensionful scale}\ [\mathrm{P1}];
\]

%dsp-level:0
\[
\qquad \text{§Cal-BM gains the consolidated register and}\ \neg\text{a parameter}\ [\mathrm{BM.143}(4)].
\]

%dsp-level:0
\[
(5)\quad\textbf{the rescope record (each row re-scoped at its own locus, none retracted, replaced or renumbered):}\qquad\begin{array}{l|l}\text{row} & \text{the rescope}\\\hline\mathrm{BM.88}(5),(8)\text{'s envelope zero and}\ N\text{-row} & \textbf{re-scoped:}\ \text{the vacuum-gauge read against the derived self-lapse}\ [\mathrm{BM.132}(5);\ \mathrm{BM.133}(6)]\\\hline\mathrm{BM.104}(5),(6) & \textbf{completed:}\ \neg\exists\text{-row scoped to its four grounds; the selection derived on two new criteria}\ [\mathrm{BM.133}(5)]\\\hline\mathrm{BM.111}(3)\text{'s}\ (\mathrm{R8}) & \textbf{discharged}\ [\mathrm{BM.137}(7)]\\\hline\mathrm{BM.115}(1)\text{'s threshold} & \textbf{substituted:}\ \mathcal T(k,E)\ \text{explicit in the derived lapse}\ [\mathrm{BM.132}(5)]\\\hline\mathrm{BM.118}(7),\ \mathrm{BM.130}(7)\text{'s}\ (\mathrm{R7}'),(\mathrm{R9}'),(\mathrm{R9}''),(\mathrm{R10}) & \textbf{discharged or reclassified}\ [\mathrm{BM.136}(6);\ \mathrm{BM.142}(8);\ \mathrm{BM.143}(3),(6)]\\\hline\mathrm{BM.122}(6),(7)\text{'s bar and census} & \textbf{resolved ground by ground; the}\ \neg\text{formed column discharged}\ [\mathrm{BM.142}(1)\text{--}(5)]\\\hline\mathrm{BM.129}(6),\ \mathrm{BM.130}(1)\text{'s per-class value row} & \textbf{discharged at the floor and six-state loci}\ [\mathrm{BM.136}(4)]\\\hline\mathrm{Def\text{-}BM.8}(8)\text{'s gauge-coupling row} & \textbf{discharged by}\ \mathrm{Def\text{-}BM.12}\ [\mathrm{BM.135}(5)]\\\hline\mathrm{Def\text{-}BM.16}(5)\text{'s sector clause} & \textbf{declared}\ [\text{A-BM.13};\ \mathrm{BM.139}(5)]\\\hline\text{item 21(b)'s coupling row} & \textbf{re-scoped:}\ \text{the reparametrization read is formed; the absolute-frequency datum stays calibration-species}\ [\mathrm{BM.131}(7)]\\\hline\mathrm{FT.76},\ \mathrm{FT.77}\text{'s conditional clauses} & \textbf{converted to declared data at their points of use}\ [\text{A-BM.11};\ \mathrm{BM.139}(4)]\\\hline\text{every clause of}\ \mathrm{BM.1}\text{--}\mathrm{BM.130} & \textbf{no retraction, no replacement, no renumbering}\ [\text{machine: the clause census}]\end{array}
\]

%dsp-level:0
\[
(6)\quad\textbf{the standing rules:}
\]

%dsp-level:0
\[
\qquad \neg\text{a new dimensionful scale}\ [\mathrm{P1}];
\]

%dsp-level:0
\[
\qquad \neg\text{a measured value,}\ \neg\text{a numerical coupling, mass, width, lifetime, cross-section or branching fraction}\ [\mathrm{P5};\ \text{machine: zero external tokens, zero decimal numerals, zero}\ [\mathrm{EXT}]\ \text{markers in this tier}];
\]

%dsp-level:0
\[
\text{no measured external constant is identified with any corpus quantity}\ [\mathrm{P5}];
\]

%dsp-level:0
\[
\qquad \text{species names are confined to the identification layers}\ [\mathrm{Def\text{-}BM.3}(7);\ \mathrm{Def\text{-}BM.6}(7);\ \mathrm{Def\text{-}BM.14}(8);\ \text{machine: the species scan of this tier's}\ \textbf{Thm}\ \text{clauses}];
\]

%dsp-level:0
\[
\text{topology labels are internal bookkeeping}\ [\mathrm{Def\text{-}H.32}(4)];
\]

%dsp-level:0
\[
\qquad \text{amplitude}\ne\text{mass}\ [\mathrm{FT.79}(5)];
\]

%dsp-level:0
\[
\qquad \text{rung}\ne\text{scale}\ [\mathrm{D0.8}];
\]

%dsp-level:0
\[
\alpha_B:\ \neg\text{added,}\ \neg\text{derived, barred}\ [\mathrm{BM.61}(8)];
\]

%dsp-level:0
\[
\qquad \neg\text{a Yukawa-type object}\ [\mathrm{BM.88}(1);\ \mathrm{Def\text{-}BM.11}(7)];
\]

%dsp-level:0
\[
\qquad \neg\text{a seesaw-type mechanism}\ [\mathrm{BM.96}(6)];
\]

%dsp-level:0
\[
\qquad \neg\text{a right-ray matter line}\ [\mathrm{BM.96}(4),(5);\ \mathrm{BM.96}(6)];
\]

%dsp-level:0
\[
\neg\text{an external mixing matrix}\ [\mathrm{FT.145}(3);\ \mathrm{BM.75}(6)];
\]

%dsp-level:0
\[
\qquad \neg\text{an imported ratio}\ [\mathrm{BM.30}(5);\ \mathrm{BM.37}(6)];
\]

%dsp-level:0
\[
\qquad \neg\text{a cosmological coordinate}\ [\mathrm{Def\text{-}BM.17}(6)(\mathrm i);\ \mathrm{BM.118}(6)];
\]

%dsp-level:0
\[
\neg\text{a theorem of this part imports a HistLog disposition as premise}\ [\text{machine: the pointer scan of this tier's}\ \textbf{Req}\ \text{lines}];
\]

%dsp-level:0
\[
(7)\quad\textbf{the ledger balanced (the final state):}\qquad\begin{array}{l|l}\text{ledger row} & \text{the final state}\\\hline\text{items 16, 17, 19} & \textbf{closed}\ [\mathrm{BM.139}(6);\ \mathrm{BM.134}(7);\ \mathrm{BM.138}(7)]\\\hline\text{Q10, Q11} & \textbf{closed}\ —\ \text{Q11 at every}\ n\ [\mathrm{BM.135}(7);\ \mathrm{BM.136}(3),(6)]\\\hline\text{deficits 9, 10} & \textbf{closed / re-certified}\ [\mathrm{BM.135}(4)]\\\hline\text{deficit 11} & \textbf{stays closed}\ [\mathrm{BM.113}(7);\ \mathrm{BM.136}]\\\hline\text{deficit 12} & \textbf{closed as adjudicated}\ [\mathrm{BM.140}(7)]\\\hline\text{deficits 2, 3, 4, 6} & \textbf{stand}\ —\ \text{bar, selection and calibration species, named in}\ (8)\ [\mathrm{BM.143}(6),(7)]\\\hline\text{residues}\ (\mathrm{R1})\text{--}(\mathrm{R10}) & \textbf{every row dispositioned}\ [\mathrm{BM.143}(6)]\\\hline\text{item 7} & \textbf{closed as an equation with its declared slot, unchanged}\ [\mathrm{A\text{-}BM.9};\ \mathrm{BM.98};\ \mathrm{BM.137}(4)]\\\hline\text{items 5, 12, 13, 20, 21} & \textbf{cross-era rows updated at the ledger, named for Part 5 where this part's share ends}\ [\mathrm{BM.143}(7)]\end{array}
\]

%dsp-level:0
\[
(8)\quad\textbf{the era closure:}
\]

%dsp-level:0
\[
\qquad \boxed{\ \textbf{the Baryomesonic era is closed.}\ \ \text{Its ledger is balanced, its books are shelved.}\ }
\]

%dsp-level:0
\[
\textbf{what was derived across the era:}
\]

%dsp-level:0
\[
\qquad \text{the sealed composites and their observability switch; the junction grammar and the}\ 4\mathbb Z\ \text{phase; the Noether layer from the Janus topology;}
\]

%dsp-level:0
\[
\text{the calibration lock and the neutral-class row; the flicker ensembles and the path-integrated mass; the CPT sectors and the unit-class mirror rows; the additive charge read, now gauged;}
\]

%dsp-level:0
\[
\text{the metric-energy layer with the dispersion, the lapse and the zero-point bound; the decay vertex closed with its kinematics; the EM self-interaction and the full-scope splitting with the exact ratio;}
\]

%dsp-level:0
\[
\text{the constituent transport splitting and the six-state assembly; the organic bridge with its priced residue; the chiral assignment issued; the mass-ratio census in four incommensurable currencies;}
\]

%dsp-level:0
\[
\text{the EM self-energy correction with the running coupling, now complete; the flicker integral and the rest-mass Q-value; the six objective resolutions with the curvature composition, now selected;}
\]

%dsp-level:0
\[
\text{the antimatter and resonance layer with the meson seal and the spin--statistics closure; the annihilation channel, the concrete spin and the fermion identification;}
\]

%dsp-level:0
\[
\text{and this tier's duality, lapse, resolvent, censuses, ensemble, mechanism, declarations, ceilings, residue, amplitude and cleaning;}
\]

%dsp-level:0
\[
\textbf{what remains for Part 5:}
\]

%dsp-level:0
\[
\qquad (\mathrm{P5.1})\text{--}(\mathrm{P5.10})\ \text{of}\ \mathrm{BM.143}(7)\ —\ \text{the EM binding channel first, the lepton mass read beside it};
\]

%dsp-level:0
\[
\textbf{the open items that are}\ \neg\textbf{baryomesonic:}
\]

%dsp-level:0
\[
\qquad \text{item 7's value side}\ (\mathrm{C\text{-}BM.1});
\]

%dsp-level:0
\[
\qquad \text{items 12, 21(d)}\ \text{(the selection)};
\]

%dsp-level:0
\[
\qquad \text{item 13's finite-}\mathfrak h\ \text{loci};
\]

%dsp-level:0
\[
\qquad \text{item 21(a),(b),(e),(i)};
\]

%dsp-level:0
\[
\qquad \text{deficits 2, 3, 4, 6}\ \text{(bar and calibration species)};
\]

%dsp-level:0
\[
\qquad \text{the W-species chiral gauging};
\]

%dsp-level:0
\[
\qquad \text{the}\ n\ge4\ \text{sealing cells};
\]

%dsp-level:0
\[
(9)\quad\textbf{the cost and the gain:}
\]

%dsp-level:0
\[
\qquad \textbf{the cost:}\ \ \text{four declarations, all conventions or interfaces, none adding a parameter};
\]

%dsp-level:0
\[
\qquad \text{zero scales fixed};
\]

%dsp-level:0
\[
\qquad \text{zero measured values};
\]

%dsp-level:0
\[
\qquad \textbf{the gain:}\ \ \text{the lapse, the composition, the resolvent, the phase space, the amplitudes and the cross-sections; the charge census gauged; the statistics per class; the exchange structure at every}\ n;
\]

%dsp-level:0
\[
\ \text{the two-chain ensemble with its formed binding; the seed-tensoring mechanism; the dodecahedron explained; the Jensen residue computed; every residue dispositioned};
\]

%dsp-level:0
\[
\boxed{\ \textbf{the era's every slotted objective is addressed: derived or adjudicated; every void residue is closed as void; every phantom gap is reclassified; the calibration items are logged; the books are shelved}\ }
\]

%dsp-level:0
\[
\qquad [\text{machine: the objective census of}\ (1)\ \text{against the tier's directives}].
\]

\textbf{Req.} D0.8, D3.4 (Part 1); CT3.22 (Part 2); P1, P3, P5 (Parts 1–2); Def-H.32(4), FT.79(5), FT.145(3) (Part 3); Def-BM.3(7), Def-BM.6(7), Def-BM.11(7), Def-BM.14(8), Def-BM.16(5), Def-BM.17(6), Def-BM.20, BM.1–BM.143 (this part, at the clauses cited in (1)–(8)); A-BM.7, A-BM.9, A-BM.10, A-BM.11, A-BM.12, A-BM.13 (§II.4, logged before use); the clause census of BM.1–BM.130 against retraction, replacement and renumbering, the standing-rule token scans, the objective census [machine: \texttt{checks/\allowbreak{}verify\_\allowbreak{}H\_\allowbreak{}bm\_\allowbreak{}closing\_\allowbreak{}tier.py} §L–§M].

\textbf{Pf.}

%dsp-level:0
\[
(1)\text{--}(3):\ \text{the assembly of the derived and adjudicated rows of}\ \mathrm{BM.131}\text{--}\mathrm{BM.143}\ [\text{machine: §A--§K green}];
\]

%dsp-level:0
\[
\qquad \text{the inventory rows re-listed from their own tags}\ [\text{machine}]
\]

%dsp-level:0
\[
(4):\ \text{the packaging census: four addenda, zero P4 members, zero scales}\ [\text{machine}];
\]

%dsp-level:0
\[
\qquad (5):\ \text{the rescope rows verified at their own loci and the clause census of}\ \mathrm{BM.1}\text{--}\mathrm{BM.130}\ \text{against retraction, replacement, renumbering}\ [\text{machine}]
\]

%dsp-level:0
\[
(6):\ \text{the standing-rule scans: external tokens, decimal numerals,}\ [\mathrm{EXT}]\ \text{markers, registry tokens, the barred-object rows}\ (\alpha_B,\ \text{Yukawa-type, seesaw-type, right-ray matter line, mixing matrix, imported ratio, cosmological-coordinate rows}),
\]

%dsp-level:0
\[
\ \text{and the}\ \textbf{Thm}\text{-clause species scan}\ [\text{machine}]
\]

%dsp-level:0
\[
(7),
\]

%dsp-level:0
\[
(8),
\]

%dsp-level:0
\[
(9):\ \text{the ledger census against the frozen residue tables and the HistLog closure pointers}\ [\text{machine: the row-by-row match}];
\]

%dsp-level:0
\[
\qquad \text{the era-closure assembly}.\ \blacksquare
\]

\textbf{Phys.} The closing tier closes. Fourteen directives came in and every one left with a verdict: the duality derived, the lapse a function, the composition selected, the resolvent and phase space constructed, the charge read censused and gauged, the permutation action shown to be the colour vertex's slot read, the exchange structure extended to every chain count, the meson's ensemble formed with a strictly negative binding, the carrier mechanism exposed as a finite-group twirl, four conventions declared, the dodecahedron explained as the tightest convex vertex gap of the Platonic series, the Jensen residue an exact rational function, the supply vertex nilpotent and the cross-sections symbolic and exact, the void closed and the phantoms reclassified. The era's ledger is balanced row by row: every residue the baryomesonic era ever named now carries a disposition, and the rows that remain open are named as what they are — calibration slots, selection questions, cross-era constructions — none of them baryomesonic structure. The Baryomesonic era is closed; the Atomic era's requirements are written down; the books are shelved.

\medskip
\subsection*{\texorpdfstring{B.1.35\quad BM.152 — The status of the Continuum-Limit Tier and the closure of the foundational era}{BM.152 B.1.35}}
\addcontentsline{toc}{subsection}{B.1.35 BM.152}
\phantomsection\label{tag:BM-152}
{\small \textit{Record.}\quad Source: Part 4 (Part4\_Baryomesonic\_era.md); tier: Tier BM.24 — The Continuum-Limit Tier: the Force Operator over the Graph, the Continuum-Limit Function, the Metric–Energy–Continuum Map, the Breaking Threshold, the Locked-Sector Complement, the Universality, the Reclassification of the Continuum Conditions, and the Closure of the Foundational Era; canonical tag \texttt{BM.152}. Reproduced from the revision-1 canon; a cross-reference stub marks the former locus in the main corpus.}
\medskip
\textbf{Thm.}

%dsp-level:0
\[
(1)\quad\textbf{the objectives and their verdicts:}\qquad\begin{array}{l|l|l}\text{objective} & \text{verdict} & \text{locus}\\\hline\text{the force operator over the graph} & \textbf{derived}\ —\ \text{the census complete, the action composed of derived maps} & [\mathrm{BM.145}]\\\hline\text{the continuum-limit function} & \textbf{derived}\ —\ \text{five closed forms, four exact domain conditions, the}\ \mathfrak h\to0\ \text{row} & [\mathrm{BM.146}]\\\hline\text{the metric–energy–continuum map} & \textbf{derived}\ —\ \text{the dictionary bijection, the pure-trace closed form, the inverse read} & [\mathrm{BM.147}]\\\hline\text{the breaking threshold} & \textbf{derived}\ —\ E_{\mathrm{crit}}=2r/\mathfrak h\ \text{in three exact forms, the envelope corollary} & [\mathrm{BM.148}]\\\hline\text{the locked-sector complement} & \textbf{derived}\ —\ \text{the two faces of one capacity, the complement theorem} & [\mathrm{BM.149}]\\\hline\text{the universality} & \textbf{derived}\ —\ \text{one operator, one arena, one action, three fingerprints} & [\mathrm{BM.150}]\\\hline\mathrm{CT2.6}(\mathrm{ii})\text{--}(\mathrm{iii}) & \textbf{reclassified as derived}\ —\ \text{the extra-axiomatic list is empty} & [\mathrm{BM.151}]\\\hline\text{the foundational-era closure} & \textbf{issued} & [(7)]\end{array}
\]

%dsp-level:0
\[
(2)\quad\textbf{the inventory gained:}\qquad\begin{array}{l|l}\text{new object} & \text{its exact content}\\\hline\hat{\mathcal F}\ \text{as a composed map} & (\Gamma_{\mathcal R},\text{registrations})\mapsto(M,g_{\mu\nu},T_{\mu\nu},\operatorname{spec}):\ \ \text{the eight-step table}\ [\mathrm{BM.145}(3)]\\\hline\mathcal C(E,r) & \bigl(g_{\mu\nu},N,\mathcal T,w_{\mathrm{mat}},\Gamma\bigr)\ \ \text{in closed form, formed on}\ (\mathrm C1)\text{--}(\mathrm C4)\ [\mathrm{BM.146}(3),(4)]\\\hline g_{\mu\nu}(E,r) & \eta_{\mu\nu}+\frac{\mathfrak hE}{2r}\mathbb 1_4+O(\cdot):\ \ \text{pure trace, conformal mode alone}\ [\mathrm{BM.147}(2)]\\\hline E_{\mathrm{crit}},S_{\mathrm{crit}},k_{\mathrm{crit}} & \frac{2r}{\mathfrak h},\quad\frac{2r\ln2}{\mathfrak h},\quad\frac{E^2\ln2}2\ [\mathrm{BM.148}(1)]\\\hline\text{the complement row} & \text{chart}\ \neg\text{formed}\ \wedge\ \text{ledger exact}\ \wedge\ \text{membrane-governed}\ [\mathrm{BM.149}(7)]\\\hline\text{the three fingerprints} & (T\text{-channel},\ \text{mask slot},\ \text{coupling profile})\ [\mathrm{BM.150}(3)]\\\hline\text{the coefficient bound} & |r_a|\le\sqrt3;\qquad|\Delta|\le2\sqrt3;\qquad\mathrm{TV}\le2\sqrt3\,m\ [\mathrm{BM.151}(2)]\\\hline\text{the empty list} & \text{the extra-axiomatic inputs of}\ \text{A §0.2a}:\ \ \varnothing\ [\mathrm{BM.151}(3)]\end{array}
\]

%dsp-level:0
\[
(3)\quad\textbf{the derived negatives (each with its ground):}\qquad\begin{array}{l|l}\text{negative} & \text{ground}\\\hline\neg\text{an aspect or block outside the census} & \mathrm{A8};\ \mathrm{CT0.12};\ \mathrm{CT0.13};\ \mathrm{FT.83}(2)\text{--}(3)\\\hline\neg\text{a higher-curvature term} & \mathrm{CT2.7}:\ \text{the Lovelock closure in}\ D=4\\\hline\neg\text{an accessible}\ D^-\ \text{pole} & \mathrm{CT4.2}:\ \text{the one-way membrane}\\\hline\neg\text{a deleted ledger entry on absorption} & \mathrm{CT4.3}:\ \text{locked},\ \neg\text{destroyed}\\\hline\neg\text{a self-formed cell at its own envelope} & \mathrm{BM.148}(4):\ E^2\ln2>2\ \ \forall E\ge2\\\hline\neg\text{a numerical}\ E_{\mathrm{crit}}\ \text{or}\ \mathfrak h & \mathrm{P1};\ \mathrm{P5}:\ \text{both stay symbolic}\\\hline\neg\text{a nonlinear joint solution} & \mathrm{BM.43}(5)\ \text{re-cited},\ \neg\text{extended}\\\hline\neg\text{an extra-axiomatic input} & \mathrm{BM.151}(3):\ \text{the list is empty}\\\hline\neg\text{a force-dependent limit} & \mathrm{BM.150}(7):\ \text{three fingerprints only}\\\hline\neg\text{imported vocabulary} & \text{the tier's objects are the corpus's own}\ [\text{machine: the token scan}]\\\hline\neg\text{a cosmological coordinate} & \mathrm{Def\text{-}BM.17}(6)(\mathrm i);\ \mathrm{BM.118}(6)\end{array}
\]

%dsp-level:0
\[
(4)\quad\textbf{the declarations this tier adds: none.}
\]

%dsp-level:0
\[
\qquad \textbf{the declaration-class input dissolves:}\ \ \text{A §0.2a's extra-axiomatic list}\ [\mathrm{BM.151}(3)];
\]

%dsp-level:0
\[
\neg\text{a P4 member},
\]

%dsp-level:0
\[
\qquad \neg\text{a new dimensionful scale}\ [\mathrm{P1}],
\]

%dsp-level:0
\[
\qquad \text{§Cal-BM unchanged}\ [\mathrm{BM.143}(4)\text{'s register stands}].
\]

%dsp-level:0
\[
(5)\quad\textbf{the rescope record (each row re-scoped at its own locus, none retracted):}\qquad\begin{array}{l|l}\text{row} & \text{the rescope}\\\hline\text{A §0.2a's list sentence} & \textbf{amended in place: the list is empty; the reclassification markers cite this tier}\ [\mathrm{BM.151}(3)]\\\hline\mathrm{CT2.6}\text{'s condition labels and Req row} & \textbf{amended in place: reclassified as derived, conclusions untouched}\ [\mathrm{BM.151}(3)]\\\hline\mathrm{BM.8}(4)\text{'s}\ \text{Rmk}\ \text{and every point-of-use declaration} & \textbf{frozen as historical scope rows; they stand as scope-inheritance markers}\ [\mathrm{BM.151}(3)]\\\hline\mathrm{FT.156},\ \mathrm{BM.11}(3),\ \mathrm{BM.34}\ \text{and every}\ \mathrm{CT2.6}\text{-inheriting tag} & \textbf{unchanged: they inherit the same conditions, now derived at the same scope}\\\hline\text{every clause of}\ \mathrm{BM.1}\text{--}\mathrm{BM.144} & \textbf{no retraction, no replacement, no renumbering}\ [\text{machine: the clause census}]\end{array}
\]

%dsp-level:0
\[
(6)\quad\textbf{the standing rules:}
\]

%dsp-level:0
\[
\qquad \neg\text{a new dimensionful scale}\ [\mathrm{P1}];
\]

%dsp-level:0
\[
\qquad \neg\text{a measured value and}\ \neg\text{a numerical coupling, mass, width, lifetime or cross-section}\ [\mathrm{P5};\ \text{machine: zero external tokens, zero decimal numerals, zero}\ [\mathrm{EXT}]\ \text{markers in this tier}];
\]

%dsp-level:0
\[
\text{no measured external constant is identified with any corpus quantity}\ [\mathrm{P5}];
\]

%dsp-level:0
\[
\qquad \text{species names are confined to the identification layers}\ [\mathrm{Def\text{-}BM.3}(7);\ \mathrm{Def\text{-}BM.6}(7);\ \mathrm{Def\text{-}BM.14}(8);\ \text{machine: the species scan of this tier's}\ \textbf{Thm}\ \text{clauses}];
\]

%dsp-level:0
\[
\text{topology labels are internal bookkeeping}\ [\mathrm{Def\text{-}H.32}(4)];
\]

%dsp-level:0
\[
\qquad \text{amplitude}\ne\text{mass}\ [\mathrm{FT.79}(5)];
\]

%dsp-level:0
\[
\qquad \text{rung}\ne\text{scale}\ [\mathrm{D0.8}];
\]

%dsp-level:0
\[
\alpha_B:\ \neg\text{added},\ \neg\text{derived, barred}\ [\mathrm{BM.61}(8)];
\]

%dsp-level:0
\[
\qquad \neg\text{a Yukawa-type object};
\]

%dsp-level:0
\[
\qquad \neg\text{a seesaw-type mechanism}\ [\mathrm{BM.96}(6)];
\]

%dsp-level:0
\[
\qquad \neg\text{a right-ray matter line};
\]

%dsp-level:0
\[
\qquad \neg\text{an external mixing matrix}\ [\mathrm{FT.145}(3)];
\]

%dsp-level:0
\[
\neg\text{an imported ratio}\ [\mathrm{BM.30}(5)];
\]

%dsp-level:0
\[
\qquad \neg\text{a cosmological coordinate}\ [\mathrm{Def\text{-}BM.17}(6)(\mathrm i)];
\]

%dsp-level:0
\[
\qquad \neg\text{black-hole physics at any boundary locus}\ [\mathrm{P3};\ \mathrm{BM.148}(6)];
\]

%dsp-level:0
\[
\qquad \neg\text{a theorem of this part imports a HistLog disposition as premise}\ [\text{machine: the pointer scan of this tier's}\ \textbf{Req}\ \text{lines}];
\]

%dsp-level:0
\[
(7)\quad\textbf{the foundational-era closure:}
\]

%dsp-level:0
\[
\qquad \boxed{\ \textbf{Era 1 — the foundational corpus of Parts 1--4 (the qubit, bosonic, fermion and baryomesonic eras) — is closed.}\ }
\]

%dsp-level:0
\[
\textbf{the balance row:}
\]

%dsp-level:0
\[
\qquad \text{the four eras are balanced:}\ \ \text{the axiom surface}\ (\mathrm{A1}\text{--}\mathrm{A9},\ \mathrm{C1}\text{--}\mathrm{C12},\ \mathrm{D0.1}\text{--}\mathrm{D4.3},\ \mathrm{P1}\text{--}\mathrm{P5})\ \ \text{is the entire non-derived surface};
\]

%dsp-level:0
\[
\qquad \text{the extra-axiomatic list is empty}\ [\mathrm{BM.151}(3)];
\]

%dsp-level:0
\[
\text{the declared surface is exactly: the interface definitions with their priced non-derived clauses}\ [\mathrm{Def\text{-}H.35}\text{--}\mathrm{Def\text{-}H.43};\ \mathrm{Def\text{-}BM.20}],
\]

%dsp-level:0
\[
\ \ \text{the theorem-local dynamical scopes}\ [\mathrm{L0};\ \mathrm{CT4.5a};\ \mathrm{L4}],
\]

%dsp-level:0
\[
\ \ \text{the materialization scopes}\ [\mathrm{M1}/\mathrm{M2};\ \mathrm{Def\text{-}H.20}],
\]

%dsp-level:0
\[
\ \ \text{the identification and convention declarations}\ [\text{A-BM.1--A-BM.13}],
\]

%dsp-level:0
\[
\ \ \text{and the calibration registers}\ [\text{P4};\ \mathrm{Cal\text{-}F};\ \text{§Cal-BM}];
\]

%dsp-level:0
\[
\textbf{the single free dimensionful scale is}\ \mathfrak h\ [\mathrm{P1}];
\]

%dsp-level:0
\[
\qquad \textbf{the ledger is balanced: every residue dispositioned}\ [\mathrm{BM.143}(6);\ \mathrm{BM.144}(7)];
\]

%dsp-level:0
\[
(8)\quad\textbf{what Part 5 (the Atomic era) receives from the foundational corpus:}
\]

%dsp-level:0
\[
\qquad (\mathrm{P5.1})\text{--}(\mathrm{P5.10})\ \ \text{of}\ \mathrm{BM.143}(7)\ \ \text{re-cited as the demand ledger};
\]

%dsp-level:0
\[
\textbf{and the foundational additions:}
\]

%dsp-level:0
\[
\qquad \text{the unconditional chart layer}\ [\mathrm{BM.151}];
\]

%dsp-level:0
\[
\qquad \text{the capacity threshold of any bound-state construction}\ [\mathrm{BM.148}];
\]

%dsp-level:0
\[
\qquad \text{the locked-sector complement}\ [\mathrm{BM.149}];
\]

%dsp-level:0
\[
\qquad \text{the universal limit behind every scattering read}\ [\mathrm{BM.150};\ \mathrm{BM.134};\ \mathrm{BM.142}];
\]

%dsp-level:0
\[
(9)\quad\textbf{the open items that are}\ \neg\textbf{foundational-era structure:}
\]

%dsp-level:0
\[
\qquad \text{item 7's value side}\ (\mathrm{C\text{-}BM.1});
\]

%dsp-level:0
\[
\qquad \text{items 12, 21(d)}\ \ \text{(the selection)};
\]

%dsp-level:0
\[
\qquad \text{item 13's finite-}\mathfrak h\ \ \text{loci};
\]

%dsp-level:0
\[
\qquad \text{item 21(a),(b),(e),(i)};
\]

%dsp-level:0
\[
\qquad \text{deficits 2, 3, 4, 6}\ \ \text{(bar and calibration species)};
\]

%dsp-level:0
\[
\qquad \text{the W-species chiral gauging};
\]

%dsp-level:0
\[
\qquad \text{the}\ n\ge4\ \ \text{sealing cells};
\]

%dsp-level:0
\[
\qquad (\mathrm{R4}')\text{'s ensemble refinement};
\]

%dsp-level:0
\[
\qquad \text{the full nonlinear joint solve}\ [\mathrm{BM.43}(5)];
\]

%dsp-level:0
\[
(10)\quad\textbf{the era closure:}
\]

%dsp-level:0
\[
\qquad \boxed{\ \textbf{the foundational era is closed; its books are shelved; the corpus stands on its axioms alone}\ }
\]

%dsp-level:0
\[
\qquad [\text{machine: the objective census of}\ (1)\ \text{against the tier's directives}].
\]

\textbf{Req.} A1–A9 (Part 1); C1–C13 (Part 1); D0.1–D4.3 (Parts 1–2); P1–P5 (Parts 1–2); CT0.\textit{, CT1.}, CT2.\textit{, CT3.}, CT4.\textit{ (Parts 1–2, at the clauses cited in BM.145–BM.151); L0–L5 (Part 1); Def-H.}, FT.\textit{ (Part 3, at the clauses cited); Def-BM.}, BM.1–BM.151 (this part, at the clauses cited in (1)–(9)); A-BM.1–A-BM.13 (§II.4, logged before use); the clause census of BM.1–BM.144 against retraction, replacement and renumbering, the standing-rule token scans, the objective census [machine: \texttt{checks/\allowbreak{}verify\_\allowbreak{}H\_\allowbreak{}continuum\_\allowbreak{}limit\_\allowbreak{}era1\_\allowbreak{}closure.py} §H].

\textbf{Pf.}

%dsp-level:0
\[
(1)\text{--}(3):\ \text{the assembly of the derived and adjudicated rows of}\ \mathrm{BM.145}\text{--}\mathrm{BM.151}\ [\text{machine: §A--§G green}];
\]

%dsp-level:0
\[
\qquad \text{the inventory rows re-listed from their own tags}\ [\text{machine}]
\]

%dsp-level:0
\[
(4):\ \text{the packaging census: zero declarations issued, one declaration-class input retired}\ [\text{machine}];
\]

%dsp-level:0
\[
\qquad (5):\ \text{the amendment pins and the clause census of}\ \mathrm{BM.1}\text{--}\mathrm{BM.144}\ [\text{machine}]
\]

%dsp-level:0
\[
(6):\ \text{the standing-rule scans: external tokens, decimal numerals,}\ [\mathrm{EXT}]\ \text{markers, registry tokens, the barred-object rows, the imported-vocabulary scan, and the}\ \textbf{Thm}\text{-clause species scan}\ [\text{machine}]
\]

%dsp-level:0
\[
(7)\text{--}(10):\ \text{the balance census against}\ \text{A §0.2a's amended text, the declared-surface inventory, and the ledger pointers}\ [\text{machine: the row-by-row match}];
\]

%dsp-level:0
\[
\qquad \text{the era-closure assembly}.\ \blacksquare
\]

\textbf{Phys.} The capstone lands. The force operator acts over the graph through eight derived steps, and its continuum limit is one function of two ledger variables with four exact domain conditions and five closed forms. The metric is the dictionary, the energy determines it through a pure-trace linearized solve, the continuum is the quotient target, and the map's capacity is exact: two $r$ over $\mathfrak h$. Past capacity nothing is lost — the ledger retains what the chart cannot read, behind a one-way membrane with maximal locked entropy. The limit is universal because the operator, the arena and the action are one, and force identity lives in three derived fingerprints and nowhere else. And the two conditions that the corpus had carried as its only extra-axiomatic imports fall out of its own structure: the ladder witnesses the richness, the spectral bound witnesses the regularity. The list is empty. The axiom surface is the entire non-derived surface. Era 1 — the foundational corpus of four eras — is closed, its ledger balanced, its books shelved, and the atomic era's requirement list written.

\medskip
\subsection*{\texorpdfstring{B.1.36\quad BM.30 — The mass-route adjudication: three exact ratios, one read choice (catalog entry)}{BM.30 B.1.36}}
\addcontentsline{toc}{subsection}{B.1.36 BM.30}
\phantomsection\label{tag:BM-30-cat}
{\small \textit\{Catalog.\}\quad The full theorem remains in the main corpus as Theorem 4.25 (Part 4); its derived clauses are consumed by later derivations. The hypothesis-level ruling is recorded here verbatim:}
\begin{quote}\small \end{quote}
\medskip
\subsection*{\texorpdfstring{B.1.37\quad BM.31 — The boundary-data comparison and the grammar verdict (catalog entry)}{BM.31 B.1.37}}
\addcontentsline{toc}{subsection}{B.1.37 BM.31}
\phantomsection\label{tag:BM-31-cat}
{\small \textit\{Catalog.\}\quad The full theorem remains in the main corpus as Theorem 4.26 (Part 4); its derived clauses are consumed by later derivations. The hypothesis-level ruling is recorded here verbatim:}
\begin{quote}\small \end{quote}
\medskip
\subsection*{\texorpdfstring{B.1.38\quad BM.45 — The assembled reaction: the conservation table, the naming row, the final read, and the adjudication (catalog entry)}{BM.45 B.1.38}}
\addcontentsline{toc}{subsection}{B.1.38 BM.45}
\phantomsection\label{tag:BM-45-cat}
{\small \textit\{Catalog.\}\quad The full theorem remains in the main corpus as Theorem 4.38 (Part 4); its derived clauses are consumed by later derivations. The hypothesis-level ruling is recorded here verbatim:}
\begin{quote}\small \end{quote}
\medskip
\subsection*{\texorpdfstring{B.1.39\quad BM.49 — The beta-decay re-adjudication with the correction, and the endpoint-to-average refutation (catalog entry)}{BM.49 B.1.39}}
\addcontentsline{toc}{subsection}{B.1.39 BM.49}
\phantomsection\label{tag:BM-49-cat}
{\small \textit\{Catalog.\}\quad The full theorem remains in the main corpus as Theorem 4.42 (Part 4); its derived clauses are consumed by later derivations. The hypothesis-level ruling is recorded here verbatim:}
\begin{quote}\small The four thirds remains what the Tier BM.8 adjudication made it: the ratio of twenty neutral sign-assignments to fifteen charged ones, a price of the charge filter on the ensemble measure.\end{quote}
\medskip
\subsection*{\texorpdfstring{B.1.40\quad BM.63 — The decay re-adjudication on the refined read: the closure survives, with an exact margin (catalog entry)}{BM.63 B.1.40}}
\addcontentsline{toc}{subsection}{B.1.40 BM.63}
\phantomsection\label{tag:BM-63-cat}
{\small \textit\{Catalog.\}\quad The full theorem remains in the main corpus as Theorem 4.54 (Part 4); its derived clauses are consumed by later derivations. The hypothesis-level ruling is recorded here verbatim:}
\begin{quote}\small \end{quote}
\medskip
\subsection*{\texorpdfstring{B.1.41\quad BM.94 — The flicker integral: the snapshot read's category error, the materialized distribution, and the two crystallization weights (catalog entry)}{BM.94 B.1.41}}
\addcontentsline{toc}{subsection}{B.1.41 BM.94}
\phantomsection\label{tag:BM-94-cat}
{\small \textit\{Catalog.\}\quad The full theorem remains in the main corpus as Theorem 4.79 (Part 4); its derived clauses are consumed by later derivations. The hypothesis-level ruling is recorded here verbatim:}
\begin{quote}\small \end{quote}
\medskip
\subsection*{\texorpdfstring{B.1.42\quad BM.97 — The rest mass is the pole mass: the species adjudication, the rest-mass Q-value, and the exact critical locus (catalog entry)}{BM.97 B.1.42}}
\addcontentsline{toc}{subsection}{B.1.42 BM.97}
\phantomsection\label{tag:BM-97-cat}
{\small \textit\{Catalog.\}\quad The full theorem remains in the main corpus as Theorem 4.82 (Part 4); its derived clauses are consumed by later derivations. The hypothesis-level ruling is recorded here verbatim:}
\begin{quote}\small \end{quote}
\medskip
\subsection*{\texorpdfstring{B.1.43\quad BM.117 — The groupings, the ceilings, and the octet adjudication (catalog entry)}{BM.117 B.1.43}}
\addcontentsline{toc}{subsection}{B.1.43 BM.117}
\phantomsection\label{tag:BM-117-cat}
{\small \textit\{Catalog.\}\quad The full theorem remains in the main corpus as Theorem 4.100 (Part 4); its derived clauses are consumed by later derivations. The hypothesis-level ruling is recorded here verbatim:}
\begin{quote}\small \end{quote}
\medskip
\subsection*{\texorpdfstring{B.1.44\quad BM.141 — The Jensen residue: the exact closed form, its range, its zero locus, the two-chain companion, and the width adjudication (catalog entry)}{BM.141 B.1.44}}
\addcontentsline{toc}{subsection}{B.1.44 BM.141}
\phantomsection\label{tag:BM-141-cat}
{\small \textit\{Catalog.\}\quad The full theorem remains in the main corpus as Theorem 4.119 (Part 4); its derived clauses are consumed by later derivations. The hypothesis-level ruling is recorded here verbatim:}
\begin{quote}\small \end{quote}
\medskip
\section*{\texorpdfstring{B.2 Refuted Hypotheses}{B.2 Refuted Hypotheses}}
\addcontentsline{toc}{section}{B.2 Refuted Hypotheses}
\markboth{B.2 Refuted Hypotheses}{B.2 Refuted Hypotheses}
\subsection*{\texorpdfstring{B.2.1\quad The two distance claims of the derivation record, resolved}{B.2.1 distance claims}}
\addcontentsline{toc}{subsection}{B.2.1 distance claims}
{\small \textit{Record.}\quad Source: Part 3 (FT.48 block); moved from the FT.48 block (the framing reference to the derivation record is retained here as provenance; the main-text stub is neutral).}
\medskip
\textbf{Rmk (the two distance claims, resolved).} "The geometric distance between stable jumps is the winding/integer $k$" holds \textit{for closed materialized loops} [FT.22, FT.34(2)–(3)] and \textit{fails for open one-way geodesics} — an open path has no winding and no derived quantization (FT.22 is closed-chain; class-(iii) edges stay continuous ). Clause (2) is the correct two-sided replacement: the integer ledger distance and the metric proper length are distinct reads of the same connector, Cauchy–Schwarz-bound, equal-in-square only at uniformity. The "mass $=$ geometric distance" identification is likewise not claimed: $m_{\mathrm{hid}}$ counts bits, $d_{12}$ counts quanta; their bridge is the FT.33 ratio under Def-H.37(1), not an equality.

\medskip
\subsection*{\texorpdfstring{B.2.2\quad FT.73 — The bridge audit: dielectric non-existence, the field-degree no-go, and the surviving species}{FT.73 B.2.2}}
\addcontentsline{toc}{subsection}{B.2.2 FT.73}
\phantomsection\label{tag:FT-73}
{\small \textit{Record.}\quad Source: Part 3 (Part3\_Fermion\_era.md); tier: Tier F.24 — The Flicker Phase Structure, the Cycle Grammar of Covers, the Dust Equation of State, Face Decomposition, and the Bridge Audit; canonical tag \texttt{FT.73}. Reproduced from the revision-1 canon; a cross-reference stub marks the former locus in the main corpus.}
\medskip
\textbf{Thm.} On the proposed bridge program (a $D^-$-medium/drag mechanism equating $m_{\mathrm{hid}}$ with $m_{\mathrm{pole}}$, with $\alpha$ emerging as a screening factor): $(1)$ \textbf{Dielectric non-existence (inventory theorem).} A dielectric medium requires a polarization-response read (induced content vs applied read). The locked sector's complete inventory is: single possibility space [D4.1], absorption at rate $\propto O(n)$ [D4.2, CT4.1], one-way [CT4.2], retained and inaccessible [CT4.3, D1.11], no gauge structure [CT4.4], unresolvable [FT.3], never observable [FT.20], no slot realization at $n\ge4$ [FT.59]. \textbf{No response read exists}; $\varepsilon_{D^-}$ and "$D^-$ vacuum polarization" have no defined ledger referent, and any dynamics supplying one contradicts CT4.2/D1.11 — inadmissible from this part: no construction of Part 3 overturns a Parts 1–2 tag. The vacuum-side corpus data are instead: $\rho_\circ=\tfrac14 I$ [C11], $D^+=\Lambda\mathbb I$ with $w=-1$ [D4.3, CT1.25], $D^-$ dust with $w=0$ [CT4.5], and the entropy triangle's exhaustiveness [CT4.7, CT4.10] — the "equation of state of the vacuum" is computed, frozen, and two-sided. $(2)$ \textbf{Field-degree no-go (drag refuted).} Under Def-H.37(2) the pole-side factors at the five floors are $\{2,\ 2-\sqrt2,\ 1,\ 2-2\cos\tfrac\pi{10},\ 2-\sqrt3\}$ [FT.54(4), FT.63(6)], with minimal polynomials: $2\cos\tfrac\pi{10}$: $y^4-5y^2+5$ (irreducible over $\mathbb Q$: no rational root, no $\mathbb Z$-quadratic factor [machine]); $2-2\cos\tfrac\pi{10}$: $x^4-8x^3+19x^2-12x+1$ (irreducible [machine; FT.63 anchor re-derived]). Hence the five factors generate $\mathbb Q(\sqrt2,\sqrt3,\cos\tfrac\pi{10})$, of degree $16$ over $\mathbb Q$: writing $y=2\cos\tfrac\pi{10}$, $\mathbb Q(y)/\mathbb Q$ is Galois of group $\mathbb Z_4$ (generator $y\mapsto\sqrt5/y$; machine: $y\cdot\tfrac{\sqrt5}y=\sqrt5=2y^2-5$, $(\tfrac{\sqrt5}y)^2=\tfrac{5-\sqrt5}2$), with unique quadratic subfield $\mathbb Q(\sqrt5)$; and $\sqrt5\notin\mathbb Q(\sqrt2,\sqrt3)$ [$\sigma_2/\sigma_3$-averaging forces $\sqrt5\in\mathbb Q$, machine] — so the two quartics intersect in $\mathbb Q$ and $[\mathbb Q(\sqrt2,\sqrt3,y):\mathbb Q]=4\cdot4=16$. Any bridge formula $m^2_{\mathrm{pole}}=F(n_1,\dots,n_m;c_1,\dots,c_k)$ with $F$ rational over its arguments, the $n_i$ \textbf{integer-valued combinatorial ledger reads} (including $\rho_{D^-}=\sum O(n)$, lock counts, covering counts, $Z^\pm$, $K$-reads, $m_{\mathrm{hid}}\mathfrak h/\ln2$, $g(V)$, $\beta_1$) and finitely many declared constants $c_j\in K$ places every pole value in $K$: covering the five computed floors forces $[K:\mathbb Q]\ge16$. The $\Theta_R$-read (Def-H.37(2)) attains all five factors with $K=\mathbb Q$: \textbf{curvature coupling is the unique surviving bridge species within the computed read classes} [consistent with the form-uniqueness sharpening of FT.54(1): state-independence forces $m^2_{\mathrm{pole}}=F(\Theta_R)$]. $(3)$ \textbf{Dual split and the exclusion of quadratic couplings.} The dual pairs share the informational mass — $m_{\mathrm{hid}}$: $24{=}24$ (cube/octa), $60{=}60$ (dodeca/icosa) [FT.44(2), FT.63(5)] — while the pole side splits both [FT.63(6)]: cube/octa by $(2-\sqrt2)/1\notin\mathbb Q$; dodeca/icosa because $2-2\cos\tfrac\pi{10}$ has degree $4$ over $\mathbb Q$ while $2-\sqrt3$ has degree $2$ (a rational ratio would place the quartic inside $\mathbb Q(\sqrt3)$). Stronger: even a \textit{single declared quadratic} coupling cannot serve — all constants in one quadratic field $K_2$ would force $(2-\sqrt2)/(2-\sqrt3)=4-2\sqrt2+2\sqrt3-\sqrt6\in K_2$ [machine: exact in the $\{1,\sqrt2,\sqrt3,\sqrt6\}$ basis], which lies in no quadratic field ($\mathbb Q(\sqrt2,\sqrt3)$ has degree $4$ over $\mathbb Q$ and over $K_2$). Drag-type proposals ($m\propto$ locked-content reads, e.g. $\rho_{D^-}$: integers $896$ vs $2976$ on the cube/octa pair — the wrong splitting entirely: rational, against an irrational pole ratio) are refuted; the $\Theta_R$-read tracks both splittings natively. $(4)$ \textbf{The $\alpha$ flag.} EM.20 derives $\alpha_{\mathrm{EM}}^{(2)}=\pi$ under Comm-min+Eff — a pure number, independent of $\mathfrak h$, $G_N$, $\kappa$, and every P4 parameter [Eff: historical scope token — Part 3 §0.2(ii); HistLog §I.4]. The proposed "$\alpha\approx1/137$ as the medium's screening factor" contradicts the frozen calibration layer and presupposes the medium refuted in (1): recorded as an inadmissible ask; no screening computation exists or is constructed. $(5)$ \textbf{Adjudication summary (proposal transcript, audit pattern).} \textit{Confirmed (already derived)}: per-edge binary freedom [A4, D0.1, Def-H.25]; the constraint surface $=$ the cap census [CT0.11; FT.68(1)–(3), FT.69(1)]; photon $=$ trivial chain [D3.7, frozen]; "mass $=$ time-averaged real edges" $=$ the $m_{\mathrm{eff}}$ reads [FT.68(1), FT.70(5)]; irrational mass hierarchy $=$ defect radicals [FT.63(6)]; triangle–triangle composition laws $=$ the bowtie/double-triangle/face cells [FT.49, FT.50]; locked content "sits there and gravitates" $=$ CT4.5 dust. \textit{Corrected}: the octahedron "forbidden zone" [FT.70(4)]; "triangle minimal winding $k{=}1$" — the floor fermion is the digon ($k{=}1$, $j{=}\tfrac12$ [FT.28]); the triangle's minimal on-shell winding is $k{=}3$ ($\Theta{=}6$, $\kappa{\equiv}2$; $\kappa{\equiv}1$ gives $\Theta{=}3$, off-shell [FT.22, FT.32]); "tetrahedron $=3$ quarks $+$ singlet seal" — tetrahedron dark sealing is meson-only (even-$\varepsilon$ [FT.63(7)]); the first $\varepsilon$-phase is the face floor [FT.49]; "everything is triangles" [FT.72(1)]; "$3$ colours $=$ triangle edges" — colour is grounded at rung $3$ [D3.1, CT3.x] with the $\{2,3\}$-partition rule [FT.34(5)]; the triangle's $3$ edges supply the first $\varepsilon$-type \textit{correspondence}, not the derivation. \textit{Amplitude-layer verdict on the whole picture}: closed triangles cancel standalone [FT.67(3a)] and never occur inside covers [FT.69(3)(iv)] — the triangle-composition picture is spectral-layer content [Def-H.16, M1/M2], legitimate there and untouched by this tier; the two layers meet only at item 13's open bridge. \textit{Inadmissible}: every particle/cosmological identification (electron/proton/pion/$\rho$/dark matter/black hole/orbitals) — items 5/12 discipline: no calibration data, no identifications, recorded not adopted.

\textbf{Req.} FT.3, FT.20, FT.22, FT.28, FT.32, FT.44, FT.49, FT.50, FT.54, FT.57, FT.59, FT.63, FT.64, FT.67, FT.69–FT.72; Def-H.16; A8; D1.11, D3.1, D3.7, D4.1, D4.2, D4.3; CT1.25 (Part 2); CT3.17 (Part 2); CT4.1–CT4.7, CT4.10 (Part 2); EM.20 (Part 2); C11; Eff/Comm-min (inherited, declared at use); field-degree and irreducibility arguments [math].

\textbf{Pf.}

%dsp-level:0
\[
(1):\ \text{inventory list exhaustive over Parts 1–2's}\ D^-\ \text{tags}\ [\text{D4.1/D4.2};\ \mathrm{CT4.1\text{–}4.7};\ \mathrm{FT.3/FT.20/FT.59}];
\]

%dsp-level:0
\[
\quad\text{response read}\ \Rightarrow\ \text{accessible re-emission}\ \Rightarrow\ \neg\mathrm{CT4.2}\wedge\neg\mathrm{D1.11}
\]

%dsp-level:0
\[
(2):\ \text{irreducibility: rational-root test} + \mathbb Z\text{-quadratic ansatz}\ (bd{=}5;\ a{+}c{=}0;\ \ldots)\ \text{no solution}\ [\text{machine}];
\]

%dsp-level:0
\[
\quad\sqrt5\notin\mathbb Q(\sqrt2,\sqrt3):\ \sqrt5
\]

%dsp-level:0
\[
= a{+}b\sqrt2{+}c\sqrt3{+}d\sqrt6
\]

%dsp-level:0
\[
\sigma_2:\ b{=}d{=}0;\ \sigma_3:\ c{=}0;\ a^2{=}5\ \text{impossible in}\ \mathbb Q;
\]

%dsp-level:0
\[
\qquad [\mathbb Q(\sqrt2,\sqrt3):\mathbb Q]=4,\ [\mathbb Q(\cos\tfrac\pi{10}):\mathbb Q]=4,\ \cap=\mathbb Q\ \implies\ \text{compositum degree}\ 16
\]

%dsp-level:0
\[
F\ \mathbb Q\text{-rational in integer reads} + c_j\in K:\ \text{value field}\subseteq K;
\]

%dsp-level:0
\[
\qquad \text{five floors}\implies K\supseteq\mathbb Q(\sqrt2,\sqrt3,\cos\tfrac\pi{10})
\]

%dsp-level:0
\[
\mathbb Q(y)\ \text{Galois}\ \mathbb Z_4:\ y\mapsto\sqrt5/y:\ \text{unique quadratic subfield}\ \mathbb Q(\sqrt5)\ [\text{math};\ \text{machine}:\ 2y^2-5=\sqrt5];
\]

%dsp-level:0
\[
\quad\sqrt5\notin\mathbb Q(\sqrt2,\sqrt3)\implies\text{intersection}\ \mathbb Q\implies\text{degree}\ 16
\]

%dsp-level:0
\[
(3):\ (2{-}\sqrt2)(2{+}\sqrt3)
\]

%dsp-level:0
\[
= 4{-}2\sqrt2{+}2\sqrt3{-}\sqrt6\ \text{[machine, basis}\ \{1,\sqrt2,\sqrt3,\sqrt6\}\text{]}:\ \text{all components}
\]

%dsp-level:0
\[
\ne 0\implies\notin\ \text{any quadratic subfield}
\]

%dsp-level:0
\[
(4):\ \mathrm{EM.20}\ \text{verbatim};\qquad(5):\ \text{each row cited at its tag}.\ \blacksquare
\]

\textbf{Rmk (effect on Gap 7).} The no-go does not derive Def-H.37(1)/Def-H.37(2) — the transfer remains an interface clause — but it \textit{exhausts the alternative species within computed reads}: a bridge exists in the inventory's language only as a curvature-coupling ($\Theta_R$-type) formula. Gap 7's residue narrows accordingly: derive or replace the Def-H.37(1)-family transfer itself; the search space for replacements no longer includes combinatorial/drag/dielectric candidates.

\textbf{Phys.} Ask the ledger for a mechanism turning hidden bits into pole mass by dragging through the dark sediment, and the answer is a clean negative in two movements. The sediment cannot polarize: it absorbs, retains, and stays silent — that silence is a theorem, not a gap. And no counting formula can reach the pole side: the five solids' pole factors live in number fields of degree up to sixteen, while every combinatorial read is an integer. What survives is the one bridge the interface layer already fixes — mass from the angle defect of wrapped geometry — now alone in its class.

\vspace{2pt}\noindent\rule{\textwidth}{0.4pt}\vspace{2pt}

\medskip
\subsection*{\texorpdfstring{B.2.3\quad FT.77 — The SM-hierarchy audit: a bounded-hierarchy negative result}{FT.77 B.2.3}}
\addcontentsline{toc}{subsection}{B.2.3 FT.77}
\phantomsection\label{tag:FT-77}
{\small \textit{Record.}\quad Source: Part 3 (Part3\_Fermion\_era.md); tier: Tier F.26 — The Platonic Pole Spectrum: Exact Multiplets, the Bounded-Hierarchy Ceiling, and the SM-Hierarchy Audit; canonical tag \texttt{FT.77}. Reproduced from the revision-1 canon; a cross-reference stub marks the former locus in the main corpus.}
\medskip
\textbf{Thm.} \textbf{External-data clause.} The three hierarchies compared below — charged leptons $1{:}207{:}3477$, up-type quarks $1{:}600{:}80000$, down-type quarks $1{:}20{:}900$ — are \textit{supplied measured data} (2–3 significant figures). They enter no derivation, no definition and no corpus result; A §0.2 admits no imported measured value, and this document identifies no particles (Open Gaps 5/12). This clause is an \textbf{adjudication cell} in the FT.73(5) pattern: its output is a verdict about a proposed identification, not a physical claim. $(1)$ \textbf{Audit method, maximally generous to the hypothesis.} The comparison pool is \textit{every} nonzero pole value of \textit{every} solid at \textit{every} weight of the full periodic ladder [FT.76(3)] — $(2,4,3,10,6)$ distinct values per solid, $25$ in all, $16$ distinct across the five solids, giving $120$ pairwise ratios and $\binom{16}{3}=560$ triples; no winding, no variant and no iterate is excluded — and ratios are taken in \textbf{mass} ($m=\sqrt{D}\,\hat\mu'$), since the ladder supplies $m^2$. A match requires an \textit{aligned triple}: three pool values whose normalized ratios reproduce a family's two ratios simultaneously. $(2)$ \textbf{The ceiling obstruction.} The pool's mass ratios are bounded above by $\sqrt{4/D_{\mathrm{dodeca}}}=6.392453\ldots$ [FT.76(3)] while the smallest top-ratio among the three families is $900$. Hence \textbf{no aligned triple can match any SM family}; the best achievable discrepancy factors are

%dsp-level:0
\[
\times140.79\ (\text{down-type}),\qquad\times543.92\ (\text{charged leptons}),\qquad\times12514.76\ (\text{up-type}),
\]

exhaustive over all $560$ triples (machine). Pairwise, the closest single ratio is the pool maximum against the families' minimum ($m_\tau/m_\mu=16.797$): $\times2.63$ — within an order of magnitude, and structurally void: it pairs the two extremes and extends to no triple. $(3)$ \textbf{The dimension trap, disclosed.} Comparing $m^2$-ratios to SM \textit{mass} ratios — dimensionally wrong — produces the only near-coincidences available: $20.431729$ vs $m_s/m_d=20$ ($\times1.0216$); $40.863458$ vs $m_b/m_s=45$ ($\times1.1012$); $14.928203$ vs $m_\tau/m_\mu=16.797$ ($\times1.1252$). All three are rejected: the ladder computes $m^2_{\mathrm{pole}}$, and Def-H.37(2) supplies no square-root reassignment. $(4)$ \textbf{Structural failures, independent of any fitting.} (a) The tetrahedron's multiplet is fully degenerate [FT.76(5a)] — it cannot carry a three-member family in either variant. (b) Even-$k$ multiplets contain an exactly massless weight [FT.76(5b)] — no charged SM fermion is massless, so the cube in (A) and the octahedron in (B) are excluded as family carriers outright. (c) Only two multiplets in the entire audit have three distinct masses — dodecahedron $k_f=5$ (A) and icosahedron $k_{vf}=5$ (B) — with mass ratios $1{:}2.902113{:}4.520147$ and $1{:}2.732051{:}3.732051$; their best discrepancies are $\times199.11$ and $\times241.15$ (down-type), $\times769$ and $\times932$ (leptons), $\times17699$ and $\times21436$ (up-type). (d) Five solids against three generations, with both dual pairs degenerate in $m_{\mathrm{hid}}$/$\Pi$ [FT.44(1)] and split only on the pole side [FT.76(8)]: no solid-to-generation assignment is available without calibration data the corpus does not supply. (e) Two of the four readings of $\Theta_R$ named by FT.63(4)/(6) give identically zero pole mass [FT.76(1)]. $(5)$ \textbf{Quantitative distance.} Reaching even the least hierarchical family ($1{:}900$ in mass) with the ceiling $\max=4\hat\mu'^2$ would require a lightest defect value $D\le4/810000=1/202500=4.938\times10^{-6}$, i.e. — exactly, since $D=4\sin^2(\delta_v/4)$ — a cone defect $\delta_v\le4\arcsin\tfrac1{900}=4.4444\times10^{-3}$ rad $=0.2546^\circ$. The smallest Platonic defect is $\delta_v=\pi/5=36^\circ$: the available defects are a factor $45\pi=141.37$ too \textbf{large} in angle and $19822$ too large in $D$ — equivalently, the Platonic mass span is $140.79$ times too \textit{small}. The two shortfalls agree to $0.5\%$ because $m\propto\delta_v$ at small angles. The obstruction is one of scale, not of pattern. $(6)$ \textbf{Verdict.} The hypothesis "Platonic vertex defects $=$ SM family hierarchies" \textbf{fails at the pole-mass layer}, conditionally on Def-H.37(2)/Def-H.38 and the vertex-figure convention, and — by the ceiling argument of (2) — unconditionally on the winding choice and on which solid is assigned to which family. No match is claimed, no parameter is fitted, and no identification is made. What survives as corpus content: the exact bounded multiplets of FT.76, the ceiling theorem, and the sharpened form of Gap 7's residue — \textbf{any replacement for the Def-H.37(2) transfer must supply a dynamical range of at least three orders of magnitude if it is ever to reach observed mass hierarchies; the defect species, on the Platonic skeletons, supplies at most $6.4\times$.}

\textbf{Req.} FT.44, FT.60, FT.63, FT.73, FT.76; Def-H.37(2), Def-H.38; A §0.2 (policy); Open Gaps 5/12/16(e); external measured hierarchies (comparison only, marked); exhaustive triple search, ratio bounds [machine].

\textbf{Pf.}

%dsp-level:0
\[
(1):\ \text{pool}
\]

%dsp-level:0
\[
= \bigcup_S\{D(m,S)\ne0:\ 0\le m<P_S/2\}:\ (2,4,3,10,6)\ \text{per solid},
\]

%dsp-level:0
\[
\ 25\ \text{in all},
\]

%dsp-level:0
\[
\ 16\ \text{distinct across solids}
\]

%dsp-level:0
\[
(2):\ \max_{\text{pool}}\ \text{mass ratio}
\]

%dsp-level:0
\[
= \sqrt{4/D_{\mathrm{dodeca}}}
\]

%dsp-level:0
\[
= 6.392453\ldots<16.797
\]

%dsp-level:0
\[
= \min\ \text{SM ratio}:\ \text{every pairwise discrepancy}\ge\times2.63
\]

%dsp-level:0
\[
\text{triples}:\ \text{normalized top ratio}
\]

%dsp-level:0
\[
\le 6.392453:\ \text{discrepancy}
\]

%dsp-level:0
\[
\ge 900/6.392453
\]

%dsp-level:0
\[
= 140.7910:\ \text{exhaustive over all}\ 560\ [\text{machine}]
\]

%dsp-level:0
\[
(3):\ \text{direct evaluation of the stated pairs};
\]

%dsp-level:0
\[
\qquad \text{rejection}:\ [m^2_{\mathrm{pole}}]\ \text{is a squared mass by Def-H.37(2)}\ \text{— no reassignment declared}
\]

%dsp-level:0
\[
(4a)\text{–}(4c):\ \text{multiplet cardinalities and ratios from FT.76}(2),(7)\ [\text{machine}];
\]

%dsp-level:0
\[
\qquad (4d):\ \mathrm{FT.44}(1),(2);
\]

%dsp-level:0
\[
\qquad (4e):\ \mathrm{FT.76}(1)
\]

%dsp-level:0
\[
(5):\ D\le4/900^2=1/202500;\qquad D=4\sin^2(\delta/4)\le1/202500\iff\sin(\delta/4)\le\tfrac1{900}\iff\delta\le4\arcsin\tfrac1{900}
\]

%dsp-level:0
\[
x\le\arcsin x
\]

%dsp-level:0
\[
\le x+x^3\ (0<x\le\tfrac12):\ \delta
\]

%dsp-level:0
\[
\in [\tfrac1{225},\ \tfrac1{225}+\tfrac{4}{729\cdot10^6}]
\]

%dsp-level:0
\[
= [4.4444\times10^{-3},\ 4.44445\times10^{-3}]\ \text{rad}\ [\text{certified}]
\]

%dsp-level:0
\[
D_{\mathrm{dodeca}}\cdot202500=19822.11\ldots;\qquad(\pi/5)\,/\,4\arcsin\tfrac1{900}=45\pi=141.3716\ldots;\qquad\sqrt{19822.11}=140.791\ \checkmark
\]

%dsp-level:0
\[
(6):\ \text{conjunction of}\ (2)\text{–}(5).\ \blacksquare
\]

\textbf{Phys.} Take the five solids, give each the interface defect formula (Def-H.37(2)) and the interface symmetric-power carrier (Def-H.38), and ask whether the numbers look like a family of quarks or leptons. They do not, and the reason is not a bad choice of solid or winding: the formula can only ever return values between zero and four times one scale, because a cone defect is an angle and angles are bounded. A hierarchy of one to a thousand needs a dynamical range the geometry does not have. The honest output is therefore a boundary — the defect transfer, on these skeletons, is a spectrum of order one, and any future version of it that hopes to speak about observed masses must bring its range from somewhere else.

\vspace{2pt}\noindent\rule{\textwidth}{0.4pt}\vspace{2pt}

\medskip
\subsection*{\texorpdfstring{B.2.4\quad FT.87 — The tumbling-model adjudication (external proposal, row by row)}{FT.87 B.2.4}}
\addcontentsline{toc}{subsection}{B.2.4 FT.87}
\phantomsection\label{tag:FT-87}
{\small \textit{Record.}\quad Source: Part 3 (Part3\_Fermion\_era.md); tier: Tier F.29 — Gravitational Reads in Three States, the Flicker Ripple, and the Tumbling-Model Adjudication; canonical tag \texttt{FT.87}. Reproduced from the revision-1 canon; a cross-reference stub marks the former locus in the main corpus.}
\medskip
\textbf{Thm.} \textbf{External-data and status clause.} The proposal adjudicated here ("tumbling/rotor masses, Regge trajectories, shell magic numbers, weak $\pm2$ parity shift, cover class $=$ generation, $\ln Z$ as Yukawa, calibration from a measured lepton mass") is \textit{external}; its SM comparisons are marked [EXT]; nothing below is adopted, and no particle is identified. Verdicts: $(1)$ \textbf{The rotor term is not new physics — it is Def-H.37(2) leading order.} $4\sin^2(m\Theta_R/4)=\tfrac{m^2\Theta_R^2}{4}\bigl(1-\tfrac{x^2}{3}+\tfrac{2x^4}{45}+\cdots\bigr)$, $x=m\Theta_R/4$ (both Taylor coefficients machine-exact in $\mathbb Q$), so the proposed slope $1/R^2$ \textbf{is} the declared coefficient $\hat\mu'^2\Theta_R^2/4$ [FT.54/FT.60(5), EW.33 recovered as its Taylor term]. Consequences: the rotor inherits Def-H.37(2)s conditional status, and it inherits \textbf{FT.76(3)'s ceiling} — the sine saturates at $4$, so no winding ladder of this species can grow without bound. $(2)$ \textbf{Two dimension traps (flagged, as the proposal's own rules require).} (a) $m^2=m^2_{\mathrm{static}}+\hbar^2k^2/(2I)$ adds an \textit{energy} ($\mathfrak q^{1/2}$) to a \textit{mass squared} ($\mathfrak q$); (b) $k^2/R^2$ with $R$ an edge count is $\mathfrak q^0$, not $\mathfrak q$ — repaired only as $k^2/(R^2\mathfrak h^2)$; (c) $I=\Delta S\cdot R^2$ has dimension $\mathfrak q^{-1}$ where a moment of inertia is $\mathfrak q^{-1/2}$. Same species as FT.77's $m^2$-vs-$m$ and FT.79's amplitude-vs-mass traps. $(3)$ \textbf{"Regge verified linear for all 12 shapes" is a tautology.} $m^2=m^2_{\mathrm{static}}+(4/R^2)j^2$ with $j=|k|/2$ is linear in $j^2$ \textit{by algebra}, so the check cannot fail. The derived read is \textbf{not} linear in $j^2$: its successive slopes in $j^2$ vary (cube $1.4142,0,-0.4714$; octa $1.5,-0.75,0$; dodeca $0.3633,0.2939,0.1959$; icosa $0.8660,0.4330,0$), and the tetrahedron's ladder is \textbf{exactly flat} ($4\sin^2(m\pi/4)=2$ for every odd $m$ [FT.76(5a)]). $(4)$ \textbf{The quadratic form is a poor approximation where it is used.} At the documented uniform-minimal windings the exact-to-quadratic ratio is $0.8106$ (cube, dodeca, icosa: $18.9\%$ error), $0.4053$ (octa: $59.5\%$), $0.0901$ (tetra: $91.0\%$); the exact values reproduce FT.76(2)'s top weights $(2,2,4,2,2)$ — an independent cross-check of both sides. $(5)$ \textbf{Shell "magic numbers" collide with a proven no-go.} Extra-stable rotational levels would require a weight-level selection criterion, and FT.66(4)'s dilemma plus FT.61/FT.66(2)'s blindness theorems prove that \textbf{no weight/ledger criterion can select} among equal-record stationary variants on the Platonic curvature class; nuclear magic numbers are external data with no corpus referent. The proposal's stability islands are not constructible. $(6)$ \textbf{The moment-of-inertia table is exact but embedding-dependent.} $I=KR^2$ reproduces $(0.8889,2.0000,\text{ratio }2.25)$, $(1.5585,4.5000,2.8873)$, $(6.2882,18.0000,2.8625)$ for $R^2=(\tfrac13,\tfrac38,\tfrac34)$ (machine). The \textbf{ratios are embedding-free} — $R^2$ cancels and the ratio equals $E_{\mathrm{int}}/\langle K\rangle$, FT.87's own within-shape state ratio — but $I$ itself needs a Euclidean circumradius, an \textit{embedding} datum the corpus does not supply (Def-H.21/FT.47/FT.54 give an intrinsic piecewise-flat cone metric; \texttt{circumradius} counts $0$). The "cold dark matter" conclusion is already derived without it ($w=0$ dust [CT4.5]; $T^{\mathrm{tr}}=0$; not accessible [D1.11/CT4.2/CT4.3]; FT.86(9)) and the cosmological name is inadmissible [FT.73(5), items 5/12]. $(7)$ \textbf{Composite tumbling: the spin side is derived, the mass side is not.} Winding additivity $k(\Gamma)=k_1+k_2+d_{12}$ [FT.34(3)], the cone's CG range $|k_1-k_2|\le k\le k_1+k_2$ [FT.52(3)] and maximal-weight composition [FT.60(6)] support the aligned/anti-aligned structure, and under Def-H.38 three seeds give $\mathrm{Sym}^3$ ($j=\tfrac32$) plus mixed-symmetry $j=\tfrac12$ — the proposal's "aligned $j=3/2$ / mixed $j=1/2$" pair \textbf{is} corpus-consistent. But the splitting $\Delta m^2=k^2/R^2$ is the undeclared rotor of (1)–(2); the ratio $(3^2-1^2)/(2^2-0^2)=2$ against [EXT] $\approx1.1$ is a factor $1.8$ discrepancy, not the "factor-of-2 match" claimed; and the state labels are identifications. $(8)$ \textbf{The weak $\pm2$ shift.} No tag supplies $\varepsilon_W$, and "parity" collides with the corpus's $\mathbb Z_2$ \textit{ledger} parity [FT.64's cut-space classes] — spacetime parity has no read (\texttt{left-handed}, $P_L$, $P_R$, $\gamma_5$: count $0$). What \textit{is} legitimate: $\chi=\operatorname{sign}(k)$ is a \textbf{C-odd orientation read}, the coordinate that FT.78(6)'s chiral cover pairs lack — the missing ledger named there, of FT.40's $q_{\mathrm{core}}$ species. A $\chi$-dependent splitting would split those pairs; promoting it to a mass read needs a declared projector (Open item 20, deficit 1) and is not adopted. $(9)$ \textbf{The Yukawa replacement is refuted, on its own numbers.} $4^{|\mathcal C|}$ for $|\mathcal C|=1,2,3$ is $(4,16,64)$, \textbf{not} the claimed $(8,1056,1109312)$; those are $(8,\ Z_{\mathrm{byp}}(\text{tetra})-1,\ Z_{\mathrm{byp}}(\text{cube})-1)=(8,1056,1109312)$ — total bypass amplitudes of \textbf{three different skeletons}, contradicting "the same topological core wrapped in 1, 2 or 3 layers" (and $8$ matches no cover-census value). $\ln Z$ spans $1:3.35:6.70$ against [EXT] $1:207:3477$ — short by $\times519$, so the "right order of magnitude" claim is refuted by the proposal's own arithmetic. $\ln Z$ is also \textbf{ill-defined} as an entropy: $Z$ is a signed sum $(-1)^{\sum\kappa}4^{|\mathcal C|}$ [FT.67(2)] vanishing on most patterns, while the corpus's ledger entropy is $S_{\mathfrak L}=n_{\mathcal R}\ln2$ [CT1.29]. Using $\ln Z$ as a mass is precisely the amplitude-to-mass smuggle the proposal's own rule 4 forbids [FT.79(5)]. $(10)$ \textbf{Foreign labels.} \texttt{Theorem\ 11}, \texttt{equivalence\ principle}, \texttt{N10}, \texttt{OP8}, \texttt{OP10}, \texttt{K-machine}, \texttt{moment\ of\ inertia}, \texttt{rigid\ rotor}, \texttt{Regge}, \texttt{tumbling}, \texttt{magic\ number}, \texttt{shell\ model}, \texttt{circumradius}: all count $0$ in Parts 1–2. Where such a label names something real, the real object is cited instead (OP10's involutions $\to$ FT.88; "mass lives in the trace" has no referent; the equivalence principle is FT.85(3)'s non-constructibility result). $(11)$ \textbf{The calibration phase is refused.} Fixing $\mathfrak h$ from a measured lepton mass violates P1 (one free dimensionful scale, fixed once), P5 ($\neg\exists$ quantity identified with a measured external constant) and the proposal's own rule 3. Not executed. $(12)$ \textbf{Verdict.} The tumbling model supplies \textbf{no mechanism that escapes the two ceilings}: its mass term is Def-H.37(2)s small-angle regime (bounded by FT.76(3)), its size datum is undeclared, its shell stability is excluded by FT.66(4), and its generation mechanism is an amplitude read the rules forbid. Open item 20's deficits 1 (chirality), 2 (Yukawa) and 3 (generations) are \textbf{not closed} by it; the genuinely new constructible content is FT.88's involution algebra and the C-odd orientation coordinate of clause (8).

\textbf{Req.} FT.20, FT.22, FT.34, FT.40, FT.44, FT.47, FT.52, FT.54, FT.60, FT.61, FT.64, FT.66, FT.67, FT.68, FT.70, FT.73, FT.76–FT.79, FT.85, FT.86; Def-H.21, Def-H.24, Def-H.30; D2.6, D2.7; CT0.9, CT0.11 (Part 1); CT1.12, CT1.29 (Part 2); CT4.5 (Part 2); EW.33; EM.9; Def-H.37(2), Def-H.38, P1, P5; Taylor series in $\mathbb Q$, exhaustive enumeration, dimensional analysis in $\mathfrak q$-exponents [math; machine]. External proposal and [EXT] comparisons marked.

\textbf{Pf.}

%dsp-level:0
\[
(1):\ \sin x
\]

%dsp-level:0
\[
= \sum_j(-1)^jx^{2j+1}/(2j)!
\]

%dsp-level:0
\[
\implies 4\sin^2x
\]

%dsp-level:0
\[
= 4x^2-\tfrac43x^4+\tfrac{8}{45}x^6-\cdots:\ \text{coefficients machine-exact in}\ \mathbb Q
\]

%dsp-level:0
\[
\text{slope identification}:\ \tfrac{m^2\Theta_R^2}{4}\hat\mu'^2\equiv\tfrac{m^2}{R^2}\implies R^{-2}=\hat\mu'^2\Theta_R^2/4\ [\text{no new datum}]
\]

%dsp-level:0
\[
(2):\ [m^2]=\mathfrak q,\ [\hbar^2k^2/(2I)]=\mathfrak q^{1/2},\ [k^2/R^2]=\mathfrak q^0,\ [\Delta S R^2]=\mathfrak q^{-1}\ne[\ I\ ]=\mathfrak q^{-1/2}\ [\mathrm{D2.9}]
\]

%dsp-level:0
\[
(3):\ j=|k|/2\implies(4/R^2)j^2=k^2/R^2:\ \text{the ansatz itself};
\]

%dsp-level:0
\[
\qquad \text{derived read}:\ \text{slopes in}\ j^2\ \text{non-constant}\ [\text{machine}]
\]

%dsp-level:0
\[
(4):\ 4\sin^2(k_f\delta_v/4)=(2,2,4,2,2)\ [\mathrm{FT.76}(2)];
\]

%dsp-level:0
\[
\quad 4x^2=(22.207,2.467,9.870,2.467,2.467):\ \text{ratios as stated}
\]

%dsp-level:0
\[
(5):\ \mathrm{FT.66}(4)\ \text{dilemma}+\mathrm{FT.61}/\mathrm{FT.66}(2)\ \text{blindness}\implies\neg\exists\ \text{weight-level selection criterion}
\]

%dsp-level:0
\[
(6):\ I=KR^2:\ \text{exact Fractions};\qquad I_{S2}/I_{S1}=E_{\mathrm{int}}/\langle K\rangle:\ R^2\ \text{cancels}
\]

%dsp-level:0
\[
(7):\ \mathrm{FT.34}(3),\mathrm{FT.52}(3),\mathrm{FT.60}(6);\qquad2\otimes2\otimes2=4\oplus2\oplus2:\ j\in\{\tfrac32,\tfrac12\}\ [\mathrm{Def-H.38}:\ \mathrm{Sym}^3]
\]

%dsp-level:0
\[
(9):\ 4^{|\mathcal C|}=(4,16,64);
\]

%dsp-level:0
\[
\quad 1056=1057-1,\ 1109312=1109313-1\ [\mathrm{FT.68}(1),(3)];\quad\ln(1109312)/\ln8=6.6937
\]

%dsp-level:0
\[
(10),(11):\ \text{inventory scan};\qquad\mathrm{P1}/\mathrm{P5}\ \text{verbatim};\qquad(12):\ \text{conjunction}.\ \blacksquare
\]

\textbf{Phys.} The proposal is a good instinct aimed at the wrong layer. Rotation as a mass generator is exactly the right shape of idea — and the corpus already contains its leading term, hiding inside the declared defect formula, where it is bounded by four and therefore cannot build a hierarchy. What survives the audit is smaller and more useful: winding adds like the proposal says, three seeds compose into the spin pair it guessed, and the left/right tumble it wanted is a real coordinate the ledger already has — the orientation sign that pairs cover classes. What does not survive is the shortcut: reading a partition function as a Yukawa coupling, or fixing the one free scale from a measured mass.

\vspace{2pt}\noindent\rule{\textwidth}{0.4pt}\vspace{2pt}

\medskip
\subsection*{\texorpdfstring{B.2.5\quad FT.90 — The resolution ladder: no-hair from the channel bound, and the chiral merger refuted}{FT.90 B.2.5}}
\addcontentsline{toc}{subsection}{B.2.5 FT.90}
\phantomsection\label{tag:FT-90}
{\small \textit{Record.}\quad Source: Part 3 (Part3\_Fermion\_era.md); tier: Tier F.30 — The Koide Adjudication, the Resolution Ladder, and the Minimal Bridge Declaration; canonical tag \texttt{FT.90}. Reproduced from the revision-1 canon; a cross-reference stub marks the former locus in the main corpus.}
\medskip
\textbf{Thm.} Under FT.3/FT.4/FT.15 with C10/CT0.6: $(1)$ \textbf{The no-hair theorem is already proved, and it is stronger than a limit.} FT.4: $I_o(v)\le\min(\delta h(\gamma),4\delta\tau(\gamma))\ln2$, and this bound \textit{is the entire externally available description} — $v,v'$ with the same $(\delta h,\delta\tau)$ are observationally identical. FT.3: for \textbf{every} $\Sigma''\in[\gamma]$, $I_{o\leftarrow\Sigma''}(\gamma)=I(\delta h(\gamma),\delta\tau(\gamma))$ exactly and totally, \textit{including whether $\Sigma''$ contains $D^-$-configured nodes}. So the collapse of internal structure is not a large-$\delta h$ effect: it holds at every $\delta h$. \textbf{Correction to the submitted framing:} the budget \textit{grows} with $\delta h$, so merging is tightest at \textit{small} $\delta h$; what is $\delta h$-independent is the invisibility of the internals. $(2)$ \textbf{The resolution ladder, exact.} With the saturating schedule $\delta\tau=\lceil\delta h/4\rceil$ (at most $4$ hops per tick [FT.1/FT.38, CT0.6, EM.5–7], attained by photon schedules [FT.36(2)]), the budget is $\min(\delta h,4\delta\tau)=\delta h$ bits, so

%dsp-level:0
\[
|\mathrm{Res}(\delta h)|\le2^{\delta h}:\qquad\delta h=1,2,3,4,5\ \mapsto\ 2,4,8,16,32\ \text{alternatives},
\]

and at $\delta h=4$ the bit budget equals $C_{\mathrm{cap}}=16$ registers exactly [CT0.6]. C10's warning is load-bearing: the \textbf{register cost} of $|X|$ alternatives is $|X|$, their \textbf{Shannon content} $\log_2|X|$ — the two budgets are different quantities and must not be identified. $(3)$ \textbf{Exact thresholds for the computed class counts.} The hop distance at which the budget first covers a class count:

%dsp-level:0
\[
\begin{array}{c|cccccc}
\text{target} & \text{tetra configs} & \text{tetra classes} & \text{tetra covers} & \text{cube configs} & \text{cube classes} & \text{cube covers}\\
\hline
\# & 1230 & 70 & 7 & 1\,580\,092 & 33\,564 & 264\\
\min\delta h & 11 & 7 & 3 & 21 & 16 & 9
\end{array}
\]

with $P_8$'s $42\,687$ admissible configurations at $\delta h\ge16$. Merge factors at $\delta h=4$ (budget $16$): tetrahedron $1230/16=615/8$, cube $1\,580\,092/16=395\,023/4$. $(4)$ \textbf{What survives at every resolution.} FT.15's invariant tuple $\Pi_v(C)=(q,S,m_{\mathrm{hid}})$ — invariance under internal rearrangement given a shared signature $(\delta h,\delta\tau,q,\Theta_{\mathrm{int}})$ — together with the aspect set $\mathcal A(\Sigma;C)$ and $\dim\mathcal K_\Sigma\le18$ [FT.20/FT.26]. Precisely: $\mathrm{Res}(\infty)$ does \textbf{not} equal the tuple because resolution is not the limit that removes the internals; FT.3 removes them at \textit{every} resolution, and the ladder of (2) bounds how many \textit{classes} a channel can separate. $(5)$ \textbf{The chiral merger is NOT forced by FT.15.} Sign-dependence of the tuple, component by component: $S=j=|k|/2$ [FT.52/FT.55, FT.60] — \textbf{sign-blind}; $m_{\mathrm{hid}}=E_{\mathrm{int}}\ln2/\mathfrak h$ [FT.16/FT.44(4)] — \textbf{sign-blind}; $q=q_{\mathrm{core}}=\tfrac12(\sigma_{e_1}-\sigma_{e_k})\mapsto-q_{\mathrm{core}}$ under $\omega$ [FT.40, D2.6] — \textbf{sign-sensitive}. Moreover FT.15's \textit{hypothesis} already requires equal $q$, so the theorem never merges states of differing $q$. Hence a Dirac-type merger of a chiral pair would require \textbf{declaring $q$ unobservable} (or a coarser invariant): a new declaration, not a consequence. This is the same C-odd coordinate identified in FT.78(6)/FT.87(8), and no projector exists to split or merge on it [FT.84(1) deficit 1; FT.87(8)]. $(6)$ \textbf{The "information cost of generations" question: premise void, legitimate version exact.} The submitted $(8,1056,1109312)$ and $\ln Z$ gaps rest on FT.87(9)'s refuted identification (those are $(8,Z^{\rm tetra}_{\mathrm{byp}}-1,Z^{\rm cube}_{\mathrm{byp}}-1)$, three different skeletons) and $\ln$ of a \textit{signed} sum is not an entropy [FT.67(2); the corpus entropy is $S_{\mathfrak L}=n_{\mathcal R}\ln2$, CT1.29] — an amplitude-to-entropy smuggle, void. Legitimate and derived: distinguishing $N$ classes costs $N$ registers and $\log_2N$ bits [C10], so $N=3$ costs $3$ registers / $\log_23$ bits and is available at $\delta h\ge2$. \textbf{The obstacle to generations is the absence of a label, not channel capacity} [item 20 deficit 3]. The Compton-like length $1/m_{\mathrm{hid}}=\mathfrak h/(K\ln2)$ exists [D2.9], but the threshold $\delta h\ge\lceil\log_2\#\text{classes}\rceil$ depends only on the class count, so \textbf{no Compton-wavelength relation arises}.

\textbf{Req.} FT.1, FT.3, FT.4, FT.15, FT.16, FT.20, FT.26, FT.36, FT.38, FT.40, FT.44, FT.52, FT.55, FT.60, FT.67, FT.78, FT.79, FT.86, FT.87; Def-H.1, Def-H.9; D0.6, D0.7, D2.6, D2.9; CT0.6, CT1.29 (\textbf{B/C}); C10; exact bisection and enumeration [math; machine].

\textbf{Pf.}

%dsp-level:0
\[
(1):\ \mathrm{FT.4}\ \text{Thm verbatim (bound + totality of the description)};
\]

%dsp-level:0
\[
\qquad \mathrm{FT.3}:\ I_{o\leftarrow\Sigma''}=I(\delta h,\delta\tau)\ \forall\Sigma''\in[\gamma]\ \text{exact and total}
\]

%dsp-level:0
\[
\text{direction of the budget}:\ \min(\delta h,4\delta\tau)\ \text{non-decreasing in}\ \delta h\implies\text{merging tightest at small}\ \delta h
\]

%dsp-level:0
\[
(2):\ \delta\tau=\lceil\delta h/4\rceil\implies4\delta\tau\ge\delta h\implies\min=\delta h;\qquad2^4=16=C_{\mathrm{cap}}\ [\mathrm{CT0.6}];\qquad\mathrm{C10}:\ |X|\ne\log_2|X|
\]

%dsp-level:0
\[
(3):\ \min\{\delta h:2^{\delta h}\ge N\}
\]

%dsp-level:0
\[
= \lceil\log_2N\rceil:\ \lceil\log_21230\rceil
\]

%dsp-level:0
\[
= 11,
\]

%dsp-level:0
\[
\ \lceil\log_270\rceil
\]

%dsp-level:0
\[
= 7,
\]

%dsp-level:0
\[
\ \lceil\log_27\rceil
\]

%dsp-level:0
\[
= 3,
\]

%dsp-level:0
\[
\ \lceil\log_21580092\rceil
\]

%dsp-level:0
\[
= 21,
\]

%dsp-level:0
\[
\ \lceil\log_233564\rceil
\]

%dsp-level:0
\[
= 16,
\]

%dsp-level:0
\[
\ \lceil\log_2264\rceil\qquad = 9
\]

%dsp-level:0
\[
(4):\ \mathrm{FT.15}\ \text{invariance}+\mathrm{FT.20}'\text{s criteria}+\mathrm{FT.26}(3)'s\ \le18\ \text{bound}
\]

%dsp-level:0
\[
(5):\ j=|k|/2\ [\mathrm{FT.52}/\mathrm{FT.55}];\qquad q_{\mathrm{core}}\mapsto-q_{\mathrm{core}}\ [\mathrm{FT.40}];\qquad\mathrm{FT.15}\ \text{hypothesis includes equal}\ q
\]

%dsp-level:0
\[
(6):\ \mathrm{C10}:\ \text{register cost}=|X|;\quad\lceil\log_23\rceil=2;\qquad1/m_{\mathrm{hid}}=\mathfrak h/(K\ln2):\ [\mathfrak q^{-1/2}]=\text{length}\ [\mathrm{D2.9}].\ \blacksquare
\]

\textbf{Phys.} The corpus already contains a no-hair theorem, and it is blunt: whatever sits behind the channel — including whether the far side is a locked, saturated tower — the observer receives a number of bits fixed by the channel's own hop and tick counts, and nothing else. Resolution is a budget that doubles per hop, so the ladder is exact: two alternatives at one hop, sixteen at four (which is precisely the register capacity of a single channel quantum), and seven hops to tell apart the seventy shapes of a tetrahedron. What no amount of resolution recovers is handedness of the internal kind: the tuple forgets the sign of the winding in its spin and mass entries but remembers it in its charge entry, so left and right stay distinct until somebody declares the charge unreadable — and that is a decision, not a theorem.

\vspace{2pt}\noindent\rule{\textwidth}{0.4pt}\vspace{2pt}

\medskip
\subsection*{\texorpdfstring{B.2.6\quad FT.94 — The chiral mass spectrum: not constructible as chiral, exact as a statistic, and the fourth hierarchy negative}{FT.94 B.2.6}}
\addcontentsline{toc}{subsection}{B.2.6 FT.94}
\phantomsection\label{tag:FT-94}
{\small \textit{Record.}\quad Source: Part 3 (Part3\_Fermion\_era.md); tier: Tier F.31 — The Native Clifford Census, the Weak-Chirality Pincer, and the Reversal Census of the Covers; canonical tag \texttt{FT.94}. Reproduced from the revision-1 canon; a cross-reference stub marks the former locus in the main corpus.}
\medskip
\textbf{Thm.} Under Def-H.25/Def-H.28, at (S1)–(S4): $(1)$ \textbf{The typed tuple is not constructible.} No read of the form $(I_{\mathrm{topo}},I_{\mathrm{grav}},I_{\mathrm{Higgs}},I_{\mathrm{Weak}},I_{\mathrm{strong}})$ exists in the inventory: no pair-channel block carries a block-local quantum [FT.30's verdict: the geometric quantum is a \textit{joint} $S$-even read mixing all ten metric coordinates across $G,H,D^+,W$]; per-aspect bit counts are well-posed exactly for the registration-type ledgers [FT.41's trichotomy: charge $=2$ end bits; colour $=\log_2$(sealing grammar); gravity $=L{-}1$ free integers, unbounded, $\neg$bit-type (C10); Higgs $=$ the continuous P4 datum $h_0$]; $I_{\mathrm{Higgs}}$ has no information-type referent at all (deficit 8: $\dim F_H=1$). The directive's own F5 concedes the point; the citation trail is exact. $(2)$ \textbf{"Chiral" splits no mass read.} By FT.93(3) no weak-chiral grade exists on any carrier; and cover-class reversal chirality is invisible to both mass ansätze — $K$, $|\mathcal C|$, $\Omega$, amp, norm weight, $k_{\max}$, mask dimension and cycle signature are identical on every exchanged pair [FT.78(6), re-verified per pair: tetra $1/1$, cube $111/111$, machine]. The per-class spectrum inside the declared inventory is therefore exactly FT.79(1)'s ladder — tetra $m_A\in\{0,2,4,6\}\ln2/\mathfrak h$, cube $\{0,2,4,6,8,10,12\}\ln2/\mathfrak h$ (re-derived) — with \textbf{zero spectral poles per class} [FT.80(3)]; the minimal bridge form [FT.91] remains constructed data and \textbf{not adopted} (this directive adopts nothing). $(3)$ \textbf{The computable core of C3, as a flagged statistic.} The norm-weighted average of the Ansatz-A read with $|A|^2$ census weights ($|A|^2=16^{|\mathcal C|}$ per configuration; class sum $\Omega\cdot16^{|\mathcal C|}$):

%dsp-level:0
\[
\langle m_A^2\rangle_{\mathrm{tetra}}=\frac{886272}{35201}\Bigl(\frac{\ln2}{\mathfrak h}\Bigr)^2\approx25.1775\Bigl(\frac{\ln2}{\mathfrak h}\Bigr)^2;
\]

%dsp-level:0
\[
\qquad \langle m_A^2\rangle_{\mathrm{cube}}=\frac{22146188288}{268241835}\Bigl(\frac{\ln2}{\mathfrak h}\Bigr)^2\approx82.5605\Bigl(\frac{\ln2}{\mathfrak h}\Bigr)^2
\]

exact rationals. \textbf{Flag:} this is a census statistic over classes (amplitude enters only as a weight — F4 respected: no amplitude is a mass, FT.79(5)); it is not a particle-mass read, and no pole interpretation is available without a bridge-map interface of the FT.91 species (retired candidate) [FT.80(2)/FT.91]. $(4)$ \textbf{C5: the fourth clean negative.} The widest $m_A^2$ range on any computed skeleton is $(12/2)^2=36$ against the [EXT] charged-lepton target $3477^2=12\,089\,529$ — a shortfall of $\times335\,820.25=12089529/36$ exactly; a priori, spanning $1{:}R$ requires $V\ge4R/3$ [FT.79(2)]: $V\ge4636$ for $R=3477$, against $V=8$. The verdicts of FT.77/FT.79(4)/FT.87(9) stand unmodified; the amplitude near-coincidences ($4^{|\mathcal C|}$ vs $207/3477$; the triple $(8,1056,1109312)=(8,Z^{\rm tetra}_{\rm byp}{-}1,Z^{\rm cube}_{\rm byp}{-}1)$) are rejected again as the same dimension trap [FT.79(5), FT.87(9)].

\textbf{Req.} Def-H.25, Def-H.28; FT.30, FT.41, FT.67, FT.69, FT.76(3), FT.78, FT.79, FT.80, FT.87(9), FT.91, FT.93; C10; CT0.11 (Part 1); exact class enumeration and Fraction arithmetic [math; machine]. External hierarchy marked at (4).

\textbf{Pf.}

%dsp-level:0
\[
(1):\ \mathrm{FT.30}\ \text{verdict row} + \mathrm{FT.41}(c)\ \text{trichotomy}\ [\text{cited; no computation adds a referent}]
\]

%dsp-level:0
\[
(2):\ \mathrm{FT.93}(3)\ \text{(no grade)};\qquad\text{pair-field identity census}\ [\text{machine: }112\ \text{pairs}\times9\ \text{fields}]
\]

%dsp-level:0
\[
(3):\ \langle m_A^2\rangle
\]

%dsp-level:0
\[
= \frac{\sum_{\mathcal O}\Omega_{\mathcal O}16^{|\mathcal C_{\mathcal O}|}(2|\mathcal C_{\mathcal O}|)^2}{\sum_{\mathcal O}\Omega_{\mathcal O}16^{|\mathcal C_{\mathcal O}|}}:\ \text{tetra over }7\text{ classes},
\]

%dsp-level:0
\[
\ \text{cube over }264\ [\text{machine, exact}]
\]

%dsp-level:0
\[
(4):\ \max m_A^2/\min m_A^2\big|_{\ne0}=(12/2)^2=36;
\]

%dsp-level:0
\[
\qquad 3477^2/36=12089529/36;
\]

%dsp-level:0
\[
\qquad V\ge\tfrac43\cdot3477=4636\ [\mathrm{FT.79}(2)]\ \blacksquare
\]

\textbf{Phys.} Asked for a chiral mass spectrum, the corpus answers in three voices. The typed-information tuple does not exist: force blocks hold no private quanta, gravity keeps an unbounded integer ledger, and the Higgs block keeps a continuous dial — nothing to type. The chirality that does exist on the covers cannot touch any mass: reversal flips arrows and leaves every mass-relevant count untouched, one hundred twelve exchanged pairs and not one field moved. And the numbers themselves — an exact average of twenty-five and a fifth on the tetrahedron, eighty-two and a half on the cube, in the only units the corpus has — sit under a ceiling of thirty-six in squared-mass range, five and a half orders of magnitude short of the lepton hierarchy they were asked to reproduce. Fourth time, same negative, same exact reason.

\vspace{2pt}\noindent\rule{\textwidth}{0.4pt}\vspace{2pt}

\medskip
\subsection*{\texorpdfstring{B.2.7\quad BM.126 — The spin-dependence of the Janus crossing: refuted on the four standing grounds, with the derived domain row}{BM.126 B.2.7}}
\addcontentsline{toc}{subsection}{B.2.7 BM.126}
\phantomsection\label{tag:BM-126}
{\small \textit{Record.}\quad Source: Part 4 (Part4\_Baryomesonic\_era.md); tier: Tier BM.22 — The Annihilation Channel, Concrete Spin, and Fermion Identification: the Janus Tunnel, the Zig-Zag World Line, and the Spin-Dependent Transport Defect; canonical tag \texttt{BM.126}. Reproduced from the revision-1 canon; a cross-reference stub marks the former locus in the main corpus.}
\medskip
\textbf{Thm.}

%dsp-level:0
\[
(1)\quad\textbf{the symbol census of the crossing weight:}
\]

%dsp-level:0
\[
\qquad x_+x_-=4\varepsilon_q(1-4\varepsilon_q),
\]

%dsp-level:0
\[
\qquad x_+=1-4\varepsilon_q,
\]

%dsp-level:0
\[
\qquad x_-=4\varepsilon_q
\]

%dsp-level:0
\[
\qquad [\mathrm{BM.19}(2);\ \mathrm{BM.120}(5)]:
\]

%dsp-level:0
\[
\boxed{\ \text{the alphabet of the crossing weight is}\ \{\varepsilon_q\}\ \ \textbf{alone}\ —\ \neg k,\ \neg j,\ \neg\Theta,\ \neg m,\ \neg D_S,\ \neg\Theta_R,\ \neg\hat\mu'^2\ }
\]

%dsp-level:0
\[
\qquad [\text{machine: the symbol scan of the three monomials}];
\]

%dsp-level:0
\[
(2)\quad\textbf{the refutation on the four standing grounds:}
\]

%dsp-level:0
\[
\qquad \textbf{a spin-dependent crossing condition would be a map}\ \varepsilon_q\mapsto\varepsilon_q(j)\ \text{or}\ x_\pm\mapsto x_\pm(k):
\]

%dsp-level:0
\[
(\mathrm O1)\quad\textbf{disjoint alphabets:}\ \text{the measure's alphabet is}\ \{\varepsilon_q\}\ \text{and the defect alphabet is}\ \{\delta_v,\ \hat\mu'^2\},
\]

%dsp-level:0
\[
\ \text{with}\ \{\varepsilon_q\}\cap\{\delta_v,\hat\mu'^2\}\qquad = \varnothing
\]

%dsp-level:0
\[
\qquad [\mathrm{BM.90}(5);\ \text{machine: the set intersection}];
\]

%dsp-level:0
\[
(\mathrm O2)\quad\textbf{no count}\to\textbf{strength coupling:}\ \neg\exists\ \text{a derived map turning an integer count into a coupling strength}
\]

%dsp-level:0
\[
\qquad [\mathrm{BM.36}(3);\ \mathrm{BM.61}(8)(\mathrm{O1});\ \mathrm{BM.90}(3);\ \mathrm{BM.105}(7)];
\]

%dsp-level:0
\[
(\mathrm O3)\quad\textbf{no loop:}\ \neg\exists\ \text{a loop, a}\ \beta\text{-function, a counterterm or a Yukawa-type object in the inventory}
\]

%dsp-level:0
\[
\qquad [\mathrm{BM.88}(1);\ \mathrm{EW.15};\ \mathrm{FT.145}(3);\ \mathrm{BM.61}(8)];
\]

%dsp-level:0
\[
(\mathrm O4)\quad\textbf{no second scale:}\ \exists!\ \text{free dimensionful scale}\ \mathfrak h\qquad[\mathrm{P1};\ \mathrm{A}\ \text{§0.3}];
\]

%dsp-level:0
\[
\boxed{\ \textbf{hypothesis (W7)}\ —\ \text{a spin-dependent Janus crossing}\ \varepsilon_q(j)\ \text{or a spin-gated tunnel}\ —\ \textbf{is refuted as a derived coupling, on}\ (\mathrm O1)\text{--}(\mathrm O4)\ }
\]

%dsp-level:0
\[
(3)\quad\textbf{the precedent refutations apply verbatim:}
\]

%dsp-level:0
\[
\qquad \textbf{hypothesis (W)}\ —\ \text{a defect-conditioned retention probability reweighting}\ m_{\mathrm{hid}}\ \text{across the six states}\ —\ \textbf{is refuted on four grounds}
\]

%dsp-level:0
\[
\qquad [\mathrm{BM.90}(5)];
\]

%dsp-level:0
\[
\textbf{hypothesis (W6)}\ —\ \text{six state-dependent thresholds}\ \varepsilon_i(D_S)\ —\ \textbf{is refuted on four grounds}
\]

%dsp-level:0
\[
\qquad [\mathrm{BM.105}(7)]:
\]

%dsp-level:0
\[
\boxed{\ \textbf{(W7) is the same species as (W) and (W6): a state-indexed threshold, and it falls to the same four grounds}\ }
\]

%dsp-level:0
\[
\qquad [\text{machine: the alphabet census of}\ (1)\ \text{against}\ \mathrm{BM.105}(7)\text{'s}];
\]

%dsp-level:0
\[
(4)\quad\textbf{the weight is branch-blind and spin-blind, exactly:}
\]

%dsp-level:0
\[
\qquad w_{\mathrm{mat}}=w(\varepsilon_q)\ \ \text{with}\ \ \neg\exists\ \text{a derived map}\ D_S\mapsto\varepsilon
\]

%dsp-level:0
\[
\qquad [\mathrm{BM.90}(5);\ \mathrm{BM.105}(7);\ \text{machine}]:
\]

%dsp-level:0
\[
\frac13\le w_{\mathrm{mat}}\le1,
\]

%dsp-level:0
\[
\qquad w_{\mathrm{mat}}(0)=1,
\]

%dsp-level:0
\[
\qquad w_{\mathrm{mat}}=\frac1{1+u+u^2}\ \ \text{on the marked pair specialization}
\]

%dsp-level:0
\[
\qquad [\mathrm{BM.94}(7);\ \mathrm{BM.105}(3);\ \text{machine}];
\]

%dsp-level:0
\[
\boxed{\ \Gamma_{\mathrm{tr}}=\frac{x_+x_-}{1-x_+x_-}\in\Bigl[0,\tfrac13\Bigr]\ \ \textbf{is a function of}\ \varepsilon_q\ \textbf{alone}\ }\qquad[\mathrm{BM.44}(2);\ \mathrm{BM.116}(2);\ \text{machine}];
\]

%dsp-level:0
\[
(5)\quad\textbf{what is derived instead — the domain row, not a prohibition:}
\]

%dsp-level:0
\[
\qquad \textbf{(i) the crossing is a move on the pair register, and the register exists for every pair:}\quad\mathcal J\cong\mathbb C^2\ \ \text{with}\ \neg\text{a spin index in its basis}
\]

%dsp-level:0
\[
\qquad [\mathrm{Def\text{-}BM.18}(1);\ (1);\ \text{machine}];
\]

%dsp-level:0
\[
\textbf{(ii) the statistics read is formed exactly on bound chains:}
\]

%dsp-level:0
\[
\qquad \Psi(c)=\mathrm{Sym}^{|k|}(\mathcal H_{\mathsf v})\ \ \text{requires a materialized closed bound chain on the nonnegative cone}
\]

%dsp-level:0
\[
\qquad [\mathrm{Def\text{-}H.38}\ \text{Scope};\ \mathrm{D1.5};\ \mathrm{FT.22}];
\]

%dsp-level:0
\[
\textbf{transit and trivial chains are excluded:}\quad\text{the}\ \mathrm{D3.7}\ \text{trivial chain carries}\ \neg\text{a}\ j
\]

%dsp-level:0
\[
= |k|/2\ \text{read and}\ \neg\text{an}\ \varepsilon\ \text{read; its helicity is a transit-side read with no identification}
\]

%dsp-level:0
\[
\qquad [\mathrm{D3.7};\ \mathrm{EM.8};\ \mathrm{Def\text{-}H.38}\ \text{Scope}];
\]

%dsp-level:0
\[
\textbf{(iii) the materialization restriction is a count, not a spin:}
\]

%dsp-level:0
\[
\qquad |E_{\mathcal R}|\ge2
\]

%dsp-level:0
\[
\qquad [\mathrm{Def\text{-}H.4};\ \mathrm{FT.144}(4);\ \mathrm{BM.13}(3);\ \mathrm{BM.70}(2):\ \text{the two-bit floor}];
\]

%dsp-level:0
\[
\boxed{\ \textbf{the derived dichotomy is bound against transit}\ [\mathrm{D1.5}],\ \neg\textbf{fermion against boson}\ —\ \text{and}\ \mathrm{BM.106}(1)\ \textbf{refutes a spin-parity theorem on the kind}\ }
\]

%dsp-level:0
\[
(6)\quad\textbf{the correlation the question seeks, stated exactly:}
\]

%dsp-level:0
\[
\qquad \text{the Janus register materializes chains, and}\ \varepsilon\ \text{is formed exactly on the bound chains it materializes:}
\]

%dsp-level:0
\[
\boxed{\ \varepsilon\ \textbf{is defined on the crossing's output and}\ \neg\textbf{on its rate}\ —\ \text{the crossing is spin-blind and the statistics is a read of what crossed}\ }
\]

%dsp-level:0
\[
\textbf{the transit object is the complement:}
\]

%dsp-level:0
\[
\qquad \neg\text{a winding carrier}\implies\neg\varepsilon\ \ \text{while}\ \ \text{the register's moves still apply to it}
\]

%dsp-level:0
\[
\qquad [(5)(\mathrm{ii});\ \text{machine: the two scope rows are disjoint}];
\]

%dsp-level:0
\[
(7)\quad\textbf{status:}
\]

%dsp-level:0
\[
\qquad \textbf{refuted as a derived coupling, on four grounds:}\ (2),\ (3);
\]

%dsp-level:0
\[
\qquad \textbf{derived:}\ (1),\ (4),\ (5),\ (6);
\]

%dsp-level:0
\[
\textbf{the verdict on the directive's crossing condition:}\qquad\begin{array}{l|l}j\in\mathbb Z+\tfrac12\implies\text{the tunnel is allowed} & \neg\textbf{formed: the weight carries no spin symbol}\ (1),\ \text{and}\ (\mathrm O1)\text{--}(\mathrm O4)\ \text{bar the coupling}\ (2)\\ j\in\mathbb Z\implies\text{the tunnel is forbidden} & \neg\textbf{formed: same grounds};\ \ \textbf{the derived row is}\ (5)\text{'s bound/transit domain and}\ (6)\text{'s read locus}\end{array}
\]

%dsp-level:0
\[
\neg\text{a declaration},
\]

%dsp-level:0
\[
\quad\neg\text{a P4 addition},
\]

%dsp-level:0
\[
\quad\neg\text{a new dimensionful scale}\ [\mathrm{P1}],
\]

%dsp-level:0
\[
\quad\neg\text{a measured value}\ [\mathrm{P5}],
\]

%dsp-level:0
\[
\quad\neg\text{a numerical coupling}\ [\mathrm{EM.25}(7)],
\]

%dsp-level:0
\[
\quad\neg\text{a species identification}\ [\mathrm{Def\text{-}BM.3}(7)];
\]

%dsp-level:0
\[
\alpha_B:\quad\neg\text{added},\quad\neg\text{derived},\quad\text{barred}\ —\ \mathrm{BM.61}(8)\text{'s}\ (\mathrm O1)\text{--}(\mathrm O4)\ \text{stand}\ [\text{machine: the token scan of this tier}];
\]

%dsp-level:0
\[
\neg\text{a HistLog disposition imported as premise}\ —\ \text{Q7 and deficits}\ 2,\ 3\ \text{are cited as pointers}.
\]

\textbf{Req.} A4, D0.4, D0.8, D1.5, D2.9 (Part 1); CT0.8, EM.8, EM.25(7), EW.15, D3.7 (Part 1, Part 2); P1, P4, P5 (Parts 1–2); Def-H.4, Def-H.38, FT.22, FT.144(4), FT.145(3) (Part 3); Def-BM.3(7), Def-BM.18(1), BM.13(3), BM.19(2), BM.36(3), BM.44(2), BM.61(8), BM.70(2), BM.88(1), BM.90(3), BM.90(5), BM.94(7), BM.105(3), BM.105(7), BM.106(1), BM.116(2), BM.120(5) (this part); exact symbol-set censuses and intersections, exhaustive rational grid censuses, exact rational function rows [math; machine: \texttt{checks/\allowbreak{}verify\_\allowbreak{}H\_\allowbreak{}janus\_\allowbreak{}tunnel\_\allowbreak{}concrete\_\allowbreak{}spin.py} §D].

\textbf{Pf.}

%dsp-level:0
\[
(1):\ \text{the three monomials}\ 4\varepsilon_q(1-4\varepsilon_q),\ 1-4\varepsilon_q,\ 4\varepsilon_q\ \text{scanned: the only symbol is}\ \varepsilon_q\ [\text{machine}]
\]

%dsp-level:0
\[
(2):\ (\mathrm O1)\ \{\varepsilon_q\}\cap\{\delta_v,\hat\mu'^2\}=\varnothing\ [\mathrm{BM.90}(5);\ \text{machine: the set intersection}];
\]

%dsp-level:0
\[
\qquad (\mathrm O2)\ \mathrm{BM.36}(3),\ \mathrm{BM.61}(8)(\mathrm{O1}),\ \mathrm{BM.90}(3),\ \mathrm{BM.105}(7)\ \text{re-cited}
\]

%dsp-level:0
\[
(\mathrm O3)\ \mathrm{BM.88}(1)\text{'s inventory row re-cited}\ [\text{machine: the token census}];
\]

%dsp-level:0
\[
\qquad (\mathrm O4)\ \mathrm{P1},\ \text{A §0.3}\ \text{re-cited}
\]

%dsp-level:0
\[
(3):\ \mathrm{BM.90}(5)\text{'s}\ (\mathrm W)\ \text{and}\ \mathrm{BM.105}(7)\text{'s}\ (\mathrm{W6})\ \text{refutations re-cited};
\]

%dsp-level:0
\[
\qquad \text{(W7)'s form}\ \varepsilon_q(j)\ \text{is a state-indexed threshold, the species of}\ (\mathrm{W6})\text{'s}\ \varepsilon_i(D_S)\ [\text{machine: the form census}]
\]

%dsp-level:0
\[
(4):\ \mathrm{BM.94}(7),\ \mathrm{BM.105}(3)\text{'s bounds re-cited}\ [\text{machine}];
\]

%dsp-level:0
\[
\qquad \Gamma_{\mathrm{tr}}=x_+x_-/(1-x_+x_-)\ \text{with}\ x_\pm\ \text{functions of}\ \varepsilon_q\ \text{alone}\ [(1);\ \text{machine}]
\]

%dsp-level:0
\[
(5):\ (\mathrm i)\ \mathrm{Def\text{-}BM.18}(1)\text{'s basis is}\ \{|\mathcal V\rangle,|\mathcal R\rangle\}\ \text{with}\ \neg\text{a spin index}\ [\text{machine: the symbol census}];
\]

%dsp-level:0
\[
\qquad (\mathrm{ii})\ \mathrm{Def\text{-}H.38}\ \text{Scope},\ \mathrm{D3.7},\ \mathrm{EM.8}\ \text{re-cited}
\]

%dsp-level:0
\[
(\mathrm{iii})\ \mathrm{Def\text{-}H.4},\ \mathrm{FT.144}(4),\ \mathrm{BM.13}(3),\ \mathrm{BM.70}(2)\ \text{re-cited};
\]

%dsp-level:0
\[
\qquad \mathrm{BM.106}(1)\text{'s two-parity floor census re-cited}\ [\kappa=1\implies j=\tfrac12\ \text{and}\ \kappa=2\implies j=1\ \text{on one kind}]
\]

%dsp-level:0
\[
(6):\ (5)(\mathrm{ii})\text{'s scope row and}\ (1)\text{'s alphabet row conjoined:}\ \varepsilon\ \text{is defined on}\ \Psi(c),
\]

%dsp-level:0
\[
\ \text{which requires a bound chain, while}\ x_\pm\ \text{are defined on}\ \mathcal J,
\]

%dsp-level:0
\[
\ \text{which does not}\ [\text{machine: the two symbol sets are disjoint}]
\]

%dsp-level:0
\[
(7):\ \text{the two-row adjudication and the standing-rule scans}\ [\text{machine}].\ \blacksquare
\]

\textbf{Phys.} The crossing does not know the spin, and that is a theorem rather than an omission. The two weights are polynomials in one symbol — the crystallization threshold — and a scan of all three monomials finds no winding, no spin, no defect angle and no defect scale anywhere in them. A spin-gated tunnel would therefore be a map from an integer count to a coupling strength, which is exactly the species of coupling the corpus has already refused twice, on four grounds that apply verbatim: the alphabets are disjoint, no count-to-strength map is derived, there is no loop or beta function or Yukawa-type object to generate one, and there is only one free scale. So the hypothesis is named and refuted rather than quietly ignored. What replaces it is a domain statement, and it is sharper than a prohibition would have been. The register exists for every pair, but the statistics read requires a winding carrier, and a winding carrier requires a materialized closed bound chain on the nonnegative cone — so the transit and trivial chain classes, which is where the corpus puts its radiation, carry neither a spin label nor an exchange sign, while the register's moves still apply to them. The derived dichotomy is bound against transit, not fermion against boson, and the earlier spin-parity refutation on the kind stands beside it untouched. The clean way to say the whole result is that the exchange sign is defined on what crossed and not on the rate at which it crossed.

\medskip
\subsection*{\texorpdfstring{B.2.8\quad BM.140 — The ceiling study: the Platonic defect census, the dodecahedron selection theorem, the five-currency range table, and the adjudication of deficit 12}{BM.140 B.2.8}}
\addcontentsline{toc}{subsection}{B.2.8 BM.140}
\phantomsection\label{tag:BM-140}
{\small \textit{Record.}\quad Source: Part 4 (Part4\_Baryomesonic\_era.md); tier: Tier BM.23 — The Closing Tier: the Haar–Poisson Duality and the Lapse, the Selected Curvature Composition, the Momentum-Space Resolvent, the Charge and Statistics Censuses, the Two-Chain Ensemble, the Seed-Tensoring Mechanism, the Convention Declarations, the Ceiling Study, the Jensen Residue, the Pair Amplitude, and the Era Closure; canonical tag \texttt{BM.140}. Reproduced from the revision-1 canon; a cross-reference stub marks the former locus in the main corpus.}
\medskip
\textbf{Thm.}

%dsp-level:0
\[
(1)\quad\textbf{the Platonic census from the convexity constraint:}
\]

%dsp-level:0
\[
\qquad \text{a regular solid}\ \{p,q\}\ \text{exists}\iff q\,\frac{(p-2)\pi}p<2\pi\iff(p-2)(q-2)<4
\]

%dsp-level:0
\[
\qquad [\text{math: the vertex angle sum}];
\]

%dsp-level:0
\[
\boxed{\ \text{the integer solutions}\ p,q\ge3\ \text{are exactly five:}\quad(3,3),\ (4,3),\ (3,4),\ (5,3),\ (3,5)\ }
\]

%dsp-level:0
\[
\qquad [\text{machine: the exhaustive scan}\ p,q\le12];
\]

%dsp-level:0
\[
\frac{\delta_v}{\pi}=2-\frac{q(p-2)}p=\Bigl(1,\ \frac12,\ \frac23,\ \boxed{\frac15},\ \frac13\Bigr)
\]

%dsp-level:0
\[
\qquad \text{on}\ (\text{tetra},\ \text{cube},\ \text{octa},\ \text{dodeca},\ \text{icosa})
\]

%dsp-level:0
\[
\qquad [\text{machine: exact Fractions}];
\]

%dsp-level:0
\[
(2)\quad\textbf{the dodecahedron selection theorem (why the dodecahedron bounds):}
\]

%dsp-level:0
\[
\qquad \text{at fixed valence}\ q
\]

%dsp-level:0
\[
= 3:\quad\delta_v
\]

%dsp-level:0
\[
= \frac{6\pi}p-\pi\ \ \text{is strictly decreasing in}\ p,
\]

%dsp-level:0
\[
\ \text{and convexity terminates the series at}\ p
\]

%dsp-level:0
\[
= 5\ \bigl(p=6:\ 3\cdot\frac{4\pi}6=2\pi\ \text{flat}\bigr);
\]

%dsp-level:0
\[
\text{at fixed face}\ p=3:\quad\delta_v=2\pi-\frac{q\pi}3\ \ \text{is strictly decreasing in}\ q,\ \text{terminating at}\ q=5\ \bigl(q=6:\ \text{flat}\bigr):
\]

%dsp-level:0
\[
\qquad \frac\pi3>\frac\pi5;
\]

%dsp-level:0
\[
\boxed{\ \delta_v^{\mathrm{dodeca}}=\frac\pi5\ \ \textbf{is the minimum over the five, and the minimum is unique}\ }
\]

%dsp-level:0
\[
\qquad [\text{machine: the exhaustive comparison of the five exact Fractions}];
\]

%dsp-level:0
\[
\textbf{the structural reading:}
\]

%dsp-level:0
\[
\qquad \text{the defect shrinks as the face angle}\ \frac{(p-2)\pi}p\ \text{approaches}\ \pi\ \text{and the valence stays minimal; the pentagon is the largest Platonic face, so the three-pentagon vertex is the tightest convex vertex gap the series contains};
\]

%dsp-level:0
\[
(3)\quad\textbf{the defect factors and their exact order:}
\]

%dsp-level:0
\[
\qquad D_S=4\sin^2\frac{\delta_v}4=2-2\cos\frac{\delta_v}2
\]

%dsp-level:0
\[
\qquad [\mathrm{FT.76}(2);\ \mathrm{Def\text{-}H.37}(2);\ \text{machine: the half-angle identity on each row}]:
\]

%dsp-level:0
\[
D_S=\Bigl(2,\ \ 2-\sqrt2,\ \ 1,\ \ 2-\mathfrak u,\ \ 2-\sqrt3\Bigr),
\]

%dsp-level:0
\[
\qquad \mathfrak u=2\cos\frac\pi{10},
\]

%dsp-level:0
\[
\qquad \mathfrak u^4-5\mathfrak u^2+5=0
\]

%dsp-level:0
\[
\qquad [\mathrm{BM.114}(3);\ \text{machine: each value exact in}\ \mathbb M];
\]

%dsp-level:0
\[
\boxed{\ 2-\mathfrak u\ <\ 2-\sqrt3\ <\ 2-\sqrt2\ <\ 1\ <\ 2\ <\ 4\ }
\]

%dsp-level:0
\[
\qquad [\text{machine: the dyadic interval order in}\ \mathbb M;\ \text{certificates}\ \mathfrak u^2=\frac{5+\sqrt5}2,\ \sqrt5<\mathfrak u\ \text{on}\ 5<\mathfrak u^2]
\]

%dsp-level:0
\[
(4)\quad\textbf{the ceiling and its exact algebra:}
\]

%dsp-level:0
\[
\qquad m^2_{\mathrm{pole}}=\hat\mu'^2D_S(m)\in\bigl[0,\ 4\hat\mu'^2\bigr];
\]

%dsp-level:0
\[
\qquad \text{the widest intra-solid ratio is}\ \frac4{D_S},\ \ \text{maximal for the dodecahedron}
\]

%dsp-level:0
\[
\qquad [\mathrm{FT.76}(3);\ \mathrm{BM.125}(2);\ \text{machine re-run}];
\]

%dsp-level:0
\[
x:=\frac4{2-\mathfrak u}:
\]

%dsp-level:0
\[
\qquad \boxed{\ x^4-48x^3+304x^2-512x+256=0\ }
\]

%dsp-level:0
\[
\qquad [\mathrm{FT.76}(3)\ \text{re-derived};\ \text{machine:}\ u=2-\tfrac4x\ \text{substituted into}\ \mathfrak u^4-5\mathfrak u^2+5=0\ \text{and cleared}];
\]

%dsp-level:0
\[
\sqrt x<7
\]

%dsp-level:0
\[
\qquad \Bigl[x<49\iff\mathfrak u<\frac{94}{49}\iff\sqrt5<\frac{5667}{2401}:\quad5\cdot2401^2=28824005<32114889=5667^2\Bigr]
\]

%dsp-level:0
\[
\qquad [\mathrm{BM.79}(10)\text{'s certificate re-derived};\ \text{machine}];
\]

%dsp-level:0
\[
(5)\quad\textbf{the five-currency ceiling census (assembled, each row re-certified):}\qquad\begin{array}{l|l|l}\text{currency} & \text{the exact range row} & \text{locus}\\\hline\text{pole/defect} & m^2_{\mathrm{pole}}\in[0,4\hat\mu'^2];\quad\lambda^{(n)}\le4n;\quad\text{ratio}\le\sqrt{4/(2-\mathfrak u)}<7 & [\mathrm{FT.76}(3);\ \mathrm{BM.117}(4);\ \mathrm{BM.125}(2)]\\\text{cover} & K=2|\mathcal C|\le\tfrac{3V}2\ (\text{bypass}),\ \le2V\ (\text{locked});\quad\text{cube ratios}\le6 & [\mathrm{FT.79}(2)]\\\text{form factor} & M_{\mathrm{form}}\le c_wV,\qquad c_w=(6,36,1,6,24) & [\mathrm{FT.97}(5)]\\\text{transport} & \text{the channel ratios}\ z_0/z_1:\quad q>139,\quad\nu\approx272,\quad\ell\approx3477\ \ (\text{enclosed}) & [\mathrm{BM.80}(6);\ \mathrm{FT.113};\ \mathrm{FT.119};\ \mathrm{FT.143}(3)]\\\text{informational} & m_{\mathrm{hid}}=E_{\mathrm{int}}\ln2/\mathfrak h:\quad\text{span}\ 30\ \text{on the computed-cell ladder, unbounded on the polygon ladder} & [\mathrm{BM.115}(2);\ \mathrm{FT.45};\ \mathrm{FT.47}]\end{array}
\]

%dsp-level:0
\[
\textbf{the transport rows against the pole ceiling:}
\]

%dsp-level:0
\[
\qquad \boxed{\ \sqrt{4/(2-\mathfrak u)}<7<139<\frac{z_0^{(q)}}{z_1^{(q)}}<\frac{z_0^{(\nu)}}{z_1^{(\nu)}}<\frac{z_0^{(\ell)}}{z_1^{(\ell)}}\ }
\]

%dsp-level:0
\[
\qquad [\mathrm{BM.79}(10);\ \mathrm{BM.80}(6);\ \text{machine: the exact rational enclosures re-derived — the trigonometric channels by Taylor bounds with rational remainders, the square roots by Newton certificates}];
\]

%dsp-level:0
\[
\textbf{the comparison is arithmetic,}\ \neg\textbf{a bridge:}
\]

%dsp-level:0
\[
\qquad \text{the currencies are incommensurable}\ \bigl(\ln1280/\ln960\notin\mathbb Q\bigr)\ \text{and no derived equation converts them}
\]

%dsp-level:0
\[
\qquad [\mathrm{BM.83}(2)\text{--}(4);\ \mathrm{BM.37}(2)];
\]

%dsp-level:0
\[
(6)\quad\textbf{the per-class rows (all}\ \nu\textbf{-,}\ \ell\textbf{- and q-channel classes censused):}
\]

%dsp-level:0
\[
\qquad \textbf{the}\ \nu\ \textbf{line:}\quad m_\nu^{\mathrm{rest}\,2}=\hat\mu_\nu^2\bigl(\ell_c^{(\nu)}\bigr)^2G(c(i_\nu)),
\]

%dsp-level:0
\[
\qquad \delta_\nu=\tfrac18,
\]

%dsp-level:0
\[
\qquad G_2>G_1>A_0>G_0,
\]

%dsp-level:0
\[
\qquad A_0-B_0\le G\le A_0+B_0,
\]

%dsp-level:0
\[
\qquad m_\nu^{\mathrm{rest}}>0\ \text{strictly}
\]

%dsp-level:0
\[
\qquad [\mathrm{BM.96}(1)\text{--}(3);\ \text{the SOS sandwiches re-cited}];
\]

%dsp-level:0
\[
\textbf{the}\ \ell\ \textbf{and}\ q\ \textbf{channels:}
\]

%dsp-level:0
\[
\qquad r_i=1+\sqrt2\cos\bigl(\delta+\tfrac{2\pi i}3\bigr),
\]

%dsp-level:0
\[
\qquad \delta_\ell=\tfrac29,
\]

%dsp-level:0
\[
\qquad \delta_q=\tfrac{\delta_\ell}3=\tfrac{2}{27},
\]

%dsp-level:0
\[
\qquad 3\delta_q=\delta_\ell
\]

%dsp-level:0
\[
\qquad [\mathrm{FT.113};\ \mathrm{FT.119}(1),(2);\ \mathrm{FT.143}(1)];
\]

%dsp-level:0
\[
\text{each channel's cubic is its own:}
\]

%dsp-level:0
\[
\qquad 2(r-1)^3-3(r-1)-\sqrt2\cos3\delta=0
\]

%dsp-level:0
\[
\qquad [\mathrm{FT.143}(2);\ \mathrm{FT.113};\ \text{machine: the ring identity per channel}];
\]

%dsp-level:0
\[
\textbf{the attached six-state splitting:}
\]

%dsp-level:0
\[
\qquad \Delta_g=\bigl(2,\ \sqrt2-1,\ 2\cos\tfrac\pi{10}-\sqrt3\bigr)\hat\mu'^2,
\]

%dsp-level:0
\[
\qquad 0<\Delta_g<4\hat\mu'^2\ \ \text{on every floor}
\]

%dsp-level:0
\[
\qquad [\mathrm{BM.75}(4);\ \text{machine: the bound census in}\ \mathbb M];
\]

%dsp-level:0
\[
\textbf{the six-state bridge:}\qquad\mathcal R_s=\frac{D_S}{4E^2},\qquad\frac{\mathcal R_{\max}}{\mathcal R_{\min}}=\frac{50}{1-\mathfrak u}\qquad[\mathrm{BM.82}(1),(3)\ \text{re-cited}];
\]

%dsp-level:0
\[
\textbf{every class row lies inside its own currency's range and outside every other's identification:}
\]

%dsp-level:0
\[
\qquad \neg\text{a cross-currency equality is formed}\ [\mathrm{BM.83}(4);\ \mathrm{BM.98}(6)];
\]

%dsp-level:0
\[
(7)\quad\textbf{deficit 12 adjudicated — closed:}
\]

%dsp-level:0
\[
\qquad (\mathrm a)\ \ \text{the ceiling census is complete: every formed mass read of the corpus carries an exact range row}\ [(5)];
\]

%dsp-level:0
\[
(\mathrm b)\ \ \text{the bounding solid is explained: the dodecahedron's}\ \pi/5\ \text{is the unique minimal Platonic defect}\ [(2)];
\]

%dsp-level:0
\[
(\mathrm c)\ \ \text{the SM-mass identification remains}\ \neg\text{formed on the standing bars:}\ \neg\text{a datum}\ [\mathrm{P5}],
\]

%dsp-level:0
\[
\ \text{and the range failures re-certified row-wise}\ [\mathrm{FT.77};\ \mathrm{FT.79}(4)\text{--}(6);\ \mathrm{FT.84}(2)];
\]

%dsp-level:0
\[
\boxed{\ \textbf{deficit 12 is closed as adjudicated: the ceilings are derived, the selection is explained, and what stands is a calibration bar,}\ \neg\textbf{a structural gap}\ }
\]

%dsp-level:0
\[
\qquad [\text{item 20's bracket updated at the ledger}];
\]

%dsp-level:0
\[
(8)\quad\textbf{status:}
\]

%dsp-level:0
\[
\qquad \textbf{derived:}\ (1)\text{--}(6);
\]

%dsp-level:0
\[
\qquad \textbf{adjudicated:}\ (7);
\]

%dsp-level:0
\[
\qquad \neg\text{a declaration},
\]

%dsp-level:0
\[
\qquad \neg\text{a P4 addition},
\]

%dsp-level:0
\[
\qquad \neg\text{a measured value and}\ \neg\text{an}\ [\mathrm{EXT}]\ \text{datum in any clause}\ [\mathrm{P5};\ \mathrm{FT.73}(5)\ \text{discipline}].
\]

\textbf{Req.} D2.9 (Part 1); P1, P3, P5 (Parts 1–2); Def-H.37(2), FT.45, FT.47, FT.76(1), FT.76(2), FT.76(3), FT.77, FT.79(2), FT.79(4), FT.97(5), FT.113, FT.119, FT.143(1), FT.143(2), FT.143(3), FT.84(2) (Part 3); Def-BM.14(3), Def-BM.17, Def-BM.19(3), BM.37(2), BM.75(4), BM.79(10), BM.80(6), BM.82(1), BM.82(3), BM.83(2), BM.83(4), BM.96, BM.98(6), BM.114(3), BM.115(2), BM.117(4), BM.125(2) (this part); Census FT.63, FT.76 (row-wise); the exhaustive Platonic admissibility scan, exact Fraction order, the dyadic interval embedding of $\mathbb M$, Newton square-root certificates, Taylor enclosures with rational remainders, resultant-free minimal-polynomial derivation by cleared substitution [math; machine: \texttt{checks/\allowbreak{}verify\_\allowbreak{}H\_\allowbreak{}bm\_\allowbreak{}closing\_\allowbreak{}tier.py} §I].

\textbf{Pf.}

%dsp-level:0
\[
(1):\ (p-2)(q-2)<4\ \text{with}\ p,q\ge3\ \text{integers:}\ p-2,q-2\ge1\ \text{and}\ (p-2)(q-2)\in\{1,2,3\}:
\]

%dsp-level:0
\[
\quad (1,1),(2,1),(1,2),(3,1),(1,3)\mapsto\ \text{the five pairs}\ [\text{machine: the scan}];
\]

%dsp-level:0
\[
\qquad \delta_v=2\pi-\frac{q(p-2)\pi}p\ \text{evaluated exactly}\ [\text{machine}]
\]

%dsp-level:0
\[
(2):\ \text{the monotonicity rows:}\ \frac{6\pi}p-\pi\ \text{decreasing in}\ p\in\{3,4,5\};
\]

%dsp-level:0
\[
\qquad 2\pi-\frac{q\pi}3\ \text{decreasing in}\ q\in\{3,4,5\};
\]

%dsp-level:0
\[
\qquad \min\bigl(1,\tfrac12,\tfrac23,\tfrac15,\tfrac13\bigr)=\tfrac15\ \text{unique}\ [\text{machine: exact Fraction comparison}]
\]

%dsp-level:0
\[
(3):\ 4\sin^2\frac{\delta_v}4=2-2\cos\frac{\delta_v}2\ [\text{the half-angle identity}];
\]

%dsp-level:0
\[
\qquad \text{the five values at}\ \frac{\delta_v}2\in\bigl\{\frac\pi4,\frac\pi8,\frac\pi6,\frac\pi{20},\frac\pi{12}\bigr\}:
\]

%dsp-level:0
\[
\quad 2,\ 2-\sqrt2,\ 1,\ 2-\mathfrak u,\ 2-\sqrt3\ [\text{machine: each in}\ \mathbb M];
\]

%dsp-level:0
\[
\qquad \text{the order by dyadic intervals}\ [\text{machine}]
\]

%dsp-level:0
\[
(4):\ u=2-\frac4x\ \text{into}\ u^4-5u^2+5=0,\ \text{cleared by}\ x^4:
\]

%dsp-level:0
\[
\quad (2x-4)^4-5x^2(2x-4)^2+5x^4=x^4-48x^3+304x^2-512x+256\ [\text{machine: the expansion in}\ \mathbb Z[x]];
\]

%dsp-level:0
\[
x<49
\]

%dsp-level:0
\[
\iff 4<49(2-\mathfrak u)
\]

%dsp-level:0
\[
\iff 49\mathfrak u<94\iff\mathfrak u^2<\frac{8836}{2401}\iff\frac{5+\sqrt5}2<\frac{8836}{2401}\iff\sqrt5<\frac{17672-12005}{2401}
\]

%dsp-level:0
\[
= \frac{5667}{2401}\ [\text{machine: the integer certificate}]
\]

%dsp-level:0
\[
(5),
\]

%dsp-level:0
\[
(6):\ \text{the frozen rows re-cited and the enclosures re-derived:}\ \cos\delta,
\]

%dsp-level:0
\[
\ \sin\delta\ \text{by alternating Taylor sums with rational remainder bounds at}\ \delta\in\bigl\{\frac18,\frac29,\frac{2}{27}\bigr\};
\]

%dsp-level:0
\[
\qquad \sqrt2,\sqrt6\ \text{by Newton certificates};
\]

%dsp-level:0
\[
\qquad \text{the ratio brackets by monotone interval division}\ [\text{machine}];
\]

%dsp-level:0
\[
\qquad (7),(8):\ \text{assembly}.\ \blacksquare
\]

\textbf{Phys.} The last unexplained selection in the mass ledger is explained: the dodecahedron bounds the pole spectrum because its vertex defect is the unique minimum of the Platonic series, and the series terminates there for a reason of pure convexity — three pentagons is the tightest vertex gap five regular polygons can leave, since a fourth polygon type would have to be the hexagon, and three hexagons are flat. The whole ceiling structure then assembles into one table: five currencies, each with its exact range row, mutually incommensurable, with the transport channels' hierarchies lying provably above the widest Platonic ratio and the identification bars standing. Every neutrino, lepton and quark class is censused inside its own currency. Deficit 12 closes as adjudicated: the ceilings are theorems, the bounding solid is derived, and what remains against external masses is the calibration bar the corpus has carried since its first audit — not a missing structure.

\medskip
\subsection*{\texorpdfstring{B.2.9\quad FT.68 — Joint skeletons: tetrahedron, the polygon ladder, the cube; the layer dichotomy; flicker (catalog entry)}{FT.68 B.2.9}}
\addcontentsline{toc}{subsection}{B.2.9 FT.68}
\phantomsection\label{tag:FT-68-cat}
{\small \textit\{Catalog.\}\quad The full theorem remains in the main corpus as Theorem 3.56 (Part 3); its derived clauses are consumed by later derivations. The hypothesis-level ruling is recorded here verbatim:}
\begin{quote}\small (d) \textbf{Verdict on "the path integral fuzzes the flickering 2D shape into a higher-tier solid"}: flickering thread gas — \textit{yes}, exact and strictly positive; fuzzing into solids — \textit{no} at the amplitude layer **[scope: a literal-domain (S2) verdict — in the chain-local domains every computed Platonic floor and every even polygon floor carries strictly positive amplitude (tetra\end{quote}
\medskip
\subsection*{\texorpdfstring{B.2.10\quad FT.79 — The combinatorial ceiling and the cover-hierarchy audit (second clean negative) (catalog entry)}{FT.79 B.2.10}}
\addcontentsline{toc}{subsection}{B.2.10 FT.79}
\phantomsection\label{tag:FT-79-cat}
{\small \textit\{Catalog.\}\quad The full theorem remains in the main corpus as Theorem 3.63 (Part 3); its derived clauses are consumed by later derivations. The hypothesis-level ruling is recorded here verbatim:}
\begin{quote}\small \end{quote}
\medskip
\subsection*{\texorpdfstring{B.2.11\quad BM.57 — The continuation-symmetry census: the vertex asymmetry is two, and it splits nothing (catalog entry)}{BM.57 B.2.11}}
\addcontentsline{toc}{subsection}{B.2.11 BM.57}
\phantomsection\label{tag:BM-57-cat}
{\small \textit\{Catalog.\}\quad The full theorem remains in the main corpus as Theorem 4.48 (Part 4); its derived clauses are consumed by later derivations. The hypothesis-level ruling is recorded here verbatim:}
\begin{quote}\small And it splits nothing.\end{quote}
\medskip
\section*{\texorpdfstring{B.3 No-Go Theorems}{B.3 No-Go Theorems}}
\addcontentsline{toc}{section}{B.3 No-Go Theorems}
\markboth{B.3 No-Go Theorems}{B.3 No-Go Theorems}
\subsection*{\texorpdfstring{B.3.1\quad CT4.4a — The realization domain of $\iota_n$: non-existence at $n\ge4$, persistence refuted}{CT4.4a B.3.1}}
\addcontentsline{toc}{subsection}{B.3.1 CT4.4a}
\phantomsection\label{tag:CT4-4a}
{\small \textit{Record.}\quad Source: Part 2 (Part2\_Bosonic\_era.md); tier: §G.I — The Saturation Lock; canonical tag \texttt{CT4.4a}. Reproduced from the revision-1 canon; a cross-reference stub marks the former locus in the main corpus.}
\medskip
\textbf{Thm.} $(1)$ \textbf{Non-existence.} The aspect realization $\iota_n:\mathfrak u(n)\to\mathcal K$ of D2.12 exists exactly for $n\le3$:

%dsp-level:0
\[
n \ge 4
\]

%dsp-level:0
\[
\implies D_{\mathrm{tick}}(n)
\]

%dsp-level:0
\[
= 2^{n+1}
\]

%dsp-level:0
\[
\ge 32>16
\]

%dsp-level:0
\[
= C_{\mathrm{cap}}\implies\neg\text{channel-admissible}\ [\mathrm{CT0.5},\ \mathrm{CT0.6},\ \mathrm{CT0.11}]\implies\text{the rung-}n\ \text{slot algebra is}\ \neg\text{realized}\ [\mathrm{CT4.4}\ \text{Pf clause (b), via CT3.20}]
\]

and the locked sector receiving rung-$\ge4$ content carries no gauge structure — no generators, no internal directions [CT4.4; D4.1: $F_{D^-}$ a single possibility space; D4.2: absorption at rate $\propto O(n)$]. Hence the silence of D2.12 at $n\ge4$ is a non-existence consequence, $\neg$an underdetermination (scope note added at D2.12): the unique derived read of rung-$n\ge4$ slot content is $D^-$-absorption at the ledger level — locked, retained [CT4.3], never accessible [D1.11]. $(2)$ \textbf{No Lie-homomorphic substitute.} Any Lie-homomorphic map $\iota_n$, $n\ge4$, vanishes on $\mathfrak{su}(n)$: $\mathfrak{su}(n\ge4)$ is simple with $\dim=n^2-1\ge15$ [math]; block-wise realization semantics [D2.12, D3.2] require the isomorphism class to lie inside a single aspect block [CT0.12/CT0.13 tables]; the maximal aspect block dimension is $8$ [the $\mathfrak{su}(3)$ colour block; CT3.8] $<15$; the $J^\pm$ spans are $\neg$Lie-closed ($[J_i,J_j]=i\epsilon_{ijk}W_k\notin\operatorname{span}\{J\}$ [CT1.6, CT1.9]); the $H$, EM and $D^+$ blocks are $1$-dim abelian (a simple nonabelian algebra has no nonzero abelian homomorphic image [math]); the $G$ block has $\dim5<15$ [CT0.12]; the $D^-$ block: clause (1). The ambient commutator algebra $\mathfrak u(4)$ [CT1.9] contains an abstract $\mathfrak{su}(4)$, but no aspect block does, and the force sum of A8 has no $\mathfrak{su}(4)$ summand. $(3)$ \textbf{Persistence refuted.} "Strong-hosting of $\mathfrak{su}(3)$ subalgebras at $n\ge4$" is impossible by (1)+(2). The activation persistence of CT0.13 at $n\ge4$ is \textit{configuration-level}: three-slot sealing substructures remain available as data — attachments [D3.8], invariant contractions [CT3.17], the completeness projector [D2.14] — not as $\iota$-images. \textbf{Hosting dies at rung 3; sealing persists at every rung.}

\textbf{Req.} A8; D1.11, D2.12, D3.2, D3.8, D4.1, D4.2; CT0.5, CT0.6, CT0.11, CT0.12, CT0.13, CT1.6, CT1.9 (Part 1); CT3.8, CT3.17, CT3.20 (Part 2); CT4.2, CT4.3, CT4.4; dimension arithmetic, simplicity of $\mathfrak{su}(n\ge4)$, the classification of low-dimensional simple algebras [math; machine: \texttt{checks/\allowbreak{}verify\_\allowbreak{}merge\_\allowbreak{}tierA.py} §C].

\textbf{Pf.}

%dsp-level:0
\[
(1):\ \mathrm{CT0.5}:\ D_{\mathrm{tick}}(n)=2^{n+1};
\]

%dsp-level:0
\[
\quad \mathrm{CT0.6}:\ C_{\mathrm{cap}}=16;
\]

%dsp-level:0
\[
\quad n\ge4:\ 2^{n+1}\ge32>16\implies O(n)=2^{n+1}-16>0\ [\mathrm{CT4.1}]\implies\text{locked}\ [\mathrm{CT4.2}/\mathrm{CT4.3}]
\]

%dsp-level:0
\[
\mathrm{CT4.4}:\ F_{D^-}\ \text{contributes}\ \neg\text{generator to any accessible algebra};
\]

%dsp-level:0
\[
\qquad \mathrm{CT3.20}:\ N=3\ \text{forced at the accessible rungs}\implies\neg\exists\ \mathfrak{su}(n\ge4)\ \text{aspect block}
\]

%dsp-level:0
\[
(2):\ \dim\mathfrak{su}(n)=n^2-1\ge15\ (n\ge4)\ [\text{math}];
\]

%dsp-level:0
\[
\qquad \text{aspect block dimensions}\ \{3,3,5,3,1,1\}\ [\mathrm{CT0.12}]\ \cup\ \{8\}\ [\mathrm{CT3.8}]\implies\max
\]

%dsp-level:0
\[
= 8<15\implies\neg\exists\ \text{injective block homomorphism}
\]

%dsp-level:0
\[
[J_1,J_2]=iW_3\notin\operatorname{span}\{J_i\}\ [\mathrm{CT1.6};\ \text{machine over}\ \mathbb Q(i)]\implies J\text{-spans}\ \neg\text{Lie-closed};
\]

%dsp-level:0
\[
\qquad \dim1\ \text{blocks abelian}\implies\text{simple nonabelian image}=0\ [\text{math}]
\]

%dsp-level:0
\[
(3):\ \text{by (1)+(2): no hosted image at}\ n\ge4;
\]

%dsp-level:0
\[
\qquad \mathrm{CT0.13}:\ \mathcal A(n\ge4)\ \text{is defined by activation/absorption data},\ \neg\text{by}\ \iota_n\text{-images}.\ \blacksquare
\]

\textbf{Phys.} The gauge ladder has a hard derived ceiling: nothing above colour ever hosts — no hidden $\mathfrak{su}(4)$, no grand-unification block waiting at high rung. Oversubscribed content is absorbed without structure, and what survives as "colour persists at $n\ge4$" is the sealing \textit{grammar} (which configurations are admissible), never the Lie algebra.

\textbf{Rmk (positive side, $n\le3$).} From the same frozen sources: the centre of $\mathfrak u(n)$ maps to $\mathcal K^{(\mathrm{EM})}$ [D2.12]; at $n(w)=2$: $\mathfrak{su}(2)$ hosts into $\operatorname{span}\{W_i\}$ — the $W_i$ close $\mathfrak{su}(2)$ [CT1.4]; at $n(w)=3$: $\mathfrak{su}(3)$ hosts onto $\mathcal K^{(\mathrm{strong})}$ with $8=8$ [D2.12, CT3.7, CT3.8]; the exchange generators host \textit{literally} as matrix units ($\tfrac12(\lambda_1+i\lambda_2)=E_{12}$ and cyclic [CT3.9's basis; machine: \texttt{checks/\allowbreak{}verify\_\allowbreak{}merge\_\allowbreak{}tierB.py} §G]). The hosting ceiling is therefore exact: nontrivial hosting $\iff$ contact exchange degree $\le3$ (the explicit-realization vocabulary stays outside the frozen corpus).

\medskip
\subsection*{\texorpdfstring{B.3.2\quad FT.59 — The $\iota_{n\ge4}$ specification: non-existence, persistence refuted, hosting ceiling}{FT.59 B.3.2}}
\addcontentsline{toc}{subsection}{B.3.2 FT.59}
\phantomsection\label{tag:FT-59}
{\small \textit{Record.}\quad Source: Part 3 (Part3\_Fermion\_era.md); tier: Tier F.19 — The Three Tails: Realization Non-Existence, the Multiplicity Map, and the Selection Obstruction; canonical tag \texttt{FT.59}. Reproduced from the revision-1 canon; a cross-reference stub marks the former locus in the main corpus.}
\medskip
\textbf{Thm.} $(1)$ \textbf{Non-existence (the specification question is void).} D2.12 defines $\iota_n:\mathfrak u(n)\to\mathcal K$ as the \textit{aspect realization} of the rung-$n$ slot algebra. For $n\ge4$ no such realization exists: slot algebras are realized only for $n\le3$ [CT4.4 Pf clause (b), via CT3.20/CT0.11: $n\ge4\implies\neg$channel-admissible $\implies$ slot algebra $\neg$realized], and the locked sector that receives rung-$\ge4$ content carries no gauge structure — no generators, no internal directions [CT4.4; D4.1: $F_{D^-}$ a single possibility space; D4.2: absorption at rate $\propto O(n)$]. Hence D2.12's silence at $n\ge4$ is a non-existence consequence, not an underdetermination: \textbf{the unique derived read of rung-$n\ge4$ slot content is $D^-$-absorption at the ledger level} (locked, retained [CT4.3], never observable [FT.3, FT.20, D1.11]), while the observable reads at such vertices proceed through data that never required $\iota_n$: registration signs $\sigma_e$ and chain-local pass-through reads $f_v$ [CT0.9, D3.3, FT.21(i)], core charges [FT.40], and sealing configuration data [D3.8, CT3.17, Def-H.10]. $(2)$ \textbf{Lie-theoretic backup (if a linear $\iota$ were posited anyway).} Any Lie-homomorphic $\iota_n$, $n\ge4$, has $\iota_n(\mathfrak{su}(n))=0$: (a) observable side — $\mathfrak{su}(n\ge4)$ simple with $\dim=n^2-1\ge15>9=\max\dim$ observable factor [FT.58(3): every Lie-type-respecting hosting into $\mathcal K^{\mathrm{obs}}=9\oplus9$ is trivial]; (b) non-observable blocks admit no image either: $J^\pm$ spans are not Lie-closed ($[J_i,J_j]=i\epsilon_{ijk}W_k\notin\operatorname{span}\{J\}$ [CT1.6, CT1.9; machine over $\mathbb Q(i)$]) and a simple image inside $\operatorname{span}\{J_i\}$ is abelian-or-zero; $H$, $\mathrm{EM}$, $D^+$ blocks are $1$-dim abelian (simple nonabelian $\mathfrak{su}(n\ge4)$ has no nonzero abelian homomorphic image [math]); $G$ has $\dim5<15$; $D^-$ has no gauge structure [CT4.4]. Block-wise realization semantics [D2.12/D3.2: the isomorphism class must lie in an aspect block or the pair blocks of CT0.12] excludes spread embeddings: no single block contains an $\mathfrak{su}(n\ge4)$ isomorphism class ($\dim\le9<15$ in every block; $\mathfrak{su}(4)\not\subseteq\mathfrak{su}(3)$ [math: classification]). The centre read survives only where $\iota_n$ exists ($n\le3$): at $n\ge4$ the EM row runs on registration data, as in (1). $(3)$ \textbf{Persistence refuted (FT.58(1)'s conditional branch resolved).} "Strong-hosting of $\mathfrak{su}(3)$ subalgebras at $n(w)\ge4$" is impossible: $\mathfrak{su}(3)_{ijk}\subset\mathfrak{su}(n\ge4)$ [machine: explicit matrix-unit embedding, bracket-closed inside the $15$-dim traceless span, CT1.9] together with (1)/(2) forces any hosted image of a slot $\mathfrak{su}(3)$ subalgebra at $n\ge4$ to vanish. CT0.13's "rung-3 structure persists" at $n\ge4$ is therefore necessarily \textbf{activation-level, not $\iota$-image-level}: strong $\in\mathcal A(n\ge4)$ means three-slot sealing substructures remain available as \textit{configuration} data — attachments [D3.8], invariant contractions [CT3.17], the $\mathcal C$-projector [Def-H.10], the grammar $g$ [FT.41(b)] — not that $\iota_n$ hosts anything. \textbf{Hosting dies at rung 3; sealing persists at every rung.} $(4)$ \textbf{Corollaries.} (a) \textit{Face floor settled}: $\Sigma_\triangle$'s corolla contact has $n(D)=6$ [FT.49(5)]: no realization $\iota_6$ exists; its $30$-dim exchange sector [$\mathfrak u(6)=36=1+35$; strictly new $=30$: FT.49(8), FT.41(d)] is wholly unhosted — kernel under Lie semantics, $D^-$-absorbed ledger content under CT4.4's, observationally identical [FT.20: neither contributes an observable]. The first computed cell in this regime. (b) \textit{Hosting ceiling, exact}: a corolla exchange algebra hosts nontrivially iff $N(\Sigma)\le n(w)\le3$: $N=2\to W$-hosted, $N=3\to$ strong-hosted literally [FT.58(1)]; every $n(w)\ge4$ corolla and every separated $N\ge2$ contact [FT.58(2)] is unhosted. FT.58(4)'s capacity ceiling $\mathfrak{su}(3)$ is now also the \textit{realization} ceiling. (c) \textit{Def-H.22(a) scope correction}: the corolla formula $\iota_\Sigma^{\mathrm{cor}}=\iota_{n(w)}|_{\mathfrak u(N)}$ is defined exactly for $n(w)\le3$; at $n(w)\ge4$ it is void, replaced by the absorption read of (1) — patched in place. The same domain restriction applies to Def-H.13's per-edge factor $\iota_{n(w_e)}$: a boundary edge whose external endpoint has rung $\ge4$ contributes no $\iota$-factor — its read is registration-plus-sealing data. Among computed cells only the face floor exercises the case (all six boundary edges end at $D$, $n(D)=6$); every other computed floor has environment contacts at $n(w_e)\in\{1,2\}$. (d) \textit{Two ledgers, no identification}: the algebra-kernel ledger ($\dim\mathfrak{su}(n)=n^2-1$ unhosted) and the register-absorption ledger ($O(n)=2^{n+1}-16$ locked registers/tick [CT0.5, CT0.6, D4.2]) are different counts, both growing in $n$; neither converts into the other.

\textbf{Req.} D2.12, D3.1, D3.2, D3.3, D3.8; D4.1, D4.2; A7; CT0.5, CT0.6, CT0.9, CT0.11, CT0.13, CT0.14 (Part 1); CT1.4, CT1.6, CT1.7, CT1.9 (Part 1); CT3.17, CT3.20 (Part 2); CT4.1, CT4.3, CT4.4 (Part 2); FT.3, FT.20, FT.21, FT.40, FT.41(b),(d), FT.49, FT.58; Def-H.10, Def-H.13, Def-H.22; simplicity of $\mathfrak{su}(N\ge4)$, absence of nonzero abelian images of simple algebras, $\mathfrak{su}(4)\not\subseteq\mathfrak{su}(3)$ [math].

\textbf{Pf.}

%dsp-level:0
\[
(1):\ \mathrm{CT4.4\ Pf}(b):\ \text{gauge structure arises only from pair blocks or slot algebras}\ \mathfrak u(n);
\]

%dsp-level:0
\[
\quad\text{slot algebras realized only for}\ n
\]

%dsp-level:0
\[
\le 3\ [\mathrm{CT3.20}/\mathrm{CT0.11}:\ 2^{n+1}>16\Rightarrow\neg\text{admissible}]
\]

%dsp-level:0
\[
n \ge 4:\ \text{slot content}\to F_{D^-}\ [\mathrm{D4.1}:\ \oplus_{n\ge4}F_n\ \text{receives}\ \forall\ \text{configs at rung}\ \ge4];
\]

%dsp-level:0
\[
\quad\mathrm{CT4.4}:\ F_{D^-}\ \text{scalar sink},
\]

%dsp-level:0
\[
\ \neg\exists\ \text{gauge generators};
\]

%dsp-level:0
\[
\quad\text{absorption rate}\propto O(n)\ [\mathrm{D4.2}];
\]

%dsp-level:0
\[
\ \text{locked}\ \neg\text{accessible}\ [\mathrm{D1.11},\mathrm{FT.3}]
\]

%dsp-level:0
\[
\text{observable reads without}\ \iota_n:\ \sigma_e\ [\mathrm{CT0.9}];
\]

%dsp-level:0
\[
\quad f_v\ \text{chain-local}\ [\mathrm{D3.3},\ \mathrm{FT.21}(i)\ \text{telescoping}];
\]

%dsp-level:0
\[
\quad q_{\mathrm{core}}\ [\mathrm{FT.40}];\quad\text{sealing}\ [\mathrm{D3.8},\mathrm{CT3.17},\mathrm{Def\text{-}H.10}]
\]

%dsp-level:0
\[
(2a):\ \mathrm{FT.58}(3)\ \text{verbatim};
\]

%dsp-level:0
\[
\qquad (2b):\ [J_1,J_2]=iW_3\notin\operatorname{span}_{\mathbb R}\{J_i\}\ [\mathrm{CT1.6};\ \text{machine}:\ \operatorname{rk}_{\mathbb Q(i)}\{J_1,J_2,J_3,W_3\}=4]
\]

%dsp-level:0
\[
\text{simple}\ \mathfrak{su}(n\ge4):\ \text{hom. image}
\]

%dsp-level:0
\[
= 0\ \vee\ \cong\mathfrak{su}(n)\ [\ker=\text{ideal};\ \text{math}];
\]

%dsp-level:0
\[
\quad\dim\ \text{of every aspect block}
\]

%dsp-level:0
\[
\le 9<15\ [\mathrm{CT0.12},\mathrm{CT0.14}:\ 3,3,5,3,1,1,1,8];
\]

%dsp-level:0
\[
\quad\text{abelian blocks}: 0\ \text{image}\ [\text{math}]
\]

%dsp-level:0
\[
\text{spread embeddings}:\ \mathrm{D3.2}\ \text{semantics}
\]

%dsp-level:0
\[
= \text{isomorphism class in a block};
\]

%dsp-level:0
\[
\ \mathfrak{su}(n\ge4)\ \text{in no block}\ (2a),\qquad (2b);
\]

%dsp-level:0
\[
\qquad \text{pair-total reading excluded by}\ \mathrm{CT4.4}\ (\text{Rmk below})
\]

%dsp-level:0
\[
(3):\ \mathfrak{su}(3)_{123}\hookrightarrow\mathfrak{su}(4)\ \text{by matrix units}\ [\text{machine}:\ 8\text{-dim},\ \text{bracket-closed inside}\ \mathfrak{su}(4)=\operatorname{span}\{\Sigma,\hat E^{ij},W_i,J_i,K_i\},\ \dim15\ [\mathrm{CT1.9}]];
\]

%dsp-level:0
\[
\quad\iota(\mathfrak{su}(3)_{123})\subseteq\iota(\mathfrak{su}(n))\qquad = 0
\]

%dsp-level:0
\[
\mathrm{CT0.13}:\ \mathcal A(n\ge4)
\]

%dsp-level:0
\[
= \mathcal A(3)\cup\{D^-\}\ \text{is an activation statement}\ [\mathrm{D1.12}];
\]

%dsp-level:0
\[
\quad\text{sealing data at}\ n
\]

%dsp-level:0
\[
\ge 4:\ \mathrm{D3.8}\ \text{attachments},
\]

%dsp-level:0
\[
\ \mathrm{CT3.17}\ \text{contractions},
\]

%dsp-level:0
\[
\ \mathrm{Def\text{-}H.10}\ \mathcal C,
\]

%dsp-level:0
\[
\ \mathrm{FT.41}(b)\ g:\ \text{configuration ledger},
\]

%dsp-level:0
\[
\ \neg\iota\text{-image}
\]

%dsp-level:0
\[
(4a):\ n(D)
\]

%dsp-level:0
\[
= 6\ [\mathrm{FT.49}(5)];
\]

%dsp-level:0
\[
\ (1)\Rightarrow\neg\exists\,
\]

%dsp-level:0
\[
\iota_6;
\]

%dsp-level:0
\[
\quad\dim\mathfrak u(6)
\]

%dsp-level:0
\[
= 36
\]

%dsp-level:0
\[
= 1+35;
\]

%dsp-level:0
\[
\quad\text{exchange}
\]

%dsp-level:0
\[
= 30\ [\mathrm{FT.49}(8)];
\]

%dsp-level:0
\[
\quad\text{observational identity}:\ \mathrm{FT.20}\ (D^-\ \text{never}),
\]

%dsp-level:0
\[
\ \mathrm{FT.58}(3)\ (\text{Lie}: 0)
\]

%dsp-level:0
\[
(4b):\ \mathrm{FT.58}(1)\ (n(w)\in\{2,3\})+(1)\text{--}(3)\ (n(w)\ge4)+\mathrm{FT.58}(2)\ (\text{separated});
\]

%dsp-level:0
\[
\qquad (4c):\ \mathrm{Def\text{-}H.22}(a)\ \text{composition requires}\ \iota_{n(w)}\ [\mathrm{D2.12}]\implies\text{domain}\ n(w)\le3
\]

%dsp-level:0
\[
(4d):\ n^2-1\ \text{vs}\ 2^{n+1}-16:\ \text{different growth, different registers}\ [\mathrm{CT0.5}/\mathrm{CT0.6}\ \text{vs}\ \dim\mathfrak{su}(n)].\ \blacksquare
\]

\textbf{Rmk (the D3.2 "pair blocks" clause, disambiguated).} D3.2's realization test admits isomorphism classes lying "in $\iota_m(\mathfrak u(m))$ \textit{or in the pair blocks of CT0.12}". Read against the pair \textit{total} — $\mathfrak H(\mathbb C^4)\cong\mathfrak u(4)$ [CT1.9], which does contain $\mathfrak{su}(4)$ as a subalgebra — a hypothetical rung-4 slot collective could be called "already realized". That reading is foreclosed by CT4.4, which derives \textit{non-realization} of slot algebras at $n\ge4$ from channel inadmissibility, and by CT0.13, whose derived value $\mathcal A(n\ge4)=\mathcal A(3)\cup\{D^-\}$ [D1.12: derived, $\neg$stipulated] contains no $\mathfrak{su}(4)$-type aspect for a spread image to be read as: the masked force operator is block-diagonal [FT.26(2), CT2.3/CT2.4 pattern], so a spread image would have no aspect whose block carries it. The ambiguity is recorded here and resolved by citation, not by choice.

\textbf{Rmk (what survives at the locked rungs).} Nothing observable is lost \textit{by this theorem} — the observable content of a rung-$n\ge4$ vertex was already exactly the registration/sealing ledger [FT.20's criteria never invoke $\iota_n$ for EM/$G$/$H$/strong rows; FT.21(i)]. What is settled is the collective boundary algebra: beyond three slots (or three boundary contacts at a corolla) it has no home in the aspect inventory — the kernel \textit{is} the coarse-graining [FT.58(3) Rmk], now with the mechanism named: absorption without structure [CT4.4].

\textbf{Phys.} Ask what the fourth slot's algebra maps to and the corpus answers: nothing — not because the map is unknown, but because a vertex that deep in saturation has no realized slot algebra at all; its content falls into a structureless dark sink at an exponential rate while its registration signs and sealing patterns, the only reads any observer ever had, continue unchanged. Colour's three is a ceiling twice over: capacity forbids a fourth accessible slot, and realization forbids a fourth slot's algebra even where the slots exist. The tetrahedron-face floor, with six boundary edges corolled at a rung-6 vertex, is the first computed object whose entire thirty-dimensional collective boundary structure is provably dark.

\vspace{2pt}\noindent\rule{\textwidth}{0.4pt}\vspace{2pt}

\medskip
\subsection*{\texorpdfstring{B.3.3\quad FT.117 — Confinement blindness: the one-body no-go on colour-sealed singlets}{FT.117 B.3.3}}
\addcontentsline{toc}{subsection}{B.3.3 FT.117}
\phantomsection\label{tag:FT-117}
{\small \textit{Record.}\quad Source: Part 3 (Part3\_Fermion\_era.md); tier: Tier F.38 — The Colour Channel: The $\textbackslash{}varepsilon$-Selected $C\_3$, Confinement Blindness, and the Triplet-Attached Mixing Angle; canonical tag \texttt{FT.117}. Reproduced from the revision-1 canon; a cross-reference stub marks the former locus in the main corpus.}
\medskip
\textbf{Thm.} Let $H\in\mathfrak H(\mathbb C^3)$ be any Hermitian slot operator, and let $\varepsilon_{ijk}$, $\delta^i{}_j$ be the two sealing tensors [CT3.17]:

%dsp-level:0
\[
(1)\quad\textbf{baryon singlet:}\qquad\bigl(H\otimes I_3\otimes I_3+I_3\otimes H\otimes I_3+I_3\otimes I_3\otimes H\bigr)_{(ijk),(lmn)}\,\varepsilon_{lmn}=\mathfrak T(H)\,\varepsilon_{ijk};
\]

%dsp-level:0
\[
\text{equivalently}\quad\sum_l\bigl(H_{il}\varepsilon_{ljk}+H_{jl}\varepsilon_{ilk}+H_{kl}\varepsilon_{ijl}\bigr)=\mathfrak T(H)\,\varepsilon_{ijk}
\]

%dsp-level:0
\[
\qquad [\text{all}\ 27\ \text{components, exact};\ \mathrm{CT3.16}\ \text{at finite generators}];
\]

%dsp-level:0
\[
(2)\quad\textbf{covariant specialization:}\ \ H=aI_3+zP+z^{*}P^{T}\ [\mathrm{FT.116}(3)]:
\]

%dsp-level:0
\[
\qquad \mathfrak T(H)=3a\implies\text{the one-body baryon eigenvalue is}\ 3a,\ \textbf{independent of}\ (|z|,\arg z);
\]

%dsp-level:0
\[
\text{hence the}\ C_3\text{-doublet content}\ B
\]

%dsp-level:0
\[
= 2|z|,
\]

%dsp-level:0
\[
\ \theta
\]

%dsp-level:0
\[
= \arg z\ \text{— the entire family-splitting datum of FT.111(2)}\text{— is annihilated on the baryon singlet};
\]

%dsp-level:0
\[
(3)\quad\textbf{meson singlet:}\qquad\bigl(H\otimes I_3-I_3\otimes H^{T}\bigr)\,\delta=H\cdot I_3-I_3\cdot H=0\qquad\forall H\ (\textbf{including the trace part});
\]

%dsp-level:0
\[
(4)\quad\textbf{visibility locus:}\ \ [H,E_{12}]
\]

%dsp-level:0
\[
\ne 0\ \text{for}\ H
\]

%dsp-level:0
\[
= 0\cdot I_3+(1+i)P+(1-i)P^{T}:\ \text{the same covariant data are nonzero on the adjoint channel — but the adjoint is not a sealed singlet, so it carries no pole}\ [\mathrm{CT3.18};\ \mathrm{CT3.22}:\ \text{isolated colour has norm}\ 0];
\]

%dsp-level:0
\[
(5)\quad\textbf{invariant census:}
\]

%dsp-level:0
\[
\qquad \dim\bigl(\mathbb C^3\otimes\mathbb C^3\otimes\mathbb C^3\bigr)^{\mathfrak{sl}(3)}=1\ \text{spanned by}\ \varepsilon_{ijk};
\]

%dsp-level:0
\[
\qquad \dim\bigl(\mathbb C^3\otimes\bar{\mathbb C}^3\bigr)^{\mathfrak{sl}(3)}=1\ \text{spanned by}\ \delta^i{}_j;
\]

%dsp-level:0
\[
(6)\quad\textbf{Casimir spectrum:}\ \ C:
\]

%dsp-level:0
\[
= \tfrac12\bigl(\sum_{ij}\rho(E_{ij})\rho(E_{ji})-\tfrac13\rho(\textstyle\sum_iE_{ii})^2\bigr)\ \text{on}\ \mathbb C^3\otimes\mathbb C^3\otimes\mathbb C^3\ \text{satisfies}\ C(C-3)(C-6)\qquad = 0,
\]

%dsp-level:0
\[
\operatorname{spec}C=\{0,3,6\}\ \text{with multiplicities}\ (1,16,10)\ \text{from}\ (\operatorname{tr}C,\operatorname{tr}C^2)=(108,504):
\]

%dsp-level:0
\[
\qquad 27=1\oplus(8\oplus8)\oplus10;
\]

%dsp-level:0
\[
\qquad C\big|_{\mathbb C^3}=\tfrac43\ (\text{normalization pin});
\]

%dsp-level:0
\[
(7)\quad\textbf{verdict (obstruction, derived):}\ \ \text{no colour-sealed object — baryon or meson — can carry a one-body}\ C_3\text{-covariant mass splitting: the baryon reads only the singlet coefficient}\ 3a,
\]

%dsp-level:0
\[
\ \text{the meson reads nothing.}\ \text{The three-line spectrum of the channel is readable only on the triplet-attached chain itself}\ [\mathrm{D3.8}],
\]

%dsp-level:0
\[
\ \text{whose colour factor is present but unsealed.}
\]

\textbf{Req.} CT3.15, CT3.16, CT3.17, CT3.18, CT3.22 (Part 2); D3.8, D3.9 (Parts 1–2); FT.116(3) (this tier); FT.111(2) (\textbf{Tier F.37}, the $(B,\theta)$ datum); $\mathfrak{sl}(3)$ invariant theory, exact nullspace and trace arithmetic, quadratic-Casimir spectral decomposition [math; machine: \texttt{checks/\allowbreak{}verify\_\allowbreak{}H\_\allowbreak{}colour\_\allowbreak{}channel.py} §B].

\textbf{Pf.}

%dsp-level:0
\[
(1):\ \text{entrywise over}\ \mathbb Q(i)\ \text{on all}\ 27\ \text{components for}\ 15\ \text{exact covariant witnesses and}\ 8\ \text{deterministic Hermitian witnesses};
\]

%dsp-level:0
\[
\ \text{zero residual}\ [\text{machine}]
\]

%dsp-level:0
\[
\text{structural reason}:\ \varepsilon\ \text{spans}\ \wedge^3\mathbb C^3,
\]

%dsp-level:0
\[
\ \text{the determinant representation};
\]

%dsp-level:0
\[
\ \mathfrak{sl}(3)\ \text{acts on it by}\ \mathfrak T(X),
\]

%dsp-level:0
\[
\ \text{so the one-body sum acts by}\ \mathfrak T(H)\ \text{and the traceless part annihilates}\ [\mathrm{CT3.7}:\ \mathfrak T(T_a)=0;\ \mathrm{CT3.16}]
\]

%dsp-level:0
\[
(2):\ \mathfrak T(aI_3+zP+z^{*}P^{T})=3a+z\,\mathfrak T(P)+z^{*}\,\mathfrak T(P^{T})=3a\qquad[\mathfrak T(P)=\operatorname{tr}P_{(123)}=0;\ \text{machine}]
\]

%dsp-level:0
\[
(3):\ \delta=I_3\ \text{as a matrix};
\]

%dsp-level:0
\[
\qquad (H\otimes I-I\otimes H^{T})\,\mathrm{vec}(I)=\mathrm{vec}(H\cdot I-I\cdot H)=\mathrm{vec}(H-H)=0\ [\text{machine, all witnesses}]
\]

%dsp-level:0
\[
(4):\ [H,E_{12}]\ \text{computed entrywise for}\ H=(1+i)P+(1-i)P^{T}:\ \text{nonzero}\ [\text{machine}];
\]

%dsp-level:0
\[
\qquad \text{adjoint}\notin\operatorname{Im}\mathcal P_M^{(3)}\ [\mathrm{CT3.18}];\ \text{norm}\ 0\ [\mathrm{CT3.22};\ \mathrm{CT2.8}]
\]

%dsp-level:0
\[
(5):\ \rho_3(X)
\]

%dsp-level:0
\[
= X\otimes I\otimes I+I\otimes X\otimes I+I\otimes I\otimes X,
\]

%dsp-level:0
\[
\ \ \rho_M(X)
\]

%dsp-level:0
\[
= X\otimes I-I\otimes X^{T}\ \text{on the}\ 8\ \text{rational}\ \mathfrak{sl}(3)\ \text{generators}\ \{E_{ij}\}_{i\ne j}\cup\{E_{11}-E_{22},E_{22}-E_{33}\}
\]

%dsp-level:0
\[
27-\operatorname{rk}\rho_3=1,\quad9-\operatorname{rk}\rho_M=1\ [\text{exact}\ \mathbb Q\ \text{ranks};\ \text{machine}];
\]

%dsp-level:0
\[
\qquad \text{null vectors proportional to}\ \varepsilon_{ijk},\ \delta^i{}_j\ \text{componentwise}\ [\text{machine}]
\]

%dsp-level:0
\[
(6):\ C(C-3I_{27})(C-6I_{27})=0\ \text{entrywise exact}\ [\text{machine}];
\]

%dsp-level:0
\[
\qquad m_6=\tfrac{\operatorname{tr}C^2-3\operatorname{tr}C}{18}=10,\quad m_3=\tfrac{\operatorname{tr}C-6m_6}{3}=16,\quad m_0=27-m_3-m_6=1\ [\text{exact Fractions}]
\]

%dsp-level:0
\[
C\big|_{\mathbb C^3}=\tfrac12(3I_3-\tfrac13I_3)=\tfrac43I_3\ [\text{normalization pin};\ \text{machine}].
\]

%dsp-level:0
\[
\qquad (7):\ \text{conjunction of}\ (2),(3),(4)\ \text{with}\ \mathrm{CT3.18}/\mathrm{CT3.22}.\ \blacksquare
\]

\textbf{Phys.} Confinement does more than hide colour; it erases the very datum a family needs. A sealed baryon sees the sum of the three slot operators, and that sum is the trace — one number, the singlet coefficient — with the phase and the spread identically gone. A sealed meson sees the difference, which vanishes outright, trace included. The only place the cyclic data survive is the adjoint channel, and the adjoint is precisely what the mask refuses a pole. So the hadron-level branch of the question ledger is closed not by a choice of reading but by a theorem about the two invariant tensors: baryons and mesons cannot be the family that falls out of a one-body $C_3$ mechanism. What survives is the unsealed object — the chain with its triplet attachment — and that is where the arena count of FT.118 is performed. The decomposition $27=1\oplus8\oplus8\oplus10$ records the same fact from the representation side: a single singlet line in the three-triple product, a single singlet line in the triplet-antitriplet product, and everything else confined.

\vspace{2pt}\noindent\rule{\textwidth}{0.4pt}\vspace{2pt}

\medskip
\subsection*{\texorpdfstring{B.3.4\quad FT.124 — The prior obstructions do not reach the orientation space}{FT.124 B.3.4}}
\addcontentsline{toc}{subsection}{B.3.4 FT.124}
\phantomsection\label{tag:FT-124}
{\small \textit{Record.}\quad Source: Part 3 (Part3\_Fermion\_era.md); tier: Tier F.39 — Chirality: The Orientation Doublet of the Bound Chain and Its Exchange-Form Grade; canonical tag \texttt{FT.124}. Reproduced from the revision-1 canon; a cross-reference stub marks the former locus in the main corpus.}
\medskip
\textbf{Thm.} The obstructions of Tiers F.26–F.28–Tiers F.31–F.34 and Tier F.38 are theorems about their own loci; the orientation-space chirality lives on a different space and evades each one — seven loci, census exact:

%dsp-level:0
\[
(1)\quad\textbf{the pincer}\ [\mathrm{FT.93}(3)]:\ \ \text{the classification}\ \{M:[M,W_i]=0\}
\]

%dsp-level:0
\[
= \operatorname{span}_{\mathbb C}\{\Pi_{\mathcal R},\Pi_{\mathcal V}\},
\]

%dsp-level:0
\[
\ \{M:\{M,W_i\}=0\}
\]

%dsp-level:0
\[
= \mathbb C\cdot\Pi_{\mathcal V}\ \text{is a classification inside}\ \operatorname{End}(\mathbb C^4_{\mathrm{pair}})\ \text{for the pair-register weak action};
\]

%dsp-level:0
\[
\mathcal H_\omega(c)\not\subseteq\operatorname{End}(\mathbb C^4_{\mathrm{pair}});
\]

%dsp-level:0
\[
\qquad \text{the extended weak action is inequivalent}\ \bigl(\tfrac34I_4\ \text{vs}\ 2\Pi_{\mathcal R}\ [\mathrm{FT.123}(1)]\bigr);
\]

%dsp-level:0
\[
\qquad I\otimes P_{L/R}\ \text{commutes with}\ W_i\otimes I_\omega\ \text{by tensor-factor disjointness, not by membership in either pincer class};
\]

%dsp-level:0
\[
(2)\quad\textbf{the per-frame refutation}\ [\mathrm{FT.93}(4)]:\ \ \text{it concerns pair-register frame projectors}\ P^{(\mathcal F)}_\pm\ (\operatorname{rank}2\ \text{vs the}\ 1\text{-dimensional weak kernel}\ \wedge^2);
\]

%dsp-level:0
\[
\ \ \text{on the extended carrier the weak kernel is}\ \{0\}\ (\sigma_i\ \text{invertible};\ \text{machine}):
\]

%dsp-level:0
\[
P^{\mathrm{car}}_{L/R}\ \text{are not frame projectors, and their halves are weak-blind, not weak-inactive — the refutation is silent here};
\]

%dsp-level:0
\[
(3)\quad\textbf{the trace/rank obstruction}\ [\mathrm{FT.101}(3)]:\ \ \text{data re-derived:}\ \mathfrak T\Pi_{\mathcal R}
\]

%dsp-level:0
\[
= 3
\]

%dsp-level:0
\[
\ne 1
\]

%dsp-level:0
\[
= \mathfrak T\Pi_{\mathcal V};
\]

%dsp-level:0
\[
\ \text{no traceless involution in}\ \operatorname{span}\{\Pi_{\mathcal R},\Pi_{\mathcal V}\}\ (\operatorname{tr}=3a+b\in\{\pm4,\pm2\},\ \text{machine});
\]

%dsp-level:0
\[
\hat\omega\ \text{is an honest matrix involution on}\ \mathcal H_\omega\ —\ \text{not a Hilbert-space implementation of the A4 branch flip, which stays label-level; FT.101}(3)\ \text{stands untouched on its locus};
\]

%dsp-level:0
\[
(4)\quad\textbf{S-parity misalignment}\ [\mathrm{FT.93}(6)]:\ \ \text{catalog re-run (machine):}\ 30\ \text{frames},
\]

%dsp-level:0
\[
\ \gamma_{\mathcal F}
\]

%dsp-level:0
\[
= \pm\ \text{one Pauli product each, all}\ 15\ \text{sign-classes attained}\ \times2;
\]

%dsp-level:0
\[
\ \text{every class fails the commuting side for some}\ W_i\ \textbf{and}\ \text{the anticommuting side for some}\ W_i;
\]

%dsp-level:0
\[
\text{exactly}\ 3\ \text{classes are}\ S\text{-aligned}\ (\pm\sigma_i\otimes\sigma_i);
\]

%dsp-level:0
\[
\qquad [\mathbb S\otimes I_2,\ I_4\otimes\hat\chi]
\]

%dsp-level:0
\[
= 0\ (8\times8\ \text{machine}):\ \text{the chiral grade is}\ S\text{-invariant by disjointness — no misalignment can arise};
\]

%dsp-level:0
\[
(5)\quad\textbf{the zero-expectation theorem}\ [\mathrm{FT.95}(3)]:\ \ \langle\psi|\hat\chi|\psi\rangle
\]

%dsp-level:0
\[
= 0\ \text{for every}\ \hat\omega\text{-symmetric}\ \psi\ (\text{machine})\ —\ \text{recorded, not an obstruction:}\ \hat\chi\ \text{is a per-chain eigenvalue datum},
\]

%dsp-level:0
\[
\ \text{and the corpus's own precedent is FT.87}(8)\text{'s}\ \operatorname{sign}(k),
\]

%dsp-level:0
\[
\text{a legitimate}\ C\text{-odd coordinate under the same vanishing theorem; a chiral read requires an oriented chain (boundary datum), not an ensemble average}\ [\text{Def-H.33 closing Rmk}];
\]

%dsp-level:0
\[
(6)\quad\textbf{confinement blindness}\ [\mathrm{FT.117};\ \mathrm{CT3.22}]:\ \ \hat\chi\ \text{has no colour-slot argument:}\ [H_{\mathrm{slot}}\otimes I_2,\ I_3\otimes\hat\chi]
\]

%dsp-level:0
\[
= 0\ \text{for every Hermitian slot operator}\ H\ (6\times6\ \text{machine witness});
\]

%dsp-level:0
\[
\ \text{sealed or not, the orientation factor persists};
\]

%dsp-level:0
\[
(7)\quad\textbf{the}\ Q\textbf{-phase-free bound}\ [\mathrm{FT.119}(4)]:\ \ Q
\]

%dsp-level:0
\[
= \tfrac23\ \text{is the seal-side normalization content}\ (A=1,B=\sqrt2);
\]

%dsp-level:0
\[
\ \hat\chi\ \text{is a chain-level grade: the}\ C_3\ \text{structures of Tier F.37 live on the seal triplet and are untouched}.
\]

\textbf{Req.} FT.93, FT.95, FT.101 (Tiers F.31–F.34); CT0.7, CT1.2, Cor CT0.7a (Part 1); FT.117, FT.119 (Tier F.38); FT.87(8) (Tiers F.29–F.30); FT.123 (this tier); CT3.22 (Part 2); D0.10, D1.4 (Parts 1–2); exact nullspace solves over $\mathbb Q(i)$, exhaustive frame enumeration, Kronecker disjointness identities [math; machine: \texttt{checks/\allowbreak{}verify\_\allowbreak{}H\_\allowbreak{}chirality.py} §C].

\textbf{Pf.}

%dsp-level:0
\[
(1):\ \text{pincer classes re-solved exactly: commuting-side nullspace}\ \dim_{\mathbb Q}4\ (=2\ \text{over}\ \mathbb C)
\]

%dsp-level:0
\[
= \operatorname{span}\{\Pi_{\mathcal R},\Pi_{\mathcal V}\};
\]

%dsp-level:0
\[
\ \text{anticommuting side}\ \dim_{\mathbb Q}2
\]

%dsp-level:0
\[
= \mathbb C\Pi_{\mathcal V};
\]

%dsp-level:0
\[
\ \operatorname{tr}(3a+b)
\]

%dsp-level:0
\[
\ne 0\ \text{on all four sign pairs}\ [\text{machine}]
\]

%dsp-level:0
\[
\text{inequivalence:}\ \textstyle\sum_i(\tfrac{\sigma_i}2\otimes I)^2=\tfrac34I_4\ne2\Pi_{\mathcal R}=\sum_iW_i^2\ [\text{machine}];
\]

%dsp-level:0
\[
\qquad \text{disjointness: FT.123}(1)\ \text{clause verbatim}
\]

%dsp-level:0
\[
(2):\ \bigcap_i\ker(\sigma_i\otimes I_2)=\{0\}:\ \text{exact}\ 24\times8\ \text{real system, rank}\ 8\ [\text{machine}];
\]

%dsp-level:0
\[
\qquad \operatorname{rank}P^{\mathrm{car}}_{L/R}=2\ [\mathrm{FT.123}\ \text{Pf}]
\]

%dsp-level:0
\[
(3):\ \Pi_{\mathcal R}
\]

%dsp-level:0
\[
= \tfrac12(I_4+\mathbb S),
\]

%dsp-level:0
\[
\ \Pi_{\mathcal V}
\]

%dsp-level:0
\[
= \tfrac12(I_4-\mathbb S),
\]

%dsp-level:0
\[
\ \mathbb S
\]

%dsp-level:0
\[
= \tfrac12(E_0+\Sigma):\ \operatorname{tr}
\]

%dsp-level:0
\[
= (3,1)\ [\mathrm{CT0.7}(1)\ \text{re-derived};\ \text{machine}];
\]

%dsp-level:0
\[
\qquad M=a\Pi_{\mathcal R}+b\Pi_{\mathcal V}:\ M^2=\mathrm{id}\implies a,b\in\{\pm1\}\implies\operatorname{tr}M\in\{\pm4,\pm2\}\ne0
\]

%dsp-level:0
\[
(4):\ \binom{15}4=1365\ \text{candidate}\ 4\text{-sets}\to30\ \text{pairwise-anticommuting frames}\ [\text{machine}];
\]

%dsp-level:0
\[
\qquad \gamma_{\mathcal F}^2=I_4,\ \operatorname{tr}\gamma_{\mathcal F}=0,\ \gamma_{\mathcal F}^\dagger=\gamma_{\mathcal F}\ \forall\ \text{frames};
\]

%dsp-level:0
\[
\qquad \gamma_{\mathcal F}=\pm\ \text{a single Pauli product (exact match census)}
\]

%dsp-level:0
\[
15\ \text{sign-classes},\ \text{each from exactly}\ 2\ \text{frames};
\]

%dsp-level:0
\[
\qquad \forall\ \text{class}:\ \exists i:[\gamma,W_i]\ne0\ \wedge\ \exists j:\{\gamma,W_j\}\ne0\ [\text{machine, entrywise}];
\]

%dsp-level:0
\[
\qquad S\text{-aligned classes}:\ [\gamma,\mathbb S]=0\iff\gamma=\pm\sigma_i\otimes\sigma_i:\ \text{exactly}\ 3\ [\text{machine}]
\]

%dsp-level:0
\[
[\mathbb S\otimes I_2,\ I_4\otimes\hat\chi]=(\mathbb S I_4)\otimes(I_2\hat\chi)-(I_4\mathbb S)\otimes(\hat\chi I_2)=0\ [8\times8\ \text{entrywise};\ \text{machine}]
\]

%dsp-level:0
\[
(5):\ \psi_\pm=(|c\rangle\pm|\omega(c)\rangle)/\sqrt2:\ \langle\psi_\pm|\hat\chi|\psi_\pm\rangle=0\ [\text{exact};\ \text{machine}];
\]

%dsp-level:0
\[
\qquad \hat\chi|L\rangle=+|L\rangle\ \text{remains sharp}\ [\mathrm{FT.122}(3)]
\]

%dsp-level:0
\[
(6):\ (H\otimes I_2)(I_3\otimes\hat\chi)
\]

%dsp-level:0
\[
= H\otimes\hat\chi
\]

%dsp-level:0
\[
= (I_3\otimes\hat\chi)(H\otimes I_2)\ [\text{machine on a generic Hermitian witness},\ 6\times6]
\]

%dsp-level:0
\[
(7):\ \text{FT.119}(4)\ \text{is an identity on the seal-triplet spectrum}\ [\mathrm{FT.111}(3)\ \text{normalization}];
\]

%dsp-level:0
\[
\ \hat\chi\ \text{acts on}\ \mathcal H_\omega:\ \text{disjoint bookkeeping}.\ \blacksquare
\]

\textbf{Phys.} Each prior obstruction was a correct theorem about a different room. The pincer and the frame refutation are theorems about the pair register under \textit{its} weak action; the trace/rank obstruction is about conjugation on that same register; the $S$-parity census is about pair-register pseudoscalars; the zero-expectation theorem is about ensemble averages of $C$-odd reads; confinement blindness is about one-body colour-slot operators; the $Q$ bound is about the seal-side $C_3$ spectrum. The chirality operator acts on the chain's own orientation doublet, indexed by the chain, tensor-disjoint from all of them — and the census above re-proved every obstruction at its own locus before showing the new grade walks past. Nothing was overturned; the locus moved.

\vspace{2pt}\noindent\rule{\textwidth}{0.4pt}\vspace{2pt}

\medskip
\subsection*{\texorpdfstring{B.3.5\quad BM.10 — The canonical pair: the finite-dimension obstruction, the discrete Weyl pre-limit, and the three exact bounds}{BM.10 B.3.5}}
\addcontentsline{toc}{subsection}{B.3.5 BM.10}
\phantomsection\label{tag:BM-10}
{\small \textit{Record.}\quad Source: Part 4 (Part4\_Baryomesonic\_era.md); tier: Tier BM.4 — Noether's Theorem from the Janus Seal Topology: the Discrete Conservation Law, the Continuum Promotion, the Stress-Energy of a Sealed Composite, and the Canonical-Pair Obstructions; canonical tag \texttt{BM.10}. Reproduced from the revision-1 canon; a cross-reference stub marks the former locus in the main corpus.}
\medskip
\textbf{Thm.}

%dsp-level:0
\[
(1)\quad\textbf{the 16-matrix cannot carry the pair (obstruction):}
\]

%dsp-level:0
\[
\qquad \mathfrak T[A,B]=\mathfrak T(AB)-\mathfrak T(BA)=0\ \ \forall A,B\in\operatorname{End}(\mathbb C^4)
\]

%dsp-level:0
\[
\qquad [\text{machine: all}\ 256\ \text{pairs of the}\ 16\text{-element CT1.2 basis}];
\]

%dsp-level:0
\[
\mathfrak T(i\,\mathbb1_4)=4i\ne0:
\]

%dsp-level:0
\[
\qquad [\hat x,\hat p]=i\,\mathbb1\ \text{has no solution in}\ \operatorname{End}(\mathbb C^4)
\]

%dsp-level:0
\[
\qquad [\text{Wielandt–Wintner};\ \text{machine}];
\]

%dsp-level:0
\[
\mathfrak T(i\,\mathbb1_d)=d\,i\ne0\ \ \forall d\ge1\ [\text{machine},\ d\le8]:
\]

%dsp-level:0
\[
\qquad \neg\exists\ \text{a finite-dimensional Heisenberg pair at any dimension};
\]

%dsp-level:0
\[
\text{the same obstruction holds in every finite spin representation}\ (\dim=2,3,4,5)\ [\text{machine: every commutator traceless,}\ \mathfrak T(i\mathbb1)=\dim\cdot i\ne0];
\]

%dsp-level:0
\[
\text{the arena dimension is derived,}\ \neg\text{chosen:}\ \dim\mathfrak H(\mathbb C^4)
\]

%dsp-level:0
\[
= 16
\]

%dsp-level:0
\[
= (7,3,3,3)
\]

%dsp-level:0
\[
= C_{\mathrm{cap}}
\]

%dsp-level:0
\[
= \dim\mathcal K_1\ [\mathrm{CT1.1},\ \mathrm{CT1.2},\ \mathrm{FT.142}(2);\ \text{machine: exact rank}\ 16];
\]

%dsp-level:0
\[
(2)\quad\textbf{what the arena does carry — the discrete Weyl pair (the pre-limit):}
\]

%dsp-level:0
\[
\qquad Z|k\rangle=\omega^k|k\rangle,\quad X|k\rangle=|k{+}1\!\!\!\pmod N\rangle:
\]

%dsp-level:0
\[
\qquad ZX=\omega XZ,\quad Z^N=X^N=\mathbb1;
\]

%dsp-level:0
\[
N=4:\ \omega=i\ \text{over}\ \mathbb Q(i)\ [\text{machine}];
\]

%dsp-level:0
\[
\qquad N
\]

%dsp-level:0
\[
= 16:\ \omega
\]

%dsp-level:0
\[
= \zeta_{16}\ \text{exact in the cyclotomic ring}\ \mathbb Q[z]/(z^8{+}1)
\]

%dsp-level:0
\[
= \mathbb Q[z]/(\Phi_{16}):\ \zeta_{16}^8
\]

%dsp-level:0
\[
= -1,
\]

%dsp-level:0
\[
\ \zeta_{16}^4\ \text{of order}\ 4,
\]

%dsp-level:0
\[
\ \operatorname{ord}(\zeta_{16})
\]

%dsp-level:0
\[
= 16\ [\text{machine: all}\ 256\ \text{entries}];
\]

%dsp-level:0
\[
[Z,X]=(\omega-1)XZ:
\]

%dsp-level:0
\[
\qquad \text{the finite commutator is a root-of-unity multiple of a product}\ —\ \text{torsion},\ \neg\ \text{the scalar}\ i\,\mathbb1;
\]

%dsp-level:0
\[
\textbf{the abelian counterpart is exact at every step size:}
\]

%dsp-level:0
\[
\qquad [\hat x,D_\ell]=-T_\ell\ [\mathrm{BM.8}(1);\ \text{machine}]:
\]

%dsp-level:0
\[
\qquad \text{the Heisenberg relation is the}\ \ell\to0\ \text{read of the hop algebra on the chart},
\]

%dsp-level:0
\[
\ \neg\text{a relation inside the finite register};
\]

%dsp-level:0
\[
(3)\quad\textbf{the pair on the chart and the Robertson bound:}
\]

%dsp-level:0
\[
\qquad \hat x^i:=\text{multiplication by}\ \chi\text{'s coordinate}\ [\mathrm{CT2.6}],
\]

%dsp-level:0
\[
\qquad \hat p_i:=-i\nabla_i\ \ (\text{self-adjoint w.r.t.}\ \sqrt{-g}\,d^4x)\ [\text{math}];
\]

%dsp-level:0
\[
[\hat x^i,\hat p_j]=i\,\delta^i_j\ \ (\hbar=1\ [\mathrm{D2.9}])
\]

%dsp-level:0
\[
\qquad [\text{Stone–von Neumann};\ \text{machine in the quotient}\ \mathbb Q(i)\langle x,p\rangle/([x,p]-i)];
\]

%dsp-level:0
\[
\Delta A\,\Delta B\ge\tfrac12\bigl|\langle[A,B]\rangle\bigr|
\]

%dsp-level:0
\[
\qquad [\text{Robertson; Cauchy–Schwarz};\ \text{math}]\implies\boxed{\ \Delta x\,\Delta p\ge\tfrac12\ };
\]

%dsp-level:0
\[
\qquad \Delta t\,\Delta E\ge\tfrac12\ \text{in the form of}\ (4);
\]

%dsp-level:0
\[
\qquad \Delta L_z\,\Delta(\sin\theta)\ge\tfrac12|\langle\cos\theta\rangle|\ \text{in the form of}\ (5);
\]

%dsp-level:0
\[
\textbf{saturation, exact:}\quad A:=\hat x+i\hat p,\quad A^\dagger=\hat x-i\hat p:
\]

%dsp-level:0
\[
\qquad [A,A^\dagger]=2,
\]

%dsp-level:0
\[
\qquad 4\hat x^2=(A+A^\dagger)^2,
\]

%dsp-level:0
\[
\qquad 4\hat p^2=-(A-A^\dagger)^2,
\]

%dsp-level:0
\[
\qquad [A^\dagger A,A^\dagger]=2A^\dagger\ [\text{machine}];
\]

%dsp-level:0
\[
\langle\hat x^2\rangle_n=\langle\hat p^2\rangle_n=n+\tfrac12
\]

%dsp-level:0
\[
\qquad [\text{machine: exact rationals}]\implies\Delta x\,\Delta p=n+\tfrac12\ge\tfrac12,
\]

%dsp-level:0
\[
\qquad \text{equality}\iff n=0;
\]

%dsp-level:0
\[
\qquad \text{the width grows strictly with}\ n;
\]

%dsp-level:0
\[
(4)\quad\textbf{the}\ (t,E)\ \textbf{pair — no time operator, the Mandelstam–Tamm form:}
\]

%dsp-level:0
\[
\qquad \neg\exists\,
\]

%dsp-level:0
\[
\hat T\ \text{self-adjoint with}\ [\hat T,\hat H]
\]

%dsp-level:0
\[
= -i\ \text{at finite dimension}\ [\text{machine:}\ \mathfrak T(-i\mathbb1_2)=-2i\ne0];
\]

%dsp-level:0
\[
\text{a self-adjoint}\ \hat T\ \text{conjugate to a semibounded}\ \hat H\ \text{would make}\ \operatorname{spec}(\hat H)=\mathbb R\ [\text{Pauli};\ \text{math}];
\]

%dsp-level:0
\[
\qquad \text{and the corpus's own arrow gives}\ \neg\exists\,\tau^{-1}\ [\mathrm{A6};\ \mathrm{FT.156}(5);\ \mathrm{FT.144}(4)]:
\]

%dsp-level:0
\[
\textbf{the admissible read is Mandelstam–Tamm:}\qquad\Delta t:=\frac{\Delta A}{\bigl|\tfrac{d}{dt}\langle A\rangle\bigr|}\implies\Delta t\,\Delta E\ge\tfrac12\qquad[\text{math}];
\]

%dsp-level:0
\[
\textbf{exact witness:}\quad A=\tfrac12\sigma_x,\ H=\tfrac12\sigma_z,\ \psi\propto(1,i):
\]

%dsp-level:0
\[
\qquad [\sigma_z,\sigma_x]=2i\sigma_y,\quad\langle A\rangle=\langle H\rangle=0,\quad\langle\sigma_y\rangle=1\ [\text{machine}];
\]

%dsp-level:0
\[
(\Delta A)^2=(\Delta E)^2=\tfrac14,
\]

%dsp-level:0
\[
\qquad [H,A]=\tfrac i2\sigma_y\implies\tfrac{d}{dt}\langle A\rangle=-\tfrac12,
\]

%dsp-level:0
\[
\qquad (\Delta t)^2(\Delta E)^2=\tfrac14\implies\Delta t\,\Delta E=\tfrac12
\]

%dsp-level:0
\[
\qquad [\text{machine: saturation}];
\]

%dsp-level:0
\[
(5)\quad\textbf{the}\ (\theta,L)\ \textbf{pair — no angle operator, the Carruthers–Nieto form:}
\]

%dsp-level:0
\[
\qquad \neg\exists\,\hat\theta\ \text{with}\ [\hat\theta,\hat L_z]=i\ \text{in any finite spin representation}\ [(1)\ \text{re-applied};\ \text{machine}];
\]

%dsp-level:0
\[
\text{on the circle}\ \hat L_z=-i\partial_\theta,\ \ \hat L_z e_m=m\,e_m:
\]

%dsp-level:0
\[
\qquad [\hat L_z,\sin\theta]=-i\cos\theta,
\]

%dsp-level:0
\[
\qquad [\hat L_z,\cos\theta]=i\sin\theta
\]

%dsp-level:0
\[
\qquad [\text{machine: exact on the interior of the truncated Fourier basis},\ |m|\le4;\ \text{boundary}=\text{truncation scope}];
\]

%dsp-level:0
\[
\textbf{Robertson then gives the exact pair:}
\]

%dsp-level:0
\[
\qquad \Delta L_z\,\Delta(\sin\theta)\ge\tfrac12|\langle\cos\theta\rangle|,
\]

%dsp-level:0
\[
\qquad \Delta L_z\,\Delta(\cos\theta)\ge\tfrac12|\langle\sin\theta\rangle|
\]

%dsp-level:0
\[
\qquad [\text{math}];
\]

%dsp-level:0
\[
\textbf{the naive bound is refuted exactly:}\quad\psi=e_0\ (\text{an}\ \hat L_z\ \text{eigenstate}):
\]

%dsp-level:0
\[
\qquad \langle\hat L_z\rangle=(\Delta L_z)^2=0,\quad\langle\cos\theta\rangle=\langle\sin\theta\rangle=0,\quad(\Delta\cos\theta)^2=(\Delta\sin\theta)^2=\tfrac12\ [\text{machine}];
\]

%dsp-level:0
\[
(\Delta\theta)(\Delta L_z)\ge\tfrac12\ \text{would require}\ 0\ge\tfrac14\ \text{in squared form}:\ \textbf{false};
\]

%dsp-level:0
\[
\qquad \text{the CN bound reads}\ 0\ge0:\ \text{saturated}\ [\text{machine}];
\]

%dsp-level:0
\[
\qquad \text{the CN inequality holds on exact rational witnesses}\ [\text{machine}];
\]

%dsp-level:0
\[
\textbf{at the seal there is no angle at all:}\quad\text{the rotation group of the}\ \varepsilon\text{-seal is the finite}\ C_3
\]

%dsp-level:0
\[
= A_3\ [\mathrm{BM.3}(1);\ \mathrm{FT.116}]:
\]

%dsp-level:0
\[
\qquad |\widehat{C_3}|=3=|C_3|\ [\text{machine}]\ —\ \text{three one-dimensional characters, no continuous parameter};
\]

%dsp-level:0
\[
\text{the angle}\ \theta\ \text{exists only after}\ \mathrm{BM.8}(3)\text{'s promotion, where}\ u(\theta)\ \text{is defined for every}\ \theta\ \text{and the CN form applies};
\]

%dsp-level:0
\[
(6)\quad\textbf{the width law:}
\]

%dsp-level:0
\[
\qquad \Delta x\,\Delta p\ge\tfrac12\ \ (\hbar=1)\implies\Delta x\ge\frac1{2\,\Delta p},
\]

%dsp-level:0
\[
\qquad V=(\Delta x)^3\ge\frac1{8(\Delta p)^3}
\]

%dsp-level:0
\[
\qquad [\text{machine:}\ \Delta x\,\Delta p=\tfrac12\ \text{on rational witnesses};\ (1/(2\Delta p))^3=1/(8(\Delta p)^3)];
\]

%dsp-level:0
\[
\rho=\frac{m_{\mathrm{hid}}}{V}\le8\,m_{\mathrm{hid}}(\Delta p)^3<\infty\ \ \forall\Delta p>0\ [\mathrm{BM.9}(5)\text{–}(6)]:
\]

%dsp-level:0
\[
\qquad \text{the stress-energy of a sealed composite is a finite-width distribution};
\]

%dsp-level:0
\[
\text{the width is the scale at which the coefficient fields of}\ \mathrm{D2.1}\ \text{become}\ \mathrm{CT2.6}(\mathrm{iii})\text{-regular of bounded variation}\ —\ \text{hypothesis content at the point of use,}\ \neg\text{a derived smoothness}.
\]

\textbf{Rmk (obstruction discipline).} Clauses (1), (4) and (5) are negative theorems of the corpus's own kind [FT.53, FT.59, FT.93(3) pattern]: they locate the domain in which a structure can exist rather than asserting the structure. The canonical pair, the time operator and the angle operator are all chart-side objects [CT2.6]; none of them is a $16$-matrix object, and no clause of this tier imports one into the register.

\textbf{Req.} A3, A6 (Part 1); D2.9 (Parts 1–2); CT0.9, CT1.1, CT1.2 (Part 1); CT2.6 (Part 2; conditions (ii)–(iii) at the point of use — A §0.2a); Def-BM.3, BM.3(1), BM.7(4), BM.8(1), BM.8(3), BM.9(5), BM.9(6) (this part); FT.116, FT.142(2), FT.144(4), FT.156(5) (Part 3); the Wielandt–Wintner trace obstruction, Stone–von Neumann, Robertson's inequality, the Pauli semiboundedness argument, the Mandelstam–Tamm bound, exact cyclotomic arithmetic in $\mathbb Q[z]/(z^8{+}1)$, exact noncommutative quotient arithmetic in $\mathbb Q(i)\langle x,p\rangle/([x,p]-i)$ [math; machine: \texttt{checks/\allowbreak{}verify\_\allowbreak{}H\_\allowbreak{}noether.py} §D].

\textbf{Pf.}

%dsp-level:0
\[
(1):\ \mathfrak T(AB)=\mathfrak T(BA)\ \text{in every finite dimension}\ [\text{math}]\implies\mathfrak T[A,B]=0;
\]

%dsp-level:0
\[
\qquad \mathfrak T(i\mathbb1_d)=di\ne0\ [\text{machine}];
\]

%dsp-level:0
\[
\qquad \text{the}\ 16\text{-basis pairs: all}\ 256\ \text{traces zero}\ [\text{machine over}\ \mathbb Q(i)];
\]

%dsp-level:0
\[
\qquad \text{spin representations}\ \dim=2,3,4,5:\ \text{the full matrix-unit basis}\ [\text{machine}]
\]

%dsp-level:0
\[
(2):\ Z,
\]

%dsp-level:0
\[
X\ \text{entrywise}: (ZX)_{ij}
\]

%dsp-level:0
\[
= \omega^{i}\delta_{j,i+1},
\]

%dsp-level:0
\[
\ (XZ)_{ij}
\]

%dsp-level:0
\[
= \omega^{j}\delta_{j,i+1}
\]

%dsp-level:0
\[
= \omega^{i-1}\delta_{j,i+1}
\]

%dsp-level:0
\[
\implies ZX
\]

%dsp-level:0
\[
= \omega XZ\ [\text{machine at}\ N=4\ \text{over}\ \mathbb Q(i);\ N=16\ \text{over}\ \mathbb Q[z]/(z^8{+}1)];
\]

%dsp-level:0
\[
\qquad \Phi_{16}=z^8+1\ [\text{math}]\implies z^8=-1,\ z^{16}=1,\ z^k\ne1\ (1\le k<16)\ [\text{machine}]
\]

%dsp-level:0
\[
(3):\ [x,p]=i\ \text{by construction of the quotient}\ [\text{machine}];
\]

%dsp-level:0
\[
\qquad A=\hat x+i\hat p\implies[A,A^\dagger]=[x,p]\cdot(-i)+[p,x]\cdot i=2\ [\text{machine}];
\]

%dsp-level:0
\[
\qquad 4x^2-(A+A^\dagger)^2=0,\ 4p^2+(A-A^\dagger)^2=0\ [\text{machine: normal-ordered reduction}\ p\,x=x\,p-i]
\]

%dsp-level:0
\[
\text{number states}: \langle A^\dagger A\rangle
\]

%dsp-level:0
\[
= 2n,
\]

%dsp-level:0
\[
\ \langle AA^\dagger\rangle
\]

%dsp-level:0
\[
= 2n+2,
\]

%dsp-level:0
\[
\ \langle A^2\rangle
\]

%dsp-level:0
\[
= \langle A^{\dagger2}\rangle
\]

%dsp-level:0
\[
= 0\implies\langle x^2\rangle
\]

%dsp-level:0
\[
= \langle p^2\rangle
\]

%dsp-level:0
\[
= \tfrac{4n+2}4
\]

%dsp-level:0
\[
= n+\tfrac12\ [\text{machine: exact Fractions}];
\]

%dsp-level:0
\[
\qquad \text{Robertson's constant}\ \tfrac12\ [\text{math}]
\]

%dsp-level:0
\[
(4):\ \text{trace obstruction re-applied at}\ d=2\ [\text{machine}];
\]

%dsp-level:0
\[
\qquad \text{Pauli:}\ e^{i\hat T s}\hat He^{-i\hat T s}=\hat H+s\implies\operatorname{spec}\hat H\ \text{is}\ \mathbb R\text{-translated}\ [\text{math}];
\]

%dsp-level:0
\[
\qquad \text{witness}: \psi
\]

%dsp-level:0
\[
= (1,i),
\]

%dsp-level:0
\[
\ \langle\psi|\psi\rangle
\]

%dsp-level:0
\[
= 2\ [\text{exact, no radicals}];
\]

%dsp-level:0
\[
\ \text{the four expectation values and the commutator as displayed}\ [\text{machine over}\ \mathbb Q(i)]
\]

%dsp-level:0
\[
(5):\ [L_z,\cos\theta]e_m
\]

%dsp-level:0
\[
= \tfrac{(m{+}1)e_{m+1}+(m{-}1)e_{m-1}}2-m\tfrac{e_{m+1}+e_{m-1}}2
\]

%dsp-level:0
\[
= \tfrac{e_{m+1}-e_{m-1}}2
\]

%dsp-level:0
\[
= i\sin\theta\,
\]

%dsp-level:0
\[
e_m\ [\text{exact}];
\]

%dsp-level:0
\[
\qquad [L_z,\sin\theta]e_m
\]

%dsp-level:0
\[
= \tfrac{(m{+}1)e_{m+1}-(m{-}1)e_{m-1}}{2i}-m\tfrac{e_{m+1}-e_{m-1}}{2i}
\]

%dsp-level:0
\[
= \tfrac{e_{m+1}+e_{m-1}}{2i}
\]

%dsp-level:0
\[
= -i\cos\theta\,
\]

%dsp-level:0
\[
e_m\ [\text{machine}]
\]

%dsp-level:0
\[
\text{eigenstate witness}: \langle\cos^2\theta\rangle=\langle e_0|\tfrac{e_2+2e_0+e_{-2}}4\rangle=\tfrac12\ [\text{machine}];
\]

%dsp-level:0
\[
\qquad (\Delta L_z)^2(\Delta\sin\theta)^2-\tfrac14\langle\cos\theta\rangle^2=0\ge0;
\]

%dsp-level:0
\[
\qquad \text{the naive form in squared reads}\ 0\ge\tfrac14:\ \text{refuted}\ [\text{exact}]
\]

%dsp-level:0
\[
|\widehat{C_3}|
\]

%dsp-level:0
\[
= 3\ [\text{machine}:\ \text{the three characters}\ \chi_k(g)=\omega^k,\ k\in\{0,1,2\},\ \text{each of exact order dividing}\ 3]
\]

%dsp-level:0
\[
(6):\ \Delta x\,
\]

%dsp-level:0
\[
\Delta p
\]

%dsp-level:0
\[
= \tfrac12\ \text{on the rational witnesses}\ (3/5,5/6),
\]

%dsp-level:0
\[
\ (1,1/2),
\]

%dsp-level:0
\[
\ (7/2,1/7),
\]

%dsp-level:0
\[
\ (11/13,13/22)\ [\text{machine}];
\]

%dsp-level:0
\[
\qquad \text{cubing}: (1/(2\Delta p))^3=1/(8(\Delta p)^3)\ [\text{machine}];
\]

%dsp-level:0
\[
\qquad \mathrm{BM.9}(5)\ \text{re-cited for}\ \rho V=m_{\mathrm{hid}}.\ \blacksquare
\]

\textbf{Phys.} The uncertainty principle is not in the sixteen matrices, and the corpus proves it cannot be: every commutator in a finite dimension is traceless while the target $i\mathbb1$ is not, so no position–momentum pair lives in the register, in any spin representation, or anywhere else finite. What the register does carry is the discrete predecessor — a clock and a shift whose commutation phase is a root of unity, order four on the pair channel and order sixteen on the arena, torsion at every finite size. The genuine pair appears one level up, on the continuum chart, where multiplication and covariant differentiation are unbounded and Robertson's inequality delivers $\Delta x\,\Delta p\ge\tfrac12$ with the constant attained exactly by the ground state and the width growing as $n+\tfrac12$. The other two conjugate pairs come with their own obstructions and their own correct forms: there is no self-adjoint time operator — the tick has no inverse and a semibounded generator forbids one — so the time–energy bound holds in Mandelstam–Tamm form, saturated by an exact two-level witness; and there is no angle operator, so the angular bound holds in the Carruthers–Nieto form, while the naive $\Delta\theta\,\Delta L\ge\tfrac12$ is refuted outright by any angular-momentum eigenstate. At the seal itself the rotation group is the cyclic three, which has three characters and no continuous angle at all: the angle is born in the promotion. What the width then buys is the composite's extent — a strictly positive volume, a finite density, and a stress-energy tensor that is a distribution rather than a spike.

\medskip
\subsection*{\texorpdfstring{B.3.6\quad FT.55 — The no-go and the surviving spin realization (catalog entry)}{FT.55 B.3.6}}
\addcontentsline{toc}{subsection}{B.3.6 FT.55}
\phantomsection\label{tag:FT-55-cat}
{\small \textit\{Catalog.\}\quad The full theorem remains in the main corpus as Theorem 3.47 (Part 3); its derived clauses are consumed by later derivations. The hypothesis-level ruling is recorded here verbatim:}
\begin{quote}\small \end{quote}
\medskip
\subsection*{\texorpdfstring{B.3.7\quad FT.66 — The Bott correction: witness re-evaluation, reconstructability, and the inadmissibility of absolute selection (catalog entry)}{FT.66 B.3.7}}
\addcontentsline{toc}{subsection}{B.3.7 FT.66}
\phantomsection\label{tag:FT-66-cat}
{\small \textit\{Catalog.\}\quad The full theorem remains in the main corpus as Theorem 3.54 (Part 3); its derived clauses are consumed by later derivations. The hypothesis-level ruling is recorded here verbatim:}
\begin{quote}\small \end{quote}
\medskip
\subsection*{\texorpdfstring{B.3.8\quad FT.82 — Gauge-block reconstruction (audit pass B): the algebra closes, the coupling structure does not (catalog entry)}{FT.82 B.3.8}}
\addcontentsline{toc}{subsection}{B.3.8 FT.82}
\phantomsection\label{tag:FT-82-cat}
{\small \textit\{Catalog.\}\quad The full theorem remains in the main corpus as Theorem 3.65 (Part 3); its derived clauses are consumed by later derivations. The hypothesis-level ruling is recorded here verbatim:}
\begin{quote}\small \end{quote}
\medskip
\subsection*{\texorpdfstring{B.3.9\quad BM.61 — $\alpha_B$: the inventory, the coupling-blindness theorem, and the obstruction (catalog entry)}{BM.61 B.3.9}}
\addcontentsline{toc}{subsection}{B.3.9 BM.61}
\phantomsection\label{tag:BM-61-cat}
{\small \textit\{Catalog.\}\quad The full theorem remains in the main corpus as Theorem 4.52 (Part 4); its derived clauses are consumed by later derivations. The hypothesis-level ruling is recorded here verbatim:}
\begin{quote}\small \end{quote}
\medskip
\section*{\texorpdfstring{B.4 Traps and Phantom Gaps}{B.4 Traps and Phantom Gaps}}
\addcontentsline{toc}{section}{B.4 Traps and Phantom Gaps}
\markboth{B.4 Traps and Phantom Gaps}{B.4 Traps and Phantom Gaps}
\subsection*{\texorpdfstring{B.4.1\quad FT.125 — The deficit ledger: deficit 1 closed at the projector level; deficit 1$'$ opened}{FT.125 B.4.1}}
\addcontentsline{toc}{subsection}{B.4.1 FT.125}
\phantomsection\label{tag:FT-125}
{\small \textit{Record.}\quad Source: Part 3 (Part3\_Fermion\_era.md); tier: Tier F.39 — Chirality: The Orientation Doublet of the Bound Chain and Its Exchange-Form Grade; canonical tag \texttt{FT.125}. Reproduced from the revision-1 canon; a cross-reference stub marks the former locus in the main corpus.}
\medskip
\textbf{Thm.} Under Def-H.33 and FT.122–FT.124 (and the chain-carrier lift of FT.126):

%dsp-level:0
\[
(1)\quad\textbf{Item 20 deficit 1 (chirality projector) — CLOSED at the projector level}\ (\text{under Def-H.34}).\ \ P_L,
\]

%dsp-level:0
\[
P_R\ \text{are corpus-constructible on}\ \mathcal H_\omega(c);
\]

%dsp-level:0
\[
\ \text{they satisfy every defining property}\ —\ \text{idempotent, orthogonal, complete, rank}\ 1\ [\mathrm{FT.122}(2)],
\]

%dsp-level:0
\[
\ \text{weak-commuting on the extended carrier}\ [\mathrm{FT.123}(1)],
\]

%dsp-level:0
\[
\ \hat\omega\text{-exchanged}\ [\mathrm{FT.123}(2)]\ —\ \text{and D2.6's vocabulary clause is an operator identity}\ [\mathrm{FT.122}(4)].
\]

%dsp-level:0
\[
\textbf{Residue:}\ \text{item 21(i)'s frame section}\ e\mapsto\mathcal F_e,
\]

%dsp-level:0
\[
\ \text{its coupling to the EW.32 transport, and a global Dirac-type operator on the emergent geometry remain unconstructed — Tier F.39 supplies the extended carrier and its grade, not the section};
\]

%dsp-level:0
\[
(2)\quad\textbf{Item 20 deficit 5 (gauge anomalies) — blocker RELOCATED, not removed.}\ \ \mathrm{FT.81}(4)\ \text{names three prerequisites: a chiral projector, a Weyl symbol, a hypercharge-assignment read. Tier F.39 discharges only the first. The anomaly sums}\ \bigl(\sum Y,\ \sum Y^3,\ [\mathrm{SU}(3)]^2\mathrm U(1),\ [\mathrm{SU}(2)]^2\mathrm U(1)\bigr)\ \text{are formable for any}\ \textbf{declared}\ \text{chiral assignment; no corpus-internal statement is possible until deficits}\ 1'\ (+2/3)\ \text{close};
\]

%dsp-level:0
\[
(3)\quad\textbf{Item 20 deficit 6 (the neutrino row) — PARTIALLY UNBLOCKED, precisely:}\ \ \operatorname{im}P_R\ \text{is a constructible right-handed Weyl-type carrier}\ [\mathrm{Def\text{-}H.33}];
\]

%dsp-level:0
\[
\ \text{species content remains absent (deficits 2/3 stand) — whether the corpus spectrum contains such a state is a separate, still-unposed question};
\]

%dsp-level:0
\[
(4)\quad\textbf{NEW deficit 1' (chiral matter assignment).}\ \ \text{Which species sits in which gauge representation on which chiral ray is not derived. The SM pattern}\ —\ \text{left}\ \mathrm{SU}(2)_L\ \text{doublets, right singlets}\ [\mathrm{FT.81}(2)]\ —\ \text{requires a rule correlating the}\ \hat\chi\text{-grade with the gauge representation and with the weak coupling itself}\ [\mathrm{FT.123}\ \text{Phys: the grade here is weak-blind}].\ \text{The algebraic}\ L/R\ \text{naming is fixed by the same rule}\ [\mathrm{FT.122}(6);\ \text{Def-H.33 canonicity Rmk}];
\]

%dsp-level:0
\[
(5)\quad\textbf{Item 20 deficit 11 (spin–statistics) — UNCHANGED.}\ \ \text{A single-chain chiral grade is not an exchange operator on two-chain states; FT.80}(5)\text{'s inventory still contains no exchange operator. Equating the grade's}\ \mathbb Z_2\ \text{with the permutation sign's}\ \mathbb Z_2\ \text{is a kind confusion, barred by the FT.89}(2)\ \text{discipline}\ —\ \text{the same species as the torsor/ray distinction of FT.122}(7);
\]

%dsp-level:0
\[
(6)\quad\textbf{the rest of the ledger:}\ \ \text{deficits}\ 2\ (\text{Yukawa}),
\]

%dsp-level:0
\[
\ 3\ (\text{generations/flavour}),
\]

%dsp-level:0
\[
\ 4\ (\text{CKM/PMNS})\ \text{remain open, now downstream of}\ 1'\ \text{rather than of}\ 1;
\]

%dsp-level:0
\[
\qquad \text{deficits}\ 7,
\]

%dsp-level:0
\[
8,
\]

%dsp-level:0
\[
9,
\]

%dsp-level:0
\[
10,
\]

%dsp-level:0
\[
12\ \text{unchanged}\ [\mathrm{FT.124}(6)\text{–}(7):\ \text{the mass ceilings are properties of the pole-side and colour-side reads}];
\]

%dsp-level:0
\[
(7)\quad\textbf{the reversal accounting:}\ \ \text{every prior obstruction tag remains frozen and valid on its own locus}\ [\mathrm{FT.124},\ \text{clause by clause}];
\]

%dsp-level:0
\[
\ \text{no result of Tiers F.1–F.11–Tier F.38 is overturned — the closure is by relocation, not by refutation}.
\]

\textbf{Req.} FT.122, FT.123, FT.124 (this tier); FT.79, FT.80, FT.81, FT.84 (Tiers F.26–F.28); FT.87(8), FT.89(2) (Tiers F.29–F.30); FT.93(7) (Tiers F.31–F.34); FT.117, FT.119 (Tier F.38); item 20, item 21(i), Def-H.38, Def-H.34; exact bookkeeping; machine: \texttt{checks/\allowbreak{}verify\_\allowbreak{}H\_\allowbreak{}chirality.py} §D.

\textbf{Pf.}

%dsp-level:0
\[
(1):\ \mathrm{FT.122}(1)\text{–}(5)+\mathrm{FT.123}(1)\text{–}(2)\ \text{discharge the defining properties}\ [\text{machine}\ \S\mathrm{A}\text{–}\S\mathrm{B}];
\]

%dsp-level:0
\[
\qquad \text{residue}:\ \text{item 21(i)'s enumeration stands — no frame section is constructed anywhere in Tier F.39}
\]

%dsp-level:0
\[
(2):\ \mathrm{FT.81}(4)\ \text{verbatim:}\ \text{``no chiral projector, no Weyl symbol and no hypercharge-assignment read''};
\]

%dsp-level:0
\[
\qquad \text{Tier F.39 supplies the first; the second and third are deficit}\ 1'/2/3\ \text{content}
\]

%dsp-level:0
\[
(3):\ \operatorname{im}P_R=\mathbb C|R\rangle\ [\mathrm{FT.122}(2)\text{–}(3):\ \text{rank}\ 1];\qquad\text{species: deficits 2/3 unchanged}
\]

%dsp-level:0
\[
(4):\ \mathrm{FT.81}(2)\ \text{records the SM chiral pattern as external data};
\]

%dsp-level:0
\[
\qquad \mathrm{FT.123}\ \text{Phys}:\ [W_i\otimes I_\omega,\ I\otimes P_{L/R}]=0\ \text{is blindness, not coupling};
\]

%dsp-level:0
\[
\qquad \neg\exists\ \text{assignment rule in Tiers F.1–F.11–Tier F.39}
\]

%dsp-level:0
\[
(5):\ \text{grade}\ \mathbb Z_2=\{\pm\hat\chi\}\ \text{on}\ \mathcal H_\omega(c)\ (\text{one chain});
\]

%dsp-level:0
\[
\qquad \text{exchange sign}\ \mathbb Z_2\ \text{would act on}\ \mathcal H_\omega(c_1)\otimes\mathcal H_\omega(c_2)\ \text{by}\ c_1\leftrightarrow c_2\ —\ \text{no such operator in the inventory}\ [\mathrm{FT.80}(5)];
\]

%dsp-level:0
\[
\qquad \text{kind confusion}\ [\mathrm{FT.89}(2);\ \mathrm{FT.122}(7)\ \text{precedent}]
\]

%dsp-level:0
\[
(6),(7):\ \mathrm{FT.124}\ \text{clauses verbatim};\qquad\text{no contradiction with any frozen tag}.\ \blacksquare
\]

\textbf{Phys.} Chirality exists in the corpus: an exact, canonical-up-to-naming $\mathbb Z_2$ grade on the chain's orientation doublet, with projectors that commute with the weak action, are exchanged by orientation reversal, and walk past all seven prior obstructions because none of them was ever about this space. The ledger keeps its discipline: one item closed at precisely the level achieved — the projector, not the coupling, not the section, not the species; one sharper item opened in its place; nothing else moved. That is the shape the deficit-list discipline expects from a real closure.

\vspace{2pt}\noindent\rule{\textwidth}{0.4pt}\vspace{2pt}

\medskip
\subsection*{\texorpdfstring{B.4.2\quad BM.143 — The house cleaning: the void residues closed, the phantom gaps reclassified, the calibration register consolidated, and the Part-5 requirements named}{BM.143 B.4.2}}
\addcontentsline{toc}{subsection}{B.4.2 BM.143}
\phantomsection\label{tag:BM-143}
{\small \textit{Record.}\quad Source: Part 4 (Part4\_Baryomesonic\_era.md); tier: Tier BM.23 — The Closing Tier: the Haar–Poisson Duality and the Lapse, the Selected Curvature Composition, the Momentum-Space Resolvent, the Charge and Statistics Censuses, the Two-Chain Ensemble, the Seed-Tensoring Mechanism, the Convention Declarations, the Ceiling Study, the Jensen Residue, the Pair Amplitude, and the Era Closure; canonical tag \texttt{BM.143}. Reproduced from the revision-1 canon; a cross-reference stub marks the former locus in the main corpus.}
\medskip
\textbf{Thm.}

%dsp-level:0
\[
(1)\quad(\mathrm{R5})\ \textbf{closed as void — the interior multipole content}\ \iota\ \textbf{is not a missing construction:}
\]

%dsp-level:0
\[
\qquad \textbf{four derived grounds:}
\]

%dsp-level:0
\[
(\mathrm i)\quad\text{the trace-only law:}\ \ \text{every Hermitian slot operator on a sealed singlet reads only its trace}\ —\ \mathfrak T(H)
\]

%dsp-level:0
\[
= 3a\ \ \text{on the covariant class};
\]

%dsp-level:0
\[
\ \text{the meson reads}\ 0\ \ \text{including the trace part}
\]

%dsp-level:0
\[
\qquad [\mathrm{BM.1}(2);\ \mathrm{FT.117}(1)\text{--}(3)];
\]

%dsp-level:0
\[
(\mathrm{ii})\quad\text{the void theorem:}\ \ \iota\ \ \text{is a functional of an interior charge density, and}\ \neg\exists\ \text{an interior density in the inventory:}\ \ \text{the question ``shell or ball?'' has no corpus referent}
\]

%dsp-level:0
\[
\qquad [\mathrm{BM.87}(3)];
\]

%dsp-level:0
\[
(\mathrm{iii})\quad\text{the end signs carry no position datum:}\ \ q_{\mathrm{core}}
\]

%dsp-level:0
\[
= \tfrac12(\sigma_s-\sigma_t)\ \ \text{is endpoint data; a neutral-class internal-multipole self-energy is}\ \neg\text{derived at the monopole order}
\]

%dsp-level:0
\[
\qquad [\mathrm{D3.4};\ \mathrm{BM.46}(2)];
\]

%dsp-level:0
\[
(\mathrm{iv})\quad\text{the envelope discipline:}\ \ \text{the coarse-graining reads are chart envelopes,}\ \neg\text{ledger distances}
\]

%dsp-level:0
\[
\qquad [\mathrm{Def\text{-}BM.4}(5);\ \mathrm{BM.13}];
\]

%dsp-level:0
\[
\textbf{the information-flux row:}
\]

%dsp-level:0
\[
\qquad \text{sealing restricts every read to traces, ends and junctions: one dangling attachment zeroes the colour sector}\ [\mathrm{FT.19}];
\]

%dsp-level:0
\[
\qquad \neg\text{a triplet pole, physical norm}\ 0\ [\mathrm{CT3.22}];
\]

%dsp-level:0
\[
\boxed{\ \textbf{the internal constituent states of a sealed composite are}\ \neg\textbf{knowable by theorem, and}\ (\mathrm{R5})\ \textbf{is closed as void: coarse graining bounds the information flux, and no construction could exist inside the bound}\ }
\]

%dsp-level:0
\[
(2)\quad(\mathrm{R2})\ \textbf{reclassified — the loop realization is a phantom gap:}
\]

%dsp-level:0
\[
\qquad \text{the named object was the vertex-figure locus}\ \Theta_R=\delta_v\ \ \text{as a ledger configuration}
\]

%dsp-level:0
\[
\qquad [\mathrm{BM.67}(9),(10);\ \mathrm{BM.82}(6);\ \mathrm{Def\text{-}BM.19}(7);\ \mathrm{Census\ FT.63}(4)\ \mathrm{Rmk}];
\]

%dsp-level:0
\[
\textbf{the reclassification:}
\]

%dsp-level:0
\[
\qquad \text{it is registration data of the same species as the edge phases}\ —\ \text{its alphabet is derived and its per-configuration value is determined at crystallization}
\]

%dsp-level:0
\[
\qquad [\text{A §0.2a};\ \mathrm{CT0.9};\ \mathrm{Def\text{-}EM.20};\ \mathrm{EM.20};\ \mathrm{BM.127}(2):\ \widehat{ds}^{\,2}(e)=\kappa_e\ \text{registration data}];
\]

%dsp-level:0
\[
\textbf{every physical read is derived without the value:}
\]

%dsp-level:0
\[
\qquad \text{spin}\ j=\frac{|k|}2\ [\mathrm{Def\text{-}H.38};\ \mathrm{BM.113}(3)];
\]

%dsp-level:0
\[
\qquad \text{parity}\ \varepsilon=(-1)^{|k|}\ [\mathrm{BM.113}(2);\ \mathrm{BM.127}(1)];
\]

%dsp-level:0
\[
\qquad \text{concrete spin}\ \Theta=4j\ [\mathrm{BM.124}(1)];
\]

%dsp-level:0
\[
\text{angular momentum from the metric:}\ \ \partial_\mu M^{\mu\nu\lambda}=0,
\]

%dsp-level:0
\[
\qquad p^2=p_r^2+\frac{L^2}{r^2},
\]

%dsp-level:0
\[
\qquad \text{the}\ \ell\text{-selection census}
\]

%dsp-level:0
\[
\qquad [\mathrm{BM.42}(1)\text{--}(3);\ \mathrm{BM.9}(3);\ \mathrm{BM.11}];
\]

%dsp-level:0
\[
\text{the defect ladder is uniform over the realizations}\ [\mathrm{Def\text{-}BM.19}(3);\ \mathrm{BM.124}(3)];
\]

%dsp-level:0
\[
\qquad \text{and}\ \text{A-BM.11}\ \text{converts the convention-tagged clauses to declared data}\ [\mathrm{BM.139}(4)];
\]

%dsp-level:0
\[
\boxed{\ (\mathrm{R2})\ \textbf{is closed as}\ \neg\textbf{a gap: no structure is missing, no selection principle is absent — the realization is the ledger's own registration act, and the laws hold uniformly over it}\ }
\]

%dsp-level:0
\[
(3)\quad(\mathrm{R10})\ \textbf{reclassified — the connector winding is registration data:}
\]

%dsp-level:0
\[
\qquad d=\Theta(P)\ge0
\]

%dsp-level:0
\[
\qquad [\mathrm{BM.2}(2);\ \mathrm{BM.129}(3);\ \mathrm{FT.34}(3):\ \Theta(P)=\Theta(\bar P)];
\]

%dsp-level:0
\[
\textbf{the alphabet is derived}\ (d\in\mathbb Z_{\ge0};\ \ \Theta\in2\mathbb Z\ \ \text{on shell}\ [\mathrm{FT.22}]);
\]

%dsp-level:0
\[
\qquad \textbf{the parity law is derived}\ \bigl(\varepsilon_M=(-1)^d\ \ \text{with a closed connector}\ [\mathrm{BM.129}(3)]\bigr);
\]

%dsp-level:0
\[
\qquad \textbf{the computed cells are derived}\ (\text{an open connector carries}\ S=0\implies d=0\implies\varepsilon_M=+1\ [\mathrm{Def\text{-}H.32}(4);\ \mathrm{BM.129}(3)]);
\]

%dsp-level:0
\[
\boxed{\ (\mathrm{R10})\ \textbf{is closed as registration-species: the per-class value is fixed at crystallization, exactly as}\ (\mathrm{R2})\text{'s}\ }
\]

%dsp-level:0
\[
(4)\quad\textbf{the calibration register consolidated (}\ \text{§Cal-BM}\ \textbf{):}
\]

%dsp-level:0
\[
\qquad \textbf{three items are calibration data,}\ \neg\textbf{structural gaps:}
\]

%dsp-level:0
\[
\begin{array}{l|l|l}\text{slot} & \text{the datum} & \text{the locus}\\\hline(\mathrm{C\text{-}BM.1}) & \rho_{\mathrm B}=\hat\mu'^2\mathfrak h^2/(\ln2)^2\in\mathfrak q^0:\ \ \text{the bridge ratio; the slot declared, the value}\ \neg\text{issued} & [\mathrm{A\text{-}BM.9};\ \mathrm{BM.69}(5);\ \mathrm{BM.98};\ (\mathrm{R1})\ \text{reclassified}]\\\hline(\mathrm{C\text{-}BM.2}) & e_0:\ \ \text{the numerical value of the electric coupling; one calibration datum decides the branch} & [\mathrm{EM.25}(7);\ \mathrm{BM.88}(8);\ \mathrm{BM.62}(12);\ \mathrm{BM.84}(9);\ \mathrm{BM.133}(6)]\\\hline(\mathrm{C\text{-}BM.3}) & \mu_\nu,\ \mu,\ \rho_\ell,\ \rho_\nu:\ \ \text{the channel scales, inherited from}\ \mathrm{Cal\text{-}F},\ \neg\text{declared in this part} & [\mathrm{FT.145}(3);\ \mathrm{BM.83}(5),(7);\ \text{§Cal-BM's inherited row}]\end{array}
\]

%dsp-level:0
\[
\textbf{the parameter count is unchanged:}
\]

%dsp-level:0
\[
\qquad \neg\text{a new free parameter};
\]

%dsp-level:0
\[
\qquad \mathfrak h\ \ \text{remains the single dimensionful scale}
\]

%dsp-level:0
\[
\qquad [\mathrm{P1}];
\]

%dsp-level:0
\[
\qquad \neg\text{a value is issued by any row}\ [\mathrm{P5}];
\]

%dsp-level:0
\[
\textbf{the register is a consolidation,}\ \neg\textbf{an addition:}
\]

%dsp-level:0
\[
\qquad \text{§Cal-BM's ``None introduced by this part'' row stands as a parameter count}\ [\mathrm{A\text{-}BM.9}\text{'s clause re-cited}];
\]

%dsp-level:0
\[
(5)\quad(\mathrm{R4}')\ \textbf{reclassified — the pair correlations are void for this era:}
\]

%dsp-level:0
\[
\qquad \textbf{the uniformity census:}\ \ \text{every ensemble-dependent derived clause of this part holds either measure-free or uniformly over the formed reads:}
\]

%dsp-level:0
\[
w_{\mathrm{mat}}=\tfrac13\textstyle\sum_ia_i\in\bigl[\tfrac13,1\bigr]\ \ \text{measure-free}
\]

%dsp-level:0
\[
\qquad [\mathrm{BM.105}(2)\text{--}(4)];
\]

%dsp-level:0
\[
\qquad \text{the ensemble-independent bar}
\]

%dsp-level:0
\[
\qquad [\mathrm{BM.91}(4)\text{--}(8)];
\]

%dsp-level:0
\[
\qquad \text{the uniform closure}
\]

%dsp-level:0
\[
\qquad [\mathrm{BM.92}(1),(2)];
\]

%dsp-level:0
\[
\text{the two-sided parameter-free bar}
\]

%dsp-level:0
\[
\qquad [\mathrm{BM.99}(1)\text{--}(4)];
\]

%dsp-level:0
\[
\qquad \text{the marked pair specialization is}\ \neg\text{adopted as the measure}
\]

%dsp-level:0
\[
\qquad [\mathrm{BM.94}(4),(9):\ \text{marked}];
\]

%dsp-level:0
\[
[\text{machine: the clause scan of the ensemble-dependent rows of}\ \mathrm{BM.19},\ \mathrm{BM.21},\ \mathrm{BM.89},\ \mathrm{BM.91},\ \mathrm{BM.92},\ \mathrm{BM.94},\ \mathrm{BM.99},\ \mathrm{BM.105},\ \mathrm{BM.137}\ \text{against their cited scopes}];
\]

%dsp-level:0
\[
\boxed{\ (\mathrm{R4}')\ \textbf{is reclassified as a Part-5 ensemble-refinement datum: no baryomesonic clause requires correlations beyond per-pair independence}\ }
\]

%dsp-level:0
\[
(6)\quad\textbf{the era's final residue census (the ledger balanced):}\qquad\begin{array}{l|l|l}\text{residue} & \text{the final disposition} & \text{the locus}\\\hline(\mathrm{R1})\ \ \rho_{\mathrm B}\text{'s value} & \textbf{reclassified: calibration slot, value}\ \neg\text{issued} & [(4);\ \mathrm{A\text{-}BM.9}]\\\hline(\mathrm{R2})\ \ \text{the loop realization} & \textbf{reclassified: phantom gap — registration data} & [(2);\ \mathrm{A\text{-}BM.11}]\\\hline(\mathrm{R3})\ \ \text{the flavour composition} & \textbf{split: the winding assignment derived; the occupancy is the generation-selection question, named for Part 5} & [\mathrm{BM.136}(4),(5)]\\\hline(\mathrm{R4}')\ \ \text{the pair correlations} & \textbf{reclassified: Part-5 ensemble datum, void for BM} & [(5)]\\\hline(\mathrm{R5})\ \ \iota\ \text{and the multipoles} & \textbf{closed as void} & [(1)]\\\hline(\mathrm{R6})\ \ \text{the Jensen residue} & \textbf{closed: computed} & [\mathrm{BM.141}]\\\hline(\mathrm{R7}')\ \ \text{the exchange at}\ n\ge4 & \textbf{closed: constructed at every}\ n & [\mathrm{BM.136}(3)]\\\hline(\mathrm{R8})\ \ \text{the 2-chain sector masses} & \textbf{closed: computed, with the ensemble binding formed} & [\mathrm{BM.137}(3),(4)]\\\hline(\mathrm{R9}')\ \ \text{the cross-section} & \textbf{closed: constructed, symbolic} & [\mathrm{BM.142}(5)]\\\hline(\mathrm{R9}'')\ \ \text{the radiation legs} & \textbf{closed at the vertex level: the supply vertex constructed} & [\mathrm{BM.142}(2)]\\\hline(\mathrm{R10})\ \ \text{the connector winding} & \textbf{reclassified: registration data} & [(3)]\\\hline\mathsf g,\ N & \textbf{closed: selected and derived} & [\mathrm{BM.133}(4);\ \mathrm{BM.132}(7)]\\\hline e_0,\ \mu_\nu & \textbf{reclassified: calibration} & [(4)]\\\hline m_\nu\ \text{absolute},\ i_\nu,\ \text{the mixing matrix} & \textbf{stand barred: deficits 6, 4 — calibration and bar species,}\ \neg\text{BM constructions} & [\mathrm{BM.96}(3);\ \mathrm{BM.77}(2);\ \mathrm{FT.145}(3);\ \mathrm{BM.75}(6)]\\\hline\text{the generation count as a selection} & \textbf{stands: cross-era selection species, named for Part 5} & [\mathrm{BM.60}(1)]\\\hline\text{the two mass ceilings} & \textbf{closed as adjudicated: deficit 12} & [\mathrm{BM.140}(7)]\\\hline\text{the}\ \vee^2\ \text{sector choice} & \textbf{declared} & [\mathrm{BM.139}(5);\ \text{A-BM.13}]\end{array}
\]

%dsp-level:0
\[
(7)\quad\textbf{what Part 5 (the Atomic era) requires, and what it receives:}
\]

%dsp-level:0
\[
\qquad \textbf{received from this part (the supply rows):}
\]

%dsp-level:0
\[
\begin{array}{l|l}\text{the requirement} & \text{the supplying loci}\\\hline\text{the unit-class baryon} & \text{the unit-class row}\ [\mathrm{BM.26};\ \mathrm{BM.28}(4)];\quad m_n/m_p<\frac{21}{20}\ \text{and the sub-critical order}\ [\mathrm{BM.84};\ \mathrm{BM.91}(2);\ \mathrm{BM.99}];\quad\text{the decay closure}\ [\mathrm{BM.45};\ \mathrm{BM.63};\ \mathrm{BM.71};\ \mathrm{BM.92}]\\\hline\text{the}\ \ell\text{-channel carrier} & \text{the}\ \ell\text{-channel naming row}\ [\mathrm{Def\text{-}BM.6}(1);\ \mathrm{BM.23}(5):\ \text{the forced carrier}];\quad\text{the symbolic mass row}\ [\mathrm{BM.43}(2)];\quad\text{the channel ratios}\ [\mathrm{FT.113}]\\\hline\text{the EM binding} & \text{the two-body potential and the hydrogenic rows}\ [\mathrm{BM.48}(1),(2)];\quad\text{the binding category bar}\ [\mathrm{BM.48}(4)];\quad\text{the complete running coupling}\ [\mathrm{BM.88}(4);\ \mathrm{BM.132};\ \mathrm{BM.133}(6)]\\\hline\text{the scattering reads} & \text{the resolvent, the phase space and the cross-sections}\ [\mathrm{Def\text{-}BM.20};\ \mathrm{BM.134};\ \mathrm{BM.142}]\\\hline\text{the composite grammar} & \text{the seals, the statistics and the exchange structure at every}\ n\ [\mathrm{Def\text{-}BM.1};\ \mathrm{Def\text{-}BM.16};\ \mathrm{BM.113};\ \mathrm{BM.129};\ \mathrm{BM.136}]\\\hline\text{the thresholds} & \text{the materialization condition with the derived lapse}\ [\mathrm{Def\text{-}BM.17}(3);\ \mathrm{BM.115};\ \mathrm{BM.132}(5)]\end{array}
\]

%dsp-level:0
\[
\textbf{to construct in Part 5 (the demand rows):}
\]

%dsp-level:0
\[
\qquad (\mathrm{P5.1})\ \ \text{the EM binding channel}\ —\ \neg\exists\ \text{an EM sealing grammar in the inventory; the lepton--baryon bound state is a dynamics-of-}Q\ \text{object}
\]

%dsp-level:0
\[
\qquad [\mathrm{BM.48}(3);\ \mathrm{BM.54}(3);\ \mathrm{BM.6}(5)];
\]

%dsp-level:0
\[
(\mathrm{P5.2})\ \ \text{the lepton mass read:}\ E_{\mathrm{int}}^{(\ell)},
\]

%dsp-level:0
\[
\ E_{\mathrm{int}}^{(\nu)}\ \ \neg\text{computed at any floor};
\]

%dsp-level:0
\[
\ \ m_r\ \text{of a lepton--baryon pair requires it}
\]

%dsp-level:0
\[
\qquad [\mathrm{BM.43}(2);\ \mathrm{BM.48}(2);\ \text{or the}\ (\mathrm{C\text{-}BM.1})\ \text{slot}];
\]

%dsp-level:0
\[
(\mathrm{P5.3})\ \ \text{the dynamics rates and the tick-rate/field coupling}\ —\ \text{emergent time is a named prerequisite}
\]

%dsp-level:0
\[
\qquad [\mathrm{BM.6}(5);\ \mathrm{FT.131}(3);\ \mathrm{BM.131}(7)\text{'s reparametrization read is the formed piece}];
\]

%dsp-level:0
\[
(\mathrm{P5.4})\ \ \text{the frame section and a global Dirac-type operator, both unconstructed}
\]

%dsp-level:0
\[
\qquad [\mathrm{BM.6}(8);\ \mathrm{Def\text{-}H.34};\ \mathrm{FT.93}];
\]

%dsp-level:0
\[
(\mathrm{P5.5})\ \ \text{the generation/flavour selection}
\]

%dsp-level:0
\[
\qquad [(\mathrm{R3})\text{'s occupancy half}];
\]

%dsp-level:0
\[
\qquad (\mathrm{P5.6})\ \ \text{item 16(c)'s remaining cells}
\]

%dsp-level:0
\[
\qquad [\mathrm{BM.139}(6)];
\]

%dsp-level:0
\[
(\mathrm{P5.7})\ \ (\mathrm{R4}')\text{'s ensemble refinement}\ [(5)];
\]

%dsp-level:0
\[
\qquad (\mathrm{P5.8})\ \ \text{the}\ n\ge4\ \text{sealing-grammar cells}
\]

%dsp-level:0
\[
\qquad [\mathrm{BM.31}(2);\ \mathrm{BM.33};\ \mathrm{BM.136}(3)];
\]

%dsp-level:0
\[
\qquad (\mathrm{P5.9})\ \ \text{the finite-}\mathfrak h\ \text{residue at the clean-bifurcation loci}
\]

%dsp-level:0
\[
\qquad [\mathrm{FT.66}(3)];
\]

%dsp-level:0
\[
(\mathrm{P5.10})\ \ \text{the embedding datum, if a rotor read is ever formed}\qquad[\mathrm{BM.6}(7);\ \mathrm{FT.87}(6)];
\]

%dsp-level:0
\[
(8)\quad\textbf{the Census annex row:}
\]

%dsp-level:0
\[
\qquad \text{the Census is frozen by its charter}\ —\ \text{the}\ 25\ \text{Category-B cells and their consumption block are data, and the era's final state is a ledger record}
\]

%dsp-level:0
\[
\qquad [\text{the separation charter};\ \mathrm{BM.5}(2);\ \text{machine: the cell inventory unchanged}];
\]

%dsp-level:0
\[
\textbf{the era's closing census is}\ (6)\text{'s table, carried at the ledger}\;
\]

%dsp-level:0
\[
(9)\quad\textbf{status:}
\]

%dsp-level:0
\[
\qquad \textbf{derived:}\ (1)\text{--}(3)\text{'s ground rows},\ (5)\text{'s uniformity census};
\]

%dsp-level:0
\[
\qquad \textbf{adjudicated:}\ (1)\text{--}(6);
\]

%dsp-level:0
\[
\qquad \textbf{named:}\ (7);
\]

%dsp-level:0
\[
\qquad \neg\text{a declaration beyond the four of}\ \mathrm{BM.139},
\]

%dsp-level:0
\[
\qquad \neg\text{a P4 addition},
\]

%dsp-level:0
\[
\qquad \neg\text{a measured value}\ [\mathrm{P5}].
\]

\textbf{Req.} A4 (Part 1); D0.8, D3.4 (Part 1); CT0.9, CT3.22 (Parts 1–2); Def-EM.20, EM.20, EM.25(7) (Part 2); P1, P5 (Parts 1–2); Def-H.4, Def-H.32(4), Def-H.34, Def-H.38, FT.22, FT.66(3), FT.93, FT.113, FT.117(1), FT.131(3), FT.145(3) (Part 3); Def-BM.1, Def-BM.4(5), Def-BM.6(1), Def-BM.16, Def-BM.17(3), Def-BM.19(3), BM.1(2), BM.2(2), BM.5(2), BM.6(5), BM.6(7), BM.6(8), BM.9(3), BM.11, BM.13, BM.19, BM.21, BM.23(5), BM.26, BM.28(4), BM.31(2), BM.33, BM.37(2), BM.42, BM.43(2), BM.45, BM.46(2), BM.48, BM.54(3), BM.60(1), BM.62(12), BM.63, BM.67(9), BM.69(5), BM.71, BM.75(6), BM.77(2), BM.82(6), BM.83(5), BM.84, BM.87(3), BM.88(4), BM.88(8), BM.89, BM.91, BM.92, BM.94(4), BM.94(9), BM.96(3), BM.98, BM.99, BM.105, BM.113, BM.115, BM.124(1), BM.124(3), BM.127, BM.129(3), BM.131(7), BM.132(5), BM.133(6), BM.134, BM.136, BM.137, BM.139, BM.141, BM.142, Def-BM.20, A-BM.9, A-BM.11 (this part); Census FT.63 (row-wise); the clause-scope inventory scan, the document pins [machine: \texttt{checks/\allowbreak{}verify\_\allowbreak{}H\_\allowbreak{}bm\_\allowbreak{}closing\_\allowbreak{}tier.py} §L].

\textbf{Pf.}

%dsp-level:0
\[
(1)\text{--}(3):\ \text{the ground rows re-cited at their own loci with the registration-species citations}\ [\text{A §0.2a};\ \mathrm{BM.127}(2)];
\]

%dsp-level:0
\[
\qquad \text{machine: the pin census}
\]

%dsp-level:0
\[
(4):\ \text{the three slots assembled from}\ \mathrm{A\text{-}BM.9},\ \mathrm{EM.25}(7),\ \mathrm{FT.145}(3),\ \mathrm{BM.83}(5),(7);
\]

%dsp-level:0
\[
\qquad \text{the count row:}\ \neg\text{a new symbol is introduced}\ [\text{machine: the packaging census}];
\]

%dsp-level:0
\[
(5):\ \text{the clause scan: each ensemble-dependent row of the cited tags carries a measure-free bound, a uniformity clause, or the marked specialization}\ [\text{machine}];
\]

%dsp-level:0
\[
(6)\text{--}(9):\ \text{the census assembly against the frozen residue tables of}\ \mathrm{BM.86}(6),
\]

%dsp-level:0
\[
\ \mathrm{BM.93}(6),
\]

%dsp-level:0
\[
\ \mathrm{BM.100}(6),
\]

%dsp-level:0
\[
\ \mathrm{BM.107}(7),
\]

%dsp-level:0
\[
\ \mathrm{BM.118}(7),
\]

%dsp-level:0
\[
\ \mathrm{BM.130}(7)\ \text{and the ledger pointers}\ [\text{machine: the row-by-row discharge census}].\ \blacksquare
\]

\textbf{Phys.} The era's books are cleaned before they are shelved. One residue closes as void: the interior multipole content was never a missing object — the trace-only law, the absent interior density, the positionless end signs and the envelope discipline together make the internal charge resolution of a sealed composite unknowable by theorem, and coarse graining is exactly what bounds the flux. Two residues turn out to be phantom gaps: the loop realization and the connector winding are registration data of the same species as the edge phases, determined at crystallization, with every physical law — spin, parity, angular momentum from the metric, the defect ladder — derived uniformly over all realizations. Three items move to the calibration register with their values unissued and their parameter count unchanged, and the pair-correlation residue is reclassified after a uniformity census shows no clause of the era ever needed it. The final table balances every residue the era named against a disposition, and the atomic era's shopping list is written: six supply rows it receives and ten demand rows it must build, the binding channel first.

\medskip
\subsection*{\texorpdfstring{B.4.3\quad FT.84 — The deficit table and the two-ceiling synthesis (audit pass D)}{FT.84 B.4.3}}
\addcontentsline{toc}{subsection}{B.4.3 FT.84}
\phantomsection\label{tag:FT-84}
{\small \textit{Record.}\quad Source: Part 5 (Census.md); tier: The cells; canonical tag \texttt{FT.84}. Reproduced from the revision-1 canon; a cross-reference stub marks the former locus in the main corpus.}
\medskip
\textbf{Thm.} Under Def-H.29 and FT.81–FT.83: $(1)$ \textbf{The ranked deficit list} (load-bearing order; class: (i) missing declaration, (ii) structural no-go already proven, (iii) genuinely open):

%dsp-level:0
\[
\begin{array}{c|l|l|l}
\# & \text{deficit} & \text{corpus status} & \text{class}\\
\hline
1 & \text{chirality (Weyl/projector)} & \text{no read; }\omega\text{-vocabulary only [D2.6]} & \text{(iii) open — root deficit}\\
2 & \text{Yukawa sector / fermion masses} & \text{no read; no P4 parameter} & \text{(iii) open, constrained by P1/P4}\\
3 & \text{three generations / flavour} & \text{no read } & \text{(iii) open}\\
4 & \text{CKM/PMNS mixing, CP phase} & \text{no read} & \text{(iii) open; needs 1 and 3}\\
5 & \text{gauge-anomaly structure} & \text{not poseable} & \text{(ii) blocked by 1}\\
6 & \text{neutrino species/mass/mixing} & \text{no read} & \text{(iii) open; needs 1 and 3}\\
7 & \Lambda_{\mathrm{QCD}}\ \text{/ any second scale} & \text{no read} & \text{(ii) no-go [A §0.3, P1, P5]}\\
8 & \text{Higgs as }4\text{ real dof} & \dim F_H=1 & \text{(ii) no-go [CT0.12 allocation]}\\
9 & \text{fractional confined charges} & \text{ternary read} & \text{(ii) partial no-go [EM.9, CT3.22]}\\
10 & |q|=2\ \text{observable states} & \text{single-chain read} & \text{(iii) open; needs a multi-chain read}\\
11 & \text{spin-statistics} & \text{external} & \text{(iii) open [FT.80(5)]}\\
12 & \text{SM masses (any species)} & \text{two ceilings proven} & \text{(ii) no-go [FT.76(3), FT.79(2)]}
\end{array}
\]

$5$ structural no-goes or blocked items (each with a cited no-go), $7$ open. Deficit $1$ is the root of $4,5,6$ and $11$; deficits $7$–$9$ and $12$ are derived impossibilities within the current inventory. $(2)$ \textbf{The two-ceiling synthesis.} Both mass audits failed for the \textit{same} structural reason at different loci: the pole side is bounded because a cone defect is an angle ($m^2_{\mathrm{pole}}\in[0,4\hat\mu'^2]$, widest ratio $6.39$ [FT.76(3)]), and the cover side is bounded because a mass counted in crystallized edges is bounded by the cap census ($K\le3V/2$, $|\mathcal C|\le3V/4$, widest ratio $6$ on the cube [FT.79(2)]). Deficit $2$ names the deeper cause: \textbf{there is no Yukawa sector}, so no mechanism in the inventory generates fermion mass at all — the two ceilings bound the \textit{reads that exist}, and the SM's masses live in a sector the corpus does not contain. The mass negatives of FT.77/FT.79 are consequences of deficit $2$, not evidence against the reads themselves. $(3)$ \textbf{What is reproduced, stated positively} (each with its citation): the gauge block dimensions and their sum $8+3+1=12$ [CT0.14, CT0.12, CT1.1]; $\mathfrak{su}(2)$ and $\mathfrak{su}(3)$ closure with correct roots, normalizations and no ninth gluon [FT.82(2)–(3); CT3.8, CT3.13]; the derived colour count $3$ [CT3.20]; the full electroweak mixing structure with an exactly massless photon and $m_W=m_Z\cos\theta_W$, $1/e^2=1/g^2+1/g'^2$ [EW.20/23/24/35]; the vev's doublet-and-singlet quantum numbers [EW.20, EW.32]; confinement as a theorem with norm-zero isolated triplets [CT3.22, CT3.18, FT.19]; the baryon/meson sealing grammar [D3.9, CT3.15, CT3.17, FT.34(5)]; integer observable charge with the meson set exactly equal to the ternary read [EM.9, D3.4, FT.81(3)]; the spin ladder $j=|k|/2$ with $\mathrm{Sym}^k$ carriers [CT0.7a, EW.32, FT.52–FT.60, Def-H.38]; a derived Newton constant from the single scale [P3, CT2.9]; the fermion/boson \textit{content-type} split [D1.5, Def-H.4]; and, beyond the SM, a dark sector with exact equations of state [D4.1–D4.3, CT1.25, CT4.1, CT4.5, FT.71]. $(4)$ \textbf{Verdict.} The corpus is \textbf{not} a Standard Model completion and does not claim to be: it is a derivation of the gauge-and-geometry skeleton from one scale and sixteen registers, with the chiral matter sector, the flavour sector and every dimensionful mass scale left outside the inventory. Its completeness fraction on the audited entry list is $13/24$, concentrated as $4/4$ on gauge species and $0.19$ on the fermion sector. No identification is made; Open Gaps 5/12 are unchanged and now carry a sharper form: \textit{the mass question is not open but blocked from two sides (deficit $2$ above, ceilings below), while the gauge question is substantially answered.} $(5)$ \textbf{Consequences.} (a) Any further mass-read search must first extend P4 with a flavour/Yukawa sector or accept that fermion masses have no read — searching for the SM hierarchy inside the existing reads is exhausted by FT.77/FT.79 and this clause. (b) The large-$V$ route of FT.79(2) is now demoted: even at $V\ge1200$ the resulting ladder would be an \textit{even-integer} spectrum with no chirality and no flavour, so it could match a hierarchy's \textit{range} and still not its \textit{structure}. (c) The highest-value open direction is deficit $1$: a chirality read would unblock $4,5,6,11$ at once, and D2.6's $\omega$-action is the only existing hook — whether $\omega$ (or the FT.78(6) cover chirality) can be promoted from vocabulary to a projector is a well-posed question with no answer in the current inventory.

\textbf{Req.} Def-H.29; FT.11, FT.12, FT.19, FT.20, FT.26, FT.27, FT.30, FT.34, FT.41, FT.52–FT.60, FT.71, FT.76–FT.83; A §0.2, §0.3; D1.5, D2.6, D3.4, D3.9, D4.1–D4.3; CT0.11–CT0.14 (Part 1); CT2.9, CT3.15, CT3.17, CT3.18, CT3.20, CT3.22 (\textbf{C/F}); CT4.1, CT4.4, CT4.5 (Part 2); EM.9, EM.20 (Part 2); EW.8, EW.13, EW.20, EW.23, EW.24, EW.27, EW.32, EW.33, EW.35 (Part 2); P1–P5, Def-H.38; Open Gaps 5/7/12/13/16; external SM parameter and hierarchy data (marked) [machine].

\textbf{Pf.}

%dsp-level:0
\[
(1):\ \text{each row's class follows from its cited tag}:\ \text{(ii) rows cite a theorem of non-existence or an allocation argument}
\]

%dsp-level:0
\[
\text{(iii) rows cite an inventory scan}\ [\mathrm{FT.81}(5)]\ \text{plus the absence of any blocking theorem};
\]

%dsp-level:0
\[
\qquad 1\Rightarrow4,5,6,11:\ \text{each requires a chiral read}
\]

%dsp-level:0
\[
(2):\ \mathrm{FT.76}(3):\ m^2_{\mathrm{pole}}\le4\hat\mu'^2,\ \max\ \text{ratio}\ \sqrt{4/D_{\mathrm{dodeca}}}=6.3925;
\]

%dsp-level:0
\[
\qquad \mathrm{FT.79}(2):\ K=2|\mathcal C|\le3V/2\implies\max\ \text{ratio}\ 6\ (V=8)
\]

%dsp-level:0
\[
\text{deficit}\ 2:\ \text{Yukawa count}\ 0\ [\text{scan}]\implies\neg\exists\ \text{fermion-mass mechanism};
\]

%dsp-level:0
\[
\qquad \text{the ceilings bound existing reads, not the SM's sector}
\]

%dsp-level:0
\[
(3):\ \text{conjunction of the cited tags, each machine-verified in FT.81–FT.83};
\]

%dsp-level:0
\[
\qquad (4):\ (1)+(3)\ \text{with}\ 13/24
\]

%dsp-level:0
\[
(5a):\ \mathrm{FT.77}(2),\ \mathrm{FT.79}(4)\ \text{exhaustive audits};
\]

%dsp-level:0
\[
\qquad (5b):\ \text{even-integer ladder}\ [\mathrm{FT.79}(1)];
\]

%dsp-level:0
\[
\qquad (5c):\ \mathrm{D2.6},\ \mathrm{FT.78}(6).\ \blacksquare
\]

\textbf{Phys.} Asked whether it contains the Standard Model, the corpus answers with a clean split that no amount of mass-fitting would have revealed. The forces are there: eight gluons in an algebra that closes with the right roots, three weak generators whose mixture leaves one combination exactly massless, a mixing angle and an electric coupling tied together by identities rather than fits, a confinement rule that turns fractional quarks into integer-charge mesons and baryons, and gravity's constant fixed by the single scale. The matter is not: nothing is left- or right-handed, nothing comes in flavours, nothing mixes, and no second scale exists for any of it. That is why the two mass audits failed the way they did — they were reading masses out of registers that hold geometry and combinatorics, in a document with no Yukawa sector at all. The honest summary is not "the corpus fails the Standard Model" but "it derives the gauge skeleton and leaves the chiral matter sector undeclared", with one hook — the orientation-reversal involution — where a chirality read might someday be hung.

\vspace{2pt}\noindent\rule{\textwidth}{0.4pt}\vspace{2pt}

\vspace{2pt}\noindent\rule{\textwidth}{0.4pt}\vspace{2pt}

\medskip
\subsection*{\texorpdfstring{B.4.4\quad BM.37 — The mass-route artifacts: the disjoint input alphabets and the closure of the counting trap (catalog entry)}{BM.37 B.4.4}}
\addcontentsline{toc}{subsection}{B.4.4 BM.37}
\phantomsection\label{tag:BM-37-cat}
{\small \textit\{Catalog.\}\quad The full theorem remains in the main corpus as Theorem 4.31 (Part 4); its derived clauses are consumed by later derivations. The hypothesis-level ruling is recorded here verbatim:}
\begin{quote}\small \end{quote}
\medskip
\subsection*{\texorpdfstring{B.4.5\quad BM.36 — The force-structure ladder: the mask dimensions, the transport channel, and the inter-family constraint (catalog entry)}{BM.36 B.4.5}}
\addcontentsline{toc}{subsection}{B.4.5 BM.36}
\phantomsection\label{tag:BM-36-cat}
{\small \textit\{Catalog.\}\quad The full theorem remains in the main corpus as Theorem 4.30 (Part 4); its derived clauses are consumed by later derivations. The hypothesis-level ruling is recorded here verbatim:}
\begin{quote}\small \end{quote}
\medskip
\subsection*{\texorpdfstring{B.4.6\quad FT.73(5)}{FT.73(5) B.4.6}}
\addcontentsline{toc}{subsection}{B.4.6 FT.73(5)}
{\small \textit{Pointer.}\quad The standing trap pattern — comparing a derived pure number with a measured coupling constant (the $\pi$-against-$1/137$ species) — is recorded at FT.73(5), reproduced in full at B.2.2.}
\medskip
\subsection*{\texorpdfstring{B.4.7\quad FT.77(3),(5)}{FT.77(3),(5) B.4.7}}
\addcontentsline{toc}{subsection}{B.4.7 FT.77(3),(5)}
{\small \textit{Pointer.}\quad The dimension-trap disclosures (comparing $m^2$-ratios to measured mass ratios; the amplitude-span near-coincidences) travel with the full audit at B.2.3.}
\medskip
\subsection*{\texorpdfstring{B.4.8\quad FT.89(2)}{FT.89(2) B.4.8}}
\addcontentsline{toc}{subsection}{B.4.8 FT.89(2)}
{\small \textit{Pointer.}\quad The Koide category trap (reading the derived rational $\tfrac23$ as a measured mass-ratio datum) is adjudicated at B.1.13; the supersession of its classification clause by FT.158(5) is recorded there and at Theorem 3.94.}
\medskip
\subsection*{\texorpdfstring{B.4.9\quad Census FT.99(6)}{Census FT.99(6) B.4.9}}
\addcontentsline{toc}{subsection}{B.4.9 Census FT.99(6)}
{\small \textit{Pointer.}\quad The unreduced-versus-lowest-terms comparison trap is recorded, with its integrity ledger, in the Census annex (Appendix A, cell FT.99(6)(b)).}
\medskip
\end{document}